% =====================================================================
%  Analog Game Design, lezione 1 — Introduzione al corso: docenti,
%  metodo, project work, figure dell'industria e primo glossario
%  Fonte: trascrizione della prima lezione (AGD1) del corso breve di
%  Analog Game Design, seconda edizione: una lezione online su Zoom di
%  circa tre ore, di pomeriggio, tenuta insieme dai tre docenti "cuore"
%  del corso (Daniele, Ruggero e Beatrice), con nove iscritti, una
%  pausa di una decina di minuti, molte slide e una breve dimostrazione
%  dal vivo di Tabletop Simulator.
%  Pacchetti richiesti nel preambolo del libro (già presenti nel
%  main.tex): \usepackage{booktabs}, \usepackage{tikz} con
%    \usetikzlibrary{shapes.geometric,shapes.symbols,shapes.misc,
%      arrows.meta,positioning,calc,fit}
%  (le librerie vanno caricate nel preambolo, non nel capitolo).
%  Consigliati, come per gli altri capitoli: \usepackage[italian]{babel}
%  e \usepackage[T1]{fontenc} (il capitolo compila anche senza).
%
%  Collocazione nel libro: parte "Analog Game Design", primo capitolo
%  (numerazione propria 2.1). Nel file principale del libro:
%
%    \part{Analog Game Design}
%    \renewcommand{\thechapter}{2.\arabic{chapter}}
%    \setcounter{chapter}{0}
%    % =====================================================================
%  Analog Game Design, lezione 1 — Introduzione al corso: docenti,
%  metodo, project work, figure dell'industria e primo glossario
%  Fonte: trascrizione della prima lezione (AGD1) del corso breve di
%  Analog Game Design, seconda edizione: una lezione online su Zoom di
%  circa tre ore, di pomeriggio, tenuta insieme dai tre docenti "cuore"
%  del corso (Daniele, Ruggero e Beatrice), con nove iscritti, una
%  pausa di una decina di minuti, molte slide e una breve dimostrazione
%  dal vivo di Tabletop Simulator.
%  Pacchetti richiesti nel preambolo del libro (già presenti nel
%  main.tex): \usepackage{booktabs}, \usepackage{tikz} con
%    \usetikzlibrary{shapes.geometric,shapes.symbols,shapes.misc,
%      arrows.meta,positioning,calc,fit}
%  (le librerie vanno caricate nel preambolo, non nel capitolo).
%  Consigliati, come per gli altri capitoli: \usepackage[italian]{babel}
%  e \usepackage[T1]{fontenc} (il capitolo compila anche senza).
%
%  Collocazione nel libro: parte "Analog Game Design", primo capitolo
%  (numerazione propria 2.1). Nel file principale del libro:
%
%    \part{Analog Game Design}
%    \renewcommand{\thechapter}{2.\arabic{chapter}}
%    \setcounter{chapter}{0}
%    % =====================================================================
%  Analog Game Design, lezione 1 — Introduzione al corso: docenti,
%  metodo, project work, figure dell'industria e primo glossario
%  Fonte: trascrizione della prima lezione (AGD1) del corso breve di
%  Analog Game Design, seconda edizione: una lezione online su Zoom di
%  circa tre ore, di pomeriggio, tenuta insieme dai tre docenti "cuore"
%  del corso (Daniele, Ruggero e Beatrice), con nove iscritti, una
%  pausa di una decina di minuti, molte slide e una breve dimostrazione
%  dal vivo di Tabletop Simulator.
%  Pacchetti richiesti nel preambolo del libro (già presenti nel
%  main.tex): \usepackage{booktabs}, \usepackage{tikz} con
%    \usetikzlibrary{shapes.geometric,shapes.symbols,shapes.misc,
%      arrows.meta,positioning,calc,fit}
%  (le librerie vanno caricate nel preambolo, non nel capitolo).
%  Consigliati, come per gli altri capitoli: \usepackage[italian]{babel}
%  e \usepackage[T1]{fontenc} (il capitolo compila anche senza).
%
%  Collocazione nel libro: parte "Analog Game Design", primo capitolo
%  (numerazione propria 2.1). Nel file principale del libro:
%
%    \part{Analog Game Design}
%    \renewcommand{\thechapter}{2.\arabic{chapter}}
%    \setcounter{chapter}{0}
%    % =====================================================================
%  Analog Game Design, lezione 1 — Introduzione al corso: docenti,
%  metodo, project work, figure dell'industria e primo glossario
%  Fonte: trascrizione della prima lezione (AGD1) del corso breve di
%  Analog Game Design, seconda edizione: una lezione online su Zoom di
%  circa tre ore, di pomeriggio, tenuta insieme dai tre docenti "cuore"
%  del corso (Daniele, Ruggero e Beatrice), con nove iscritti, una
%  pausa di una decina di minuti, molte slide e una breve dimostrazione
%  dal vivo di Tabletop Simulator.
%  Pacchetti richiesti nel preambolo del libro (già presenti nel
%  main.tex): \usepackage{booktabs}, \usepackage{tikz} con
%    \usetikzlibrary{shapes.geometric,shapes.symbols,shapes.misc,
%      arrows.meta,positioning,calc,fit}
%  (le librerie vanno caricate nel preambolo, non nel capitolo).
%  Consigliati, come per gli altri capitoli: \usepackage[italian]{babel}
%  e \usepackage[T1]{fontenc} (il capitolo compila anche senza).
%
%  Collocazione nel libro: parte "Analog Game Design", primo capitolo
%  (numerazione propria 2.1). Nel file principale del libro:
%
%    \part{Analog Game Design}
%    \renewcommand{\thechapter}{2.\arabic{chapter}}
%    \setcounter{chapter}{0}
%    \input{AGD/AGD1_introduzione_all_analog_game_design}
%
%  (Con hyperref conviene aggiungere anche
%  \renewcommand{\theHchapter}{agd.\arabic{chapter}} nella parte, per
%  evitare ancore duplicate con i capitoli 1, 2, ... di Game Design.)
%  Tutte le etichette di questo capitolo hanno il prefisso "a1:"
%  (a1:cap:..., a1:sec:..., a1:fig:..., a1:tab:...). Il capitolo è
%  autonomo: non contiene \ref verso altri capitoli.
%  Le figure sono ricostruzioni schematiche di quanto i docenti spiegano
%  a voce e con le slide, non copie delle slide.
% =====================================================================

\chapter[Introduzione all'Analog Game Design]{Introduzione all'Analog Game
Design: il corso, le figure dell'industria e un primo linguaggio comune}
\label{a1:cap:intro}

Questa è la \textbf{prima lezione} del corso di \textbf{Analog Game Design},
cioè di progettazione di giochi \emph{analogici}: giochi da tavolo, giochi di
carte, giochi di ruolo da tavolo, giochi di ruolo dal vivo (LARP), giochi
\emph{pervasivi} ed esperienze a evento. È un corso \textbf{breve}, di circa tre
mesi (``da qui a dicembre''), separato dal corso di Game Design, e siamo alla sua
\textbf{seconda edizione}. Come avvertono subito i docenti, è ``una giornata di
introduzione'': si spiega come funziona il corso, chi sono i docenti, con quali
metodi e con quale calendario si lavorerà, e ci si presenta a vicenda. ``C'è poco
gioco analogico'', si scusano. Ma nella seconda parte la lezione comincia
comunque a costruire un \textbf{linguaggio comune}: le figure professionali
dell'industria del gioco analogico, la differenza tra \emph{game design} e
\emph{game development}, la situazione del mercato italiano rispetto a quello
estero, e un primo glossario di design (\emph{game} e \emph{play}, il framework
MDA, le parole ``vietate'', il playtest e l'iterazione). Si chiude con lo
strumento di settore per i prototipi digitali dei giochi da tavolo,
\textbf{Tabletop Simulator}, e con una lunga sessione di domande.

La lezione si svolge \textbf{online}, su Zoom, ed è condotta a più voci dai tre
docenti principali del corso: \textbf{Daniele}, narrative designer, che guida la
parte organizzativa; \textbf{Ruggero}, game designer, che guida la parte sui
fondamenti e le domande finali; \textbf{Beatrice}, progettista di giochi
pervasivi, che interviene soprattutto nella prima metà. Gli studenti iscritti
sono \textbf{nove}; uno è assente e un altro arriva a lezione iniziata. Ci sono
una pausa di una decina di minuti (si riprende alle 18:15) e, dopo le 19:15,
un'ultima mezz'ora di domande con il solo Ruggero.

Tre avvertenze di lettura. La prima: la trascrizione è automatica e storpia
molti nomi propri, compresi quelli dei docenti e dell'accademia (``Dive'',
``IVE'', ``Accademia Live'', ``Ai''); dove la ricostruzione è sicura, o molto
probabile, il capitolo usa la forma corretta e lo segnala in nota; dove non lo
è, riporta la forma della trascrizione e lo dice. La seconda: durante la pausa la
trascrizione registra, ripetuta in loop per decine di righe, una frase senza
senso che non ha nulla a che fare con la lezione, e che il capitolo ignora. La
terza: in molti passaggi più docenti parlano uno dopo l'altro senza che la
trascrizione segni il cambio di voce; il capitolo attribuisce le parole a chi le
pronuncia quando il contesto lo rende chiaro (chi parla di sé, chi viene
chiamato per nome, chi ``passa la parola''), e resta generico altrimenti.

Quanto alla data, la lezione si colloca con ogni probabilità \textbf{nella prima
metà di settembre del 2026}\footnote{La data non è mai detta. La si deduce da
alcuni indizi: il form iniziale va compilato ``entro e non oltre il 15
settembre''; i gruppi verranno annunciati il 25 settembre; le date del
calendario (sabato 3 ottobre, lunedì 12 ottobre, sabato 7 e 14 novembre,
venerdì 4 dicembre) cadono nei giorni della settimana indicati soltanto nel
2026; un docente parla di un argomento ``molto complesso nel 2026''; e fa
ancora molto caldo (``le temperature sono ancora indecenti''). È una
deduzione, non un'informazione esplicita.}: i riferimenti a ``oggi'',
``quest'estate'' e ``l'anno scorso'' vanno letti rispetto a quel momento.

Il percorso del capitolo segue quello della lezione:
\begin{enumerate}
  \item l'avvio, i \textbf{docenti} e le \textbf{presentazioni} dei partecipanti
        (sezioni da~\ref{a1:sec:avvio} a~\ref{a1:sec:partecipanti});
  \item i \textbf{materiali} e gli strumenti del corso: il form, la cartella
        condivisa, Discord (sezione~\ref{a1:sec:materiali});
  \item lo \textbf{scopo} e il \textbf{metodo} del corso, le regole di
        comportamento e il carico di lavoro (sezioni~\ref{a1:sec:scopo}
        e~\ref{a1:sec:regole});
  \item il \textbf{project work} e il \textbf{calendario} (sezione~\ref{a1:sec:project});
  \item la regola sull'\textbf{intelligenza artificiale} e la lunga discussione
        che ne segue (sezione~\ref{a1:sec:ia});
  \item le \textbf{figure professionali} dell'industria del gioco analogico,
        con un approfondimento su marketing e sales
        (sezioni~\ref{a1:sec:figure} e~\ref{a1:sec:marketing});
  \item la differenza tra \textbf{game design} e \textbf{game development}
        (sezione~\ref{a1:sec:design-development});
  \item il \textbf{mercato}: Italia ed estero (sezione~\ref{a1:sec:mercato});
  \item i \textbf{fondamenti di design}: perché serve una terminologia comune,
        \emph{game} e \emph{play}, il framework \textbf{MDA} (e il DPE), le
        \textbf{parole vietate}, chi crea, chi gioca e chi facilita, il
        \textbf{playtest} e l'\textbf{iterazione} (sezioni
        da~\ref{a1:sec:terminologia} a~\ref{a1:sec:playtest});
  \item \textbf{Tabletop Simulator} e il primo task (sezione~\ref{a1:sec:tts});
  \item le \textbf{domande finali}: il tema del project work, il bilanciamento,
        i ruoli nei team, la pubblicazione, il gioco di ruolo in Italia, le
        \emph{house rule} di Monopoly, i giochi ``comunicati male''
        (sezioni~\ref{a1:sec:domande-daniele} e~\ref{a1:sec:domande}).
\end{enumerate}

% ---------------------------------------------------------------------
\section{L'avvio della lezione}
\label{a1:sec:avvio}

La lezione comincia con qualche minuto di attesa. ``Siamo in nove in classe'', e
si aspetta che arrivino tutti; uno dei docenti scherza sul fatto che l'unica
persona che arriverà ``sempre in ritardo di venerdì'' è Gianmarco. Mancano
Giordano e Tiziana; Tiziana arriva, mentre la segreteria fa sapere poco dopo che
Giordano non potrà esserci e vedrà la registrazione. Si aspetta anche Beatrice,
la terza docente, che arriva a sua volta dopo pochi minuti.

Nell'attesa, Daniele anticipa il programma della giornata e qualche
informazione pratica:
\begin{itemize}
  \item oggi è una giornata di \textbf{introduzione}: come funziona il corso, chi
        sono i docenti, le presentazioni degli studenti, le \textbf{metodologie}
        e il \textbf{calendario}. Gli studenti avranno ``un fracco di domande'', e
        alla fine della lezione c'è ``proprio uno slot apposta'' per le domande;
  \item le lezioni sono \textbf{registrate}: di solito la registrazione viene
        consegnata dopo circa \textbf{12--24 ore} sul profilo personale della
        piattaforma dell'accademia, che va attivato con la segreteria (Tiziana,
        già studentessa, probabilmente ce l'ha già);
  \item i docenti ``cuore'' del corso sono tre (Daniele, Ruggero e Beatrice),
        ``soprattutto Ruggero''; a loro si aggiungeranno tre docenti ospiti, per
        meno ore, probabilmente una lezione a testa, ``tranne Marta, che avrà due
        lezioni quest'anno''.
\end{itemize}
E una battuta sul caldo: ``spero che voi siate carichi, anche se le temperature
sono ancora indecenti''. ``Carichi moderatamente, per quello che potete''.

Poi Daniele condivide lo schermo (``non solo le slide, ma anche un po' di altre
cose'') e apre ufficialmente il corso: ``benvenuti e benvenute nella seconda
edizione del corso di Analog Game Design. Questa è la prima lezione''. Il
percorso è di circa \textbf{tre mesi}, fino a dicembre, e durante il corso i
docenti cercheranno di far provare agli studenti ``cosa significa fare Analog
Game Design'' con un \textbf{project work}, un progetto di gruppo di cui si
spiegheranno più avanti la struttura e le richieste.

% ---------------------------------------------------------------------
\section{I docenti}
\label{a1:sec:docenti}

``Anzitutto partiamo da delle presentazioni doverose''. Ognuno dei tre docenti
principali si presenta con una slide che riporta i contatti e una selezione dei
lavori fatti.

\subsection{Daniele: narrativa, fiction e gioco di ruolo}

\textbf{Daniele}\footnote{La trascrizione riporta ``Daniel Cusetto''; nel
resto della lezione tutti lo chiamano ``Daniele''. Il cognome non è
ricostruibile con certezza.} ha già conosciuto quasi tutti gli studenti, perché
ha fatto loro il colloquio di ammissione. È \textbf{narrative designer}, o meglio
\emph{game narrative designer}; insegna nell'accademia che ospita il
corso\footnote{La trascrizione storpia il nome dell'accademia in molti modi
(``Dive'', ``IVE'', ``Accademia Live'', ``Ai'', ``AIVA''). Con ogni probabilità
si tratta dell'\textbf{AIV}, l'Accademia Italiana Videogiochi, che ha sede a
Roma: lo fanno pensare i corsi triennali di game design e narrative design citati
dai docenti e il fatto che una studentessa racconti di essersi trasferita a Roma.
È una ricostruzione, non un'informazione esplicita; nel seguito il capitolo
parla semplicemente dell'``accademia''.} ma lavora soprattutto come \textbf{brand
advisor}. Nel corso si occuperà della parte di \textbf{gioco di ruolo},
\textbf{narrazione} e \textbf{fiction}.

I suoi contatti, ``in quest'ordine'', sono email, Telegram e Discord. Ci tiene a
dare alcune indicazioni precise:
\begin{itemize}
  \item se si vuole avere la certezza di una risposta, bisogna scrivergli su
        \textbf{Telegram}, perché lì ha tutti i suoi clienti;
  \item Discord lo guarda di meno; l'\textbf{email} è per le cose ufficiali;
  \item non bisogna contattarlo altrove: non su Facebook, non su LinkedIn, non su
        Instagram, ``perché li vedo molto di meno''.
\end{itemize}
La slide mostra anche alcuni giochi, attività e brand su cui ha lavorato negli
anni, ma non si dilunga: ``nel caso ne parleremo''.

\subsection{Ruggero: game design e ricerca e sviluppo}

\textbf{Ruggero} (spesso chiamato ``Ruger'') ha fatto il \textbf{game designer}
per varie aziende e oggi è game designer e si occupa di \textbf{ricerca e
sviluppo} per un'azienda del settore\footnote{La trascrizione dice ``per MS'';
il nome dell'azienda non è ricostruibile con certezza. Più avanti Ruggero cita
tra i giochi su cui ha lavorato il gioco da tavolo ufficiale di
\emph{Cyberpunk}, \emph{Marvel United} e \emph{Status}, un gioco sulla
geopolitica.}. Ha un background ``misto'': nel suo portfolio ci sono giochi da
tavolo, giochi di ruolo, giochi analogici, LARP, videogiochi, giochi in realtà
virtuale. Si occupa anche di \textbf{divulgazione} sul gioco e di analisi di
\textbf{game studies} e game design. Più avanti, nelle domande finali, dirà di
avere ``50 e più giochi prodotti e progettati'' e di insegnare da vari anni.

Lo si può contattare liberamente via Gmail e Telegram, e cerca di rispondere
\textbf{entro 24 ore}, salvo imprevisti o eventi: ``sentitevi in libertà di
scrivermi se avete dubbi, senza problemi''. Quando arriva a parlare dei propri
giochi, la slide mostra ``un mix di varie cose, dalle più famose alle meno
famose'', ma Ruggero si imbarazza a presentarli, tra le risate dei colleghi.

\subsection{Beatrice: i giochi pervasivi}

\textbf{Beatrice} (``Bea'') si occupa di \textbf{giochi pervasivi}
(\emph{pervasive games}). ``Se non li avete mai sentiti nominare, non
preoccupatevi'': questo corso è ``uno dei pochi, forse l'unico in Italia'' che
li racconta esplicitamente. Li si vedrà nelle prossime lezioni, ma la definizione
di base è questa: i giochi pervasivi \textbf{utilizzano gli spazi e gli oggetti
della vita quotidiana e li trasformano in spazi e oggetti di gioco}.

Beatrice lavora come \textbf{freelance} e si occupa di diversi tipi di
progettazione: lavora sia con le \textbf{aziende} sia con gli \textbf{eventi dal
vivo}. La sua attività diretta si trova sotto l'etichetta \textbf{Minime
Differenze}\footnote{La trascrizione riporta il nome dell'etichetta come
``Minime Differenze'' e l'indirizzo del sito come ``mininedifferenze.com''; la
grafia esatta del sito non è verificabile dalla trascrizione.}, mentre i lavori
per le aziende sono ``sparpagliati'' sui siti delle varie aziende. Molte delle
cose mostrate nella sua slide sono attività proposte al pubblico ``con un
biglietto, con un gioco in formato evento'': un tema che si toccherà nel corso,
per capire come \textbf{comunicare} i propri giochi in modo che siano
comprensibili, accessibili per il pubblico e \textbf{vendibili}.

Ammette un ``grosso limite'' con i messaggi: è un po' lenta a rispondere, ``ma
rispondo sempre a tutto''. La si può contattare per email (cognome punto nome) o
su Telegram. È ``molto curiosa di vedere cosa produrrete''.

\subsection{I docenti ospiti}

Daniele presenta poi i tre docenti che il corso avrà ``solo per alcune
lezioni'', tutte verso \textbf{novembre}:
\begin{itemize}
  \item \textbf{Diego Ricchiuti}, direttore del corso di Game Design
        dell'accademia, farà una lezione sull'\textbf{RNG} (\emph{random number
        generation}, la casualità nei giochi): ``vi parleremo di RNG quando sarà
        necessario. Comunque avrete una lezione con lui, sappiatelo'';
  \item \textbf{Marta}\footnote{La trascrizione riporta il cognome come
        ``Cercasassi''; la grafia non è verificabile.} farà \textbf{due lezioni
        sul development}: quest'anno ``abbiamo deciso di fare doppietta'',
        perché l'anno scorso il suo intervento è stato ``molto, molto utile'';
  \item \textbf{Flavio}\footnote{La trascrizione riporta il cognome come
        ``Mortarino'' e, più avanti, ``Mortalino'', e il nome della sua casa
        editrice come ``Grampibert''; nessuno dei due è ricostruibile con
        certezza.} farà la lezione di \textbf{game production}; è un editore, e
        di solito partecipa anche alla tavola rotonda con gli editori della
        giornata finale.
\end{itemize}
Le slide sono già online e gli studenti avranno accesso alla cartella condivisa:
``non vi preoccupate di queste cose, ne parleremo a breve''. Ma prima di parlare
del corso, ``vogliamo conoscere voi''.

% ---------------------------------------------------------------------
\section{Le presentazioni dei partecipanti}
\label{a1:sec:partecipanti}

Le domande per la presentazione sono semplici, da affrontare ``a turno, senza
fretta, con calma'':
\begin{itemize}
  \item nome e occupazione (``se volete dirla'');
  \item perché si vogliono fare giochi analogici, e a che livello: che cosa ha
        portato lo studente al corso;
  \item che tipo di esperienze o conoscenze si hanno già nell'ambito, se si è già
        fatto qualcosa oppure no.
\end{itemize}
``Chi vuole partire senza nessun tipo di ritegno e di vergogne?''. Per far
parlare gli studenti guardandoli, Daniele interrompe la condivisione dello
schermo e ricorda lui le domande quando serve. Le presentazioni sono interessanti
anche per i commenti dei docenti, che anticipano diversi principi del corso.

\subsection{Tiziana}

Tiziana fa l'\textbf{infermiera}, ma nella vita ha fatto molti lavori:
impiegata, insegnante e ``tanti altri lavoretti''. Gioca ``da sempre'': il gioco
fa parte anche della sua vita familiare, perché suo fratello fa il programmatore
di videogiochi. Ha già \textbf{organizzato} giochi: una caccia al tesoro per far
conoscere il territorio, due scene con un delitto (giochi con omicidio da
risolvere), giochi per intrattenere le persone. Ha riscoperto di recente i giochi
da tavolo, arrivando a \emph{Carcassonne} con famiglia e amici; poiché è
difficile trovare persone con voglia di giocare nella propria cerchia ristretta,
ha scoperto il mondo delle \textbf{associazioni} ludiche, ``e lì è stato
un'esplosione'': prima in Veneto (è di Vicenza), poi a Roma, dove si è
trasferita.

Ha in mente un progetto ambizioso: far conoscere una certa realtà con la modalità
del \textbf{gioco di ruolo}, o comunque proporre un gioco suo.

\paragraph{Il commento: gli strumenti, non il gioco dei sogni.} Ruggero coglie
l'occasione per un'avvertenza che vale per tutti, visto che molti studenti
arrivano con idee già pronte. Nel corso si lavorerà a un progetto, ma
l'obiettivo dei docenti è dare \textbf{strumenti} da applicare anche
\emph{fuori} dal progetto, ai giochi che gli studenti vorranno fare un domani:
``che non siano esercizi per il corso, ma per voi in primis''. Quindi: ``non
rimaneteci male se non riuscirete a lavorare al vostro gioco dei sogni in questo
corso''. Anche perché ``sarebbe impossibile, in due mesi, fare il gioco dei
sogni''. Il corso darà invece la possibilità di fare un'esperienza ludica
insieme e di imparare a replicarla sulle proprie idee: ``le idee non buttatele,
anzi, il contrario: tenetele''.

\subsection{Nicola}

Nicola è uno \textbf{studente}. Sei mesi prima ha fatto un colloquio con Daniele
perché stava valutando di passare all'accademia, e ha deciso, ``nel dubbio'', di
cominciare da questo corso breve per capire come funziona la formazione;
eventualmente si iscriverà al corso di game design vero e proprio. Gli interessa
in particolare il lato dell'esperienza del giocatore, su cui è ``molto critico''
(qui usa la parola ``divertimento'', e Daniele lo ferma con un ``ne parleremo'':
è una delle parole che il corso vieta, come si vedrà nella
sezione~\ref{a1:sec:parole-vietate}).

Ha 21 anni ed è \textbf{giocatore di ruolo da circa otto}. Prima ancora di
conoscere i giochi di ruolo si era creato un sistema tutto suo, ``pessimo,
disgustoso, terribile'', un sistema a percentuale con il d100 ``terrificante'';
poi ha scoperto i giochi di ruolo da tavolo, i classici (\emph{Dungeons \&
Dragons}, \emph{Daggerheart}, \emph{Vampire}\footnote{La trascrizione riporta
``DD'', ``Daggerart'' e ``Vampire'': quasi certamente \emph{Dungeons \&
Dragons}, \emph{Daggerheart} e \emph{Vampire: The Masquerade}.}). Gioca anche a
\textbf{card game} e ad altri giochi sociali, pur avendo sempre avuto difficoltà
a trovare gruppi, ``avendo vissuto in un buco dimenticato dal mondo''. Come
giocatore è ``l'eterno narratore'' del gruppo, quello che prepara la storia e si
costruisce le regole. Come autore ha fatto soprattutto giochi su carta, ``perché
sono incapace a programmare, per adesso''.

Daniele osserva che il tema del giocatore e del suo rapporto con il gioco è ``un
argomento importante'', anzi ``un vaso di Pandora enorme'' (espressione che
tornerà molte volte nella lezione), e che il corso tratterà anche che cos'è
l'analog game design in generale e quanti tipi di prodotti esistono.

\subsection{Rita}

Rita\footnote{Daniele la chiama prima ``Nino'' (probabilmente un errore della
trascrizione) e poi la ringrazia come Rita.} ha 30 anni e si collega da una
regione del Sud, dove ``ancora si muore di caldo''. Ha studiato
all'\textbf{Accademia di Belle Arti di Napoli}, \textbf{design della
comunicazione visiva}: si occupa dell'aspetto più imprenditoriale della
progettazione artistica. Non volendo progettare per un'impresa (``fare loghi,
robe del genere''), vorrebbe applicare la propria creatività alla realizzazione
di giochi analogici.

Non è una giocatrice di ruolo e non lo è mai stata: è una giocatrice
\textbf{tradizionale}, cresciuta con Monopoly, la dama e giochi tradizionali
cinesi. Ha già \textbf{idee, prototipi e progetti} realizzati, e vuole capire
come renderli \textbf{prodotti per la vendita} e se questo è il campo adatto a
lei; magari anche confrontarsi con i colleghi. La sua progettazione è orientata
al \textbf{sociale} e all'\textbf{educativo}: ``responsabilizzare un po' le nuove
generazioni su come sta andando il mondo e su come potrebbe andare, visto che è
tutto nelle nostre mani''.

\paragraph{Il commento: siate curiosi.} ``Pane per i tuoi denti'', le risponde
Daniele, soprattutto per i LARP e i giochi pervasivi: ``vi apriremo un po' la
testa da questo punto di vista''. E aggiunge un consiglio per tutti: ``siate
curiosi'', perché imparare da altri ambiti dà moltissimo. Ognuno è lì per fare
qualcosa di preciso, ``però siate aperti anche ad altre esperienze''.

\subsection{Tancredi}

Tancredi ha 21 anni e non ha ancora grande esperienza lavorativa: qualche
``piccolo progettino'' che non è andato come voleva. Gli interessa il mondo dei
giochi, digitali e non, con un approccio \textbf{educativo} e \textbf{sociale}.
È arrivato soprattutto dai videogiochi, a cui gioca da quando era piccolo; ai
giochi da tavolo si è avvicinato solo nell'ultimo anno, grazie a una realtà
ludica, perché in famiglia ``non siamo molto da giochi''. Ha realizzato qualche
piccolo videogioco con vari engine: prima quelli senza un vero linguaggio di
programmazione, ``più a nodi'', poi qualcosa di più complicato, e ora sta
lavorando a un piccolo prototipo.

\paragraph{Il commento: niente di quello che vedrete è sprecato.} ``Non butterai
via niente di analog game design'', gli risponde Daniele, ``perché il 99 per cento
delle cose che vedremo si ritrovano quasi uno a uno nel mondo dei videogiochi''.
Tanto che nei corsi triennali dell'accademia (anche in quello di narrative
design, dove insegna Ruggero) il primo anno si parte dai \textbf{giochi da tavolo
e di ruolo}, proprio per spiegare meccaniche più complesse che servono poi per i
videogiochi. Ruggero aggiunge che vale anche il contrario: alcuni dei
\textbf{framework} di game design che si vedranno nelle lezioni tecniche vengono
dal videogioco e sono stati riapplicati all'analogico. ``Evitare la
crossmedialità qui è impossibile''. Anzi, tutte le conoscenze ``mediatiche
adiacenti'' (musica, letteratura, cinema) saranno utili.

\subsection{Lorenzo Vergata}

I Lorenzo in classe sono due; il primo a presentarsi è \textbf{Lorenzo
Vergata}, 26 anni, \textbf{dottorando al Politecnico di Milano}. Gioca ai
videogiochi da quando aveva due anni e mezzo, ma ha giocato al suo primo gioco da
tavolo ``credo nove mesi fa''. Gli fa piacere sentire che molti in classe sono
interessati ai giochi \textbf{sociali} ed \textbf{educativi}, un settore che
interessa anche a lui e che ``può essere esplorato molto'', anche con qualche
aiuto dell'intelligenza artificiale per esperienze più ricercate.

Di esperienza nell'analogico dice di avere fatto ``il minimo sindacale''. Ma
racconta che diciotto anni prima, da bambino, era ``più game designer di
adesso'': prendeva un gioco e ne cambiava le regole, anche al biliardino (``no,
ora introduciamo questa regola''). Gli piace sperimentare, creare regole e
vedere che cosa accade in sé e negli altri. Soprattutto vorrebbe ``una chiave di
lettura per poter rimanipolare la realtà'': tre settimane prima è stato a fare la
\textbf{vendemmia in Francia} e ha visto tantissime dinamiche che potrebbero
diventare un gioco da tavolo o un videogioco. Per esempio: finire contro un rovo
pieno di semi che si attaccano ai vestiti lo rallentava nella raccolta
dell'uva, mentre in un altro filare era più rapido. ``Si può fare pure una
dinamica'', e vorrebbe gli strumenti per farne giochi.

Arriva subito la battuta: una ``vendemmia gamificata'' da proporre in Oltrepò
Pavese (``io sono dell'Oltrepò Pavese, quindi devo per forza bere vino'')\footnote{La
trascrizione dice ``nell'altro papavese'' e non chiarisce chi pronunci la
battuta.}. Daniele conferma che quest'anno il corso avrà ``tante cose sul
sociale'': ``ci piace parecchio''.

\subsection{L'altro Lorenzo}

Il secondo Lorenzo (``secondo, ma non meno importante'') è uno \textbf{scrittore
di racconti brevi} e sta scrivendo il suo primo romanzo. Gli è sempre piaciuto
raccontare storie, e da più grande scriverle; gioca ai videogiochi ed è ``un
amante delle storie e della cultura pop''. Negli ultimi anni è entrato nel mondo
del \textbf{gioco di ruolo da tavolo} e vorrebbe unire le due passioni, le
storie e il gioco di ruolo, magari facendone un lavoro; gli interessa anche lo
\emph{strategic writing}. Daniele gli ricorda che ci sarà una lezione proprio
sulla \textbf{fiction}.

\subsection{Angelo}

Angelo si collega senza videocamera dal lavoro: fa il \textbf{birraio} (in un
piccolo birrificio in Brianza, a Bernareggio, come racconterà dopo la pausa) e,
con il caldo, ``vi risparmio l'orrore''. Per molto tempo era convinto che i
giochi da tavolo fossero solo i classici (Monopoly, Risiko) o i party game; poi,
quattro o cinque anni fa, ha ``sbattuto la faccia'' contro giochi davvero
appassionanti, con una specie di narrazione dietro che li rendeva ancora più
coinvolgenti, e non ha più smesso. È uno che ``ha sempre un botto di idee, ma poi
non mette in pratica assolutamente nulla'': cercava un corso che gli desse le
basi e soprattutto lo \textbf{motivasse} a provare a fare qualcosa.

\paragraph{Il commento: bisogna buttarsi.} Daniele racconta che anche l'anno
prima un suo caro amico si era iscritto proprio per questo motivo. Il corso è
breve, ma cerca di portare gli studenti ``da 0 a 100'' e di far realizzare un
\textbf{prodotto} da presentare alla fine. Il punto è che i giochi ``bisogna
farli, bisogna buttarsi''. E, venendo dalla narrazione, Daniele lo paragona allo
scrivere: si scrive e si riscrive, e allo stesso modo si fa \textbf{playtest}, si
rifà, si aggiorna. ``Non è che si fanno così''.

\subsection{Gianmarco (dopo le slide)}

Gianmarco arriva a lezione iniziata e si presenta più tardi, alla fine della
prima parte, dopo che i docenti hanno chiuso le slide organizzative e la
discussione sull'intelligenza artificiale (``Gianmarco, ti giuro, ti facciamo
presentare''). Lavora nel \textbf{team di innovazione} di un'azienda italiana.
Vuole fare giochi perché è una sua passione, che lo gratifica, e in passato ci ha
provato ``in maniera scomposta''. Cerca un \textbf{metodo} e un confronto con
persone più esperte: ``studiare solo dei manuali è molto sterile'', ``stare
chiusi nel pertugio non aiuta'', ed ``essere autodidatta ha estremamente i suoi
limiti''. Come esperienza è un giocatore e ha letto ``un po' di materiale, di
blog, di manuali'', ``niente di rilevante''.

Manca solo Giordano, che si presenterà alla lezione successiva.

\paragraph{Uno sguardo d'insieme.} La classe è molto varia: un'infermiera,
una designer della comunicazione visiva, un birraio, un dottorando, uno
scrittore, un impiegato nell'innovazione, due studenti ventunenni. Alcuni
vengono dal gioco di ruolo, altri dai giochi tradizionali, altri dai
videogiochi; quasi tutti hanno scoperto i giochi da tavolo moderni da pochi anni
o da pochi mesi. E almeno tre studenti dichiarano un interesse per i giochi
\textbf{sociali} ed \textbf{educativi}. Come si vedrà nella
sezione~\ref{a1:sec:scopo}, i docenti considerano questa varietà una risorsa.

% ---------------------------------------------------------------------
\section{I materiali e gli strumenti del corso}
\label{a1:sec:materiali}

Prima delle slide sul corso, Daniele sbriga un po' di ``clean up casalingo''.

\paragraph{Il form iniziale.} C'è un \textbf{form} da compilare ``entro e non
oltre il 15 settembre'': il QR code è nelle slide, e la segreteria lo ricorderà
su Discord. Richiede circa 10--15 minuti; si può fare subito o in giornata, ``e
ve lo togliete di torno''. È un form standardizzato per tutti i corsi (``non
siamo io e Ruggero a farvi passare la tortura del test''), diviso in due parti:
\begin{itemize}
  \item una parte di \textbf{conoscenza del candidato}: quanto tempo si ha a
        disposizione per il corso, come ci si organizza con il lavoro, che cosa
        piace giocare. È la parte che interessa ai docenti, perché li aiuta a
        organizzare le lezioni;
  \item una parte di \textbf{test logico}, che i docenti non hanno scritto (``io
        non sono responsabile del form'', ``no, no, nemmeno io'').
\end{itemize}

\paragraph{La cartella condivisa.} Tutti gli studenti vengono aggiunti alla
cartella del corso su \textbf{Google Drive}, che si trova tra i file
``condivisi con me''. Gli studenti ne sono \textbf{gestori dei contenuti} e
possono caricarci quello che vogliono. Dentro ci sono:
\begin{itemize}
  \item le \textbf{attività in classe};
  \item le \textbf{dispense};
  \item i \textbf{materiali didattici}, con le slide (quelle della lezione sono
        già lì);
  \item le cartelle dei \textbf{project work} futuri;
  \item i \textbf{task}, che ``saranno pochi, ma ci saranno'';
  \item il \textbf{calendario studenti}, ``importantissimo'', con due schede: il
        calendario delle lezioni e quello del project work.
\end{itemize}

\paragraph{Discord.} Il corso ha anche un \textbf{server Discord}, a cui tutti
sono stati invitati; ci sono i docenti (Daniele, Ruggero, Diego e altri), che si
possono taggare. Di solito è lì che la segreteria pubblica le informazioni, per
esempio se una lezione ``dovesse per caso essere cancellata''. Nelle domande
finali Ruggero aggiungerà che c'è un canale apposito per il corso, che gli
studenti possono usare liberamente per organizzarsi tra loro.

\paragraph{Le reaction di Zoom.} Durante le lezioni su Zoom si possono usare le
\textbf{reaction} (il pollice in alto, il cuoricino, la mano alzata): un modo
rapido per far capire ai docenti se si è capito, se si vuole intervenire. Nel
frattempo arrivano i primi pollici.

% ---------------------------------------------------------------------
\section{Lo scopo e il metodo del corso}
\label{a1:sec:scopo}

La prima serie di slide (che Daniele, per un lapsus, presenta come ``cos'è il
corso di pulire il pavimento di casa'') spiega che cosa il corso vuole fare e
come.

\subsection{Diventare autori analogici}

Lo \textbf{scopo del corso} è fornire le competenze per \textbf{diventare autori
analogici} e stimolare la creatività e l'\textbf{autorialità} di ciascuno: ``voi
imparerete a diventare autori di giochi analogici''. La differenza tra
\emph{game designer} e \emph{game developer} verrà chiarita più avanti (sezione~\ref{a1:sec:design-development}),
ma la parola chiave qui è \emph{autorialità}: ``vogliamo voi, vogliamo le vostre
idee''. Magari non si potrà usare l'idea del gioco della propria vita, ma si
potranno portare nelle lezioni le proprie idee e le proprie sensibilità.

Gli obiettivi elencati nella slide sono questi:
\begin{itemize}
  \item dare \textbf{maggiore consapevolezza nel design e nello sviluppo}. In
        Italia, purtroppo, ``si fa design a volte in maniera non consapevole''.
        Consapevole vuol dire avere almeno un'introduzione alle basi del design e
        dello sviluppo, e fare le cose ``con metodo, con raziocinio, con
        consapevolezza'';
  \item fornire gli \textbf{strumenti essenziali} per un design consapevole;
  \item stimolare la \textbf{sperimentazione} e la \textbf{trasversalità
        mediatica}. Non si chiede di fare ``il \emph{Root} della situazione o il
        \emph{Carcassonne} della situazione''\footnote{La trascrizione dice ``il
        lumaven della situazione o il carcassone della situazione''; il primo
        titolo non è ricostruibile con certezza. Il senso è chiaro: non si
        chiede di rifare un grande classico del gioco da tavolo.}, ma un
        project work \textbf{calibrato sulla durata del corso}, e però
        \textbf{sperimentale} e ``assolutamente transmediale'';
  \item portare \textbf{numerosi esempi} di tipologie di gioco, dall'analogico al
        tavolo, al pervasivo ``e oltre''. Lo si vedrà nella lezione sui
        \textbf{casi studio} del venerdì successivo: ``voi non avete idea di
        quante cose esistono, bellissime, sconosciutissime e favolose''. Daniele
        confessa che preparando e aggiornando le slide i colleghi gli hanno
        mostrato cose che lui stesso non sapeva nemmeno che esistessero;
  \item dare gli strumenti per \textbf{analizzare} i giochi altrui. Ruggero lo
        sottolinea: se il corso riesce nel suo intento, gli studenti sapranno
        non solo fare i propri giochi ma anche analizzare giochi già fatti, che è
        ``una delle skill importanti anche nella parte di ricerca e sviluppo''.
        Si vedranno casi studio e giochi pubblicati, sempre con l'obiettivo di
        capire \textbf{che cosa fanno e perché lo fanno}, per poi fare meglio il
        proprio lavoro;
  \item imparare a \textbf{fare giochi insieme}. Esistono gli autori che lavorano
        da soli (i \emph{solo developer}), ``ma sappiate che anche l'autore che
        lavora da solo poi si deve interfacciare con il pubblico, con un editore,
        magari con un editor, con un publisher''. ``Non lavorerete mai da soli'':
        bisogna imparare a lavorare con gli altri.
\end{itemize}

\subsection{Un corso intersezionale}

Il corso è \textbf{intersezionale}: \textbf{non ci sono lezioni dedicate a una
singola tipologia di gioco}. Non ci sarà ``la lezione sul gioco da tavolo'',
``la lezione sul gioco di ruolo'', ``la lezione sul gioco pervasivo''. Ci saranno
lezioni su temi diversi, ciascuna pensata per essere utile a \textbf{tutte e tre
le sfaccettature} del corso; alcune saranno più importanti per il gioco da
tavolo, altre per il pervasivo, altre per il gioco di ruolo, ma nessuna sarà
esclusiva. I docenti vogliono ``contaminazioni e applicazioni assolutamente
trasversali''.

\subsection{Un corso cross-generazionale}

È anche un corso \textbf{cross-generazionale}: i gradi di esperienza di gioco
quest'anno sono abbastanza allineati, ma le esperienze di vita, le età, le
estrazioni sociali sono varie. ``Ed è un pro, non è un contro. Più diversificata
è la classe, più imparate meglio anche voi''. La slide ha una piccola nota per
fasce d'età:
\begin{description}
  \item[Under 30] Concentratevi sulla vostra \textbf{crescita}. Evitate di
        confrontarvi con le persone più grandi che magari ``ce la fanno meglio''
        in alcune cose.
  \item[Over 30] Concentratevi sul \textbf{consolidamento}, ma siate
        \textbf{umili}: i docenti diranno cose che faranno ``storcere il naso''.
        ``Bisogna un po' reimparare per imparare'': bisogna sapersi adattare.
\end{description}
Ruggero contestualizza: lo dicono perché in passato hanno visto studenti che,
essendo più adulti o avendo giocato di più, pensavano che l'esperienza
\emph{di gioco} si trasformasse automaticamente in competenza \emph{di
creazione} di giochi. Non è così. E non bisogna sottovalutare chi gioca meno:
proprio perché gioca cose diverse, per risolvere un problema di design specifico
``è una ricchezza incredibile avere qualcun altro che riesce a uscire dalla
scatola''.

\subsection{Un corso pratico e adattabile}

Il corso è \textbf{pratico}: gli studenti si concentreranno su un progetto di
gruppo, il \textbf{project work}, lavorando in team. Ed è \textbf{adattabile}: il
programma è definito, ma ``si aggiorna sui vostri stimoli, priorità e
necessità''. Se durante il corso i docenti vedono che serve più una cosa
dell'altra, sono ``prontissimi'' a dare più spazio a un argomento invece che a
un altro.

\subsection{Inquiry-based learning: domandate}

Il metodo è quello dell'\textbf{inquiry-based learning}, l'apprendimento basato
sulle domande. Di solito le persone dicono ``ho una domanda, magari sarà
stupida''. ``Non esistono domande stupide. Pensate che ci sia una domanda
stupida? Fatela, alzate la mano''. Vanno bene anche le
\textbf{provocazioni}: se qualcosa non torna, è meglio chiedere. La condivisione
è fondamentale.

I docenti danno \textbf{metodi, dati e conoscenze}, ma lasciano trovare le
risposte agli studenti: ``vi lasciamo un pochettino annegare nel mare prima di
darvi la soluzione, in modo tale che la possiate trovare da soli''. Di
conseguenza \textbf{non ci saranno interrogazioni sulle definizioni}: interessa
vedere \textbf{come} gli studenti applicano ciò che imparano. Se una
definizione non viene in mente, non è un problema; l'importante è saperla usare.

\subsection{Feedback, non valutazione}

``Sì al feedback, no alla valutazione''. I docenti danno \textbf{feedback
verbali}, anche sulle attività da fare a casa (i task, che non saranno tanti).
Il corso è calibrato per \textbf{studenti lavoratori}: c'è un po' da lavorare
fuori dalle lezioni, ma in modo calibrato.

Il principio è \textbf{fallire oggi per imparare per il futuro}: \emph{fail
forward}, \emph{fail to learn}. I docenti lasciano provare e fallire almeno una
volta prima di dare la soluzione: ``il classico metodo Montessori, quello
originale, non quello travisato''.

\subsection{Le soft skill e i criteri di valutazione}

Grande attenzione va alle \textbf{soft skill}: i lavori di gruppo, lo scambio
propositivo, l'atteggiamento positivo, la propositività, la proattività.
``Lavoro di gruppo'' non significa stare ogni giorno in sala studio con gli
altri: va benissimo che un gruppo si organizzi dandosi priorità e tempi in base
alla disponibilità di ciascuno. L'importante è \textbf{collaborare}, essere
\textbf{professionali}, rispettarsi a vicenda, rispettare le idee proprie e
altrui e lasciare spazio agli altri, perché i progetti ``non saranno solo
vostri, saranno di tutte le persone del vostro gruppo''.

La valutazione del project work non riguarda solo la \textbf{documentazione}
prodotta, ma anche:
\begin{itemize}
  \item la \textbf{presentazione} finale;
  \item l'\textbf{autovalutazione};
  \item la \textbf{crescita personale};
  \item l'\textbf{atteggiamento} tenuto nel corso.
\end{itemize}
Daniele cita: ``da ognuno secondo le sue capacità, a ognuno secondo i suoi
bisogni''. Lavorare in gruppo non vuol dire che, in tre, ciascuno lavori
esattamente il 33\% del tempo. Ai docenti serve piuttosto capire \textbf{chi fa
che cosa} nel gruppo, e che il processo interno sia \textbf{trasparente}.

\paragraph{Lavorare in cloud.} Per lo stesso motivo il consiglio è di
\textbf{non lavorare in locale}, ma direttamente nelle cartelle condivise del
project work e dei task. Si è liberi di fare altrimenti, ma lavorando in cloud
si dà ai docenti la possibilità di vedere \emph{come} si lavora, il proprio
``processo mentale''.

% ---------------------------------------------------------------------
\section{Le regole di comportamento e il carico di lavoro}
\label{a1:sec:regole}

Alcune slide sono di quelle che si ripetono in ogni corso (``\emph{repetita
iuvant}'', storpiato per scherzo in ``ripetita Juventus'').

\paragraph{Comportamento.} Prestare attenzione durante le lezioni; interagire;
portare rispetto; essere proattivi; alzare la mano e fare domande.
\textbf{Tolleranza zero} per pregiudizi, maleducazione, mancanza di rispetto,
\emph{cheat} e scorciatoie. Daniele ammette che ``cheat e scorciatoie'' viene da
altri corsi dove il tema è più importante, ma Ruggero lo ferma: ``non
sottovalutate questa cosa'', perché implica anche \textbf{non far fare a ChatGPT
i task}. Daniele ha una slide apposta, ``perché è un discorso molto complesso nel
2026'' (sezione~\ref{a1:sec:ia}).

\paragraph{Le lezioni sono intense.} Le lezioni possono essere \textbf{dense},
``soprattutto le mie, sappiatelo: io sono logorroico'', avverte Daniele. Bisogna
chiedere spiegazioni se serve, prendere appunti, rivedere le registrazioni. Le
lezioni sono ``delle scintille'': modi per far esplorare i temi oltre le slide.
``Andate oltre, se riuscite''.

\paragraph{Il carico di lavoro.} Le regole sui task sono queste:
\begin{itemize}
  \item i task \textbf{sono per voi, non per noi};
  \item se si ha poco tempo, è meglio farli rapidamente piuttosto che non farli:
        ``provateci'';
  \item i docenti cercano di non chiedere più di \textbf{un'ora o due a
        settimana} oltre alle ore di lezione;
  \item non tutte le lezioni avranno un task, e il \textbf{project work} è il più
        importante di tutti.
\end{itemize}

\paragraph{I feedback dei docenti.} I feedback ``si guadagnano'' man mano che si
reagisce bene: non bisogna aspettarsi ``dei papiri'' sui task. Se i docenti
hanno voglia di dare feedback più ampi lo fanno, soprattutto se vedono che lo
studente se li merita. Ad Analog di solito non ci sono task \emph{rejected}; ma
se un docente rifiuta un task, vuol dire che si è fatto \textbf{meno del livello
minimo} richiesto dal corso. I feedback possono non essere immediati: ``siate un
pochettino impietositi dal fatto che siamo dei docenti''. E se un feedback non
arriva, o se ne vogliono di più, bisogna scrivere, ``perché siamo tre esseri
umani e quindi possiamo anche dimenticarci''.

% ---------------------------------------------------------------------
\section{Il project work e il calendario}
\label{a1:sec:project}

\subsection{Come funziona il project work}

Una prima slide con le date del project work è da aggiornare (``perdonatemi il
refuso''); Daniele la corregge al volo e rimanda alla slide successiva, più
precisa. Il funzionamento è questo:
\begin{itemize}
  \item i \textbf{gruppi li formano i docenti}. Gli studenti sono nove, quindi si
        faranno \textbf{tre gruppi da tre}: una dimensione pensata anche per il
        fatto che molti sono studenti lavoratori, così da non ritrovarsi ``in due,
        a esplodere'';
  \item ogni gruppo deve prima preparare un \textbf{pitch} del proprio gioco, da
        presentare in anticipo: un'idea, una bozza del gioco;
  \item poi si lavora al gioco \textbf{nelle pause del corso}: da qui a dicembre
        ci saranno due pause, quindi non ci sarà lezione tutti i lunedì e tutti i
        venerdì;
  \item ci saranno \textbf{review intermedie} con i \textbf{tutoraggi}, per lo più
        il \textbf{sabato mattina}, dalle 11 alle 12, con Ruggero su Google Meet.
        Servono ad aiutare i gruppi a portare avanti il project work.
\end{itemize}

\subsection{Il calendario delle lezioni}

Il \textbf{calendario} è nel calendario studenti della cartella condivisa.
Attenzione: le date \textbf{possono variare}. ``Il 90\% rimarranno così'', ma
può capitare che, per motivi di lavoro, qualcuno stia male o non possa, o che
``scompaia l'elettricità da tutte le città d'Italia''. E l'ordine delle lezioni
può cambiare se i docenti vedono che agli studenti serve prima una lezione
invece di un'altra.

Ci sono \textbf{due pause importantissime}:
\begin{itemize}
  \item dal \textbf{13 ottobre al 5 novembre}: è il periodo delle grandi fiere
        (\textbf{Essen} e \textbf{Lucca}\footnote{La trascrizione dice ``Essen,
        Luca, Cose'': si tratta con ogni probabilità della fiera Spiel di Essen,
        in Germania, e di Lucca Comics \& Games, che si tengono entrambe in
        ottobre.}), ``un periodo un po' frenetico'', e cominciano anche i corsi
        accademici, quindi la pausa permette ai docenti di respirare. Per gli
        studenti è la pausa più importante: quasi un mese in cui sviluppare la
        parte importante del gioco;
  \item dal \textbf{24 novembre all'11 dicembre}, dopo qualche lezione a novembre.
\end{itemize}

E ci sono \textbf{due date da evidenziare}:
\begin{itemize}
  \item il \textbf{25 settembre}: nei primi 5--10 minuti della lezione (che per il
        resto sarà tenuta da Beatrice) i docenti annunceranno i tre gruppi;
  \item il \textbf{12 ottobre}: il \textbf{pitch} dei giochi, con l'idea
        embrionale. Tra il 25 settembre e il 12 ottobre ogni gruppo prepara un
        breve pitch in formato slide; quel giorno i docenti lo rivedono e danno
        feedback: se il gioco va bene, se non va bene, come svilupparlo.
\end{itemize}
Dal 12 ottobre fino all'inizio di dicembre si sviluppa il gioco, con i
tutoraggi: ``sfruttate questi momenti di pausa''.

\subsection{Il calendario del project work}

La slide successiva, ``reale e aggiornata'', riporta il calendario del project
work (tabella~\ref{a1:tab:calendario} e figura~\ref{a1:fig:calendario}).
``Per ora non pensate al project work: ci penserete dal 25 settembre in poi''.

\begin{table}[htbp]
\centering
\begin{tabular}{@{}lp{9.2cm}@{}}
\toprule
\textbf{Data} & \textbf{Che cosa succede} \\
\midrule
ven 25 settembre\footnotemark & lezione: annuncio dei gruppi e del
\emph{mandate}, cioè di che cosa si dovrà fare \\
sab 3 ottobre, ore 11 & prima review e tutoraggio con Ruggero (dopo una
decina di giorni di lavoro) \\
lun 12 ottobre & pitch del gioco durante la lezione \\
13 ottobre -- 5 novembre & prima pausa: sviluppo del gioco \\
sab 7 novembre & tutoraggio \\
sab 14 novembre & tutoraggio \\
24 novembre -- 11 dicembre & seconda pausa: sviluppo del gioco \\
ven 4 dicembre, entro mezzanotte & consegna del gioco sul Drive (non c'è
lezione) \\
ven 11 dicembre & i giochi vengono visti a lezione \\
14 dicembre & giornata finale, con la tavola rotonda con gli editori \\
\bottomrule
\end{tabular}
\caption{Il calendario del project work come lo presenta Daniele. Le date
possono variare.}
\label{a1:tab:calendario}
\end{table}
\footnotetext{Daniele dice ``lunedì 25 settembre'', ma nel 2026 il 25 settembre
è un venerdì; il giorno della settimana nella tabella è corretto di
conseguenza.}

\paragraph{Perché la consegna è il 4 dicembre.} Il 4 dicembre non c'è lezione:
si deve solo consegnare il gioco, entro mezzanotte, sul Drive. La presentazione
a lezione è l'11 dicembre, ma i tre docenti hanno bisogno di \textbf{una
settimana} per rivedere il gioco in privato e valutare anche la
\textbf{documentazione}; per questo la consegna non coincide con la
presentazione. Che cosa voglia dire esattamente ``consegnare il gioco'' verrà
spiegato più avanti.

\paragraph{La giornata finale.} Il 14 dicembre c'è la \textbf{giornata finale},
con una \textbf{tavola rotonda con gli editori}: di solito viene invitato Flavio,
che è anche docente di game production, più qualche altro editore, ``vediamo
quest'anno chi riusciamo a portarvi''.

\begin{figure}[htbp]
\centering
\begin{tikzpicture}[
    font=\footnotesize,
    x=0.1cm,
    evento/.style={draw, thick, rounded corners=2pt, fill=white,
                   text width=1.8cm, align=center, inner sep=2pt},
    pausa/.style={fill=gray!20, draw=gray!60},
    tacca/.style={thick}
  ]
  % asse: giorno 0 = 25 settembre; 1 unità = 1 giorno
  \draw[pausa] (18,-0.35) rectangle (41,0.35);
  \draw[pausa] (60,-0.35) rectangle (77,0.35);
  \node[font=\scriptsize\itshape] at (29.5,0.18) {pausa};
  \node[font=\scriptsize\itshape] at (65,0.18) {pausa};
  \draw[-{Stealth[length=2.5mm]}, thick] (-3,0) -- (85,0);
  % eventi sopra
  \draw[tacca] (0,-0.15) -- (0,0.9);
  \node[evento, above] at (0,0.9) {25 set\\ gruppi e \emph{mandate}};
  \draw[tacca] (17,-0.15) -- (17,2.2);
  \node[evento, above, fill=gray!8] at (17,2.2) {12 ott\\ \textbf{pitch}};
  \draw[tacca] (70,-0.15) -- (70,0.9);
  \node[evento, above, fill=gray!8] at (70,0.9) {4 dic\\ \textbf{consegna}};
  % eventi sotto
  \draw[tacca] (8,0.15) -- (8,-0.9);
  \node[evento, below] at (8,-0.9) {3 ott\\ tutoraggio};
  \draw[tacca] (43,0.15) -- (43,-0.7);
  \draw[tacca] (50,0.15) -- (50,-0.7);
  \draw (43,-0.7) -- (50,-0.7);
  \draw[tacca] (46.5,-0.7) -- (46.5,-0.9);
  \node[evento, below] at (46.5,-0.9) {7 e 14 nov\\ tutoraggi};
  \draw[tacca] (77,0.15) -- (77,-0.9);
  \node[evento, below] at (77,-0.9) {11 dic\\ giochi a lezione};
  \draw[tacca] (80,-0.15) -- (80,2.2);
  \node[evento, above] at (80,2.2) {14 dic\\ giornata finale};
\end{tikzpicture}
\caption{Una linea del tempo del project work, ricostruita dalle slide (scala
approssimativa: dal 25 settembre al 14 dicembre). Le bande grigie sono le due
pause senza lezioni, in cui i gruppi sviluppano il gioco.}
\label{a1:fig:calendario}
\end{figure}

Finita questa parte, Daniele apre alle domande: ``Domande sul calendario? Sulle
slide che ho fatto? Considerazioni? Non abbiate timore, io non ci credo''. E la
domanda arriva, sull'unico tema che aveva lasciato da parte.

% ---------------------------------------------------------------------
\section{L'intelligenza artificiale nel corso}
\label{a1:sec:ia}

\subsection{La slide di Daniele}

Daniele aveva già presentato una slide che definisce ``importantissima''. ``Lo so
che è paradossale, me ne rendo conto. Lo so che magari mi darete del luddista,
dell'anti-tecnologia: non è così. Uso l'intelligenza artificiale per lavoro. Ma
in questo corso, vi prego, \textbf{non usate l'intelligenza artificiale}''. La
richiesta vale per le attività, per i task e \textbf{anche per la
prototipazione}. Le ragioni:
\begin{itemize}
  \item ``voglio che siate voi a fare il gioco, voglio che siate voi a imparare a
        farlo''. Daniele parla da utilizzatore: prima di ChatGPT, che ha smesso
        di usare, oggi soprattutto di Claude, di cui ``si sottovaluta quanto
        possa essere incerto, insicuro, impreciso, non all'altezza'';
  \item nella game industry, \textbf{i publisher rifiutano automaticamente} i
        pitch e i prototipi fatti con l'intelligenza artificiale. Daniele lo ha
        visto di persona in ``tantissime pitch review'' quell'estate, anche al
        \textbf{First Playable} di Firenze\footnote{La trascrizione dice ``First
        Playboy a Firenze'': si tratta del First Playable, l'evento di business
        dell'industria videoludica che si tiene a Firenze.}: anche nel mondo dei
        videogiochi i publisher che vedono pitch o prototipi fatti con l'IA
        ``rigettano automaticamente la build'', perché hanno ancora timore
        dell'IA e delle sue possibili conseguenze;
  \item ``non so come sarà il futuro'': non sa se bisognerà abituarsi a usarla
        tutto il giorno. Non è ``totalmente contrario'', purché sia usata ``con
        raziocinio''; secondo lui oggi, negli Stati Uniti, non lo è;
  \item ma in un corso che si chiama \emph{Analog Game Design} gli studenti devono
        saper fare i giochi da soli. In futuro, se vorranno farsi aiutare
        dall'IA per essere più efficienti, saranno liberissimi di farlo; durante
        il corso no, perché i docenti vogliono valutare \textbf{le loro
        competenze e capacità}.
\end{itemize}

\subsection{La domanda: che cosa vuol dire usare l'IA in un gioco analogico?}

La domanda di uno studente è precisa: che cosa si intende per ``task fatti con
l'IA''? Qual è l'uso dell'IA in un gioco analogico?

\paragraph{Ruggero: l'esperienza di chi valuta i prototipi.} Ruggero risponde
dalla sua esperienza nel settore, dove riceve prototipi da valutare per
un'eventuale pubblicazione. Distingue due usi:
\begin{enumerate}
  \item le \textbf{immagini generate}: illustrazioni, graphic design e così via.
        Per molti publisher è ``un \emph{reject} immediato, proprio immediato'';
  \item l'IA usata per il \textbf{regolamento}, per scrivere le regole e le
        meccaniche. È più difficile da individuare, ma molti publisher se ne
        accorgono comunque, perché le IA ``non sono ancora brave abbastanza da
        creare un regolamento coeso meccanicamente e scritto con un senso per il
        vostro gioco''. Quello che fanno è prendere parti di altri regolamenti,
        elaborare le strutture che vedono nei regolamenti pubblicati e
        riapplicarle. Ma molti giochi si prestano poco a un regolamento
        strutturato come quello di altri, perché \textbf{ogni gioco è un
        ecosistema meccanico a sé}. Il risultato sono regolamenti con
        riferimenti non immediati, e soprattutto con \textbf{la stessa meccanica
        descritta in modo diverso in pagine diverse}: due difetti che fanno
        perdere tempo all'editore, e in questo caso ai docenti.
\end{enumerate}
Poi aggiunge un punto ``personale, soggettivo'': i docenti sono interessati a
vedere il lavoro, la creatività, ``il frutto della vostra creazione''. E
viceversa gli studenti vogliono i commenti dei docenti: ``immaginatevi se voi
fate un gioco e noi vi diamo i commenti scritti con ChatGPT''. Altrimenti, dice
ridendo, gli studenti fanno il gioco con ChatGPT, i docenti lo danno a ChatGPT
per i commenti, ``e facciamo molto alla svelta''.

\paragraph{Daniele: perché gli studenti usano l'IA.} Daniele racconta la sua
esperienza con gli studenti, che secondo lui ricorrono all'IA in due momenti:
\begin{enumerate}
  \item quando vanno in \textbf{crisi}: in ansia da prestazione o per i tempi.
        Cercano di risolvere il problema della propria disorganizzazione
        facendolo risolvere all'IA, senza accorgersi che l'output è ``molto,
        molto scarso'';
  \item quando hanno l'impressione che l'intelligenza dell'IA sia
        \textbf{migliore} di quello che scrivono loro, o di quello che c'è in
        giro. Per Daniele è una ``mancanza di palato critico'': l'output dell'IA
        sembra buono in superficie, ma quando si va nel dettaglio e ci si fa
        delle domande ``di solito fa un po' di \emph{crumbling}'', si sbriciola.
\end{enumerate}
Lui usa Claude ogni giorno e ne conosce i limiti: ``non c'è nemmeno una frase
che io copio direttamente''. Gli fa fare il lavoro più noioso, più ripetitivo e
soprattutto \textbf{strutturato}, ``replicare delle cose come se fosse una
macro''. Lo usa per il testo e per il codice, mai per le immagini: per motivi di
\textbf{diritto d'autore}, e perché ha molti amici artisti ed è stufo di vedere
il loro lavoro ``dilaniato'' dall'IA.

Ammette che i docenti non possono controllare gli studenti; per questo
``ve lo consiglio''. E riprende quella che definisce la cosa più bella detta da
Ruggero (``ti ruberò la frase per il corso di narrative''): i docenti i feedback
li danno di persona. Daniele a volte manda 30 minuti di audio a uno studente, o
fa intere call per rivedere script, pitch e giochi: ``è il mio lavoro, mi piace
farlo e voglio farlo io, quindi mi aspetto che lo studente lo faccia lui''. Come
freelance, poi, è ``stufo marcio'' dei clienti che arrivano con qualcosa fatto
con l'IA (``guarda che bello, pubblica questo''), perché il più delle volte il
risultato è pessimo: non sanno come usarla. ``È anche una questione di
competenza''.

\paragraph{Il prototipo non deve essere bello.} Questo vale anche per il
\textbf{prototipo}, su cui si tornerà nelle lezioni di prototipazione: in un
prototipo non servono ``graficoni allucinanti''. Le grafiche bellissime fatte con
l'IA spesso \textbf{peccano di user experience}: sembrano bellissime, ma quando
il giocatore deve leggere la carta non capisce qual è il costo, qual è il
potere, qual è l'effetto.

\paragraph{Beatrice: si perde l'occasione di imparare.} Beatrice aggiunge due
motivi che valgono ``al di là dei nostri credo''. Il primo: chi usa l'IA in
sostituzione (per esempio per scrivere un testo, per accelerare i tempi o per
trovare strategie) \textbf{perde la parte in cui può imparare}. Nel game design
bisogna ``tenere insieme i pezzi'' di un regolamento, ma anche dell'esperienza
dell'utente e di come trasmetterla; i passaggi necessari (creare il regolamento,
creare le meccaniche, sperimentare, raccogliere i feedback) ``si capiscono solo
vivendoli''. E si sta pagando un corso per imparare. Il secondo, sul piano
editoriale: chi presenta a un editore un prodotto che ha come base l'IA rischia
da un lato di non avere imparato abbastanza (e di ritrovarsi ``a piedi'' sul
campo senza IA), dall'altro che l'editore se ne accorga, ed è un tema
\textbf{reputazionale}. Ci sono editori ``più illuminati di altri'', quindi
anche qui c'è un mercato da capire; ma i docenti preferiscono dire che cosa è
più sensato per il mercato come lo conoscono.

\paragraph{Ruggero: non abbiate paura del giudizio.} Ruggero vuole dare uno
``spotlight'' a un punto detto da Daniele e rimasto ``semi-ignorato'': a volte
l'uso dell'IA nasce dalla \textbf{paura del giudizio}, dal voler consegnare
qualcosa di più pulito, più strutturato possibile. ``Vi giuro, non ve ne
preoccupate''. Ruggero detesta i voti e non ragiona per voti; non valuterà
nessuno per grafiche fatte a matita. Insegna da anni, ha avuto studenti con
dislessia e discalculia: refusi ed errori non gli interessano, ``non è un corso
di editing o di proofreading''. Gli interessa che le \textbf{meccaniche siano
coerenti con lo scopo di design}. Immagini fatte a matita o con Paint ``non sono
un problema, anzi, le preferisco'' rispetto all'``ennesima immagine fatta con
ChatGPT o Midjourney''.

\paragraph{Daniele: il pensiero critico, costo e qualità.} L'ultimo intervento
è di Daniele, sul \textbf{pensiero critico}. Ragionando per assurdo (``un
pensiero di Belzebù''), immagina che un domani l'IA venga usata ovunque. I suoi
limiti, secondo lui, resterebbero due:
\begin{description}
  \item[Il costo] Tutti dicono che sta diminuendo e che diminuirà, ma non se ne
        può essere sicuri al 100\%.
  \item[La qualità] Cioè la capacità di leggere l'output. Se si dà un regolamento
        a Claude o a ChatGPT chiedendo quali problemi ha, il modello risponde con
        i problemi che, in percentuale, è più probabile che gli esseri umani
        abbiano già individuato nei dati su cui è stato addestrato. Ne derivano
        due problemi: se il problema del proprio gioco non è mai stato
        individuato prima, il modello non lo vedrà; e viceversa può segnalare
        problemi che il gioco non ha.
\end{description}
Fa l'esempio di Grammarly, che usa nella versione premium per rivedere i testi
in inglese: spesso non accetta i suoi consigli, perché se vuole scrivere un
dialogo di videogioco in \textbf{dialetto Cockney} londinese, lo deve scrivere
in Cockney. Se Grammarly dice che è sbagliato, ``sei tu che non hai capito il
motivo per cui voglio scrivere questo testo in Cockney''. La conclusione: prima
di usare un'IA per un'attività bisogna sapere \textbf{che cosa si vuole
ottenere}, qual è l'obiettivo, e quali \textbf{competenze} servono per valutare
in modo critico se l'output di un qualunque LLM è davvero valido. Daniele non
vuole dilungarsi oltre: è un argomento ``molto caldo'' per lui, che è anche
scrittore e si sente dire da tutti ``ti sostituiranno''.

Subito dopo si presenta finalmente Gianmarco (lo si è visto nella
sezione~\ref{a1:sec:partecipanti}), e poi si fa una pausa di una decina di
minuti. Daniele avverte che nella terza parte lui non ci sarà (alle sette e un
quarto ha un impegno di lavoro), e lo stesso vale per Beatrice: la parte finale
di domande e risposte sarà con il solo Ruggero, ``in ottime mani''.

% ---------------------------------------------------------------------
\section{Le figure professionali dell'industria analogica}
\label{a1:sec:figure}

Dopo la pausa (in cui si scopre che Angelo lavora in un piccolo birrificio in
Brianza e si scherza sul rischio di un \emph{drinking game}, tra la birra e la
vendemmia) la lezione riprende con le slide di Ruggero. ``Oggi è una lezione
introduttiva'': gli argomenti si vedranno ``molto dall'alto'' e verranno
approfonditi nelle lezioni successive. Ma servono già ora, per creare
\textbf{l'inizio di un linguaggio condiviso}: ci sono termini che torneranno
spesso, a cominciare da quelli del mercato e dell'industria del gioco
analogico, un'industria ``molto varia, molto caotica, anche molto mutevole'',
con tanti tipi diversi di giochi e altrettanti tipi diversi di entità che ci
lavorano.

\subsection{Una decina di etichette}

I docenti hanno raccolto ``una decina di etichette'', cioè di \textbf{ruoli} che
nell'industria sono o possono essere presenti. Due avvertenze:
\begin{itemize}
  \item nelle aziende più piccole \textbf{la stessa persona copre spesso più
        ruoli}, per forza di cose. Sarà anche il caso degli studenti, che
        lavoreranno in gruppi di tre e non di venti: è facile che ciascuno finisca
        a fare più cose. I docenti ne terranno conto, e qui entreranno in gioco
        anche competenze diverse dal game design (programmi di grafica,
        illustrazione, impaginazione, layout, \emph{proofreading}\ldots);
  \item Daniele precisa che \textbf{non si chiede agli studenti di diventare}
        illustratori o altro. Ruggero concorda, ma aggiunge che è utile sapere
        anche \emph{che cosa non è il proprio lavoro}. Se in un gruppo nessuno sa
        illustrare, ai docenti va benissimo un prototipo che mostri dove vanno
        le immagini: \textbf{placeholder}, concept, anche semplici box di testo
        con scritto ``qui vorremmo un'immagine di\ldots''. ``È importante anche
        sapere che cosa sapete fare e che cosa non sapete fare''.
\end{itemize}
La tabella~\ref{a1:tab:figure} riassume le figure; le sottosezioni seguenti le
riprendono una per una, con i commenti dei docenti.

\begin{table}[htbp]
\centering
\small
\begin{tabular}{@{}p{3.3cm}p{9.6cm}@{}}
\toprule
\textbf{Figura} & \textbf{Di che cosa si occupa} \\
\midrule
Game designer & crea le regole alla base dell'esperienza: le meccaniche, come
interagiscono tra loro e come chi gioca interagisce con esse \\
Game developer & fa da ponte tra chi fa game design e il publisher: bilancia il
prototipo, cura i modelli matematici, adatta il gioco (e il paragame) alle
necessità del publisher \\
Editor (proofreader) & revisiona e adatta i testi a partire dalla prima stesura
del designer (correzione, linguaggio inclusivo\ldots) \\
Publisher (editore) & pubblica e produce fisicamente il gioco; il suo marchio è
sulla scatola o sulla copertina \\
Sales e marketing & il lato commerciale: vendite (anche B2B) e comunicazione del
prodotto \\
Graphic designer & icone meccaniche, wireframe, UX: il lato funzionale degli
elementi visivi \\
Illustratore & le illustrazioni: il \emph{look and feel}, il world building, il
``vibe'' del gioco \\
Sensitivity reader (cultural consultant) & verifica l'accuratezza e la
rappresentazione di un tema sensibile che il team non conosce \\
Producer & coordina il progetto (quasi project management), parla con le
tipografie, verifica che i materiali siano sostenibili economicamente \\
\bottomrule
\end{tabular}
\caption{Le figure professionali dell'industria del gioco analogico presentate
nella slide di Ruggero, in sintesi.}
\label{a1:tab:figure}
\end{table}

\subsection{Game designer}

Il \textbf{game designer} (o la game designer) è chi si occupa, ``sembra banale
a dirlo'', di \textbf{creare le regole alla base dell'esperienza} che si vuole
far arrivare. Vale per tutti i tipi di giochi, non solo per quelli analogici, e
che cosa faccia concretamente in un progetto o in una serie specifica ``può
mutare tantissimo''. Nell'analogico non verrà chiesto di programmare o di
prototipare in digitale; verrà chiesto di \textbf{creare le meccaniche di
gioco}, di capire \textbf{come le meccaniche si interfacciano tra loro} e
\textbf{come le persone che giocano si interfacceranno con le meccaniche}.

\subsection{Game developer}

Il \textbf{game developer} (o la game developer) sarà l'argomento delle due
lezioni di Marta. La parola \emph{development} è nota soprattutto
dall'ambito videoludico, ma nell'analogico ha un significato molto specifico
``che non c'entra niente con la programmazione'': il game developer è
\textbf{un ponte tra chi fa game design e il publisher}. Di solito si occupa di
\textbf{bilanciare il prototipo}, di curare i \textbf{modelli matematici}, di
adattare il gioco e il \textbf{paragame}\footnote{La trascrizione dice ``il
Wix, il paragame''. Il \emph{paragame} è, in termini generali, tutto ciò che
accompagna il gioco senza esserne il sistema di regole (regolamento, scatola,
materiali, comunicazione); i docenti lo tratteranno in una lezione dedicata, e
qui il termine è solo citato.} alle necessità del publisher. La distinzione
viene ripresa e approfondita nella sezione~\ref{a1:sec:design-development}.

\subsection{Editor}

L'\textbf{editor}, a volte chiamato \textbf{proofreader} a seconda dell'azienda,
si occupa della revisione dei \textbf{testi}: non della stesura iniziale, perché
di norma, data la scala dell'industria, è chi fa game design a scrivere un
\textbf{primo draft}, e l'editor ci passa sopra. Non solo per correggere i
refusi: spesso il suo lavoro è anche \textbf{adattare} i testi, per esempio
renderli \textbf{gender neutral}. Se ne vedranno casi studio più avanti.

\subsection{Publisher}

Il \textbf{publisher}, o \textbf{editore}, è l'entità che nell'ambito
commerciale \textbf{pubblica} giochi, giochi di ruolo o giochi da tavolo. È una
figura ``più solidificata'' nel gioco da tavolo e nel gioco di ruolo da tavolo;
\textbf{nell'ambito pervasivo non esiste proprio in questa forma}, come
mostrerà Beatrice parlando dei casi studio. Se si compra un gioco da tavolo o un
gioco di ruolo in un negozio, il marchio del publisher è di solito sul retro
della copertina o sul fronte della scatola: è l'entità che, di solito,
\textbf{produce fisicamente} il gioco.

\subsection{Sales e marketing}

L'ambito \textbf{sales e marketing}, a volte affidato a un'unica persona e a
volte separato, riguarda il lato commerciale. In alcune aziende è interno al
publisher, e si occupa di parlare con le fabbriche e le tipografie, di ``far
quadrare i conti''. In altre realtà è completamente diverso, specialmente nel
\textbf{B2B} o nei giochi \textbf{didattici} e \textbf{formativi}, di cui
Daniele e Beatrice ``sapranno parlare meglio''. Il tema viene ripreso più avanti
in dettaglio (sezione~\ref{a1:sec:marketing}).

\subsection{Graphic designer e illustratore}

Ruggero presenta insieme \textbf{graphic designer} e \textbf{illustratore}
proprio perché sono ruoli molto diversi che si tende a far coincidere:
\begin{itemize}
  \item chi fa \textbf{graphic design} si occupa delle \textbf{icone meccaniche},
        del \textbf{wireframing}, della \textbf{UX}: del lato \textbf{meccanico e
        funzionale} degli elementi visivi, cioè di tutto ciò che serve
        meccanicamente nel gioco, o anche alla sua ``qualità della vita'';
  \item chi fa \textbf{illustrazione} si occupa dell'\emph{art}, delle
        illustrazioni: copre il \textbf{look and feel}, la parte di
        \textbf{world building} e il ``vibe'' che si dà al gioco.
\end{itemize}

\paragraph{Daniele: che cos'è la user experience.} Daniele interviene sulla
\textbf{user experience}, ``importantissima per me'': ha un background da
\textbf{layout artist} (impaginava magazine), un mestiere in cui la UX conta
moltissimo. Per user experience non si intende solo l'\textbf{interfaccia
grafica}, la UI, che anche nei videogiochi è disegnata da un artista. La UX è
qualcosa che può fare anche il game designer: \textbf{disegnare l'esperienza del
giocatore}, cioè ``il gesto, l'intenzione, l'intuito'' che stanno dietro
l'aspetto visivo dell'interfaccia. Per fare UX non serve saper disegnare.
L'esempio è quello della \textbf{porta}: se si progetta una porta, bisogna
indicare all'utente come si apre, come si prende la maniglia; ``quella è user
experience''. Il design della maniglia in sé, invece, è
\textbf{interfaccia}.

\paragraph{Un aneddoto sulle icone.} Ruggero aggiunge che molte persone che
fanno \textbf{graphic design ben fatto non sanno illustrare}, e spesso vale il
viceversa. Lo si nota quando le icone di un gioco sono bellissime ma non
funzionali: andando a guardare i crediti, di solito si scopre che le icone le ha
fatte la persona che ha fatto le illustrazioni. ``Comunque non andiamo adesso in
questo, perché è un altro vaso di Pandora''.

\subsection{Sensitivity reader}

Il \textbf{sensitivity reader}, o \textbf{cultural consultant} a seconda del
progetto, fa un lavoro di \emph{sensitivity reading} del gioco. È una figura
molto usata nell'ambito \textbf{formativo e didattico}. L'esempio di Ruggero: se
si fa un gioco per un ospedale su una malattia specifica, chi lo progetta non è
un medico, ma il gioco deve essere \textbf{accurato} e \textbf{rappresentativo}
delle dinamiche che quella malattia comporta. Serve una persona che conosca quel
tema, anche se non sa che cosa sia il game design. Sono figure chiamate di
solito \textbf{a progetto}, per giochi specifici, che compensano competenze che
il team non ha ``per forza di cose, perché noi non sappiamo tutto''; e a volte,
per lavoro, si devono fare giochi su temi che non si conoscono personalmente.

\paragraph{Daniele: non è un censore.} Daniele segnala un problema italiano:
alcune persone, anche online, hanno confuso il sensitivity reader con un
\textbf{censore}, ``una persona che fa censura, soprattutto in materia di
\emph{woke}''. ``Assolutamente non è così''. La parola \emph{sensitive} non
significa che il pubblico sia sensibile a un tema, ma che \textbf{l'argomento di
cui tratta il gioco è sensibile}. In un gioco sulla medicina, realistico o no,
non si possono inserire imprecisioni. Per questo esistono i sensitivity reader
``più disparati'': da quelli per la cultura indiana o afroamericana fino a quello
su ``come si allevano gli gnu in Africa''. È una persona competente su un
argomento che aiuta a essere \textbf{veritieri}. Certo, uno degli ambiti
principali in cui si usano è quello culturale, ma ``non fate l'errore di pensare
che siano dei censori''. Daniele ha lavorato con ``centinaia'' di sensitivity
reader e nessuno gli ha mai detto ``non puoi scrivere questa cosa''. Di solito
sono come degli \textbf{editor}: danno opzioni, mostrano i punti problematici,
anche quelli di cui non ci si era nemmeno accorti.

Ruggero conferma con un esempio di quando lavorava in \textbf{Danimarca} con
l'\textbf{OMS}: giochi sul \textbf{diabete di tipo 2} nei bambini, per un target
di ragazzi delle scuole medie. C'era un team dell'OMS specializzato nella
comunicazione sul diabete, per garantire che il gioco fosse accurato e ne desse
una rappresentazione adatta al target. Il sensitivity reader aiuta a rendere il
gioco più accurato e ``più tematicamente coeso''.

\subsection{Producer}

Il \textbf{producer} è una figura a volte interna al publisher. In un progetto
grande è la figura di riferimento, ``quasi'' un project manager. In un progetto
piccolo si occupa spesso anche di parlare con le \textbf{tipografie} e di fare la
\textbf{lista dei materiali}, per assicurarsi che i materiali non siano solo
funzionali al gioco ma anche \textbf{economicamente sostenibili} da produrre:
``magari è un gioco bellissimo, ma poi costa più di quanto potete permettervi, e
nessuno vi comprerà mai un card game a 200 euro''.

\subsection{Perché conoscere queste figure}

Ruggero chiude la slide ``tornando a ombrello'': perché parlarne, se gli studenti
sono lì per fare game design? Non solo per saper dare un \textbf{nome} alle
figure. Non c'è niente di male a creare un gioco solo per sé; ma chi vuole fare
giochi a livello lavorativo, perché li giochino altre persone, ``magari più
persone possibile'', \textbf{si interfaccerà con queste figure}, e sapere che
cosa fanno e come parlarci è parte importante del percorso. Nel corso se ne
toccheranno diverse: il development con Marta, la produzione con Flavio, e
lezioni sulla UX, sul paragame, sul graphic design. ``Abituatevi mentalmente:
non siete in solitaria contro il mondo''.

% ---------------------------------------------------------------------
\section{Marketing e sales}
\label{a1:sec:marketing}

Arrivato all'angolo della slide con marketing e sales, Ruggero passa la parola
ai colleghi: ``io è meglio che non mi fate parlare di marketing e sales'', anche
per esperienze personali.

\subsection{Beatrice: i giochi trattati come prodotti}

Beatrice premette che le due voci sono in parte autoesplicative. Nel mondo
capitalista in cui viviamo, i giochi sono ``strumenti molto potenti per fare
\emph{n} cose'', ma sono trattati come \textbf{prodotti}, e il corso ragionerà
anche in quest'ottica. Nel \textbf{marketing} delle industrie ludiche e
videoludiche ci sono varie professionalità che gli studenti probabilmente già
conoscono: dal \emph{copy} alla grafica, tutto ciò che riguarda il
\textbf{vendere} il gioco o l'attività di gioco.

Il lato \textbf{sales} è forse la distinzione meno nota. Come in tante aziende, è
la parte che si occupa delle vendite \textbf{B2B}, cioè tra aziende: fare
accordi per vendere prodotti ludici a diverse realtà, per esempio ai
\textbf{distributori}, o a una casa editrice che vende alle librerie, che poi
vendono al consumatore finale. Chi si occupa di sales nel mondo ludico capisce
come vendere \textbf{alle aziende}. I due mondi si parlano spesso, e soprattutto
nelle realtà piccole si trovano figure che fanno ``un po' sales, un po'
marketing''.

\subsection{Daniele: che cos'è (e che cosa non è) il marketing}

Daniele il marketing si è trovato ``purtroppo'' a farlo, e lo fa anche per
l'accademia. ``Purtroppo'' perché in Italia è mal compreso: molti lo
identificano con l'\emph{advertisement} online. Non è così: il marketing è una
\textbf{disciplina molto più antica}, e in Italia ``abbiamo molto dimenticato''
come si fa il marketing tradizionale. Alcune cose da ricordare:
\begin{enumerate}
  \item \textbf{Il marketing non dice che cosa pubblicare, ma come pubblicarlo}.
        Esistono, nei videogiochi e nei giochi da tavolo, realtà che fanno
        prodotti \textbf{marketing driven}: l'esempio sono gli
        \textbf{hyper-casual game} per smartphone, quelli con il personaggio che
        va a destra e a sinistra e l'esercito che si allarga. Possono anche
        guadagnare, ma ``sono un po' la morte della creatività''. Il marketing
        ideale \textbf{prende un prodotto pensato da altri e gli dà le maggiori
        chance di essere venduto}: cerca il pubblico ideale per il prodotto che
        si ha, senza doverlo cambiare. Anzi, il mondo videoludico, e ludico in
        generale, ``è lastricato'' di errori di dipartimenti marketing che hanno
        cambiato completamente non il gioco, ma il modo in cui l'hanno
        pubblicizzato, e il gioco è fallito.
  \item \textbf{Usare strumenti del marketing non significa fare marketing}. Game
        design e game development a volte usano strumenti che vengono dal
        marketing (gli studenti ne vedranno con Beatrice, con Daniele e con
        Marta).
  \item \textbf{Il marketing è un gerundio}: si fa. È un'azione che si compie per
        comunicare; la comunicazione è attiva, e il marketing è l'insieme delle
        attività con cui si comunica il proprio gioco, dal \emph{copy} ai social.
        Può includere il digital marketing, ``ma non è solo quello''.
  \item \textbf{Il marketing esiste anche se si crede che non esista}. In Italia
        molti non vogliono ``fare marketing'' perché lo associano al lucro. Ma
        il marketing è \textbf{lo studio di come gli esseri umani comprano e
        vendono}, ed è innato in noi: ci sono logiche di cui siamo vittime, e
        ``anche se sei il bastian contrario, l'anticapitalista'', magari si
        compra comunque su Amazon perché fa comodo. Sono meccanismi che
        esistono, ``affascinanti anche, a volte''.
\end{enumerate}
Gli studenti non dovranno fare marketing, ma dovranno ricordare alcuni di questi
concetti, soprattutto nelle lezioni con Marta, che tratta argomenti legati al
marketing. E, chiude Daniele, come game designer e game developer ``per lo più
speriamo che il marketing non ci rovini i giochi'', soprattutto nel videoludico.

% ---------------------------------------------------------------------
\section{Game design e game development}
\label{a1:sec:design-development}

Ruggero torna sulla distinzione più importante della slide, ``proseguendo con
questa strada'':
\begin{description}
  \item[Game design] è chi \textbf{progetta l'esperienza} tramite la
        \textbf{creazione delle meccaniche} di gioco.
  \item[Game development] è chi fa \textbf{mediazione} tra la persona che scrive
        il game design e la \textbf{casa editrice}.
\end{description}

\subsection{I red flag dei designer}

Ruggero lo ripete perché lo vede ogni giorno nel suo lavoro: tra giochi
pubblicati e prototipi, ne valuta ``quasi uno al giorno, se faccio un conto
annuale'', e solo una minima percentuale arriva alla pubblicazione. Molte volte
i \textbf{red flag} vengono proprio da designer che non hanno presente questa
definizione, e che cominciano a mettere \textbf{clausole sulla non modifica del
gioco}: il gioco deve restare com'è. ``Questa cosa è impossibile con un
publisher, spesso''. Il publisher, per forza di cose, potrà chiedere di cambiare
la grafica, di reimpaginare il regolamento, di cambiare i materiali perché
troppo costosi: chiederà dei \textbf{compromessi}. E la persona che, all'interno
del publisher, fa questo lavoro è il \textbf{developer}.

Sta al designer capire \textbf{su che cosa del proprio gioco si può trovare un
compromesso e su che cosa no}. Publisher diversi lavorano in modo diverso, e il
designer ha sempre l'autorità di dire ``no, non voglio pubblicare questo gioco
se queste sono le modifiche''. Ma deve sapere chi fa il lavoro di development,
perché se il suo obiettivo è la \textbf{pubblicazione tradizionale} si
interfaccerà con questa figura.

\subsection{Daniele: una distinzione che manca al videogioco}

Daniele sottolinea che questa distinzione è \textbf{tipica del mondo analogico}
ed è ``molto interessante''. In Italia la parola \emph{developer} è diventata
negli anni quasi un sinonimo di programmatore di videogiochi, o di chiunque
lavori nel videoludico. Secondo lui c'entra anche il fatto che qualche anno fa
il sito \textbf{Gamasutra} ha cambiato nome in \textbf{Game Developer}, e la
parola è diventata un ombrello per tutte le professioni del videogioco. ``Mi fa
un po' ribrezzo'': il game developer dovrebbe essere un ruolo di design, un
designer che, come nell'analogico, fa \textbf{mediazione}. Nel videogioco i
publisher hanno in effetti i loro developer interni, che adattano il gioco, per
esempio alle console e ai loro kit di sviluppo. Ma la distinzione dovrebbe
esistere anche nel videogioco, perché è ``difficilissimo'' avere la mentalità
per fare le due cose insieme: chi fa design ha un focus che raramente si può
sovrapporre a quello di chi fa mediazione.

A Daniele piace la parola \emph{develop}, \textbf{sviluppare}, nel senso di
prendere qualcosa che è stato progettato e \textbf{renderlo possibile nel mondo
reale}. Ci sono progetti che hanno bisogno di uno sviluppo che ``polishi''
l'esperienza: la semplifichi, la migliori, la renda comprensibile e la concluda,
ma soprattutto la renda \textbf{stampabile} e \textbf{fruibile} per l'utente. La
metafora crossmediale è quella della \textbf{fotografia}: chi sviluppa le
fotografie non è chi scatta la foto, ma chi la rende visibile, anche
fisicamente. ``È un po' retro come metafora, ma ci capiamo''. ``È vero, però è
giusto'', conferma Ruggero.

% ---------------------------------------------------------------------
\section{Il mercato: Italia ed estero}
\label{a1:sec:mercato}

Il corso è in italiano e si svolge in Italia, ma bisogna conoscere anche il
\textbf{mercato estero}: molti giochi progettati da autori italiani vengono
pubblicati all'estero, e conoscere il mercato estero aiuta a capire come il
proprio gioco può arrivarci, ``o viceversa, come non può arrivarci''. La
tabella~\ref{a1:tab:mercato} riassume il confronto della slide.

\begin{table}[htbp]
\centering
\small
\begin{tabular}{@{}p{6.4cm}p{6.4cm}@{}}
\toprule
\textbf{Italia} & \textbf{Estero (soprattutto Europa)} \\
\midrule
il gioco è ancora visto come un'attività \textbf{infantile}, anche se la
normalizzazione è in crescita &
forte tradizione di gioco anche in un contesto \textbf{familiare adulto},
soprattutto ``dalla Germania in su'' \\
risposta \textbf{conservativa} sul tema ludico &
\textbf{premi di prestigio} riconosciuti in tutto il mondo (Spiel des Jahres,
As d'Or) \\
forte focus sulla \textbf{localizzazione} di giochi stranieri &
si cercano soprattutto giochi \textbf{originali} \\
investimento \textbf{molto basso} sulla creazione di giochi originali interni;
pochissimi designer assunti &
\textbf{developer interni} pagati con un contratto (Francia, Germania,
Danimarca\ldots) \\
\bottomrule
\end{tabular}
\caption{Il mercato del gioco analogico in Italia e all'estero, secondo la
slide di Ruggero e i commenti di Daniele.}
\label{a1:tab:mercato}
\end{table}

\subsection{L'Italia}

\begin{itemize}
  \item In Italia il gioco è ancora visto come un'\textbf{attività infantile},
        anche se la \textbf{normalizzazione} è in crescita. C'è ancora una forte
        risposta \textbf{conservativa} sul tema ludico.
  \item C'è un forte focus sulla \textbf{localizzazione} dei giochi stranieri: in
        una qualunque ludoteca o negozio, molti dei giochi da tavolo e di ruolo
        sono di publisher o designer esteri, pubblicati in italiano da publisher
        italiani. (La localizzazione, cioè portare un titolo dall'estero in
        Italia, esiste come funzione all'interno del publisher; Ruggero non l'ha
        messa nella slide delle figure perché non è il focus del corso.)
  \item C'è, ``purtroppo, preparatevi psicologicamente'', un \textbf{investimento
        molto basso} sulla creazione di giochi originali interni alla casa
        editrice. Pochissime case editrici italiane hanno designer interni pagati
        a tempo pieno con contratto a tempo indeterminato. Daniele lo dice senza
        mezzi termini: ``il mercato è molto piccolo, conosciamo praticamente
        tutto il settore'', e le persone che in Italia fanno game design
        analogico a tempo pieno con un contratto a tempo indeterminato ``sono
        meno delle dita delle mie due mani''. Di solito la formula è quella del
        \textbf{freelance}, del contratto, delle \textbf{royalty}. Chi va da un
        publisher aspettandosi di essere assunto a tempo pieno come game
        designer deve sapere che ``non siamo purtroppo ancora in quel tipo di
        situazione''.
  \item C'è però un contraltare: in Italia, soprattutto nel gioco di ruolo ma
        anche nel gioco da tavolo, c'è stata molta \textbf{imprenditoria}, anche
        grazie a \textbf{Kickstarter}, che ha aperto le porte ad altri sistemi
        di \textbf{crowdfunding}. È un altro argomento, perché richiede un team,
        un investimento e così via.
  \item E gli autori non mancano, anzi: le occasioni di incontro tra autori
        vedono ogni anno un aumento dei possibili autori\footnote{La trascrizione
        dice ``i DAG a Parma''; il riferimento è probabilmente a IDEAG, il
        circuito di incontri tra autori di giochi italiani, ma non è
        ricostruibile con certezza.}. Bisogna ``sgomitare'' tra tanti autori,
        che magari arrivano da esperienze molto diverse.
\end{itemize}
Ruggero precisa: l'investimento basso è ``dall'alto''. Le persone che fanno e
vogliono fare giochi sono tantissime e in aumento; scarseggia l'investimento,
sia dei publisher disposti a pagare stipendi sia dei \textbf{bandi} pubblici. I
bandi esistono (Beatrice ne parlerà per i giochi pervasivi, che spesso si
muovono in quella direzione, e Ruggero cita qualcuno che sta facendo un card
game con un bando italiano, ``una cosa assurda, ma lo sta facendo''), ma sono
una rarità, soprattutto rispetto all'estero.

\subsection{L'estero}

\begin{itemize}
  \item L'estero (dipende dalla nazione, ma in generale il resto d'Europa) ha una
        forte \textbf{tradizione di gioco} anche in un contesto \textbf{familiare
        adulto}, soprattutto ``dalla Germania in su''.
  \item Ha \textbf{premi di prestigio} riconosciuti in tutto il mondo, come lo
        \textbf{Spiel des Jahres} tedesco e l'\textbf{As d'Or} francese\footnote{La
        trascrizione dice ``lo Spield AR, la SDOR''.}.
  \item Cerca soprattutto \textbf{giochi originali}. Lo si vede dal fatto che
        molti autori italiani di giochi da tavolo pubblicano in Germania o in
        Francia e \textbf{non vengono pubblicati in Italia}, ``che è tutto
        dire''.
  \item Ha \textbf{figure di development interne}: non è raro vedere publisher
        francesi, tedeschi, danesi con developer pagati internamente, a
        contratto.
\end{itemize}
``Poi la situazione è ancora più drammatica'', ammette Ruggero, ``ma per ora
vogliamo tenere il morale un po' alto. È l'introduzione. Ci sarà tempo di andare
nella tristezza''.

Prima di proseguire, Ruggero insiste sulle domande: con lui si può interrompere
anche mentre parla, alzando la mano con le reaction o semplicemente togliendo il
microfono e chiedendo. ``Non mi offenderò mai, anzi, sarò felice di vedervelo
fare e mi ricorderò positivamente di voi''.

% ---------------------------------------------------------------------
\section{Perché serve una terminologia comune}
\label{a1:sec:terminologia}

La serie di slide successiva è ``una piccolissima parte di fondamenti di
design'' (``non spaventatevi, vi vedo che vi spaventate un po'\,''). Anche qui
molti temi verranno ripresi in dettaglio, ma l'idea è arrivare alla prossima
lezione con almeno un \textbf{glossario}, un \textbf{linguaggio condiviso}: se
qualcuno dice \emph{game developer}, tutti sanno di che cosa sta parlando.

Ruggero parte da una domanda, e precisa che non è retorica: \textbf{perché ci
serve una terminologia comune?} Gli studenti rispondono:
\begin{itemize}
  \item Gianmarco, ``molto pragmatico'': perché bisogna lavorare in team e
        capirsi. Vero, nel corso e anche fuori;
  \item un'altra studentessa: per la \textbf{comunicazione}, tra di noi e con
        tutti quelli che sono interessati nel settore. Vero anche questo:
        conoscere i termini comuni permette di parlare con chi fa game design o
        development nel dettaglio, nel tecnico, nello specifico.
\end{itemize}
Ruggero aggiunge altre due ragioni.

\paragraph{Saper comunicare il proprio gioco.} La prima riguarda la
\textbf{comunicazione del gioco} verso l'esterno. ``E lo dico non tirando acqua
al mio mulino, ma anzi togliendola'': \textbf{non è spesso il buon game design
a far diffondere un gioco}. ``Vorrei che lo fosse, sarebbe bellissimo, ma
purtroppo non è così''. Ci sono giochi bellissimi che non sanno comunicare, che
non sanno dire che cosa fanno, e che quindi è difficile, da fuori, trovare
invitanti; e ci sono giochi mediocri dal punto di vista del design che si
comunicano molto bene, e che quindi conoscono in molti. Saper parlare di quello
che si fa è utile \textbf{a sé stessi}, in primis.

\paragraph{Documentare per sé stessi.} La seconda è la \textbf{documentazione}:
documenti di design, di brainstorming e così via. ``Sembrerà una battuta, ma
provate a recuperare un documento di brainstorming una settimana dopo, se non
avete scritto bene le cose che avete scritto''. Chi l'ha fatto sa già di che
cosa si parla; ``chi non l'ha ancora vissuto, forse lo vivrà''.

% ---------------------------------------------------------------------
\section{Game e play}
\label{a1:sec:game-play}

La \textbf{prima, importantissima divisione terminologica} è quella tra
\emph{game} e \emph{play}. In italiano c'è un problema: \textbf{gioco} e
\textbf{giocare} hanno la stessa radice, mentre in altre lingue, tra cui
l'inglese, sono due parole strutturalmente diverse:
\begin{description}
  \item[Game] il \textbf{gioco come artefatto}: il gioco come ecosistema di
        meccaniche, come struttura di regole, come prodotto (``tra
        virgolette''), come oggetto fisico.
  \item[Play] l'\textbf{atto}, il verbo: le persone che giocano, l'approccio
        fisico all'attività del gioco.
\end{description}
È importante perché gli studenti faranno \textbf{game design}; esistono anche
ambiti di \emph{play design}, ma qui ``creerete giochi, non giocattoli''.

\subsection{Giochi e giocattoli}

Qual è la differenza? Uno studente risponde: \textbf{la presenza di regole}.
Esatto: un gioco ha delle regole, un'attività di \emph{play} non
necessariamente. Con un giocattolo si può fare \emph{play}: c'è chi gioca con i
Lego, chi gioca a guardia e ladri o a nascondino al parco, e magari alcune di
quelle regole sono ``più o meno semi-solidificate''. Ma il \emph{play} permette
spesso un approccio non strettamente legato alle regole.

Viceversa, \textbf{un game può esistere anche se nessuno lo sta giocando}. Se si
prepara su un tavolo una scacchiera con i pezzi degli scacchi, chi passa e la
vede riconosce il gioco degli scacchi, anche se nessuno sta giocando. Se invece
si vedono dei bambini correre in un parco giochi, quei bambini \emph{stanno
giocando}, ma non stanno giocando \emph{a un gioco} specifico. ``È importante
separare questa cosa''; ci si tornerà, ``ma per ora abituatevi a questa
differenza''.

\paragraph{Approfondimento: giocattoli, giochi e regole.} La stessa distinzione
è il punto di partenza del capitolo che Jesse Schell, in \emph{The Art of Game
Design}, dedica alla definizione di gioco (\emph{The Experience Rises Out of a
Game}). Schell nota che usiamo perfino un linguaggio diverso: ``un gioco è
qualcosa a cui si gioca'' (\emph{a game is something you play}), ``un giocattolo
è qualcosa con cui si gioca'' (\emph{a toy is something you play with}); e
arriva a definire il giocattolo come ``un oggetto con cui si gioca''. Analizzando
poi la definizione di gioco di Avedon e Sutton-Smith, ne ricava che ``i giochi
hanno regole, i giocattoli no'', e che le regole sono uno degli aspetti
\emph{definitori} del gioco, insieme alla volontarietà, agli obiettivi, al
conflitto e alla possibilità di vincere o perdere. Schell osserva anche che
molti giochi sono costruiti \textbf{sopra} un giocattolo (``una palla è un
giocattolo, ma il baseball è un gioco''), e tra i suoi consigli di
prototipazione c'è proprio \emph{Build the Toy First}: costruire prima il
giocattolo, cioè qualcosa che sia già piacevole da manipolare, e poi il gioco
intorno. È una lettura utile della distinzione di Ruggero: il \emph{play} può
fare a meno delle regole, il \emph{game} no.

% ---------------------------------------------------------------------
\section{Il framework MDA (e il DPE)}
\label{a1:sec:mda}

\subsection{Meccaniche, dinamiche, estetiche}

Un'altra terminologia fondamentale è quella del framework \textbf{MDA}
(\emph{Mechanics, Dynamics, Aesthetics}; in italiano a volte ``MDE''). Viene da
un \textbf{paper} nato nell'ambito del videogioco\footnote{È il paper di Robin
Hunicke, Marc LeBlanc e Robert Zubek, \emph{MDA: A Formal Approach to Game
Design and Game Research}, del 2004. Daniele aggiunge il paper alla cartella
condivisa durante la lezione.}, che ha ``solidificato'' queste terminologie, e
divide un gioco in tre aree:
\begin{description}
  \item[Meccaniche] le \textbf{regole per come sono scritte}. ``Ogni volta che
        pescate un 2 fate dei punti'': quella è una meccanica.
  \item[Dinamiche] l'\textbf{applicazione in tempo reale} delle meccaniche. La
        dinamica è il giocatore che, pescando una carta, \textbf{spera} che esca
        un 2. Quel processo non è una regola: ``il gioco non vi dice che dovete
        sperare che esca un 2''.
  \item[Estetiche] l'insieme di quello che il gioco vuole comunicare,
        l'\textbf{esperienziale}.
\end{description}

L'esempio di Ruggero mette insieme le tre aree. In un gioco in cui pescare un 2
fa fare punti (meccanica) e in cui quindi si spera che esca un 2 (dinamica),
\textbf{che estetica si ottiene se nel mazzo di 100 carte c'è un solo 2? E se ce
ne sono 50?} L'emozione di aver pescato \emph{l'unico} 2 del mazzo è
completamente diversa da quella di pescarne uno sapendo che ce ne sono
cinquanta e che è molto probabile trovarlo. Le regole sono quasi identiche; cambia
l'esperienza.

\subsection{Il falso amico: ``estetica''}

Daniele avverte di non sottovalutare la parola \textbf{estetica}, che in italiano
è un \textbf{falso amico}: sembra indicare la parte grafica, e invece non lo è.
Si incontreranno designer che hanno studiato MDA e altri che non l'hanno
studiato, e lo si capirà proprio da questo: i secondi diranno che le
\emph{aesthetics} sono la grafica. Le \emph{aesthetics} sono la parte
\textbf{esperienziale}: come nell'esempio di Ruggero, quello che lascia al
giocatore aver pescato l'unico 2 del mazzo.

\subsection{Il DPE: design, play, experience}

Per gli italiani, aggiunge Daniele, è più facile capire l'MDA attraverso un
altro paper, quello del \textbf{DPE}\footnote{È il framework \emph{Design, Play,
Experience} di Brian Winn, proposto come estensione dell'MDA per i giochi
educativi (\emph{serious games}). Anche questo paper viene aggiunto da Daniele
alla cartella condivisa.}, che di fatto fa un parallelismo assegnando
\textbf{design, play, experience}:
\begin{itemize}
  \item le \textbf{meccaniche} sono il lato \textbf{design};
  \item le \textbf{dinamiche} sono il lato \textbf{play};
  \item le \textbf{estetiche} sono l'\textbf{experience}, l'esperienza che riceve
        chi gioca.
\end{itemize}
``Se vivete le \emph{aesthetics} come l'esperienza che vive chi sta giocando,
state a postissimo''. In questo momento Beatrice lascia la lezione (``ciao,
Bea''). I paper di MDA e DPE sono in inglese; chi vuole può dargli una lettura.

\begin{figure}[htbp]
\centering
\begin{tikzpicture}[
    font=\small,
    blocco/.style={draw, thick, rounded corners=3pt, fill=gray!12,
                   minimum width=3.2cm, minimum height=1.0cm, align=center},
    dpe/.style={font=\footnotesize\itshape, align=center},
    esempio/.style={draw, rounded corners=2pt, text width=3.3cm,
                    align=left, font=\footnotesize, inner sep=4pt},
    freccia/.style={-{Stealth[length=2.5mm]}, thick}
  ]
  \node[blocco] (m) at (0,0)   {\textbf{Meccaniche}\\ \footnotesize le regole scritte};
  \node[blocco] (d) at (4.3,0) {\textbf{Dinamiche}\\ \footnotesize le regole in azione};
  \node[blocco] (a) at (8.6,0) {\textbf{Estetiche}\\ \footnotesize l'esperienza};
  \draw[freccia] (m) -- (d);
  \draw[freccia] (d) -- (a);
  \node[dpe, above=0.15cm of m] {DPE: \textbf{design}};
  \node[dpe, above=0.15cm of d] {DPE: \textbf{play}};
  \node[dpe, above=0.15cm of a] {DPE: \textbf{experience}};
  \node[esempio, below=0.45cm of m] (em)
    {``Ogni volta che peschi un 2 fai punti''};
  \node[esempio, below=0.45cm of d] (ed)
    {pescando, il giocatore \emph{spera} che esca un 2 (nessuna regola glielo
     chiede)};
  \node[esempio, below=0.45cm of a] (ea)
    {un solo 2 su 100 carte: pescarlo è un evento;\\
     50 su 100: è quasi scontato};
  \draw[dashed] (m) -- (em);
  \draw[dashed] (d) -- (ed);
  \draw[dashed] (a) -- (ea);
  % prospettive
  \draw[freccia, gray] ([yshift=-0.3cm]em.south west) --
    node[below, font=\footnotesize\itshape, text=black]
    {il designer agisce sulle meccaniche}
    ([yshift=-0.3cm]ea.south east);
  \draw[freccia, gray] ([yshift=-1.2cm]ea.south east) --
    node[below, font=\footnotesize\itshape, text=black]
    {chi gioca incontra prima l'esperienza}
    ([yshift=-1.2cm]em.south west);
\end{tikzpicture}
\caption{Il framework MDA con l'esempio del ``2'' usato da Ruggero, e la
corrispondenza con il DPE suggerita da Daniele. Le due frecce in basso
riassumono la lettura classica del paper: il designer costruisce le meccaniche
per ottenere un'esperienza, mentre il giocatore parte dall'esperienza.}
\label{a1:fig:mda}
\end{figure}

\paragraph{Approfondimento: le due direzioni dell'MDA e le otto estetiche.} Le
frecce della figura~\ref{a1:fig:mda} riprendono un'idea centrale del paper
originale, che la lezione non esplicita: designer e giocatore guardano il gioco
da lati opposti. In \emph{Challenges for Game Designers}, Brenda Brathwaite e
Ian Schreiber la riassumono così: i designer si chiedono quale estetica
vogliono ottenere, definiscono le dinamiche che porterebbero a quella
sensazione, e poi creano le meccaniche che producono quelle dinamiche. Lo stesso
libro conferma l'avvertenza di Daniele sul falso amico: ``le \emph{aesthetics}
non si riferiscono all'aspetto del gioco, ma alla risposta emotiva'' che il
designer vuole suscitare. E offre esempi di dinamiche nel senso di Ruggero: in
uno sparatutto i punti di \emph{spawn} sono una meccanica, mentre appostarsi
accanto a uno di essi per colpire chi riappare (lo \emph{spawn camping}) è una
dinamica; negli scacchi i movimenti dei pezzi sono meccaniche, mentre le
aperture codificate sono dinamiche. Per dare un nome alle estetiche, il paper
MDA propone una tassonomia, poi detta degli ``otto tipi di \emph{fun}'', che
Schell riprende in \emph{The Art of Game Design} come \emph{LeBlanc's Taxonomy of
Game Pleasures}: \textbf{sensazione}, \textbf{fantasia}, \textbf{narrativa},
\textbf{sfida}, \textbf{compagnia} (\emph{fellowship}), \textbf{scoperta},
\textbf{espressione} e \textbf{sottomissione} (il piacere di entrare nel
``cerchio magico'' e lasciarsi andare a un passatempo). È un primo esempio di
quello che i docenti chiedono con le parole vietate della prossima sezione:
invece di dire che un gioco è ``divertente'', dire \emph{quale} esperienza
offre.

% ---------------------------------------------------------------------
\section{Le parole vietate}
\label{a1:sec:parole-vietate}

``Attenzione, adesso qua triggeriamo delle persone''. Ruggero annuncia una regola
che ``tornerà'': nel corso ci sono delle \textbf{parole vietate}. ``Non perché
siamo dei censori''; il perché si vedrà in una lezione apposita, perché
spiegarlo adesso richiederebbe un'ora. Ma Daniele vuole spiegarne ``un
minimo''.

\subsection{La regola}

In \textbf{nessun materiale scritto} consegnato ai docenti, task compresi, sono
accettate queste parole, né in inglese né in italiano:
\begin{center}
\begin{tabular}{@{}ll@{}}
\toprule
\textbf{Inglese} & \textbf{Italiano} \\
\midrule
\emph{fun} & divertimento \\
\emph{immersion} & immersione \\
\emph{entertainment} & intrattenimento \\
\bottomrule
\end{tabular}
\end{center}
``Non li voglio neanche vedere''. Daniele lo ripete con un paragone: ``se mi fate
i disegni con Paint, benissimo, top, fantastico, siamo ancora best friend. Se mi
scrivete \emph{divertimento}, io sbrocco''. È una \textbf{regola dell'accademia}:
negli stessi corsi triennali di narrative design e game design le parole sono
vietate agli studenti per tre anni; ``per voi solo per tre mesi''. Daniele
aggiungerebbe volentieri un quarto termine, \textbf{gamification}.

\subsection{Perché}

Il motivo approfondito lo darà Ruggero nella lezione apposita. Daniele dà il
suo. Ruggero e Diego Ricchiuti sono molto contrari a \emph{fun} e a
\emph{immersion}; lui ha aggiunto \emph{intrattenimento}. Sono \textbf{parole
troppo passe-partout e troppo ambigue}:
\begin{itemize}
  \item \textbf{intrattenimento} in Italia è stata usata per anni come un modo
        per \textbf{sminuire}: ``ma sì, tanto è un gioco solo di
        intrattenimento, è un romanzo solo di intrattenimento''. Ma anche se un
        prodotto è ``solo di intrattenimento'', non vuol dire che non ci si debba
        lavorare con attenzione. E poi intrattenimento ``non significa niente'':
        ci sono tanti modi in cui una cosa può intrattenere;
  \item lo stesso vale per \textbf{fun}, e per \textbf{immersion}, che ha un
        problema in più: è una parola passe-partout per cui si può dire che un
        gioco è ``immersivo'' nel senso generale del termine, ma poi non si sa
        davvero che cosa voglia dire, \emph{immersivo in che senso}.
\end{itemize}

L'obiettivo didattico è che gli studenti \textbf{vadano oltre queste parole}:
``un po' come quando alle medie vi dicono smettetela di dire che tutto è bello,
perché ci sono vari gradi di bello, ci sono altri modi di dire bello''. Bisogna
\textbf{decostruirle}: quando si sta per usarle, chiedersi \textbf{che cosa si
intende} con quella parola e cercarne una \textbf{più precisa}.

Ruggero chiede, in uno dei ``rari casi'', di \textbf{non fare domande adesso} su
questa regola: ogni anno arrivano domande e commenti, e ci saranno lezioni
apposite di tre ore, con strumenti per capire il perché. Ma da subito gli
studenti devono abituarsi a evitare questi termini, sia perché si spera che
abbiano termini più precisi, sia perché queste parole ``portano con sé dei
problemi agganciati'' che si vedranno. ``Non spaventatevi, non verrete sparati in
un vicolo cieco di notte''.

\paragraph{Una piccola gamification.} Il divieto va visto come uno
\textbf{stimolo}, ``un po' come se stessimo gamificando la vostra esperienza'',
una sorta di user experience delle lezioni. Imponendosi di non usare queste
parole, si noterà come migliora la capacità di \textbf{descrivere i giochi} e le
proprie esperienze ludiche, e di essere più precisi. È un invito a ``obbligarsi a
essere più precisi e ad avere una mentalità più critica''. La lezione di
Ruggero sull'\textbf{involvement}, promette Daniele, sarà ``fenomenale'' e
``vi aprirà la testa''.

\paragraph{Approfondimento: anche i libri faticano a definire il ``fun''.} La
difficoltà che il divieto vuole far emergere è visibile anche nei testi di game
design. Schell, in \emph{The Art of Game Design}, prova a definire il \emph{fun}
e propone ``piacere con sorprese'' (\emph{pleasure with surprises}), ammettendo
subito che è strano poter usare una parola per tutta la vita, sapendo con
certezza che cosa significa, senza riuscire a esprimerla chiaramente quando
qualcuno lo chiede; e invita a mettere alla prova la definizione con
controesempi. La sua lente sul \emph{fun} (\emph{The Lens of Fun}) si limita,
non a caso, a chiedere quali parti del gioco siano \emph{fun} e \emph{perché}.
Anche \emph{immersion} ha molte letture: in \emph{21st Century Game Design},
Bateman e Boon la descrivono, tra l'altro, come la capacità di ``perdersi''
nel mondo artificiale presentato dal gioco, legata alla sospensione
dell'incredulità e ai giochi di mimesi. Sono significati diversi sotto la stessa
parola: proprio il motivo per cui, in un documento di design, conviene dire
quale si intende.

% ---------------------------------------------------------------------
\section{Chi crea, chi gioca, chi facilita}
\label{a1:sec:ruoli-esperienza}

Un'altra distinzione importante è quella tra \textbf{chi crea}, \textbf{chi
gioca} e \textbf{chi facilita}. Spesso, quando si fa un gioco, lo si fa provare
di persona a chi gioca. Ma ``la norma dei giochi pubblicati non funziona così'':
quando si pubblica un gioco non si vede chi lo compra e chi lo gioca, e non si
sa se qualcuno, a casa propria con i propri amici, segue le regole alla lettera e
interpreta correttamente quello che il designer aveva descritto. Quindi già in
fase di progettazione bisogna tenere a mente \textbf{chi sta facendo game design
e chi sta giocando}.

E nei giochi a evento, nei giochi pervasivi, nei giochi di ruolo da tavolo e dal
vivo (LARP) c'è un terzo pezzo del puzzle, ``importantissimo'': \textbf{chi
facilita}. La domanda da farsi è: \textbf{il regolamento è scritto per chi gioca
o per chi deve facilitare il gioco?} Bisogna abituarsi a pensare a tutte le
persone coinvolte nell'esperienza.

% ---------------------------------------------------------------------
\section{Playtest e iterazione}
\label{a1:sec:playtest}

\subsection{Nessun gioco nasce perfetto}

L'ultima voce del glossario è ``un'altra cosa importantissima'':
\textbf{playtest} e \textbf{iterazione}. Il corso è \textbf{pragmatico}:
``specialmente da parte mia'', dice Ruggero, non si vedranno teorie ``messe lì
per fare accademia sul game design nelle nostre torri d'avorio''. Ci saranno
teorie e framework, ma tutti estremamente pragmatici, perché l'obiettivo non è
solo capire i giochi, ma \textbf{saperli fare e iterare}.

Playtest e iterazione sono \textbf{inevitabili} nel game design, perché
\textbf{nessun gioco nasce perfetto}, soprattutto perché le persone sono
diverse: un gioco perfetto per una persona può essere incomprensibile per
un'altra. Chi progetta giochi per altre persone deve capire, dopo averli fatti
provare, \textbf{che cosa cambiare, che cosa sistemare, che cosa aggiungere, che
cosa togliere}. Ci saranno lezioni su playtest, prototipazione, raccolta dei
feedback, metodi etnografici. Ma bisogna abituarsi già ora all'idea che
\textbf{non esiste la consegna definitiva}. A un certo punto i task vanno
consegnati per forza di cose, ma, come diceva Daniele all'inizio, alla fine del
corso gli studenti avranno un \textbf{prototipo}, ``che sarà un prototipo'', e
gli strumenti per portare avanti il lavoro anche dopo.

\subsection{L'aiuto dopo il corso}

Non solo gli strumenti: anche l'aiuto dei docenti. Daniele racconta di aver
portato a un editore, a Play\footnote{Probabilmente \emph{Play -- Festival del
Gioco} di Modena; la trascrizione non permette di esserne certi.}, uno dei
progetti realizzati l'anno precedente da due studenti, un gioco di ruolo in
solitario (``un \emph{solo game} RPG molto, molto bello''). Se gli studenti
vorranno trasformare il proprio lavoro in un prototipo vero e proprio, possono
contattarlo: lo fa ``per lavoro, ma anche per passione''. Gestisce anche una
community per aiutare i designer ``un po' depressi'' a rimettersi a progettare.

Quindi: ``non sentitevi persi se durante la fase di iterazione vi sentite
persi, è normale''.

\subsection{Il dizionario comune e il production dictionary}

Daniele riprende un punto detto ``benissimo'' da Ruggero e che secondo lui non è
stato sottolineato abbastanza: il \textbf{dizionario comune} serve anche a
\textbf{non perdersi} mentre si itera. Se da un'iterazione all'altra i
documenti usano termini completamente diversi, e non c'è un dizionario di base,
ci si perde.

Nel videogioco, una delle primissime cose che un designer fa quando entra in
un'azienda è il \textbf{production dictionary}: si annota, in un file privato,
tutte le parole che usano i colleghi, per allinearsi al dizionario comune
dell'azienda, perché lì su un prototipo possono lavorare anche cento persone.
Nel gioco da tavolo non si fa molto, perché il mondo dell'analogico è
``talmente più piccolo''; ma serve comunque imparare un \textbf{dizionario
minimo}, dai docenti e dai professionisti. E il corso ne darà una parte,
entrando nel dettaglio di alcuni termini e framework.

\paragraph{L'esempio del feedback.} Perché serva, lo mostra un esempio. Si
testa un gioco e ci si annota come feedback: ``questo gioco deve essere più
divertente''. Che cosa vuol dire? Come lo si migliora? Quella frase non dà
elementi utili. Se invece si scrive che deve essere \textbf{più adrenalinico},
\textbf{più frenetico}, \textbf{più rilassato}, \textbf{più cozy}, o
\textbf{più bilanciato matematicamente}, si sa dove andare a lavorare.

\paragraph{Approfondimento: la regola del ciclo.} Il principio ``nessun gioco
nasce perfetto'' ha un nome preciso in \emph{The Art of Game Design}: è la
\textbf{Rule of the Loop}, la regola del ciclo. ``Più volte provi e migliori il
tuo design, migliore sarà il tuo gioco''. Schell precisa che non è una lente,
cioè un punto di vista, ma ``una verità assoluta'', senza eccezioni: ci saranno
momenti in cui si cercherà di convincersi che ``questa volta il design è così
buono che non serve testarlo'', e ogni volta se ne pagherà il prezzo. Ed è
interessante la ragione per cui la regola pesa di più sul digitale: secondo
Schell il tempo e il denaro necessari per provare e correggere un sistema sono
molto maggiori nei videogiochi che nei giochi tradizionali. Per un gioco di
carte o da tavolo, con il tempo per provarlo e cambiarlo più e più volte, il
consiglio è semplicemente di farlo. Il gioco analogico è, da questo punto di
vista, la palestra ideale dell'iterazione.

% ---------------------------------------------------------------------
\section{Tabletop Simulator}
\label{a1:sec:tts}

\subsection{Il software di settore}

``Last but not least'': il corso usa un software, \textbf{Tabletop Simulator},
che si compra su \textbf{Steam} e costa tra i \textbf{10 e i 20 euro} (``aspettate
che vada in sconto, va in sconto ogni due per tre''). Daniele spiega perché:
\begin{itemize}
  \item è \textbf{il software di settore} del gioco da tavolo, quello che i
        publisher usano per i prodotti digitali, interni ed esterni;
  \item essendo su Steam, permette di fare \textbf{playtest} e di giocare con gli
        amici di Steam. Daniele stesso sta testando lì il proprio gioco;
  \item funziona anche su Mac: Windows o Mac, nessun problema.
\end{itemize}
Ci sarà una lezione apposita, di Ruggero, sulla creazione del proprio
\textbf{modulo} di gioco.

\paragraph{Che cosa non è.} Daniele chiarisce subito un ``fraintendimento
gigantesco'', soprattutto per chi ha già esperienza: Tabletop Simulator
\textbf{non richiede programmazione} e \textbf{non è un software per
trasformare i giochi da tavolo in videogiochi}. È un software che
\textbf{simula dei componenti fisici su un tavolo}: serve per giocare i board
game in digitale. Simula anche la fisica: c'è perfino il comando per
\textbf{ribaltare il tavolo} (``non usatelo ora'').

\paragraph{La partnership mancata.} Da anni Daniele prova a fare una
partnership con gli sviluppatori: due anni prima ha contattato l'azienda per
acquistare delle chiavi per gli studenti, dato che il software si usa nei corsi.
L'azienda ha risposto che non fa accordi commerciali. ``Credo sia la prima
azienda nella mia vita che dice che non vuole che io le dia dei soldi''. Gli
studenti dovranno quindi provare ad acquistarlo da sé; se i docenti riusciranno a
darglielo, bene.

\paragraph{La domanda di Gianmarco.} Gianmarco segue le lezioni da un computer
aziendale ``super chiuso'', su cui non può installare Steam: deve usare Tabletop
Simulator durante le lezioni? Ruggero risponde che per la lezione dedicata è
meglio poter seguire e rifare in diretta. Ma il software non sarà richiesto
necessariamente per il project work: un gruppo potrebbe fare un gioco pervasivo
o un gioco di ruolo, non un gioco da tavolo. Essendo un software di settore, è
comunque bene conoscerlo. Ruggero aggiunge che le versioni non ufficiali che
circolano non permettono di giocare in molti, e quindi non servono per
playtestare o giocare insieme.

\subsection{Una dimostrazione}

Daniele mostra il software aperto da Steam:
\begin{itemize}
  \item si può creare una partita \textbf{single player} o \textbf{multiplayer};
        per giocare tra compagni si crea un server multiplayer e ci si aggiunge
        come amici su Steam;
  \item si dà al server un nome, un numero di giocatori, eventualmente una
        password; creato il server, si apre la scelta dei giochi;
  \item ci sono giochi gratuiti, giochi da acquistare (DLC, non obbligatori) e
        soprattutto un \textbf{Workshop} ``ampissimo'', con decine di migliaia di
        giochi caricati dagli utenti: ``se c'è un board game che conoscete,
        quel board game è su Tabletop Simulator, giocabile dal Workshop''.
        L'esempio mostrato è l'open beta ufficiale del gioco da tavolo di
        \emph{Slay the Spire}\footnote{La trascrizione dice ``Slade Spire''.}:
        molti publisher, per farsi pubblicità, mettono il proprio gioco in demo o
        giocabile in digitale;
  \item c'è un \textbf{tutorial} iniziale che spiega come muoversi sul tavolo;
        con \emph{Esc} e \emph{Help} si trovano i comandi da tastiera.
\end{itemize}

\subsection{Il primo task}

Il primo (e per ora unico) task è questo: \textbf{entro il 10 ottobre}, giorno
della lezione di Ruggero su Tabletop Simulator, giocare \textbf{tre o quattro
volte in totale}, online, tra compagni di corso o con amici. Non si tratta di
imparare a creare moduli, ma di prendere la mano \textbf{lato giocatore}: aprire
il software, impararne i comandi, giocare. ``Per me potete anche andare a
prendere domino o mahjong'' (``mahjong è bellissimo tra l'altro'', aggiunge
Daniele, che sta facendo un videogioco sul mahjong), o il pachisi o il go. Gli
anni precedenti si chiedeva a ciascuno di proporre ai compagni un gioco da
tavolo che apprezzava. È anche un modo per conoscersi. Ruggero spiega il
perché: usarlo per giocare è facile e intuitivo, usarlo per creare moduli è più
complesso; ma chi arriva al 10 ottobre avendo già aperto il software e giocato,
troverà più intuitivo anche creare un modulo. E l'avvertimento: ``non riducetevi
all'ultimo, non fate una partita il 9 sera''.

\subsection{Tabletopia e altri vantaggi}

Tancredi racconta di aver provato a caricare un suo gioco su
\textbf{Tabletopia}. ``Questo è un vaso di Pandora'', risponde Daniele.
Tabletopia è molto usato, ma secondo lui il suo strumento di creazione dei moduli
è pessimo, ed è il motivo per cui la maggioranza dei moduli ufficiali dei
publisher è su Tabletop Simulator. Tabletopia ha un vantaggio vero, è
\textbf{gratuito}, ma ha scelte di ``qualità della vita'' poco felici.

Tabletop Simulator ha invece un vantaggio ``fenomenale'' che si tende a
sottovalutare: essendo su Steam, \textbf{il modulo si pubblica direttamente sul
Workshop di Steam}; si possono invitare gli amici a iscriversi e giocare con
loro, purché abbiano un account Steam. Rispetto ad altri software da imparare,
costa ``tra i 10 e i 20 euro, ragazzi: si parla proprio di spiccioli''. ``Io con
Tabletop Simulator ci ho playtestato la vita''.

Altre caratteristiche utili:
\begin{itemize}
  \item si possono creare \textbf{moduli non elencati} (\emph{unlisted}),
        perfetti da passare a un publisher prima che il gioco sia pubblicato:
        non sono pubblici, e nessuno può rubare i materiali;
  \item il software ha una grafica datata (``\emph{World Art 1998}'', scherza
        uno dei docenti), ma lo stanno portando su \textbf{Unity} con molti investimenti
        e aggiornamenti; ascolta la comunità e introduce funzioni utili anche per
        chi fa design;
  \item il suo server Discord è molto frequentato, e c'è un canale YouTube
        dedicato solo al design di giochi di carte\footnote{La trascrizione ne
        riporta il nome come ``Distraction Maker''; non è verificabile.} con un
        server Discord dove si fanno playtest: il 99\% è su Tabletop Simulator;
  \item molti editori, quando si propone loro un gioco da tavolo,
        \textbf{chiedono se c'è un modulo per Tabletop Simulator}.
\end{itemize}

% ---------------------------------------------------------------------
\section{Le domande a Daniele: il tema del project work}
\label{a1:sec:domande-daniele}

Daniele ha ``un minuto'' per rispondere alle domande sul corso prima di dover
andare. Uno studente chiede, ``giusto per curiosità'', se il project work avrà
un tema o sarà libero.

\begin{itemize}
  \item C'è un \textbf{mandate}, che i docenti propongono ogni anno e
        aggiornano leggermente; verrà presentato il \textbf{25 settembre}.
        ``Non facciamo spoiler''. Conterrà quello che si chiede di fare e quello
        che non si chiede: dei \textbf{must}, dei \textbf{must not} e dei
        \textbf{nice to have}.
  \item Verrà chiesto di lavorare su un'\textbf{idea originale} che non si è
        portata al corso: originale nel senso di \textbf{creata apposta}, non
        nel senso di ``mai vista prima''.
  \item L'idea va \textbf{partorita insieme}. Altri anni alcuni studenti hanno
        provato a imporre la propria idea al gruppo, ``e l'esperienza è stata
        pessima per tutti''. Quindi: \textbf{non imponete le vostre idee, fate
        brainstorming con gli altri}. Daniele, che fa ``molto spesso
        brainstorming'', darà strumenti per farlo: nessuno arriverà impreparato
        a quella lezione.
\end{itemize}

Ruggero aggiunge che il mandate è \textbf{abbastanza flessibile} perché tutti
possano fare qualcosa che li interessi: ``l'ultima cosa che vogliamo è forzarvi a
lavorare su qualcosa che non vi piace''. Bisogna lavorare su ciò che piace e per
cui si ha uno stimolo personale, anche nella \textbf{tipologia} di gioco: nessuno
sarà obbligato a fare un gioco di ruolo o un board game, e la lezione sui casi
studio mostrerà quante forme diverse di gioco analogico il corso accetta.

L'unica cosa da \textbf{non fare}, ``perché è contraddittorio per questo corso'',
sono i \textbf{giochi grandi}: ``lo chiediamo per voi''. Niente progetti
mastodontici\footnote{La trascrizione dice ``non fate delle cose da gnome'':
probabilmente il nome di un gioco da tavolo molto grande e complesso, non
ricostruibile con certezza.}, niente ``cose alla \emph{Dungeons \& Dragons}''.
Si può fare un gioco di ruolo, ma con un'esperienza ``meno corposa'': ci sono
tanti giochi di ruolo da 15 o 20 pagine, ``bellissimi, fichissimi'', che si
possono leggere online. Il consiglio è farsi una cultura anche su
\textbf{itch.io}, ``uno dei portali più importanti del gioco di ruolo
internazionale'', di cui si parlerà anche nella parte di ricerca e sviluppo.

Poi Daniele saluta (``noi ci vediamo venerdì a lezione'') e lascia la
conduzione a Ruggero.

% ---------------------------------------------------------------------
\section{Le domande finali con Ruggero}
\label{a1:sec:domande}

\subsection{Le regole delle domande}

Ruggero spiega che cosa vuol dire ``stare con lui''. Alla fine di ogni sua
lezione ci saranno \textbf{almeno 20--30 minuti di domande} libere, \emph{free
form}. Vuole che fare domande sia \textbf{normalizzato}: ``io gestisco malissimo
l'\emph{horror vacui} e divento super cringe molto alla svelta se non ci sono
domande in giro''. Sono benvenute le domande sulla lezione appena fatta, ma non
solo: il corso è crossmediale, e anche le domande che vengono da
\textbf{divagazioni}, da parti laterali, ``fatele''. Quello che sembra scontato o
banale a uno studente può non esserlo per il docente, e viceversa; e può essere
utile anche agli altri. ``Se vi sembrano stupide, fatele il doppio''. ``Magari
sembro incazzato, ma vi giuro sono rilassatissimo. Non mordo, magari abbaio un
po'\,''.

\subsection{Come approcciarsi a Steam e Tabletop Simulator}

Rita non ha mai usato né Steam né Tabletop Simulator e chiede come
approcciarsi, e se c'è modo di mettersi d'accordo tra compagni su Discord, magari
con un sondaggio sui giochi, per formare gruppi di gioco. Ruggero conferma: sul
Discord dell'accademia c'è un \textbf{canale apposito per il corso}, ed è degli
studenti. Lo possono usare per organizzarsi a giocare, per chiedere a un
compagno di prestare un gioco, per domande tecniche (come si scarica, come si
installa). ``Sarebbe proprio ideale che vi parlaste e vi organizzaste''.

Anche perché, aggiunge con un ``piccolo spoiler'' (``nessuno vi ha detto
niente''), se fossero i docenti a organizzare i gruppi di gioco, il task lo
starebbero facendo loro. Bisogna \textbf{abituarsi a organizzare le sessioni di
gioco}, perché arriveranno i task sul \textbf{playtest} e su come organizzarlo,
e a quel punto si sarà ``meno disperati''. Sembra scontato, ma saper
\textbf{organizzarsi per giocare in autonomia} ``vi sarà utilissimo come skill in
questo settore, ve lo giuro''.

\subsection{Il bilanciamento}

Angelo confessa ``una paranoia grossissima'': ci saranno lezioni sul
\textbf{bilanciamento}? ``Ma certo''. Sono proprio le lezioni di Marta sul
development: il bilanciamento è di solito un tema che tocca \textbf{i developer
più dei designer}. L'anno precedente c'era una sola lezione con Marta, e i
docenti si sono resi conto che serviva andare ancora più a fondo: per questo
quest'anno le lezioni sono due. ``Non è una paranoia, è una legittima domanda'':
Ruggero insegna da anni, vede altri corsi, parla con altri docenti e studenti, e
``mi incazzo spesso'' perché molti corsi di game design \textbf{non fanno
bilanciamento}. I docenti volevano invece un corso che coprisse tutte le parti
del game design: non solo l'idea del gioco, ma il bilanciamento, la produzione,
i playtest, l'iterazione, la scrittura dei regolamenti.

\subsection{Si può giocare con i docenti?}

Una studentessa chiede se si può giocare con i docenti. Per motivi etici,
durante il corso questo è limitato alle parti di \textbf{valutazione} dei giochi
e di \textbf{tutoraggio}. Dopo il corso, ``nessun problema, almeno con me''.

\subsection{Dove si trovano i giochi su Tabletop Simulator}

Ruggero mostra dove si trovano i giochi: nella propria libreria di Steam, nella
pagina di Tabletop Simulator, alla voce \textbf{Workshop}, dove ci sono ``più o
meno 36.000'' giochi. Una studentessa scopre intanto che l'app di Steam per
telefono è solo una \emph{companion app}: per registrarsi e giocare serve il
PC.

Per provare, chiede agli studenti un titolo:
\begin{itemize}
  \item \emph{Keep the Heroes Out}: non compare, ``non tutti i giochi ci sono,
        perché molti publisher non li mettono''; forse, essendoci già una
        versione del gioco su Steam, il modulo è stato tolto;
  \item \emph{Magic: The Gathering}: ``almeno 50 moduli'', anzi ``123 pagine'';
  \item \emph{Bang!}: tre moduli.
\end{itemize}
Bisogna fare attenzione: alcuni moduli sono in lingue specifiche (``magari il
primo è in russo''), alcuni hanno già l'espansione, altri sono solo il gioco
base. Conviene sfogliare e non fermarsi al primo risultato. E nel Workshop si
trovano anche cose molto particolari: giochi fuori produzione, giochi vecchi,
giochi ancora in prototipo.

\subsection{I ruoli nei team e le dinamiche tra le persone}

Uno studente torna sulla slide delle figure: nelle realtà piccole molti ruoli
hanno ``il doppio o triplo cappello''. Quali ruoli restano distinti, e come si
creano le sinergie tra persone con obiettivi diversi?

Ruggero premette che l'ideale ``platonico'' è un ambiente di lavoro costruttivo,
basato sulla fiducia e ``il più \emph{safe} possibile''; ma nell'industria non
sempre è così, e le strutture variano moltissimo: alcuni publisher sono molto
verticali, altri hanno ruoli già solidificati (se un publisher ha già i suoi
developer e non sta crescendo, entrarci come developer è praticamente
impossibile). ``Se prendiamo due aziende in Italia, vediamo due strutture
interne completamente diverse''. Non esiste una scorciatoia per farsi assumere
da un editore.

Poi prende ``l'esempio più comune di tutti'', la media dei giochi da tavolo
pubblicati (figura~\ref{a1:fig:relazioni}):
\begin{itemize}
  \item il \textbf{game designer} è una figura \textbf{esterna}, con un
        \textbf{contratto a royalty} con il publisher;
  \item la sua figura di riferimento è il \textbf{publisher} stesso, e la
        conversazione designer--publisher è ``molto calda'';
  \item se il publisher ha dei \textbf{developer}, il developer entra in questa
        conversazione: designer e developer collaborano molto, e internamente
        developer e publisher collaborano molto. Spesso, nei publisher piccoli,
        è il publisher a fare le veci del developer;
  \item \textbf{graphic design} e \textbf{illustrazione} tendono a parlare molto
        con il publisher e \textbf{molto poco con il designer}. A meno che non si
        tratti di uno studio indie molto piccolo, è difficile che il designer
        sappia anche solo chi è l'illustratore prima che il gioco esca.
\end{itemize}

\begin{figure}[htbp]
\centering
\begin{tikzpicture}[
    font=\small,
    ruolo/.style={draw, thick, rounded corners=3pt, fill=gray!12,
                  minimum width=2.8cm, minimum height=0.9cm, align=center},
    esterno/.style={draw, thick, dashed, rounded corners=3pt,
                    minimum width=2.8cm, minimum height=0.9cm, align=center},
    calda/.style={{Stealth[length=2.5mm]}-{Stealth[length=2.5mm]}, very thick},
    tiepida/.style={{Stealth[length=2mm]}-{Stealth[length=2mm]}, thin, dashed}
  ]
  \node[esterno] (gd) at (0,0) {Game designer\\ \footnotesize esterno, a royalty};
  \node[ruolo] (dev) at (5,1.4) {Developer\\ \footnotesize se c'è};
  \node[ruolo] (pub) at (5,-1.4) {Publisher};
  \node[ruolo] (art) at (10,-1.4) {Graphic designer\\ e illustratore};
  \draw[calda] (gd) -- (dev);
  \draw[calda] (gd) -- (pub);
  \draw[calda] (dev) -- (pub);
  \draw[calda] (pub) -- (art);
  \draw[tiepida] (gd.south) to[bend right=25]
     node[below, font=\footnotesize\itshape] {quasi nessun contatto} (art.south);
  \node[draw, dotted, rounded corners=4pt, fit=(dev)(pub)(art), inner sep=8pt,
        label={[font=\footnotesize\itshape]above right:interno al publisher (di solito)}] {};
\end{tikzpicture}
\caption{Le relazioni tipiche intorno a un gioco da tavolo pubblicato, secondo la
risposta di Ruggero. Le frecce spesse indicano una collaborazione stretta e
frequente; quella tratteggiata, un contatto quasi assente. Nei publisher piccoli
il publisher fa spesso anche da developer.}
\label{a1:fig:relazioni}
\end{figure}

\subsection{Fare da developer di sé stessi?}

Rita chiede in chat se un designer può fare anche da developer del proprio
gioco, o se Ruggero lo sconsiglia. In generale \textbf{lo sconsiglia}, per lo
stesso motivo per cui si sconsiglia in altri media: i \textbf{darling}. Il
designer ha delle meccaniche a cui è affezionato, e il developer a volte dovrà
\textbf{tagliarle}; è molto difficile, se si è ``gestori di sé stessi'', tagliare
cose che si vogliono nel gioco. ``È un po' come lasciare fare il lavoro di editing
allo scrittore che ha già scritto il libro'': può rivedersi, certo, ma difficilmente
cambierà cose che farebbe bene cambiare. Nelle lezioni con Marta si vedrà perché,
e alla fine del corso ``non avrete dubbi nemmeno voi'' nel separare designer e
developer.

Ma tutto dipende dallo \textbf{scopo del progetto}: se si fa un gioco in tre,
qualcuno quel lavoro deve farlo per forza, e magari il designer fa anche il
developer, compensando con \textbf{più playtest}. ``Il mondo attuale e il mondo
ideale sono due cose molto diverse, e lo vediamo non solo dai giochi ma dalla
politica'', ma ``non entriamo in questo vaso di Pandora''.

\subsection{Il ruolo indispensabile e la pubblicazione}

Lo studente di prima ne trae una conclusione: senza game designer il gioco non
esiste. ``Quello sicuramente''. Ma Ruggero coglie l'occasione per approfondire
un punto che forse non era passato: \textbf{fare design di un gioco per
pubblicarlo non è l'unica opzione possibile}. Un gioco non è ``meno gioco''
se non è pubblicato da un editore. È legittimo mettere un gioco su itch.io in
\textbf{print and play}, fare un gioco a evento senza venderlo, fare un
\emph{murder party}, una escape room, un board game per una persona amica:
``quell'esperienza lì non è meno gioco''. Dipende dall'obiettivo del progetto.

Ruggero ha lavorato molto anche con le \textbf{scuole}, dove il publisher non
esiste: è un gioco per le classi, e la figura più vicina è l'editoria scolastica,
``che però non fa il publisher''. Ci sono compromessi diversi sulla materialità
e su come il gioco si presenta, che si vedranno nella prototipazione e nei casi
studio. Quindi l'unica figura che deve esserci sempre è \textbf{chi fa game
design}, anche se magari non sa di chiamarsi game designer e non ha mai
sentito parlare di game design: ``questo non la rende meno game designer''.

\paragraph{Un libro consigliato.} Ruggero fa anche mentoring e gestisce gruppi
di settore, e spesso gli chiedono che cosa leggere per imparare il game design.
Non è il tipo che obbliga a leggere libri (``non vi obbligherò a leggere
niente''), ma ce n'è uno che consiglia sempre, ed è il primo che ha scelto di
leggere: \textbf{\emph{Rise of the Videogame Zinesters}} di \textbf{Anna
Anthropy}\footnote{La trascrizione dice ``Anna Antrofi'' e ``Rise of the
Videogames in Esters''.}. È un libro un po' datato, ma ancora reperibile; parla
soprattutto di videogiochi, ma ``per l'80\%'' il contenuto si applica al gioco e
al game design in generale. È un libro sulla \textbf{democratizzazione del gioco}:
il gioco come strumento per parlare di sé e come \textbf{medium legittimo}. Come
esistono film drammatici, tristi, intensi, possono esistere giochi su tematiche
delicate, e se ne vedranno nei casi studio. Molti giochi fatti con questa
direttiva non hanno un publisher, che tende a pubblicare ciò che dà un ritorno
commerciale.

\paragraph{Commerciale e personale.} ``Io spero che voi non farete giochi solo
per scopo commerciale''. Certo, servono un tetto sulla testa, il cibo in
tavola, un posto in cui dormire, e così si finisce per fare giochi che non sono
il proprio gioco dei sogni. Spesso è un mix: Ruggero ha fatto giochi commerciali
``per pagarmi l'affitto'' e giochi che voleva fare e su cui ha perso soldi,
``sapendolo proprio che sarebbero stati un fallimento''. I giochi, come i libri e
le canzoni, si fanno per molti motivi, e quelli commerciali di successo non sono
più importanti degli altri. Tra i giochi della sua slide c'erano il gioco da
tavolo ufficiale di \emph{Cyberpunk}\footnote{La trascrizione dice ``Cyberpunk
del 1977'': quasi certamente un gioco legato a \emph{Cyberpunk 2077}.} e
\emph{Marvel United}, su IP grosse e famose, ma anche giochi fatti di ``una
candela e una busta su un tavolo''. E se gli si chiede quale preferisce in
termini di contenuto, non sono quelli sui supereroi.

\subsection{Il gioco di ruolo in Italia}

Nicola, che è interessato soprattutto al gioco di ruolo e conosce già persone
del settore, ha ``una leggera paura'': che chi studia game design finisca per
fare il \textbf{freelance} a vita, ``perché non ci sarà mai un'occupazione per
bene'', o il \textbf{copywriter}. Conosce un ragazzo che ha lavorato in quel
mondo e che ora fa altro.

Ruggero risponde che la maggioranza dei giochi nel suo portfolio sono giochi di
ruolo, quindi il tema lo tocca da vicino, e che in Italia quello che dice Nicola
``è semplicemente vero''. Molti publisher piccoli hanno una persona di
riferimento che copre più ruoli: impagina, fa i playtest, fa un po' di
development. È ``un settore economicamente completamente disastrato'', per
motivi culturali che non apre: oggi in Italia è praticamente impossibile
\textbf{vivere di giochi di ruolo}. Chi ci lavora lo fa nel tempo libero, al di
fuori del lavoro principale, oppure coprendo molte altre cose: molti publisher
italiani di giochi di ruolo pubblicano anche titoli esteri e localizzano, non
fanno solo cose originali. ``È un grosso calderone''. Poi si ferma: ``questo è
un \emph{rant}'', e ``i \emph{rant} tornano sempre'', ma non vuole che gli
studenti escano dalla prima lezione ``con la disperazione nel cuore''. Ci sarà
spazio per parlarne quando avranno gli strumenti per capirlo, ``per non farvi
abbattere''.

\subsection{Che cosa bisogna sapere già}

Uno studente chiede se ci sono nozioni che dovrebbero già avere e che possono
recuperare: per esempio, sente parlare di giochi ``German'' o ``American'' e non
sa che cosa siano. Ruggero spiega l'organizzazione del corso:
\begin{itemize}
  \item le \textbf{prime due lezioni} sono di fondamenta. Questa era ``dall'alto'':
        come è strutturato il corso, le tempistiche;
  \item la \textbf{prossima lezione} è di tre ore di soli \textbf{casi studio},
        ``a manetta, a raffica'', divisi per categorie e tipologie, con
        ragionamenti su categorie e generi di gioco: lì si vedranno German,
        American ``e anche altre più utili'';
  \item nel resto del corso le terminologie specifiche sono inserite nelle
        lezioni dedicate: il \textbf{paragame} nelle lezioni sul paragame, il
        \textbf{modello matematico} nelle lezioni di Marta sul development.
\end{itemize}
Per ora agli studenti si chiede di sapere quello che si è visto oggi:
meccaniche, dinamiche ed estetiche; le figure professionali, designer e
developer; e soprattutto \textbf{game e play}.

\subsection{Game e play: i quattro quadranti}

``Specialmente \emph{game} e \emph{play}'', insiste Ruggero: è più difficile di
quello che sembra, perché l'italiano non abitua a separare le due cose, ed è
possibile averle usate per anni come sinonimi. Nella sua esperienza con chi si
avvicina al design è ``uno dei punti più critici''. Spesso i commenti che si
fanno sui giochi giocati sono commenti sul \emph{play} applicati al
\emph{game}: ``il gioco non mi è piaciuto perché, giocando con Tizio, è successo
questo''; ``il gioco mi sembra sbilanciato perché Tizio al tavolo ha fatto
questo''.

Separare \emph{game} e \emph{play} serve a capire quali giochi sono progettati
bene e quali male, e soprattutto a riconoscere i \textbf{quattro possibili
incroci} (figura~\ref{a1:fig:quadranti}): un gioco progettato bene che piace
giocare; uno progettato male che non piace; uno progettato male che piace; uno
progettato bene che non piace. ``Non è importantissimo in generale nella vita,
ci sono cose più importanti'', ma per chi fa game design lo è: ``a volte vi
chiederanno di lavorare su un gioco che non vi piace giocare, e quel gioco deve
essere fatto bene lo stesso''.

\begin{figure}[htbp]
\centering
\begin{tikzpicture}[
    font=\small,
    cella/.style={draw, thick, minimum width=4.4cm, minimum height=1.9cm,
                  text width=4.0cm, align=center}
  ]
  \node[cella, fill=gray!20] (a) at (0,0)     {\textbf{progettato bene}\\ e \textbf{mi piace} giocarlo};
  \node[cella]               (b) at (4.4,0)   {\textbf{progettato bene}\\ ma \textbf{non mi piace} giocarlo};
  \node[cella]               (c) at (0,-1.9)  {\textbf{progettato male}\\ ma \textbf{mi piace} giocarlo};
  \node[cella, fill=gray!8]  (d) at (4.4,-1.9){\textbf{progettato male}\\ e \textbf{non mi piace} giocarlo};
  \node[font=\footnotesize\itshape, above=0.15cm of a] {mi piace (\emph{play})};
  \node[font=\footnotesize\itshape, above=0.15cm of b] {non mi piace (\emph{play})};
  \node[font=\footnotesize\itshape, rotate=90, left=0.15cm of a, anchor=south]
       {bene (\emph{game})};
  \node[font=\footnotesize\itshape, rotate=90, left=0.15cm of c, anchor=south]
       {male (\emph{game})};
\end{tikzpicture}
\caption{I quattro incroci tra la qualità del \emph{game} (come è progettato) e
l'esperienza di \emph{play} (se piace giocarlo), descritti da Ruggero. Per un
designer i quadranti ``scomodi'' sono i due fuori diagonale: un gioco può
piacere pur essendo progettato male, e va saputo progettare bene anche un gioco
che non piace.}
\label{a1:fig:quadranti}
\end{figure}

\paragraph{La domanda di Tancredi.} Tancredi chiede conferma: il \emph{play}
non è il ``giocare senza regole'', ma può essere anche il giocare a un gioco? Sì:
il \emph{play} è l'\textbf{atto}, e può essere fatto con un \emph{game} o senza.
E un \emph{game} può esistere senza che nessuno lo giochi: gli scacchi non sono
meno scacchi perché in quel momento nessuno ci gioca. E se, giocando, una
persona applica le regole in modo diverso? Allora \textbf{non si sta giocando
quel gioco}.

\subsection{House rule: il caso di Monopoly}

Qui Ruggero fa un po' di \emph{foreshadowing} su una lezione di Daniele: il
discorso degli \textbf{hack}, delle \textbf{house rule}, delle \textbf{mod},
cioè di come modificare un gioco esistente. Il punto è che \textbf{un gioco
modificato da una house rule è un altro gioco}: cambiando le regole cambiano
necessariamente anche le \textbf{dinamiche} che ne derivano, e quindi
l'esperienza. Si immagini \emph{Monopoly} in cui ogni turno si possa, volendo,
prendere gratis dei soldi dalla banca: tutto il resto è identico, eppure il gioco
cambia.

E questo è talmente interiorizzato che ci sono regole che si applicano ``senza
neanche rendercene conto''. Ruggero chiede quanto duri in media una partita a
\emph{Monopoly}. Gianmarco: ``quattro ore''; altri: ``un pomeriggio''; un altro:
``finché uno dei giocatori non aggiunge la regola: ora vince chi\ldots''. In
realtà, \textbf{seguendo le regole scritte}, una partita non dura mai più di
\textbf{un'ora e venti, un'ora e mezza}. La domanda che fa sempre, e la risposta
è sempre la stessa: perché si pensa che \emph{Monopoly} duri tantissimo? Qual è
la regola che nessuno rispetta?

Uno studente indovina: \textbf{le aste}. Secondo il regolamento, se chi finisce
su un terreno non lo compra, il terreno \textbf{va all'asta} tra i giocatori, e
chiunque può aggiudicarselo. Molti invece giocano come se finire su un terreno
desse solo la \emph{facoltà} di comprarlo o di metterlo all'asta\footnote{Più
precisamente, nel regolamento ufficiale, se il giocatore che finisce su una
proprietà libera decide di non comprarla, la banca la mette all'asta e tutti i
giocatori possono fare offerte, compreso chi l'ha rifiutata. L'asta quindi non è
una scelta del giocatore, ma la conseguenza obbligata del suo rifiuto.}. Con la
regola applicata, di aste se ne fanno moltissime: all'ottavo turno ci sono già
otto terreni assegnati, al ventesimo venti, la gente va in bancarotta
velocissimamente e la partita dura un'ora e mezza. Anche le partite competitive
di \emph{Monopoly}, nota Ruggero, durano un'ora e venti, un'ora e mezza.

Uno studente obietta che giocava con le aste e durava comunque quattro ore.
``Avete dei nemici con l'analisi paralisi'', ride Ruggero; e cita un'altra house
rule tipica dell'infanzia: genitori che danno soldi per non far andare in
bancarotta. Ma \textbf{la bancarotta è parte del gioco}: se la si toglie,
\emph{Monopoly} non regge. Sul tema che \emph{Monopoly} vuole trattare, e sul
suo antenato (\emph{The Landlord's Game}), ``non apriamo adesso un altro vaso di
Pandora''.

\paragraph{Approfondimento: Monopoly, la durata e le house rule nei libri.}
L'esempio di Ruggero si ritrova quasi identico in \emph{The Art of Game Design}.
Parlando del bilanciamento tra partite brevi e lunghe (\emph{Short vs. Long}),
Schell scrive che \emph{Monopoly}, giocato secondo le regole ufficiali,
``finisce spesso in circa novanta minuti'', ma molti giocatori trovano quelle
regole troppo dure e le modificano, distribuendo premi in denaro e allentando i
vincoli sull'acquisto delle proprietà: con l'effetto collaterale di far durare la
partita molto di più, ``tipicamente tre ore, o anche di più''. Schell classifica
le \textbf{house rule} come regole create dai giocatori in risposta a un
difetto percepito dopo qualche partita. Anche \emph{Challenges for Game
Designers}, nel capitolo sull'aggiungere e togliere meccaniche, parte dallo
stesso gioco: il \emph{Parcheggio gratuito}, che per le regole originali non fa
succedere nulla, mentre quasi tutti giocano con una house rule che dà un bonus a
chi ci finisce sopra. Le due fonti confermano la lezione: una piccola regola
cambiata cambia le dinamiche, e le dinamiche cambiano l'esperienza.

\subsection{Il fine del game è il play?}

Un altro studente commenta: va bene la distinzione, ma il fine del \emph{game} è
il \emph{play}; quello che si vede è sempre l'effetto. Ruggero è d'accordo: di
solito, se si fa un gioco, lo si fa per farlo giocare, come si scrive un libro
per farlo leggere (anche se esistono, ed ``è già stato esplorato'', libri scritti
per non essere letti e giochi per zero giocatori). Ma la distinzione serve a una
cosa precisa: \textbf{separare, nel playtest, i feedback sul gioco dai feedback
sul gruppo di gioco} e sulle sue dinamiche. Se una persona esce frustratissima da
un playtest, non è necessariamente colpa del gioco; se una persona dice ``questo
è il gioco migliore del mondo'', potrebbe non essere merito del gioco ma del
gruppo. Capire a che cosa si riferisce un feedback serve a capire \textbf{dove
iterare} il prototipo.

Lo studente osserva che l'area è sfumata: alcuni giochi si basano proprio
sull'avere le persone giuste al tavolo. Un gioco di \textbf{negoziazione} con
giocatori abituati ai giochi ``multi-solitario'' diventa ingiocabile; e nei giochi
a \textbf{economia chiusa} (cita \emph{Container}\footnote{La trascrizione cita
anche un secondo titolo non ricostruibile (``Tokyo su Vigimarket'').}) un gruppo
economicamente conservativo, che spende poco e non fa girare l'economia, fa
durare la partita il doppio. Ruggero rilancia: e se il gioco dura il doppio,
questa dinamica cambia la \textbf{percezione di peso} del gioco, le energie spese,
quanto si è stanchi a fine partita. ``Tutto è collegato. \textbf{I giochi sono
ecosistemi di meccaniche}, e cambiandone una, le altre cambiano di
conseguenza''.

\paragraph{La parola vietata, alla prova.} Lo studente aggiunge che ci sono giochi
che ha trovato ``molto sbagliati, ma che mi divertono''. ``Attenzione: non si può
dire \emph{divertimento}!'', lo ferma Ruggero; nelle presentazioni aveva lasciato
correre, ``ma lì non vi avevo ancora letto la regola''. Lo studente ci prova:
``giochi che poi voglio \textbf{rigiocare}'', anche se ``sono evidentemente
sbagliati, rotti''. Ruggero apprezza, ma non vuole anticipare cose che meritano il
loro tempo: se ne parlerà nelle lezioni sui framework e soprattutto in quella di
\textbf{emotion design}, ``forse la mia lezione preferita del corso''. (Lo
studente ammette che ``forse alla fine si gioca per quello, per l'emozione''.)

\subsection{I giochi ``comunicati male''}

L'ultima domanda corposa riprende un'affermazione di inizio lezione: ci sono
giochi belli ma comunicati male. Qualche esempio? ``Mi basta girarmi per
vederne 48'', risponde Ruggero indicando gli scaffali alle sue spalle. Più che
titoli, elenca le \textbf{ragioni}:
\begin{itemize}
  \item \textbf{l'Italia localizza lentamente}. Basta guardare i premi del gioco
        dell'anno in Italia: spesso vanno a giochi usciti due o tre anni prima in
        Germania, Francia, Stati Uniti o Regno Unito, perché localizzare richiede i
        suoi tempi. Il mercato italiano è quindi molto dipendente dall'estero,
        ``molto impaurito'', e spesso aspetta che un gioco vada bene all'estero
        prima di portarlo in Italia; e alcuni giochi ``perdono il treno'';
  \item \textbf{il publisher originale chiude}: come si pubblica un gioco se non
        c'è più l'azienda con cui fare un contratto? A Ruggero è successo più
        volte;
  \item \textbf{il dipartimento commerciale italiano non sa comunicare il gioco}:
        si fa male un'edizione, la stampa vende pochissimo e non viene più
        ristampata;
  \item \textbf{la scala del mercato}: in Italia un gioco da tavolo di peso medio
        (un \emph{eurogame}, un \emph{German}) che vende 1.500 copie ``è tanto'';
        all'estero se ne vendono magari 10.000. Con 10.000 copie in inglese ci
        sono i soldi per una campagna di comunicazione: un trailer, i
        recensori, i tutorial online. In Italia tutto questo è ``ridotto all'osso'',
        e la cosa si riflette su stampa e divulgazione: anche i canali per
        informarsi su un gioco prima di comprarlo sono pochi (se ne parlerà nella
        lezione su ricerca e sviluppo). Un gioco spesso non arriva al mercato non
        perché sia brutto, ma perché l'editore non lo sa comunicare, ne fa poche
        copie e va fuori produzione;
  \item \textbf{il tempismo}: party game progettati nel 2019 e usciti nel 2020, in
        pieno Covid, quando non si poteva giocare al tavolo, sono falliti tutti, e
        gli editori ci hanno perso. ``Non potevano immaginarselo'';
  \item \textbf{i temi caldi}: dopo l'invasione russa dell'Ucraina e la guerra in
        Palestina, alcuni editori non pubblicano o non localizzano giochi
        strategici e di guerra. Ruggero ha lavorato a \textbf{\emph{Status}}, un
        gioco sulla geopolitica che tocca, tra gli altri temi, l'invasione russa
        dell'Ucraina e il genocidio in Palestina, e alcuni editori che di
        solito localizzano i loro giochi (per esempio il partner di riferimento in
        Russia) non lo hanno localizzato per motivi politici.
\end{itemize}
Lo studente voleva capire che cosa voglia dire comunicare male: non metterci
abbastanza energia, o comunicare una cosa per un'altra, ``come il trailer
sbagliato, che poi vai al cinema e ti trovi un altro film''. ``È un mix'', e la
domanda ``era complessa e meritava più spazio''.

\subsection{I giochi di carte collezionabili e i saluti}

L'ultima domanda (``posso rompere le balle un'ultima volta?'') riguarda i
\textbf{giochi di carte collezionabili}: funzionano con le stesse regole dei
giochi da tavolo? Se ne parlerà alla prossima lezione: è proprio il motivo per
cui la lezione sui casi studio dura tre ore, perché ogni genere di gioco ha le
sue dinamiche e le sue regole, e anche il modo di pubblicare cambia molto tra un
TCG, un gioco di ruolo da tavolo, un LARP e un gioco pervasivo.

Ruggero saluta: le email dei tre docenti sono all'inizio delle slide, la
registrazione arriverà, ``ci si vede questo venerdì''.

% ---------------------------------------------------------------------
\section*{Punti chiave del capitolo}
\addcontentsline{toc}{section}{Punti chiave del capitolo}

\begin{itemize}
  \item Il corso di \textbf{Analog Game Design} (seconda edizione) dura circa tre
        mesi, fino a dicembre, ed è tenuto da tre docenti principali:
        \textbf{Daniele} (narrative design, gioco di ruolo, fiction),
        \textbf{Ruggero} (game design, ricerca e sviluppo) e \textbf{Beatrice}
        (giochi pervasivi), con tre ospiti su RNG, development e game production.
  \item Lo scopo è formare \textbf{autori analogici} consapevoli: strumenti da
        usare anche fuori dal corso, non il ``gioco dei sogni''. Il corso è
        \textbf{intersezionale} (nessuna lezione dedicata a un solo tipo di
        gioco), \textbf{cross-generazionale}, \textbf{pratico e adattabile},
        basato sulle \textbf{domande} (nessuna domanda è stupida), sul
        \textbf{feedback} invece che sul voto e sul \emph{fail forward}.
  \item Il \textbf{project work} si fa in tre gruppi da tre, formati dai docenti;
        pitch il 12 ottobre, sviluppo nelle due pause, tutoraggi il sabato,
        consegna il 4 dicembre, presentazione l'11 e giornata finale con gli
        editori il 14. Si valutano documentazione, presentazione,
        autovalutazione, crescita e atteggiamento.
  \item \textbf{Niente intelligenza artificiale} nel corso: i publisher
        rifiutano pitch e prototipi con immagini generate, i regolamenti scritti
        dall'IA descrivono la stessa meccanica in modi diversi, e soprattutto chi
        la usa perde l'occasione di imparare. Un prototipo può essere fatto a
        matita.
  \item Le \textbf{figure dell'industria}: game designer, game developer,
        editor, publisher, sales e marketing, graphic designer, illustratore,
        sensitivity reader (che non è un censore) e producer. Nelle realtà
        piccole una persona copre più ruoli.
  \item Nell'analogico il \textbf{game designer} progetta l'esperienza creando
        le meccaniche; il \textbf{game developer} fa da \textbf{mediatore} con il
        publisher (bilanciamento, modelli matematici, adattamenti). Pretendere
        che il proprio gioco non venga modificato è un \emph{red flag}.
  \item Il \textbf{marketing} non decide che cosa pubblicare ma come; è
        un'azione, ed esiste anche se lo si ignora.
  \item In \textbf{Italia} il gioco è ancora visto come infantile, il mercato
        punta sulla localizzazione e investe poco su giochi e designer interni;
        all'estero (soprattutto dalla Germania in su) ci sono tradizione adulta,
        premi di prestigio e developer interni.
  \item \textbf{Game} è l'artefatto (le regole, il sistema, l'oggetto);
        \textbf{play} è l'atto di giocare. Si fanno giochi, non giocattoli.
        Separarli serve a distinguere i feedback sul gioco da quelli sul gruppo
        e a riconoscere i quattro incroci tra ``progettato bene/male'' e ``mi
        piace/non mi piace''.
  \item \textbf{MDA}: meccaniche (regole scritte), dinamiche (regole in azione),
        estetiche (esperienza, non grafica); nel \textbf{DPE}: design, play,
        experience.
  \item \textbf{Parole vietate} in ogni materiale scritto: \emph{fun},
        \emph{immersion}, \emph{entertainment} (divertimento, immersione,
        intrattenimento). Bisogna trovare parole più precise.
  \item Distinguere \textbf{chi crea, chi gioca e chi facilita}; nessun gioco
        nasce perfetto: \textbf{playtest e iterazione} sono inevitabili, e un
        \textbf{dizionario comune} aiuta a non perdersi.
  \item \textbf{Tabletop Simulator} è il software di settore per giocare (non
        creare videogiochi) i giochi da tavolo in digitale. Primo task: giocarci
        tre o quattro volte entro il 10 ottobre.
  \item Una \textbf{house rule} crea un altro gioco: \emph{Monopoly} con le aste
        applicate come da regolamento dura circa un'ora e mezza.
\end{itemize}

%
%  (Con hyperref conviene aggiungere anche
%  \renewcommand{\theHchapter}{agd.\arabic{chapter}} nella parte, per
%  evitare ancore duplicate con i capitoli 1, 2, ... di Game Design.)
%  Tutte le etichette di questo capitolo hanno il prefisso "a1:"
%  (a1:cap:..., a1:sec:..., a1:fig:..., a1:tab:...). Il capitolo è
%  autonomo: non contiene \ref verso altri capitoli.
%  Le figure sono ricostruzioni schematiche di quanto i docenti spiegano
%  a voce e con le slide, non copie delle slide.
% =====================================================================

\chapter[Introduzione all'Analog Game Design]{Introduzione all'Analog Game
Design: il corso, le figure dell'industria e un primo linguaggio comune}
\label{a1:cap:intro}

Questa è la \textbf{prima lezione} del corso di \textbf{Analog Game Design},
cioè di progettazione di giochi \emph{analogici}: giochi da tavolo, giochi di
carte, giochi di ruolo da tavolo, giochi di ruolo dal vivo (LARP), giochi
\emph{pervasivi} ed esperienze a evento. È un corso \textbf{breve}, di circa tre
mesi (``da qui a dicembre''), separato dal corso di Game Design, e siamo alla sua
\textbf{seconda edizione}. Come avvertono subito i docenti, è ``una giornata di
introduzione'': si spiega come funziona il corso, chi sono i docenti, con quali
metodi e con quale calendario si lavorerà, e ci si presenta a vicenda. ``C'è poco
gioco analogico'', si scusano. Ma nella seconda parte la lezione comincia
comunque a costruire un \textbf{linguaggio comune}: le figure professionali
dell'industria del gioco analogico, la differenza tra \emph{game design} e
\emph{game development}, la situazione del mercato italiano rispetto a quello
estero, e un primo glossario di design (\emph{game} e \emph{play}, il framework
MDA, le parole ``vietate'', il playtest e l'iterazione). Si chiude con lo
strumento di settore per i prototipi digitali dei giochi da tavolo,
\textbf{Tabletop Simulator}, e con una lunga sessione di domande.

La lezione si svolge \textbf{online}, su Zoom, ed è condotta a più voci dai tre
docenti principali del corso: \textbf{Daniele}, narrative designer, che guida la
parte organizzativa; \textbf{Ruggero}, game designer, che guida la parte sui
fondamenti e le domande finali; \textbf{Beatrice}, progettista di giochi
pervasivi, che interviene soprattutto nella prima metà. Gli studenti iscritti
sono \textbf{nove}; uno è assente e un altro arriva a lezione iniziata. Ci sono
una pausa di una decina di minuti (si riprende alle 18:15) e, dopo le 19:15,
un'ultima mezz'ora di domande con il solo Ruggero.

Tre avvertenze di lettura. La prima: la trascrizione è automatica e storpia
molti nomi propri, compresi quelli dei docenti e dell'accademia (``Dive'',
``IVE'', ``Accademia Live'', ``Ai''); dove la ricostruzione è sicura, o molto
probabile, il capitolo usa la forma corretta e lo segnala in nota; dove non lo
è, riporta la forma della trascrizione e lo dice. La seconda: durante la pausa la
trascrizione registra, ripetuta in loop per decine di righe, una frase senza
senso che non ha nulla a che fare con la lezione, e che il capitolo ignora. La
terza: in molti passaggi più docenti parlano uno dopo l'altro senza che la
trascrizione segni il cambio di voce; il capitolo attribuisce le parole a chi le
pronuncia quando il contesto lo rende chiaro (chi parla di sé, chi viene
chiamato per nome, chi ``passa la parola''), e resta generico altrimenti.

Quanto alla data, la lezione si colloca con ogni probabilità \textbf{nella prima
metà di settembre del 2026}\footnote{La data non è mai detta. La si deduce da
alcuni indizi: il form iniziale va compilato ``entro e non oltre il 15
settembre''; i gruppi verranno annunciati il 25 settembre; le date del
calendario (sabato 3 ottobre, lunedì 12 ottobre, sabato 7 e 14 novembre,
venerdì 4 dicembre) cadono nei giorni della settimana indicati soltanto nel
2026; un docente parla di un argomento ``molto complesso nel 2026''; e fa
ancora molto caldo (``le temperature sono ancora indecenti''). È una
deduzione, non un'informazione esplicita.}: i riferimenti a ``oggi'',
``quest'estate'' e ``l'anno scorso'' vanno letti rispetto a quel momento.

Il percorso del capitolo segue quello della lezione:
\begin{enumerate}
  \item l'avvio, i \textbf{docenti} e le \textbf{presentazioni} dei partecipanti
        (sezioni da~\ref{a1:sec:avvio} a~\ref{a1:sec:partecipanti});
  \item i \textbf{materiali} e gli strumenti del corso: il form, la cartella
        condivisa, Discord (sezione~\ref{a1:sec:materiali});
  \item lo \textbf{scopo} e il \textbf{metodo} del corso, le regole di
        comportamento e il carico di lavoro (sezioni~\ref{a1:sec:scopo}
        e~\ref{a1:sec:regole});
  \item il \textbf{project work} e il \textbf{calendario} (sezione~\ref{a1:sec:project});
  \item la regola sull'\textbf{intelligenza artificiale} e la lunga discussione
        che ne segue (sezione~\ref{a1:sec:ia});
  \item le \textbf{figure professionali} dell'industria del gioco analogico,
        con un approfondimento su marketing e sales
        (sezioni~\ref{a1:sec:figure} e~\ref{a1:sec:marketing});
  \item la differenza tra \textbf{game design} e \textbf{game development}
        (sezione~\ref{a1:sec:design-development});
  \item il \textbf{mercato}: Italia ed estero (sezione~\ref{a1:sec:mercato});
  \item i \textbf{fondamenti di design}: perché serve una terminologia comune,
        \emph{game} e \emph{play}, il framework \textbf{MDA} (e il DPE), le
        \textbf{parole vietate}, chi crea, chi gioca e chi facilita, il
        \textbf{playtest} e l'\textbf{iterazione} (sezioni
        da~\ref{a1:sec:terminologia} a~\ref{a1:sec:playtest});
  \item \textbf{Tabletop Simulator} e il primo task (sezione~\ref{a1:sec:tts});
  \item le \textbf{domande finali}: il tema del project work, il bilanciamento,
        i ruoli nei team, la pubblicazione, il gioco di ruolo in Italia, le
        \emph{house rule} di Monopoly, i giochi ``comunicati male''
        (sezioni~\ref{a1:sec:domande-daniele} e~\ref{a1:sec:domande}).
\end{enumerate}

% ---------------------------------------------------------------------
\section{L'avvio della lezione}
\label{a1:sec:avvio}

La lezione comincia con qualche minuto di attesa. ``Siamo in nove in classe'', e
si aspetta che arrivino tutti; uno dei docenti scherza sul fatto che l'unica
persona che arriverà ``sempre in ritardo di venerdì'' è Gianmarco. Mancano
Giordano e Tiziana; Tiziana arriva, mentre la segreteria fa sapere poco dopo che
Giordano non potrà esserci e vedrà la registrazione. Si aspetta anche Beatrice,
la terza docente, che arriva a sua volta dopo pochi minuti.

Nell'attesa, Daniele anticipa il programma della giornata e qualche
informazione pratica:
\begin{itemize}
  \item oggi è una giornata di \textbf{introduzione}: come funziona il corso, chi
        sono i docenti, le presentazioni degli studenti, le \textbf{metodologie}
        e il \textbf{calendario}. Gli studenti avranno ``un fracco di domande'', e
        alla fine della lezione c'è ``proprio uno slot apposta'' per le domande;
  \item le lezioni sono \textbf{registrate}: di solito la registrazione viene
        consegnata dopo circa \textbf{12--24 ore} sul profilo personale della
        piattaforma dell'accademia, che va attivato con la segreteria (Tiziana,
        già studentessa, probabilmente ce l'ha già);
  \item i docenti ``cuore'' del corso sono tre (Daniele, Ruggero e Beatrice),
        ``soprattutto Ruggero''; a loro si aggiungeranno tre docenti ospiti, per
        meno ore, probabilmente una lezione a testa, ``tranne Marta, che avrà due
        lezioni quest'anno''.
\end{itemize}
E una battuta sul caldo: ``spero che voi siate carichi, anche se le temperature
sono ancora indecenti''. ``Carichi moderatamente, per quello che potete''.

Poi Daniele condivide lo schermo (``non solo le slide, ma anche un po' di altre
cose'') e apre ufficialmente il corso: ``benvenuti e benvenute nella seconda
edizione del corso di Analog Game Design. Questa è la prima lezione''. Il
percorso è di circa \textbf{tre mesi}, fino a dicembre, e durante il corso i
docenti cercheranno di far provare agli studenti ``cosa significa fare Analog
Game Design'' con un \textbf{project work}, un progetto di gruppo di cui si
spiegheranno più avanti la struttura e le richieste.

% ---------------------------------------------------------------------
\section{I docenti}
\label{a1:sec:docenti}

``Anzitutto partiamo da delle presentazioni doverose''. Ognuno dei tre docenti
principali si presenta con una slide che riporta i contatti e una selezione dei
lavori fatti.

\subsection{Daniele: narrativa, fiction e gioco di ruolo}

\textbf{Daniele}\footnote{La trascrizione riporta ``Daniel Cusetto''; nel
resto della lezione tutti lo chiamano ``Daniele''. Il cognome non è
ricostruibile con certezza.} ha già conosciuto quasi tutti gli studenti, perché
ha fatto loro il colloquio di ammissione. È \textbf{narrative designer}, o meglio
\emph{game narrative designer}; insegna nell'accademia che ospita il
corso\footnote{La trascrizione storpia il nome dell'accademia in molti modi
(``Dive'', ``IVE'', ``Accademia Live'', ``Ai'', ``AIVA''). Con ogni probabilità
si tratta dell'\textbf{AIV}, l'Accademia Italiana Videogiochi, che ha sede a
Roma: lo fanno pensare i corsi triennali di game design e narrative design citati
dai docenti e il fatto che una studentessa racconti di essersi trasferita a Roma.
È una ricostruzione, non un'informazione esplicita; nel seguito il capitolo
parla semplicemente dell'``accademia''.} ma lavora soprattutto come \textbf{brand
advisor}. Nel corso si occuperà della parte di \textbf{gioco di ruolo},
\textbf{narrazione} e \textbf{fiction}.

I suoi contatti, ``in quest'ordine'', sono email, Telegram e Discord. Ci tiene a
dare alcune indicazioni precise:
\begin{itemize}
  \item se si vuole avere la certezza di una risposta, bisogna scrivergli su
        \textbf{Telegram}, perché lì ha tutti i suoi clienti;
  \item Discord lo guarda di meno; l'\textbf{email} è per le cose ufficiali;
  \item non bisogna contattarlo altrove: non su Facebook, non su LinkedIn, non su
        Instagram, ``perché li vedo molto di meno''.
\end{itemize}
La slide mostra anche alcuni giochi, attività e brand su cui ha lavorato negli
anni, ma non si dilunga: ``nel caso ne parleremo''.

\subsection{Ruggero: game design e ricerca e sviluppo}

\textbf{Ruggero} (spesso chiamato ``Ruger'') ha fatto il \textbf{game designer}
per varie aziende e oggi è game designer e si occupa di \textbf{ricerca e
sviluppo} per un'azienda del settore\footnote{La trascrizione dice ``per MS'';
il nome dell'azienda non è ricostruibile con certezza. Più avanti Ruggero cita
tra i giochi su cui ha lavorato il gioco da tavolo ufficiale di
\emph{Cyberpunk}, \emph{Marvel United} e \emph{Status}, un gioco sulla
geopolitica.}. Ha un background ``misto'': nel suo portfolio ci sono giochi da
tavolo, giochi di ruolo, giochi analogici, LARP, videogiochi, giochi in realtà
virtuale. Si occupa anche di \textbf{divulgazione} sul gioco e di analisi di
\textbf{game studies} e game design. Più avanti, nelle domande finali, dirà di
avere ``50 e più giochi prodotti e progettati'' e di insegnare da vari anni.

Lo si può contattare liberamente via Gmail e Telegram, e cerca di rispondere
\textbf{entro 24 ore}, salvo imprevisti o eventi: ``sentitevi in libertà di
scrivermi se avete dubbi, senza problemi''. Quando arriva a parlare dei propri
giochi, la slide mostra ``un mix di varie cose, dalle più famose alle meno
famose'', ma Ruggero si imbarazza a presentarli, tra le risate dei colleghi.

\subsection{Beatrice: i giochi pervasivi}

\textbf{Beatrice} (``Bea'') si occupa di \textbf{giochi pervasivi}
(\emph{pervasive games}). ``Se non li avete mai sentiti nominare, non
preoccupatevi'': questo corso è ``uno dei pochi, forse l'unico in Italia'' che
li racconta esplicitamente. Li si vedrà nelle prossime lezioni, ma la definizione
di base è questa: i giochi pervasivi \textbf{utilizzano gli spazi e gli oggetti
della vita quotidiana e li trasformano in spazi e oggetti di gioco}.

Beatrice lavora come \textbf{freelance} e si occupa di diversi tipi di
progettazione: lavora sia con le \textbf{aziende} sia con gli \textbf{eventi dal
vivo}. La sua attività diretta si trova sotto l'etichetta \textbf{Minime
Differenze}\footnote{La trascrizione riporta il nome dell'etichetta come
``Minime Differenze'' e l'indirizzo del sito come ``mininedifferenze.com''; la
grafia esatta del sito non è verificabile dalla trascrizione.}, mentre i lavori
per le aziende sono ``sparpagliati'' sui siti delle varie aziende. Molte delle
cose mostrate nella sua slide sono attività proposte al pubblico ``con un
biglietto, con un gioco in formato evento'': un tema che si toccherà nel corso,
per capire come \textbf{comunicare} i propri giochi in modo che siano
comprensibili, accessibili per il pubblico e \textbf{vendibili}.

Ammette un ``grosso limite'' con i messaggi: è un po' lenta a rispondere, ``ma
rispondo sempre a tutto''. La si può contattare per email (cognome punto nome) o
su Telegram. È ``molto curiosa di vedere cosa produrrete''.

\subsection{I docenti ospiti}

Daniele presenta poi i tre docenti che il corso avrà ``solo per alcune
lezioni'', tutte verso \textbf{novembre}:
\begin{itemize}
  \item \textbf{Diego Ricchiuti}, direttore del corso di Game Design
        dell'accademia, farà una lezione sull'\textbf{RNG} (\emph{random number
        generation}, la casualità nei giochi): ``vi parleremo di RNG quando sarà
        necessario. Comunque avrete una lezione con lui, sappiatelo'';
  \item \textbf{Marta}\footnote{La trascrizione riporta il cognome come
        ``Cercasassi''; la grafia non è verificabile.} farà \textbf{due lezioni
        sul development}: quest'anno ``abbiamo deciso di fare doppietta'',
        perché l'anno scorso il suo intervento è stato ``molto, molto utile'';
  \item \textbf{Flavio}\footnote{La trascrizione riporta il cognome come
        ``Mortarino'' e, più avanti, ``Mortalino'', e il nome della sua casa
        editrice come ``Grampibert''; nessuno dei due è ricostruibile con
        certezza.} farà la lezione di \textbf{game production}; è un editore, e
        di solito partecipa anche alla tavola rotonda con gli editori della
        giornata finale.
\end{itemize}
Le slide sono già online e gli studenti avranno accesso alla cartella condivisa:
``non vi preoccupate di queste cose, ne parleremo a breve''. Ma prima di parlare
del corso, ``vogliamo conoscere voi''.

% ---------------------------------------------------------------------
\section{Le presentazioni dei partecipanti}
\label{a1:sec:partecipanti}

Le domande per la presentazione sono semplici, da affrontare ``a turno, senza
fretta, con calma'':
\begin{itemize}
  \item nome e occupazione (``se volete dirla'');
  \item perché si vogliono fare giochi analogici, e a che livello: che cosa ha
        portato lo studente al corso;
  \item che tipo di esperienze o conoscenze si hanno già nell'ambito, se si è già
        fatto qualcosa oppure no.
\end{itemize}
``Chi vuole partire senza nessun tipo di ritegno e di vergogne?''. Per far
parlare gli studenti guardandoli, Daniele interrompe la condivisione dello
schermo e ricorda lui le domande quando serve. Le presentazioni sono interessanti
anche per i commenti dei docenti, che anticipano diversi principi del corso.

\subsection{Tiziana}

Tiziana fa l'\textbf{infermiera}, ma nella vita ha fatto molti lavori:
impiegata, insegnante e ``tanti altri lavoretti''. Gioca ``da sempre'': il gioco
fa parte anche della sua vita familiare, perché suo fratello fa il programmatore
di videogiochi. Ha già \textbf{organizzato} giochi: una caccia al tesoro per far
conoscere il territorio, due scene con un delitto (giochi con omicidio da
risolvere), giochi per intrattenere le persone. Ha riscoperto di recente i giochi
da tavolo, arrivando a \emph{Carcassonne} con famiglia e amici; poiché è
difficile trovare persone con voglia di giocare nella propria cerchia ristretta,
ha scoperto il mondo delle \textbf{associazioni} ludiche, ``e lì è stato
un'esplosione'': prima in Veneto (è di Vicenza), poi a Roma, dove si è
trasferita.

Ha in mente un progetto ambizioso: far conoscere una certa realtà con la modalità
del \textbf{gioco di ruolo}, o comunque proporre un gioco suo.

\paragraph{Il commento: gli strumenti, non il gioco dei sogni.} Ruggero coglie
l'occasione per un'avvertenza che vale per tutti, visto che molti studenti
arrivano con idee già pronte. Nel corso si lavorerà a un progetto, ma
l'obiettivo dei docenti è dare \textbf{strumenti} da applicare anche
\emph{fuori} dal progetto, ai giochi che gli studenti vorranno fare un domani:
``che non siano esercizi per il corso, ma per voi in primis''. Quindi: ``non
rimaneteci male se non riuscirete a lavorare al vostro gioco dei sogni in questo
corso''. Anche perché ``sarebbe impossibile, in due mesi, fare il gioco dei
sogni''. Il corso darà invece la possibilità di fare un'esperienza ludica
insieme e di imparare a replicarla sulle proprie idee: ``le idee non buttatele,
anzi, il contrario: tenetele''.

\subsection{Nicola}

Nicola è uno \textbf{studente}. Sei mesi prima ha fatto un colloquio con Daniele
perché stava valutando di passare all'accademia, e ha deciso, ``nel dubbio'', di
cominciare da questo corso breve per capire come funziona la formazione;
eventualmente si iscriverà al corso di game design vero e proprio. Gli interessa
in particolare il lato dell'esperienza del giocatore, su cui è ``molto critico''
(qui usa la parola ``divertimento'', e Daniele lo ferma con un ``ne parleremo'':
è una delle parole che il corso vieta, come si vedrà nella
sezione~\ref{a1:sec:parole-vietate}).

Ha 21 anni ed è \textbf{giocatore di ruolo da circa otto}. Prima ancora di
conoscere i giochi di ruolo si era creato un sistema tutto suo, ``pessimo,
disgustoso, terribile'', un sistema a percentuale con il d100 ``terrificante'';
poi ha scoperto i giochi di ruolo da tavolo, i classici (\emph{Dungeons \&
Dragons}, \emph{Daggerheart}, \emph{Vampire}\footnote{La trascrizione riporta
``DD'', ``Daggerart'' e ``Vampire'': quasi certamente \emph{Dungeons \&
Dragons}, \emph{Daggerheart} e \emph{Vampire: The Masquerade}.}). Gioca anche a
\textbf{card game} e ad altri giochi sociali, pur avendo sempre avuto difficoltà
a trovare gruppi, ``avendo vissuto in un buco dimenticato dal mondo''. Come
giocatore è ``l'eterno narratore'' del gruppo, quello che prepara la storia e si
costruisce le regole. Come autore ha fatto soprattutto giochi su carta, ``perché
sono incapace a programmare, per adesso''.

Daniele osserva che il tema del giocatore e del suo rapporto con il gioco è ``un
argomento importante'', anzi ``un vaso di Pandora enorme'' (espressione che
tornerà molte volte nella lezione), e che il corso tratterà anche che cos'è
l'analog game design in generale e quanti tipi di prodotti esistono.

\subsection{Rita}

Rita\footnote{Daniele la chiama prima ``Nino'' (probabilmente un errore della
trascrizione) e poi la ringrazia come Rita.} ha 30 anni e si collega da una
regione del Sud, dove ``ancora si muore di caldo''. Ha studiato
all'\textbf{Accademia di Belle Arti di Napoli}, \textbf{design della
comunicazione visiva}: si occupa dell'aspetto più imprenditoriale della
progettazione artistica. Non volendo progettare per un'impresa (``fare loghi,
robe del genere''), vorrebbe applicare la propria creatività alla realizzazione
di giochi analogici.

Non è una giocatrice di ruolo e non lo è mai stata: è una giocatrice
\textbf{tradizionale}, cresciuta con Monopoly, la dama e giochi tradizionali
cinesi. Ha già \textbf{idee, prototipi e progetti} realizzati, e vuole capire
come renderli \textbf{prodotti per la vendita} e se questo è il campo adatto a
lei; magari anche confrontarsi con i colleghi. La sua progettazione è orientata
al \textbf{sociale} e all'\textbf{educativo}: ``responsabilizzare un po' le nuove
generazioni su come sta andando il mondo e su come potrebbe andare, visto che è
tutto nelle nostre mani''.

\paragraph{Il commento: siate curiosi.} ``Pane per i tuoi denti'', le risponde
Daniele, soprattutto per i LARP e i giochi pervasivi: ``vi apriremo un po' la
testa da questo punto di vista''. E aggiunge un consiglio per tutti: ``siate
curiosi'', perché imparare da altri ambiti dà moltissimo. Ognuno è lì per fare
qualcosa di preciso, ``però siate aperti anche ad altre esperienze''.

\subsection{Tancredi}

Tancredi ha 21 anni e non ha ancora grande esperienza lavorativa: qualche
``piccolo progettino'' che non è andato come voleva. Gli interessa il mondo dei
giochi, digitali e non, con un approccio \textbf{educativo} e \textbf{sociale}.
È arrivato soprattutto dai videogiochi, a cui gioca da quando era piccolo; ai
giochi da tavolo si è avvicinato solo nell'ultimo anno, grazie a una realtà
ludica, perché in famiglia ``non siamo molto da giochi''. Ha realizzato qualche
piccolo videogioco con vari engine: prima quelli senza un vero linguaggio di
programmazione, ``più a nodi'', poi qualcosa di più complicato, e ora sta
lavorando a un piccolo prototipo.

\paragraph{Il commento: niente di quello che vedrete è sprecato.} ``Non butterai
via niente di analog game design'', gli risponde Daniele, ``perché il 99 per cento
delle cose che vedremo si ritrovano quasi uno a uno nel mondo dei videogiochi''.
Tanto che nei corsi triennali dell'accademia (anche in quello di narrative
design, dove insegna Ruggero) il primo anno si parte dai \textbf{giochi da tavolo
e di ruolo}, proprio per spiegare meccaniche più complesse che servono poi per i
videogiochi. Ruggero aggiunge che vale anche il contrario: alcuni dei
\textbf{framework} di game design che si vedranno nelle lezioni tecniche vengono
dal videogioco e sono stati riapplicati all'analogico. ``Evitare la
crossmedialità qui è impossibile''. Anzi, tutte le conoscenze ``mediatiche
adiacenti'' (musica, letteratura, cinema) saranno utili.

\subsection{Lorenzo Vergata}

I Lorenzo in classe sono due; il primo a presentarsi è \textbf{Lorenzo
Vergata}, 26 anni, \textbf{dottorando al Politecnico di Milano}. Gioca ai
videogiochi da quando aveva due anni e mezzo, ma ha giocato al suo primo gioco da
tavolo ``credo nove mesi fa''. Gli fa piacere sentire che molti in classe sono
interessati ai giochi \textbf{sociali} ed \textbf{educativi}, un settore che
interessa anche a lui e che ``può essere esplorato molto'', anche con qualche
aiuto dell'intelligenza artificiale per esperienze più ricercate.

Di esperienza nell'analogico dice di avere fatto ``il minimo sindacale''. Ma
racconta che diciotto anni prima, da bambino, era ``più game designer di
adesso'': prendeva un gioco e ne cambiava le regole, anche al biliardino (``no,
ora introduciamo questa regola''). Gli piace sperimentare, creare regole e
vedere che cosa accade in sé e negli altri. Soprattutto vorrebbe ``una chiave di
lettura per poter rimanipolare la realtà'': tre settimane prima è stato a fare la
\textbf{vendemmia in Francia} e ha visto tantissime dinamiche che potrebbero
diventare un gioco da tavolo o un videogioco. Per esempio: finire contro un rovo
pieno di semi che si attaccano ai vestiti lo rallentava nella raccolta
dell'uva, mentre in un altro filare era più rapido. ``Si può fare pure una
dinamica'', e vorrebbe gli strumenti per farne giochi.

Arriva subito la battuta: una ``vendemmia gamificata'' da proporre in Oltrepò
Pavese (``io sono dell'Oltrepò Pavese, quindi devo per forza bere vino'')\footnote{La
trascrizione dice ``nell'altro papavese'' e non chiarisce chi pronunci la
battuta.}. Daniele conferma che quest'anno il corso avrà ``tante cose sul
sociale'': ``ci piace parecchio''.

\subsection{L'altro Lorenzo}

Il secondo Lorenzo (``secondo, ma non meno importante'') è uno \textbf{scrittore
di racconti brevi} e sta scrivendo il suo primo romanzo. Gli è sempre piaciuto
raccontare storie, e da più grande scriverle; gioca ai videogiochi ed è ``un
amante delle storie e della cultura pop''. Negli ultimi anni è entrato nel mondo
del \textbf{gioco di ruolo da tavolo} e vorrebbe unire le due passioni, le
storie e il gioco di ruolo, magari facendone un lavoro; gli interessa anche lo
\emph{strategic writing}. Daniele gli ricorda che ci sarà una lezione proprio
sulla \textbf{fiction}.

\subsection{Angelo}

Angelo si collega senza videocamera dal lavoro: fa il \textbf{birraio} (in un
piccolo birrificio in Brianza, a Bernareggio, come racconterà dopo la pausa) e,
con il caldo, ``vi risparmio l'orrore''. Per molto tempo era convinto che i
giochi da tavolo fossero solo i classici (Monopoly, Risiko) o i party game; poi,
quattro o cinque anni fa, ha ``sbattuto la faccia'' contro giochi davvero
appassionanti, con una specie di narrazione dietro che li rendeva ancora più
coinvolgenti, e non ha più smesso. È uno che ``ha sempre un botto di idee, ma poi
non mette in pratica assolutamente nulla'': cercava un corso che gli desse le
basi e soprattutto lo \textbf{motivasse} a provare a fare qualcosa.

\paragraph{Il commento: bisogna buttarsi.} Daniele racconta che anche l'anno
prima un suo caro amico si era iscritto proprio per questo motivo. Il corso è
breve, ma cerca di portare gli studenti ``da 0 a 100'' e di far realizzare un
\textbf{prodotto} da presentare alla fine. Il punto è che i giochi ``bisogna
farli, bisogna buttarsi''. E, venendo dalla narrazione, Daniele lo paragona allo
scrivere: si scrive e si riscrive, e allo stesso modo si fa \textbf{playtest}, si
rifà, si aggiorna. ``Non è che si fanno così''.

\subsection{Gianmarco (dopo le slide)}

Gianmarco arriva a lezione iniziata e si presenta più tardi, alla fine della
prima parte, dopo che i docenti hanno chiuso le slide organizzative e la
discussione sull'intelligenza artificiale (``Gianmarco, ti giuro, ti facciamo
presentare''). Lavora nel \textbf{team di innovazione} di un'azienda italiana.
Vuole fare giochi perché è una sua passione, che lo gratifica, e in passato ci ha
provato ``in maniera scomposta''. Cerca un \textbf{metodo} e un confronto con
persone più esperte: ``studiare solo dei manuali è molto sterile'', ``stare
chiusi nel pertugio non aiuta'', ed ``essere autodidatta ha estremamente i suoi
limiti''. Come esperienza è un giocatore e ha letto ``un po' di materiale, di
blog, di manuali'', ``niente di rilevante''.

Manca solo Giordano, che si presenterà alla lezione successiva.

\paragraph{Uno sguardo d'insieme.} La classe è molto varia: un'infermiera,
una designer della comunicazione visiva, un birraio, un dottorando, uno
scrittore, un impiegato nell'innovazione, due studenti ventunenni. Alcuni
vengono dal gioco di ruolo, altri dai giochi tradizionali, altri dai
videogiochi; quasi tutti hanno scoperto i giochi da tavolo moderni da pochi anni
o da pochi mesi. E almeno tre studenti dichiarano un interesse per i giochi
\textbf{sociali} ed \textbf{educativi}. Come si vedrà nella
sezione~\ref{a1:sec:scopo}, i docenti considerano questa varietà una risorsa.

% ---------------------------------------------------------------------
\section{I materiali e gli strumenti del corso}
\label{a1:sec:materiali}

Prima delle slide sul corso, Daniele sbriga un po' di ``clean up casalingo''.

\paragraph{Il form iniziale.} C'è un \textbf{form} da compilare ``entro e non
oltre il 15 settembre'': il QR code è nelle slide, e la segreteria lo ricorderà
su Discord. Richiede circa 10--15 minuti; si può fare subito o in giornata, ``e
ve lo togliete di torno''. È un form standardizzato per tutti i corsi (``non
siamo io e Ruggero a farvi passare la tortura del test''), diviso in due parti:
\begin{itemize}
  \item una parte di \textbf{conoscenza del candidato}: quanto tempo si ha a
        disposizione per il corso, come ci si organizza con il lavoro, che cosa
        piace giocare. È la parte che interessa ai docenti, perché li aiuta a
        organizzare le lezioni;
  \item una parte di \textbf{test logico}, che i docenti non hanno scritto (``io
        non sono responsabile del form'', ``no, no, nemmeno io'').
\end{itemize}

\paragraph{La cartella condivisa.} Tutti gli studenti vengono aggiunti alla
cartella del corso su \textbf{Google Drive}, che si trova tra i file
``condivisi con me''. Gli studenti ne sono \textbf{gestori dei contenuti} e
possono caricarci quello che vogliono. Dentro ci sono:
\begin{itemize}
  \item le \textbf{attività in classe};
  \item le \textbf{dispense};
  \item i \textbf{materiali didattici}, con le slide (quelle della lezione sono
        già lì);
  \item le cartelle dei \textbf{project work} futuri;
  \item i \textbf{task}, che ``saranno pochi, ma ci saranno'';
  \item il \textbf{calendario studenti}, ``importantissimo'', con due schede: il
        calendario delle lezioni e quello del project work.
\end{itemize}

\paragraph{Discord.} Il corso ha anche un \textbf{server Discord}, a cui tutti
sono stati invitati; ci sono i docenti (Daniele, Ruggero, Diego e altri), che si
possono taggare. Di solito è lì che la segreteria pubblica le informazioni, per
esempio se una lezione ``dovesse per caso essere cancellata''. Nelle domande
finali Ruggero aggiungerà che c'è un canale apposito per il corso, che gli
studenti possono usare liberamente per organizzarsi tra loro.

\paragraph{Le reaction di Zoom.} Durante le lezioni su Zoom si possono usare le
\textbf{reaction} (il pollice in alto, il cuoricino, la mano alzata): un modo
rapido per far capire ai docenti se si è capito, se si vuole intervenire. Nel
frattempo arrivano i primi pollici.

% ---------------------------------------------------------------------
\section{Lo scopo e il metodo del corso}
\label{a1:sec:scopo}

La prima serie di slide (che Daniele, per un lapsus, presenta come ``cos'è il
corso di pulire il pavimento di casa'') spiega che cosa il corso vuole fare e
come.

\subsection{Diventare autori analogici}

Lo \textbf{scopo del corso} è fornire le competenze per \textbf{diventare autori
analogici} e stimolare la creatività e l'\textbf{autorialità} di ciascuno: ``voi
imparerete a diventare autori di giochi analogici''. La differenza tra
\emph{game designer} e \emph{game developer} verrà chiarita più avanti (sezione~\ref{a1:sec:design-development}),
ma la parola chiave qui è \emph{autorialità}: ``vogliamo voi, vogliamo le vostre
idee''. Magari non si potrà usare l'idea del gioco della propria vita, ma si
potranno portare nelle lezioni le proprie idee e le proprie sensibilità.

Gli obiettivi elencati nella slide sono questi:
\begin{itemize}
  \item dare \textbf{maggiore consapevolezza nel design e nello sviluppo}. In
        Italia, purtroppo, ``si fa design a volte in maniera non consapevole''.
        Consapevole vuol dire avere almeno un'introduzione alle basi del design e
        dello sviluppo, e fare le cose ``con metodo, con raziocinio, con
        consapevolezza'';
  \item fornire gli \textbf{strumenti essenziali} per un design consapevole;
  \item stimolare la \textbf{sperimentazione} e la \textbf{trasversalità
        mediatica}. Non si chiede di fare ``il \emph{Root} della situazione o il
        \emph{Carcassonne} della situazione''\footnote{La trascrizione dice ``il
        lumaven della situazione o il carcassone della situazione''; il primo
        titolo non è ricostruibile con certezza. Il senso è chiaro: non si
        chiede di rifare un grande classico del gioco da tavolo.}, ma un
        project work \textbf{calibrato sulla durata del corso}, e però
        \textbf{sperimentale} e ``assolutamente transmediale'';
  \item portare \textbf{numerosi esempi} di tipologie di gioco, dall'analogico al
        tavolo, al pervasivo ``e oltre''. Lo si vedrà nella lezione sui
        \textbf{casi studio} del venerdì successivo: ``voi non avete idea di
        quante cose esistono, bellissime, sconosciutissime e favolose''. Daniele
        confessa che preparando e aggiornando le slide i colleghi gli hanno
        mostrato cose che lui stesso non sapeva nemmeno che esistessero;
  \item dare gli strumenti per \textbf{analizzare} i giochi altrui. Ruggero lo
        sottolinea: se il corso riesce nel suo intento, gli studenti sapranno
        non solo fare i propri giochi ma anche analizzare giochi già fatti, che è
        ``una delle skill importanti anche nella parte di ricerca e sviluppo''.
        Si vedranno casi studio e giochi pubblicati, sempre con l'obiettivo di
        capire \textbf{che cosa fanno e perché lo fanno}, per poi fare meglio il
        proprio lavoro;
  \item imparare a \textbf{fare giochi insieme}. Esistono gli autori che lavorano
        da soli (i \emph{solo developer}), ``ma sappiate che anche l'autore che
        lavora da solo poi si deve interfacciare con il pubblico, con un editore,
        magari con un editor, con un publisher''. ``Non lavorerete mai da soli'':
        bisogna imparare a lavorare con gli altri.
\end{itemize}

\subsection{Un corso intersezionale}

Il corso è \textbf{intersezionale}: \textbf{non ci sono lezioni dedicate a una
singola tipologia di gioco}. Non ci sarà ``la lezione sul gioco da tavolo'',
``la lezione sul gioco di ruolo'', ``la lezione sul gioco pervasivo''. Ci saranno
lezioni su temi diversi, ciascuna pensata per essere utile a \textbf{tutte e tre
le sfaccettature} del corso; alcune saranno più importanti per il gioco da
tavolo, altre per il pervasivo, altre per il gioco di ruolo, ma nessuna sarà
esclusiva. I docenti vogliono ``contaminazioni e applicazioni assolutamente
trasversali''.

\subsection{Un corso cross-generazionale}

È anche un corso \textbf{cross-generazionale}: i gradi di esperienza di gioco
quest'anno sono abbastanza allineati, ma le esperienze di vita, le età, le
estrazioni sociali sono varie. ``Ed è un pro, non è un contro. Più diversificata
è la classe, più imparate meglio anche voi''. La slide ha una piccola nota per
fasce d'età:
\begin{description}
  \item[Under 30] Concentratevi sulla vostra \textbf{crescita}. Evitate di
        confrontarvi con le persone più grandi che magari ``ce la fanno meglio''
        in alcune cose.
  \item[Over 30] Concentratevi sul \textbf{consolidamento}, ma siate
        \textbf{umili}: i docenti diranno cose che faranno ``storcere il naso''.
        ``Bisogna un po' reimparare per imparare'': bisogna sapersi adattare.
\end{description}
Ruggero contestualizza: lo dicono perché in passato hanno visto studenti che,
essendo più adulti o avendo giocato di più, pensavano che l'esperienza
\emph{di gioco} si trasformasse automaticamente in competenza \emph{di
creazione} di giochi. Non è così. E non bisogna sottovalutare chi gioca meno:
proprio perché gioca cose diverse, per risolvere un problema di design specifico
``è una ricchezza incredibile avere qualcun altro che riesce a uscire dalla
scatola''.

\subsection{Un corso pratico e adattabile}

Il corso è \textbf{pratico}: gli studenti si concentreranno su un progetto di
gruppo, il \textbf{project work}, lavorando in team. Ed è \textbf{adattabile}: il
programma è definito, ma ``si aggiorna sui vostri stimoli, priorità e
necessità''. Se durante il corso i docenti vedono che serve più una cosa
dell'altra, sono ``prontissimi'' a dare più spazio a un argomento invece che a
un altro.

\subsection{Inquiry-based learning: domandate}

Il metodo è quello dell'\textbf{inquiry-based learning}, l'apprendimento basato
sulle domande. Di solito le persone dicono ``ho una domanda, magari sarà
stupida''. ``Non esistono domande stupide. Pensate che ci sia una domanda
stupida? Fatela, alzate la mano''. Vanno bene anche le
\textbf{provocazioni}: se qualcosa non torna, è meglio chiedere. La condivisione
è fondamentale.

I docenti danno \textbf{metodi, dati e conoscenze}, ma lasciano trovare le
risposte agli studenti: ``vi lasciamo un pochettino annegare nel mare prima di
darvi la soluzione, in modo tale che la possiate trovare da soli''. Di
conseguenza \textbf{non ci saranno interrogazioni sulle definizioni}: interessa
vedere \textbf{come} gli studenti applicano ciò che imparano. Se una
definizione non viene in mente, non è un problema; l'importante è saperla usare.

\subsection{Feedback, non valutazione}

``Sì al feedback, no alla valutazione''. I docenti danno \textbf{feedback
verbali}, anche sulle attività da fare a casa (i task, che non saranno tanti).
Il corso è calibrato per \textbf{studenti lavoratori}: c'è un po' da lavorare
fuori dalle lezioni, ma in modo calibrato.

Il principio è \textbf{fallire oggi per imparare per il futuro}: \emph{fail
forward}, \emph{fail to learn}. I docenti lasciano provare e fallire almeno una
volta prima di dare la soluzione: ``il classico metodo Montessori, quello
originale, non quello travisato''.

\subsection{Le soft skill e i criteri di valutazione}

Grande attenzione va alle \textbf{soft skill}: i lavori di gruppo, lo scambio
propositivo, l'atteggiamento positivo, la propositività, la proattività.
``Lavoro di gruppo'' non significa stare ogni giorno in sala studio con gli
altri: va benissimo che un gruppo si organizzi dandosi priorità e tempi in base
alla disponibilità di ciascuno. L'importante è \textbf{collaborare}, essere
\textbf{professionali}, rispettarsi a vicenda, rispettare le idee proprie e
altrui e lasciare spazio agli altri, perché i progetti ``non saranno solo
vostri, saranno di tutte le persone del vostro gruppo''.

La valutazione del project work non riguarda solo la \textbf{documentazione}
prodotta, ma anche:
\begin{itemize}
  \item la \textbf{presentazione} finale;
  \item l'\textbf{autovalutazione};
  \item la \textbf{crescita personale};
  \item l'\textbf{atteggiamento} tenuto nel corso.
\end{itemize}
Daniele cita: ``da ognuno secondo le sue capacità, a ognuno secondo i suoi
bisogni''. Lavorare in gruppo non vuol dire che, in tre, ciascuno lavori
esattamente il 33\% del tempo. Ai docenti serve piuttosto capire \textbf{chi fa
che cosa} nel gruppo, e che il processo interno sia \textbf{trasparente}.

\paragraph{Lavorare in cloud.} Per lo stesso motivo il consiglio è di
\textbf{non lavorare in locale}, ma direttamente nelle cartelle condivise del
project work e dei task. Si è liberi di fare altrimenti, ma lavorando in cloud
si dà ai docenti la possibilità di vedere \emph{come} si lavora, il proprio
``processo mentale''.

% ---------------------------------------------------------------------
\section{Le regole di comportamento e il carico di lavoro}
\label{a1:sec:regole}

Alcune slide sono di quelle che si ripetono in ogni corso (``\emph{repetita
iuvant}'', storpiato per scherzo in ``ripetita Juventus'').

\paragraph{Comportamento.} Prestare attenzione durante le lezioni; interagire;
portare rispetto; essere proattivi; alzare la mano e fare domande.
\textbf{Tolleranza zero} per pregiudizi, maleducazione, mancanza di rispetto,
\emph{cheat} e scorciatoie. Daniele ammette che ``cheat e scorciatoie'' viene da
altri corsi dove il tema è più importante, ma Ruggero lo ferma: ``non
sottovalutate questa cosa'', perché implica anche \textbf{non far fare a ChatGPT
i task}. Daniele ha una slide apposta, ``perché è un discorso molto complesso nel
2026'' (sezione~\ref{a1:sec:ia}).

\paragraph{Le lezioni sono intense.} Le lezioni possono essere \textbf{dense},
``soprattutto le mie, sappiatelo: io sono logorroico'', avverte Daniele. Bisogna
chiedere spiegazioni se serve, prendere appunti, rivedere le registrazioni. Le
lezioni sono ``delle scintille'': modi per far esplorare i temi oltre le slide.
``Andate oltre, se riuscite''.

\paragraph{Il carico di lavoro.} Le regole sui task sono queste:
\begin{itemize}
  \item i task \textbf{sono per voi, non per noi};
  \item se si ha poco tempo, è meglio farli rapidamente piuttosto che non farli:
        ``provateci'';
  \item i docenti cercano di non chiedere più di \textbf{un'ora o due a
        settimana} oltre alle ore di lezione;
  \item non tutte le lezioni avranno un task, e il \textbf{project work} è il più
        importante di tutti.
\end{itemize}

\paragraph{I feedback dei docenti.} I feedback ``si guadagnano'' man mano che si
reagisce bene: non bisogna aspettarsi ``dei papiri'' sui task. Se i docenti
hanno voglia di dare feedback più ampi lo fanno, soprattutto se vedono che lo
studente se li merita. Ad Analog di solito non ci sono task \emph{rejected}; ma
se un docente rifiuta un task, vuol dire che si è fatto \textbf{meno del livello
minimo} richiesto dal corso. I feedback possono non essere immediati: ``siate un
pochettino impietositi dal fatto che siamo dei docenti''. E se un feedback non
arriva, o se ne vogliono di più, bisogna scrivere, ``perché siamo tre esseri
umani e quindi possiamo anche dimenticarci''.

% ---------------------------------------------------------------------
\section{Il project work e il calendario}
\label{a1:sec:project}

\subsection{Come funziona il project work}

Una prima slide con le date del project work è da aggiornare (``perdonatemi il
refuso''); Daniele la corregge al volo e rimanda alla slide successiva, più
precisa. Il funzionamento è questo:
\begin{itemize}
  \item i \textbf{gruppi li formano i docenti}. Gli studenti sono nove, quindi si
        faranno \textbf{tre gruppi da tre}: una dimensione pensata anche per il
        fatto che molti sono studenti lavoratori, così da non ritrovarsi ``in due,
        a esplodere'';
  \item ogni gruppo deve prima preparare un \textbf{pitch} del proprio gioco, da
        presentare in anticipo: un'idea, una bozza del gioco;
  \item poi si lavora al gioco \textbf{nelle pause del corso}: da qui a dicembre
        ci saranno due pause, quindi non ci sarà lezione tutti i lunedì e tutti i
        venerdì;
  \item ci saranno \textbf{review intermedie} con i \textbf{tutoraggi}, per lo più
        il \textbf{sabato mattina}, dalle 11 alle 12, con Ruggero su Google Meet.
        Servono ad aiutare i gruppi a portare avanti il project work.
\end{itemize}

\subsection{Il calendario delle lezioni}

Il \textbf{calendario} è nel calendario studenti della cartella condivisa.
Attenzione: le date \textbf{possono variare}. ``Il 90\% rimarranno così'', ma
può capitare che, per motivi di lavoro, qualcuno stia male o non possa, o che
``scompaia l'elettricità da tutte le città d'Italia''. E l'ordine delle lezioni
può cambiare se i docenti vedono che agli studenti serve prima una lezione
invece di un'altra.

Ci sono \textbf{due pause importantissime}:
\begin{itemize}
  \item dal \textbf{13 ottobre al 5 novembre}: è il periodo delle grandi fiere
        (\textbf{Essen} e \textbf{Lucca}\footnote{La trascrizione dice ``Essen,
        Luca, Cose'': si tratta con ogni probabilità della fiera Spiel di Essen,
        in Germania, e di Lucca Comics \& Games, che si tengono entrambe in
        ottobre.}), ``un periodo un po' frenetico'', e cominciano anche i corsi
        accademici, quindi la pausa permette ai docenti di respirare. Per gli
        studenti è la pausa più importante: quasi un mese in cui sviluppare la
        parte importante del gioco;
  \item dal \textbf{24 novembre all'11 dicembre}, dopo qualche lezione a novembre.
\end{itemize}

E ci sono \textbf{due date da evidenziare}:
\begin{itemize}
  \item il \textbf{25 settembre}: nei primi 5--10 minuti della lezione (che per il
        resto sarà tenuta da Beatrice) i docenti annunceranno i tre gruppi;
  \item il \textbf{12 ottobre}: il \textbf{pitch} dei giochi, con l'idea
        embrionale. Tra il 25 settembre e il 12 ottobre ogni gruppo prepara un
        breve pitch in formato slide; quel giorno i docenti lo rivedono e danno
        feedback: se il gioco va bene, se non va bene, come svilupparlo.
\end{itemize}
Dal 12 ottobre fino all'inizio di dicembre si sviluppa il gioco, con i
tutoraggi: ``sfruttate questi momenti di pausa''.

\subsection{Il calendario del project work}

La slide successiva, ``reale e aggiornata'', riporta il calendario del project
work (tabella~\ref{a1:tab:calendario} e figura~\ref{a1:fig:calendario}).
``Per ora non pensate al project work: ci penserete dal 25 settembre in poi''.

\begin{table}[htbp]
\centering
\begin{tabular}{@{}lp{9.2cm}@{}}
\toprule
\textbf{Data} & \textbf{Che cosa succede} \\
\midrule
ven 25 settembre\footnotemark & lezione: annuncio dei gruppi e del
\emph{mandate}, cioè di che cosa si dovrà fare \\
sab 3 ottobre, ore 11 & prima review e tutoraggio con Ruggero (dopo una
decina di giorni di lavoro) \\
lun 12 ottobre & pitch del gioco durante la lezione \\
13 ottobre -- 5 novembre & prima pausa: sviluppo del gioco \\
sab 7 novembre & tutoraggio \\
sab 14 novembre & tutoraggio \\
24 novembre -- 11 dicembre & seconda pausa: sviluppo del gioco \\
ven 4 dicembre, entro mezzanotte & consegna del gioco sul Drive (non c'è
lezione) \\
ven 11 dicembre & i giochi vengono visti a lezione \\
14 dicembre & giornata finale, con la tavola rotonda con gli editori \\
\bottomrule
\end{tabular}
\caption{Il calendario del project work come lo presenta Daniele. Le date
possono variare.}
\label{a1:tab:calendario}
\end{table}
\footnotetext{Daniele dice ``lunedì 25 settembre'', ma nel 2026 il 25 settembre
è un venerdì; il giorno della settimana nella tabella è corretto di
conseguenza.}

\paragraph{Perché la consegna è il 4 dicembre.} Il 4 dicembre non c'è lezione:
si deve solo consegnare il gioco, entro mezzanotte, sul Drive. La presentazione
a lezione è l'11 dicembre, ma i tre docenti hanno bisogno di \textbf{una
settimana} per rivedere il gioco in privato e valutare anche la
\textbf{documentazione}; per questo la consegna non coincide con la
presentazione. Che cosa voglia dire esattamente ``consegnare il gioco'' verrà
spiegato più avanti.

\paragraph{La giornata finale.} Il 14 dicembre c'è la \textbf{giornata finale},
con una \textbf{tavola rotonda con gli editori}: di solito viene invitato Flavio,
che è anche docente di game production, più qualche altro editore, ``vediamo
quest'anno chi riusciamo a portarvi''.

\begin{figure}[htbp]
\centering
\begin{tikzpicture}[
    font=\footnotesize,
    x=0.1cm,
    evento/.style={draw, thick, rounded corners=2pt, fill=white,
                   text width=1.8cm, align=center, inner sep=2pt},
    pausa/.style={fill=gray!20, draw=gray!60},
    tacca/.style={thick}
  ]
  % asse: giorno 0 = 25 settembre; 1 unità = 1 giorno
  \draw[pausa] (18,-0.35) rectangle (41,0.35);
  \draw[pausa] (60,-0.35) rectangle (77,0.35);
  \node[font=\scriptsize\itshape] at (29.5,0.18) {pausa};
  \node[font=\scriptsize\itshape] at (65,0.18) {pausa};
  \draw[-{Stealth[length=2.5mm]}, thick] (-3,0) -- (85,0);
  % eventi sopra
  \draw[tacca] (0,-0.15) -- (0,0.9);
  \node[evento, above] at (0,0.9) {25 set\\ gruppi e \emph{mandate}};
  \draw[tacca] (17,-0.15) -- (17,2.2);
  \node[evento, above, fill=gray!8] at (17,2.2) {12 ott\\ \textbf{pitch}};
  \draw[tacca] (70,-0.15) -- (70,0.9);
  \node[evento, above, fill=gray!8] at (70,0.9) {4 dic\\ \textbf{consegna}};
  % eventi sotto
  \draw[tacca] (8,0.15) -- (8,-0.9);
  \node[evento, below] at (8,-0.9) {3 ott\\ tutoraggio};
  \draw[tacca] (43,0.15) -- (43,-0.7);
  \draw[tacca] (50,0.15) -- (50,-0.7);
  \draw (43,-0.7) -- (50,-0.7);
  \draw[tacca] (46.5,-0.7) -- (46.5,-0.9);
  \node[evento, below] at (46.5,-0.9) {7 e 14 nov\\ tutoraggi};
  \draw[tacca] (77,0.15) -- (77,-0.9);
  \node[evento, below] at (77,-0.9) {11 dic\\ giochi a lezione};
  \draw[tacca] (80,-0.15) -- (80,2.2);
  \node[evento, above] at (80,2.2) {14 dic\\ giornata finale};
\end{tikzpicture}
\caption{Una linea del tempo del project work, ricostruita dalle slide (scala
approssimativa: dal 25 settembre al 14 dicembre). Le bande grigie sono le due
pause senza lezioni, in cui i gruppi sviluppano il gioco.}
\label{a1:fig:calendario}
\end{figure}

Finita questa parte, Daniele apre alle domande: ``Domande sul calendario? Sulle
slide che ho fatto? Considerazioni? Non abbiate timore, io non ci credo''. E la
domanda arriva, sull'unico tema che aveva lasciato da parte.

% ---------------------------------------------------------------------
\section{L'intelligenza artificiale nel corso}
\label{a1:sec:ia}

\subsection{La slide di Daniele}

Daniele aveva già presentato una slide che definisce ``importantissima''. ``Lo so
che è paradossale, me ne rendo conto. Lo so che magari mi darete del luddista,
dell'anti-tecnologia: non è così. Uso l'intelligenza artificiale per lavoro. Ma
in questo corso, vi prego, \textbf{non usate l'intelligenza artificiale}''. La
richiesta vale per le attività, per i task e \textbf{anche per la
prototipazione}. Le ragioni:
\begin{itemize}
  \item ``voglio che siate voi a fare il gioco, voglio che siate voi a imparare a
        farlo''. Daniele parla da utilizzatore: prima di ChatGPT, che ha smesso
        di usare, oggi soprattutto di Claude, di cui ``si sottovaluta quanto
        possa essere incerto, insicuro, impreciso, non all'altezza'';
  \item nella game industry, \textbf{i publisher rifiutano automaticamente} i
        pitch e i prototipi fatti con l'intelligenza artificiale. Daniele lo ha
        visto di persona in ``tantissime pitch review'' quell'estate, anche al
        \textbf{First Playable} di Firenze\footnote{La trascrizione dice ``First
        Playboy a Firenze'': si tratta del First Playable, l'evento di business
        dell'industria videoludica che si tiene a Firenze.}: anche nel mondo dei
        videogiochi i publisher che vedono pitch o prototipi fatti con l'IA
        ``rigettano automaticamente la build'', perché hanno ancora timore
        dell'IA e delle sue possibili conseguenze;
  \item ``non so come sarà il futuro'': non sa se bisognerà abituarsi a usarla
        tutto il giorno. Non è ``totalmente contrario'', purché sia usata ``con
        raziocinio''; secondo lui oggi, negli Stati Uniti, non lo è;
  \item ma in un corso che si chiama \emph{Analog Game Design} gli studenti devono
        saper fare i giochi da soli. In futuro, se vorranno farsi aiutare
        dall'IA per essere più efficienti, saranno liberissimi di farlo; durante
        il corso no, perché i docenti vogliono valutare \textbf{le loro
        competenze e capacità}.
\end{itemize}

\subsection{La domanda: che cosa vuol dire usare l'IA in un gioco analogico?}

La domanda di uno studente è precisa: che cosa si intende per ``task fatti con
l'IA''? Qual è l'uso dell'IA in un gioco analogico?

\paragraph{Ruggero: l'esperienza di chi valuta i prototipi.} Ruggero risponde
dalla sua esperienza nel settore, dove riceve prototipi da valutare per
un'eventuale pubblicazione. Distingue due usi:
\begin{enumerate}
  \item le \textbf{immagini generate}: illustrazioni, graphic design e così via.
        Per molti publisher è ``un \emph{reject} immediato, proprio immediato'';
  \item l'IA usata per il \textbf{regolamento}, per scrivere le regole e le
        meccaniche. È più difficile da individuare, ma molti publisher se ne
        accorgono comunque, perché le IA ``non sono ancora brave abbastanza da
        creare un regolamento coeso meccanicamente e scritto con un senso per il
        vostro gioco''. Quello che fanno è prendere parti di altri regolamenti,
        elaborare le strutture che vedono nei regolamenti pubblicati e
        riapplicarle. Ma molti giochi si prestano poco a un regolamento
        strutturato come quello di altri, perché \textbf{ogni gioco è un
        ecosistema meccanico a sé}. Il risultato sono regolamenti con
        riferimenti non immediati, e soprattutto con \textbf{la stessa meccanica
        descritta in modo diverso in pagine diverse}: due difetti che fanno
        perdere tempo all'editore, e in questo caso ai docenti.
\end{enumerate}
Poi aggiunge un punto ``personale, soggettivo'': i docenti sono interessati a
vedere il lavoro, la creatività, ``il frutto della vostra creazione''. E
viceversa gli studenti vogliono i commenti dei docenti: ``immaginatevi se voi
fate un gioco e noi vi diamo i commenti scritti con ChatGPT''. Altrimenti, dice
ridendo, gli studenti fanno il gioco con ChatGPT, i docenti lo danno a ChatGPT
per i commenti, ``e facciamo molto alla svelta''.

\paragraph{Daniele: perché gli studenti usano l'IA.} Daniele racconta la sua
esperienza con gli studenti, che secondo lui ricorrono all'IA in due momenti:
\begin{enumerate}
  \item quando vanno in \textbf{crisi}: in ansia da prestazione o per i tempi.
        Cercano di risolvere il problema della propria disorganizzazione
        facendolo risolvere all'IA, senza accorgersi che l'output è ``molto,
        molto scarso'';
  \item quando hanno l'impressione che l'intelligenza dell'IA sia
        \textbf{migliore} di quello che scrivono loro, o di quello che c'è in
        giro. Per Daniele è una ``mancanza di palato critico'': l'output dell'IA
        sembra buono in superficie, ma quando si va nel dettaglio e ci si fa
        delle domande ``di solito fa un po' di \emph{crumbling}'', si sbriciola.
\end{enumerate}
Lui usa Claude ogni giorno e ne conosce i limiti: ``non c'è nemmeno una frase
che io copio direttamente''. Gli fa fare il lavoro più noioso, più ripetitivo e
soprattutto \textbf{strutturato}, ``replicare delle cose come se fosse una
macro''. Lo usa per il testo e per il codice, mai per le immagini: per motivi di
\textbf{diritto d'autore}, e perché ha molti amici artisti ed è stufo di vedere
il loro lavoro ``dilaniato'' dall'IA.

Ammette che i docenti non possono controllare gli studenti; per questo
``ve lo consiglio''. E riprende quella che definisce la cosa più bella detta da
Ruggero (``ti ruberò la frase per il corso di narrative''): i docenti i feedback
li danno di persona. Daniele a volte manda 30 minuti di audio a uno studente, o
fa intere call per rivedere script, pitch e giochi: ``è il mio lavoro, mi piace
farlo e voglio farlo io, quindi mi aspetto che lo studente lo faccia lui''. Come
freelance, poi, è ``stufo marcio'' dei clienti che arrivano con qualcosa fatto
con l'IA (``guarda che bello, pubblica questo''), perché il più delle volte il
risultato è pessimo: non sanno come usarla. ``È anche una questione di
competenza''.

\paragraph{Il prototipo non deve essere bello.} Questo vale anche per il
\textbf{prototipo}, su cui si tornerà nelle lezioni di prototipazione: in un
prototipo non servono ``graficoni allucinanti''. Le grafiche bellissime fatte con
l'IA spesso \textbf{peccano di user experience}: sembrano bellissime, ma quando
il giocatore deve leggere la carta non capisce qual è il costo, qual è il
potere, qual è l'effetto.

\paragraph{Beatrice: si perde l'occasione di imparare.} Beatrice aggiunge due
motivi che valgono ``al di là dei nostri credo''. Il primo: chi usa l'IA in
sostituzione (per esempio per scrivere un testo, per accelerare i tempi o per
trovare strategie) \textbf{perde la parte in cui può imparare}. Nel game design
bisogna ``tenere insieme i pezzi'' di un regolamento, ma anche dell'esperienza
dell'utente e di come trasmetterla; i passaggi necessari (creare il regolamento,
creare le meccaniche, sperimentare, raccogliere i feedback) ``si capiscono solo
vivendoli''. E si sta pagando un corso per imparare. Il secondo, sul piano
editoriale: chi presenta a un editore un prodotto che ha come base l'IA rischia
da un lato di non avere imparato abbastanza (e di ritrovarsi ``a piedi'' sul
campo senza IA), dall'altro che l'editore se ne accorga, ed è un tema
\textbf{reputazionale}. Ci sono editori ``più illuminati di altri'', quindi
anche qui c'è un mercato da capire; ma i docenti preferiscono dire che cosa è
più sensato per il mercato come lo conoscono.

\paragraph{Ruggero: non abbiate paura del giudizio.} Ruggero vuole dare uno
``spotlight'' a un punto detto da Daniele e rimasto ``semi-ignorato'': a volte
l'uso dell'IA nasce dalla \textbf{paura del giudizio}, dal voler consegnare
qualcosa di più pulito, più strutturato possibile. ``Vi giuro, non ve ne
preoccupate''. Ruggero detesta i voti e non ragiona per voti; non valuterà
nessuno per grafiche fatte a matita. Insegna da anni, ha avuto studenti con
dislessia e discalculia: refusi ed errori non gli interessano, ``non è un corso
di editing o di proofreading''. Gli interessa che le \textbf{meccaniche siano
coerenti con lo scopo di design}. Immagini fatte a matita o con Paint ``non sono
un problema, anzi, le preferisco'' rispetto all'``ennesima immagine fatta con
ChatGPT o Midjourney''.

\paragraph{Daniele: il pensiero critico, costo e qualità.} L'ultimo intervento
è di Daniele, sul \textbf{pensiero critico}. Ragionando per assurdo (``un
pensiero di Belzebù''), immagina che un domani l'IA venga usata ovunque. I suoi
limiti, secondo lui, resterebbero due:
\begin{description}
  \item[Il costo] Tutti dicono che sta diminuendo e che diminuirà, ma non se ne
        può essere sicuri al 100\%.
  \item[La qualità] Cioè la capacità di leggere l'output. Se si dà un regolamento
        a Claude o a ChatGPT chiedendo quali problemi ha, il modello risponde con
        i problemi che, in percentuale, è più probabile che gli esseri umani
        abbiano già individuato nei dati su cui è stato addestrato. Ne derivano
        due problemi: se il problema del proprio gioco non è mai stato
        individuato prima, il modello non lo vedrà; e viceversa può segnalare
        problemi che il gioco non ha.
\end{description}
Fa l'esempio di Grammarly, che usa nella versione premium per rivedere i testi
in inglese: spesso non accetta i suoi consigli, perché se vuole scrivere un
dialogo di videogioco in \textbf{dialetto Cockney} londinese, lo deve scrivere
in Cockney. Se Grammarly dice che è sbagliato, ``sei tu che non hai capito il
motivo per cui voglio scrivere questo testo in Cockney''. La conclusione: prima
di usare un'IA per un'attività bisogna sapere \textbf{che cosa si vuole
ottenere}, qual è l'obiettivo, e quali \textbf{competenze} servono per valutare
in modo critico se l'output di un qualunque LLM è davvero valido. Daniele non
vuole dilungarsi oltre: è un argomento ``molto caldo'' per lui, che è anche
scrittore e si sente dire da tutti ``ti sostituiranno''.

Subito dopo si presenta finalmente Gianmarco (lo si è visto nella
sezione~\ref{a1:sec:partecipanti}), e poi si fa una pausa di una decina di
minuti. Daniele avverte che nella terza parte lui non ci sarà (alle sette e un
quarto ha un impegno di lavoro), e lo stesso vale per Beatrice: la parte finale
di domande e risposte sarà con il solo Ruggero, ``in ottime mani''.

% ---------------------------------------------------------------------
\section{Le figure professionali dell'industria analogica}
\label{a1:sec:figure}

Dopo la pausa (in cui si scopre che Angelo lavora in un piccolo birrificio in
Brianza e si scherza sul rischio di un \emph{drinking game}, tra la birra e la
vendemmia) la lezione riprende con le slide di Ruggero. ``Oggi è una lezione
introduttiva'': gli argomenti si vedranno ``molto dall'alto'' e verranno
approfonditi nelle lezioni successive. Ma servono già ora, per creare
\textbf{l'inizio di un linguaggio condiviso}: ci sono termini che torneranno
spesso, a cominciare da quelli del mercato e dell'industria del gioco
analogico, un'industria ``molto varia, molto caotica, anche molto mutevole'',
con tanti tipi diversi di giochi e altrettanti tipi diversi di entità che ci
lavorano.

\subsection{Una decina di etichette}

I docenti hanno raccolto ``una decina di etichette'', cioè di \textbf{ruoli} che
nell'industria sono o possono essere presenti. Due avvertenze:
\begin{itemize}
  \item nelle aziende più piccole \textbf{la stessa persona copre spesso più
        ruoli}, per forza di cose. Sarà anche il caso degli studenti, che
        lavoreranno in gruppi di tre e non di venti: è facile che ciascuno finisca
        a fare più cose. I docenti ne terranno conto, e qui entreranno in gioco
        anche competenze diverse dal game design (programmi di grafica,
        illustrazione, impaginazione, layout, \emph{proofreading}\ldots);
  \item Daniele precisa che \textbf{non si chiede agli studenti di diventare}
        illustratori o altro. Ruggero concorda, ma aggiunge che è utile sapere
        anche \emph{che cosa non è il proprio lavoro}. Se in un gruppo nessuno sa
        illustrare, ai docenti va benissimo un prototipo che mostri dove vanno
        le immagini: \textbf{placeholder}, concept, anche semplici box di testo
        con scritto ``qui vorremmo un'immagine di\ldots''. ``È importante anche
        sapere che cosa sapete fare e che cosa non sapete fare''.
\end{itemize}
La tabella~\ref{a1:tab:figure} riassume le figure; le sottosezioni seguenti le
riprendono una per una, con i commenti dei docenti.

\begin{table}[htbp]
\centering
\small
\begin{tabular}{@{}p{3.3cm}p{9.6cm}@{}}
\toprule
\textbf{Figura} & \textbf{Di che cosa si occupa} \\
\midrule
Game designer & crea le regole alla base dell'esperienza: le meccaniche, come
interagiscono tra loro e come chi gioca interagisce con esse \\
Game developer & fa da ponte tra chi fa game design e il publisher: bilancia il
prototipo, cura i modelli matematici, adatta il gioco (e il paragame) alle
necessità del publisher \\
Editor (proofreader) & revisiona e adatta i testi a partire dalla prima stesura
del designer (correzione, linguaggio inclusivo\ldots) \\
Publisher (editore) & pubblica e produce fisicamente il gioco; il suo marchio è
sulla scatola o sulla copertina \\
Sales e marketing & il lato commerciale: vendite (anche B2B) e comunicazione del
prodotto \\
Graphic designer & icone meccaniche, wireframe, UX: il lato funzionale degli
elementi visivi \\
Illustratore & le illustrazioni: il \emph{look and feel}, il world building, il
``vibe'' del gioco \\
Sensitivity reader (cultural consultant) & verifica l'accuratezza e la
rappresentazione di un tema sensibile che il team non conosce \\
Producer & coordina il progetto (quasi project management), parla con le
tipografie, verifica che i materiali siano sostenibili economicamente \\
\bottomrule
\end{tabular}
\caption{Le figure professionali dell'industria del gioco analogico presentate
nella slide di Ruggero, in sintesi.}
\label{a1:tab:figure}
\end{table}

\subsection{Game designer}

Il \textbf{game designer} (o la game designer) è chi si occupa, ``sembra banale
a dirlo'', di \textbf{creare le regole alla base dell'esperienza} che si vuole
far arrivare. Vale per tutti i tipi di giochi, non solo per quelli analogici, e
che cosa faccia concretamente in un progetto o in una serie specifica ``può
mutare tantissimo''. Nell'analogico non verrà chiesto di programmare o di
prototipare in digitale; verrà chiesto di \textbf{creare le meccaniche di
gioco}, di capire \textbf{come le meccaniche si interfacciano tra loro} e
\textbf{come le persone che giocano si interfacceranno con le meccaniche}.

\subsection{Game developer}

Il \textbf{game developer} (o la game developer) sarà l'argomento delle due
lezioni di Marta. La parola \emph{development} è nota soprattutto
dall'ambito videoludico, ma nell'analogico ha un significato molto specifico
``che non c'entra niente con la programmazione'': il game developer è
\textbf{un ponte tra chi fa game design e il publisher}. Di solito si occupa di
\textbf{bilanciare il prototipo}, di curare i \textbf{modelli matematici}, di
adattare il gioco e il \textbf{paragame}\footnote{La trascrizione dice ``il
Wix, il paragame''. Il \emph{paragame} è, in termini generali, tutto ciò che
accompagna il gioco senza esserne il sistema di regole (regolamento, scatola,
materiali, comunicazione); i docenti lo tratteranno in una lezione dedicata, e
qui il termine è solo citato.} alle necessità del publisher. La distinzione
viene ripresa e approfondita nella sezione~\ref{a1:sec:design-development}.

\subsection{Editor}

L'\textbf{editor}, a volte chiamato \textbf{proofreader} a seconda dell'azienda,
si occupa della revisione dei \textbf{testi}: non della stesura iniziale, perché
di norma, data la scala dell'industria, è chi fa game design a scrivere un
\textbf{primo draft}, e l'editor ci passa sopra. Non solo per correggere i
refusi: spesso il suo lavoro è anche \textbf{adattare} i testi, per esempio
renderli \textbf{gender neutral}. Se ne vedranno casi studio più avanti.

\subsection{Publisher}

Il \textbf{publisher}, o \textbf{editore}, è l'entità che nell'ambito
commerciale \textbf{pubblica} giochi, giochi di ruolo o giochi da tavolo. È una
figura ``più solidificata'' nel gioco da tavolo e nel gioco di ruolo da tavolo;
\textbf{nell'ambito pervasivo non esiste proprio in questa forma}, come
mostrerà Beatrice parlando dei casi studio. Se si compra un gioco da tavolo o un
gioco di ruolo in un negozio, il marchio del publisher è di solito sul retro
della copertina o sul fronte della scatola: è l'entità che, di solito,
\textbf{produce fisicamente} il gioco.

\subsection{Sales e marketing}

L'ambito \textbf{sales e marketing}, a volte affidato a un'unica persona e a
volte separato, riguarda il lato commerciale. In alcune aziende è interno al
publisher, e si occupa di parlare con le fabbriche e le tipografie, di ``far
quadrare i conti''. In altre realtà è completamente diverso, specialmente nel
\textbf{B2B} o nei giochi \textbf{didattici} e \textbf{formativi}, di cui
Daniele e Beatrice ``sapranno parlare meglio''. Il tema viene ripreso più avanti
in dettaglio (sezione~\ref{a1:sec:marketing}).

\subsection{Graphic designer e illustratore}

Ruggero presenta insieme \textbf{graphic designer} e \textbf{illustratore}
proprio perché sono ruoli molto diversi che si tende a far coincidere:
\begin{itemize}
  \item chi fa \textbf{graphic design} si occupa delle \textbf{icone meccaniche},
        del \textbf{wireframing}, della \textbf{UX}: del lato \textbf{meccanico e
        funzionale} degli elementi visivi, cioè di tutto ciò che serve
        meccanicamente nel gioco, o anche alla sua ``qualità della vita'';
  \item chi fa \textbf{illustrazione} si occupa dell'\emph{art}, delle
        illustrazioni: copre il \textbf{look and feel}, la parte di
        \textbf{world building} e il ``vibe'' che si dà al gioco.
\end{itemize}

\paragraph{Daniele: che cos'è la user experience.} Daniele interviene sulla
\textbf{user experience}, ``importantissima per me'': ha un background da
\textbf{layout artist} (impaginava magazine), un mestiere in cui la UX conta
moltissimo. Per user experience non si intende solo l'\textbf{interfaccia
grafica}, la UI, che anche nei videogiochi è disegnata da un artista. La UX è
qualcosa che può fare anche il game designer: \textbf{disegnare l'esperienza del
giocatore}, cioè ``il gesto, l'intenzione, l'intuito'' che stanno dietro
l'aspetto visivo dell'interfaccia. Per fare UX non serve saper disegnare.
L'esempio è quello della \textbf{porta}: se si progetta una porta, bisogna
indicare all'utente come si apre, come si prende la maniglia; ``quella è user
experience''. Il design della maniglia in sé, invece, è
\textbf{interfaccia}.

\paragraph{Un aneddoto sulle icone.} Ruggero aggiunge che molte persone che
fanno \textbf{graphic design ben fatto non sanno illustrare}, e spesso vale il
viceversa. Lo si nota quando le icone di un gioco sono bellissime ma non
funzionali: andando a guardare i crediti, di solito si scopre che le icone le ha
fatte la persona che ha fatto le illustrazioni. ``Comunque non andiamo adesso in
questo, perché è un altro vaso di Pandora''.

\subsection{Sensitivity reader}

Il \textbf{sensitivity reader}, o \textbf{cultural consultant} a seconda del
progetto, fa un lavoro di \emph{sensitivity reading} del gioco. È una figura
molto usata nell'ambito \textbf{formativo e didattico}. L'esempio di Ruggero: se
si fa un gioco per un ospedale su una malattia specifica, chi lo progetta non è
un medico, ma il gioco deve essere \textbf{accurato} e \textbf{rappresentativo}
delle dinamiche che quella malattia comporta. Serve una persona che conosca quel
tema, anche se non sa che cosa sia il game design. Sono figure chiamate di
solito \textbf{a progetto}, per giochi specifici, che compensano competenze che
il team non ha ``per forza di cose, perché noi non sappiamo tutto''; e a volte,
per lavoro, si devono fare giochi su temi che non si conoscono personalmente.

\paragraph{Daniele: non è un censore.} Daniele segnala un problema italiano:
alcune persone, anche online, hanno confuso il sensitivity reader con un
\textbf{censore}, ``una persona che fa censura, soprattutto in materia di
\emph{woke}''. ``Assolutamente non è così''. La parola \emph{sensitive} non
significa che il pubblico sia sensibile a un tema, ma che \textbf{l'argomento di
cui tratta il gioco è sensibile}. In un gioco sulla medicina, realistico o no,
non si possono inserire imprecisioni. Per questo esistono i sensitivity reader
``più disparati'': da quelli per la cultura indiana o afroamericana fino a quello
su ``come si allevano gli gnu in Africa''. È una persona competente su un
argomento che aiuta a essere \textbf{veritieri}. Certo, uno degli ambiti
principali in cui si usano è quello culturale, ma ``non fate l'errore di pensare
che siano dei censori''. Daniele ha lavorato con ``centinaia'' di sensitivity
reader e nessuno gli ha mai detto ``non puoi scrivere questa cosa''. Di solito
sono come degli \textbf{editor}: danno opzioni, mostrano i punti problematici,
anche quelli di cui non ci si era nemmeno accorti.

Ruggero conferma con un esempio di quando lavorava in \textbf{Danimarca} con
l'\textbf{OMS}: giochi sul \textbf{diabete di tipo 2} nei bambini, per un target
di ragazzi delle scuole medie. C'era un team dell'OMS specializzato nella
comunicazione sul diabete, per garantire che il gioco fosse accurato e ne desse
una rappresentazione adatta al target. Il sensitivity reader aiuta a rendere il
gioco più accurato e ``più tematicamente coeso''.

\subsection{Producer}

Il \textbf{producer} è una figura a volte interna al publisher. In un progetto
grande è la figura di riferimento, ``quasi'' un project manager. In un progetto
piccolo si occupa spesso anche di parlare con le \textbf{tipografie} e di fare la
\textbf{lista dei materiali}, per assicurarsi che i materiali non siano solo
funzionali al gioco ma anche \textbf{economicamente sostenibili} da produrre:
``magari è un gioco bellissimo, ma poi costa più di quanto potete permettervi, e
nessuno vi comprerà mai un card game a 200 euro''.

\subsection{Perché conoscere queste figure}

Ruggero chiude la slide ``tornando a ombrello'': perché parlarne, se gli studenti
sono lì per fare game design? Non solo per saper dare un \textbf{nome} alle
figure. Non c'è niente di male a creare un gioco solo per sé; ma chi vuole fare
giochi a livello lavorativo, perché li giochino altre persone, ``magari più
persone possibile'', \textbf{si interfaccerà con queste figure}, e sapere che
cosa fanno e come parlarci è parte importante del percorso. Nel corso se ne
toccheranno diverse: il development con Marta, la produzione con Flavio, e
lezioni sulla UX, sul paragame, sul graphic design. ``Abituatevi mentalmente:
non siete in solitaria contro il mondo''.

% ---------------------------------------------------------------------
\section{Marketing e sales}
\label{a1:sec:marketing}

Arrivato all'angolo della slide con marketing e sales, Ruggero passa la parola
ai colleghi: ``io è meglio che non mi fate parlare di marketing e sales'', anche
per esperienze personali.

\subsection{Beatrice: i giochi trattati come prodotti}

Beatrice premette che le due voci sono in parte autoesplicative. Nel mondo
capitalista in cui viviamo, i giochi sono ``strumenti molto potenti per fare
\emph{n} cose'', ma sono trattati come \textbf{prodotti}, e il corso ragionerà
anche in quest'ottica. Nel \textbf{marketing} delle industrie ludiche e
videoludiche ci sono varie professionalità che gli studenti probabilmente già
conoscono: dal \emph{copy} alla grafica, tutto ciò che riguarda il
\textbf{vendere} il gioco o l'attività di gioco.

Il lato \textbf{sales} è forse la distinzione meno nota. Come in tante aziende, è
la parte che si occupa delle vendite \textbf{B2B}, cioè tra aziende: fare
accordi per vendere prodotti ludici a diverse realtà, per esempio ai
\textbf{distributori}, o a una casa editrice che vende alle librerie, che poi
vendono al consumatore finale. Chi si occupa di sales nel mondo ludico capisce
come vendere \textbf{alle aziende}. I due mondi si parlano spesso, e soprattutto
nelle realtà piccole si trovano figure che fanno ``un po' sales, un po'
marketing''.

\subsection{Daniele: che cos'è (e che cosa non è) il marketing}

Daniele il marketing si è trovato ``purtroppo'' a farlo, e lo fa anche per
l'accademia. ``Purtroppo'' perché in Italia è mal compreso: molti lo
identificano con l'\emph{advertisement} online. Non è così: il marketing è una
\textbf{disciplina molto più antica}, e in Italia ``abbiamo molto dimenticato''
come si fa il marketing tradizionale. Alcune cose da ricordare:
\begin{enumerate}
  \item \textbf{Il marketing non dice che cosa pubblicare, ma come pubblicarlo}.
        Esistono, nei videogiochi e nei giochi da tavolo, realtà che fanno
        prodotti \textbf{marketing driven}: l'esempio sono gli
        \textbf{hyper-casual game} per smartphone, quelli con il personaggio che
        va a destra e a sinistra e l'esercito che si allarga. Possono anche
        guadagnare, ma ``sono un po' la morte della creatività''. Il marketing
        ideale \textbf{prende un prodotto pensato da altri e gli dà le maggiori
        chance di essere venduto}: cerca il pubblico ideale per il prodotto che
        si ha, senza doverlo cambiare. Anzi, il mondo videoludico, e ludico in
        generale, ``è lastricato'' di errori di dipartimenti marketing che hanno
        cambiato completamente non il gioco, ma il modo in cui l'hanno
        pubblicizzato, e il gioco è fallito.
  \item \textbf{Usare strumenti del marketing non significa fare marketing}. Game
        design e game development a volte usano strumenti che vengono dal
        marketing (gli studenti ne vedranno con Beatrice, con Daniele e con
        Marta).
  \item \textbf{Il marketing è un gerundio}: si fa. È un'azione che si compie per
        comunicare; la comunicazione è attiva, e il marketing è l'insieme delle
        attività con cui si comunica il proprio gioco, dal \emph{copy} ai social.
        Può includere il digital marketing, ``ma non è solo quello''.
  \item \textbf{Il marketing esiste anche se si crede che non esista}. In Italia
        molti non vogliono ``fare marketing'' perché lo associano al lucro. Ma
        il marketing è \textbf{lo studio di come gli esseri umani comprano e
        vendono}, ed è innato in noi: ci sono logiche di cui siamo vittime, e
        ``anche se sei il bastian contrario, l'anticapitalista'', magari si
        compra comunque su Amazon perché fa comodo. Sono meccanismi che
        esistono, ``affascinanti anche, a volte''.
\end{enumerate}
Gli studenti non dovranno fare marketing, ma dovranno ricordare alcuni di questi
concetti, soprattutto nelle lezioni con Marta, che tratta argomenti legati al
marketing. E, chiude Daniele, come game designer e game developer ``per lo più
speriamo che il marketing non ci rovini i giochi'', soprattutto nel videoludico.

% ---------------------------------------------------------------------
\section{Game design e game development}
\label{a1:sec:design-development}

Ruggero torna sulla distinzione più importante della slide, ``proseguendo con
questa strada'':
\begin{description}
  \item[Game design] è chi \textbf{progetta l'esperienza} tramite la
        \textbf{creazione delle meccaniche} di gioco.
  \item[Game development] è chi fa \textbf{mediazione} tra la persona che scrive
        il game design e la \textbf{casa editrice}.
\end{description}

\subsection{I red flag dei designer}

Ruggero lo ripete perché lo vede ogni giorno nel suo lavoro: tra giochi
pubblicati e prototipi, ne valuta ``quasi uno al giorno, se faccio un conto
annuale'', e solo una minima percentuale arriva alla pubblicazione. Molte volte
i \textbf{red flag} vengono proprio da designer che non hanno presente questa
definizione, e che cominciano a mettere \textbf{clausole sulla non modifica del
gioco}: il gioco deve restare com'è. ``Questa cosa è impossibile con un
publisher, spesso''. Il publisher, per forza di cose, potrà chiedere di cambiare
la grafica, di reimpaginare il regolamento, di cambiare i materiali perché
troppo costosi: chiederà dei \textbf{compromessi}. E la persona che, all'interno
del publisher, fa questo lavoro è il \textbf{developer}.

Sta al designer capire \textbf{su che cosa del proprio gioco si può trovare un
compromesso e su che cosa no}. Publisher diversi lavorano in modo diverso, e il
designer ha sempre l'autorità di dire ``no, non voglio pubblicare questo gioco
se queste sono le modifiche''. Ma deve sapere chi fa il lavoro di development,
perché se il suo obiettivo è la \textbf{pubblicazione tradizionale} si
interfaccerà con questa figura.

\subsection{Daniele: una distinzione che manca al videogioco}

Daniele sottolinea che questa distinzione è \textbf{tipica del mondo analogico}
ed è ``molto interessante''. In Italia la parola \emph{developer} è diventata
negli anni quasi un sinonimo di programmatore di videogiochi, o di chiunque
lavori nel videoludico. Secondo lui c'entra anche il fatto che qualche anno fa
il sito \textbf{Gamasutra} ha cambiato nome in \textbf{Game Developer}, e la
parola è diventata un ombrello per tutte le professioni del videogioco. ``Mi fa
un po' ribrezzo'': il game developer dovrebbe essere un ruolo di design, un
designer che, come nell'analogico, fa \textbf{mediazione}. Nel videogioco i
publisher hanno in effetti i loro developer interni, che adattano il gioco, per
esempio alle console e ai loro kit di sviluppo. Ma la distinzione dovrebbe
esistere anche nel videogioco, perché è ``difficilissimo'' avere la mentalità
per fare le due cose insieme: chi fa design ha un focus che raramente si può
sovrapporre a quello di chi fa mediazione.

A Daniele piace la parola \emph{develop}, \textbf{sviluppare}, nel senso di
prendere qualcosa che è stato progettato e \textbf{renderlo possibile nel mondo
reale}. Ci sono progetti che hanno bisogno di uno sviluppo che ``polishi''
l'esperienza: la semplifichi, la migliori, la renda comprensibile e la concluda,
ma soprattutto la renda \textbf{stampabile} e \textbf{fruibile} per l'utente. La
metafora crossmediale è quella della \textbf{fotografia}: chi sviluppa le
fotografie non è chi scatta la foto, ma chi la rende visibile, anche
fisicamente. ``È un po' retro come metafora, ma ci capiamo''. ``È vero, però è
giusto'', conferma Ruggero.

% ---------------------------------------------------------------------
\section{Il mercato: Italia ed estero}
\label{a1:sec:mercato}

Il corso è in italiano e si svolge in Italia, ma bisogna conoscere anche il
\textbf{mercato estero}: molti giochi progettati da autori italiani vengono
pubblicati all'estero, e conoscere il mercato estero aiuta a capire come il
proprio gioco può arrivarci, ``o viceversa, come non può arrivarci''. La
tabella~\ref{a1:tab:mercato} riassume il confronto della slide.

\begin{table}[htbp]
\centering
\small
\begin{tabular}{@{}p{6.4cm}p{6.4cm}@{}}
\toprule
\textbf{Italia} & \textbf{Estero (soprattutto Europa)} \\
\midrule
il gioco è ancora visto come un'attività \textbf{infantile}, anche se la
normalizzazione è in crescita &
forte tradizione di gioco anche in un contesto \textbf{familiare adulto},
soprattutto ``dalla Germania in su'' \\
risposta \textbf{conservativa} sul tema ludico &
\textbf{premi di prestigio} riconosciuti in tutto il mondo (Spiel des Jahres,
As d'Or) \\
forte focus sulla \textbf{localizzazione} di giochi stranieri &
si cercano soprattutto giochi \textbf{originali} \\
investimento \textbf{molto basso} sulla creazione di giochi originali interni;
pochissimi designer assunti &
\textbf{developer interni} pagati con un contratto (Francia, Germania,
Danimarca\ldots) \\
\bottomrule
\end{tabular}
\caption{Il mercato del gioco analogico in Italia e all'estero, secondo la
slide di Ruggero e i commenti di Daniele.}
\label{a1:tab:mercato}
\end{table}

\subsection{L'Italia}

\begin{itemize}
  \item In Italia il gioco è ancora visto come un'\textbf{attività infantile},
        anche se la \textbf{normalizzazione} è in crescita. C'è ancora una forte
        risposta \textbf{conservativa} sul tema ludico.
  \item C'è un forte focus sulla \textbf{localizzazione} dei giochi stranieri: in
        una qualunque ludoteca o negozio, molti dei giochi da tavolo e di ruolo
        sono di publisher o designer esteri, pubblicati in italiano da publisher
        italiani. (La localizzazione, cioè portare un titolo dall'estero in
        Italia, esiste come funzione all'interno del publisher; Ruggero non l'ha
        messa nella slide delle figure perché non è il focus del corso.)
  \item C'è, ``purtroppo, preparatevi psicologicamente'', un \textbf{investimento
        molto basso} sulla creazione di giochi originali interni alla casa
        editrice. Pochissime case editrici italiane hanno designer interni pagati
        a tempo pieno con contratto a tempo indeterminato. Daniele lo dice senza
        mezzi termini: ``il mercato è molto piccolo, conosciamo praticamente
        tutto il settore'', e le persone che in Italia fanno game design
        analogico a tempo pieno con un contratto a tempo indeterminato ``sono
        meno delle dita delle mie due mani''. Di solito la formula è quella del
        \textbf{freelance}, del contratto, delle \textbf{royalty}. Chi va da un
        publisher aspettandosi di essere assunto a tempo pieno come game
        designer deve sapere che ``non siamo purtroppo ancora in quel tipo di
        situazione''.
  \item C'è però un contraltare: in Italia, soprattutto nel gioco di ruolo ma
        anche nel gioco da tavolo, c'è stata molta \textbf{imprenditoria}, anche
        grazie a \textbf{Kickstarter}, che ha aperto le porte ad altri sistemi
        di \textbf{crowdfunding}. È un altro argomento, perché richiede un team,
        un investimento e così via.
  \item E gli autori non mancano, anzi: le occasioni di incontro tra autori
        vedono ogni anno un aumento dei possibili autori\footnote{La trascrizione
        dice ``i DAG a Parma''; il riferimento è probabilmente a IDEAG, il
        circuito di incontri tra autori di giochi italiani, ma non è
        ricostruibile con certezza.}. Bisogna ``sgomitare'' tra tanti autori,
        che magari arrivano da esperienze molto diverse.
\end{itemize}
Ruggero precisa: l'investimento basso è ``dall'alto''. Le persone che fanno e
vogliono fare giochi sono tantissime e in aumento; scarseggia l'investimento,
sia dei publisher disposti a pagare stipendi sia dei \textbf{bandi} pubblici. I
bandi esistono (Beatrice ne parlerà per i giochi pervasivi, che spesso si
muovono in quella direzione, e Ruggero cita qualcuno che sta facendo un card
game con un bando italiano, ``una cosa assurda, ma lo sta facendo''), ma sono
una rarità, soprattutto rispetto all'estero.

\subsection{L'estero}

\begin{itemize}
  \item L'estero (dipende dalla nazione, ma in generale il resto d'Europa) ha una
        forte \textbf{tradizione di gioco} anche in un contesto \textbf{familiare
        adulto}, soprattutto ``dalla Germania in su''.
  \item Ha \textbf{premi di prestigio} riconosciuti in tutto il mondo, come lo
        \textbf{Spiel des Jahres} tedesco e l'\textbf{As d'Or} francese\footnote{La
        trascrizione dice ``lo Spield AR, la SDOR''.}.
  \item Cerca soprattutto \textbf{giochi originali}. Lo si vede dal fatto che
        molti autori italiani di giochi da tavolo pubblicano in Germania o in
        Francia e \textbf{non vengono pubblicati in Italia}, ``che è tutto
        dire''.
  \item Ha \textbf{figure di development interne}: non è raro vedere publisher
        francesi, tedeschi, danesi con developer pagati internamente, a
        contratto.
\end{itemize}
``Poi la situazione è ancora più drammatica'', ammette Ruggero, ``ma per ora
vogliamo tenere il morale un po' alto. È l'introduzione. Ci sarà tempo di andare
nella tristezza''.

Prima di proseguire, Ruggero insiste sulle domande: con lui si può interrompere
anche mentre parla, alzando la mano con le reaction o semplicemente togliendo il
microfono e chiedendo. ``Non mi offenderò mai, anzi, sarò felice di vedervelo
fare e mi ricorderò positivamente di voi''.

% ---------------------------------------------------------------------
\section{Perché serve una terminologia comune}
\label{a1:sec:terminologia}

La serie di slide successiva è ``una piccolissima parte di fondamenti di
design'' (``non spaventatevi, vi vedo che vi spaventate un po'\,''). Anche qui
molti temi verranno ripresi in dettaglio, ma l'idea è arrivare alla prossima
lezione con almeno un \textbf{glossario}, un \textbf{linguaggio condiviso}: se
qualcuno dice \emph{game developer}, tutti sanno di che cosa sta parlando.

Ruggero parte da una domanda, e precisa che non è retorica: \textbf{perché ci
serve una terminologia comune?} Gli studenti rispondono:
\begin{itemize}
  \item Gianmarco, ``molto pragmatico'': perché bisogna lavorare in team e
        capirsi. Vero, nel corso e anche fuori;
  \item un'altra studentessa: per la \textbf{comunicazione}, tra di noi e con
        tutti quelli che sono interessati nel settore. Vero anche questo:
        conoscere i termini comuni permette di parlare con chi fa game design o
        development nel dettaglio, nel tecnico, nello specifico.
\end{itemize}
Ruggero aggiunge altre due ragioni.

\paragraph{Saper comunicare il proprio gioco.} La prima riguarda la
\textbf{comunicazione del gioco} verso l'esterno. ``E lo dico non tirando acqua
al mio mulino, ma anzi togliendola'': \textbf{non è spesso il buon game design
a far diffondere un gioco}. ``Vorrei che lo fosse, sarebbe bellissimo, ma
purtroppo non è così''. Ci sono giochi bellissimi che non sanno comunicare, che
non sanno dire che cosa fanno, e che quindi è difficile, da fuori, trovare
invitanti; e ci sono giochi mediocri dal punto di vista del design che si
comunicano molto bene, e che quindi conoscono in molti. Saper parlare di quello
che si fa è utile \textbf{a sé stessi}, in primis.

\paragraph{Documentare per sé stessi.} La seconda è la \textbf{documentazione}:
documenti di design, di brainstorming e così via. ``Sembrerà una battuta, ma
provate a recuperare un documento di brainstorming una settimana dopo, se non
avete scritto bene le cose che avete scritto''. Chi l'ha fatto sa già di che
cosa si parla; ``chi non l'ha ancora vissuto, forse lo vivrà''.

% ---------------------------------------------------------------------
\section{Game e play}
\label{a1:sec:game-play}

La \textbf{prima, importantissima divisione terminologica} è quella tra
\emph{game} e \emph{play}. In italiano c'è un problema: \textbf{gioco} e
\textbf{giocare} hanno la stessa radice, mentre in altre lingue, tra cui
l'inglese, sono due parole strutturalmente diverse:
\begin{description}
  \item[Game] il \textbf{gioco come artefatto}: il gioco come ecosistema di
        meccaniche, come struttura di regole, come prodotto (``tra
        virgolette''), come oggetto fisico.
  \item[Play] l'\textbf{atto}, il verbo: le persone che giocano, l'approccio
        fisico all'attività del gioco.
\end{description}
È importante perché gli studenti faranno \textbf{game design}; esistono anche
ambiti di \emph{play design}, ma qui ``creerete giochi, non giocattoli''.

\subsection{Giochi e giocattoli}

Qual è la differenza? Uno studente risponde: \textbf{la presenza di regole}.
Esatto: un gioco ha delle regole, un'attività di \emph{play} non
necessariamente. Con un giocattolo si può fare \emph{play}: c'è chi gioca con i
Lego, chi gioca a guardia e ladri o a nascondino al parco, e magari alcune di
quelle regole sono ``più o meno semi-solidificate''. Ma il \emph{play} permette
spesso un approccio non strettamente legato alle regole.

Viceversa, \textbf{un game può esistere anche se nessuno lo sta giocando}. Se si
prepara su un tavolo una scacchiera con i pezzi degli scacchi, chi passa e la
vede riconosce il gioco degli scacchi, anche se nessuno sta giocando. Se invece
si vedono dei bambini correre in un parco giochi, quei bambini \emph{stanno
giocando}, ma non stanno giocando \emph{a un gioco} specifico. ``È importante
separare questa cosa''; ci si tornerà, ``ma per ora abituatevi a questa
differenza''.

\paragraph{Approfondimento: giocattoli, giochi e regole.} La stessa distinzione
è il punto di partenza del capitolo che Jesse Schell, in \emph{The Art of Game
Design}, dedica alla definizione di gioco (\emph{The Experience Rises Out of a
Game}). Schell nota che usiamo perfino un linguaggio diverso: ``un gioco è
qualcosa a cui si gioca'' (\emph{a game is something you play}), ``un giocattolo
è qualcosa con cui si gioca'' (\emph{a toy is something you play with}); e
arriva a definire il giocattolo come ``un oggetto con cui si gioca''. Analizzando
poi la definizione di gioco di Avedon e Sutton-Smith, ne ricava che ``i giochi
hanno regole, i giocattoli no'', e che le regole sono uno degli aspetti
\emph{definitori} del gioco, insieme alla volontarietà, agli obiettivi, al
conflitto e alla possibilità di vincere o perdere. Schell osserva anche che
molti giochi sono costruiti \textbf{sopra} un giocattolo (``una palla è un
giocattolo, ma il baseball è un gioco''), e tra i suoi consigli di
prototipazione c'è proprio \emph{Build the Toy First}: costruire prima il
giocattolo, cioè qualcosa che sia già piacevole da manipolare, e poi il gioco
intorno. È una lettura utile della distinzione di Ruggero: il \emph{play} può
fare a meno delle regole, il \emph{game} no.

% ---------------------------------------------------------------------
\section{Il framework MDA (e il DPE)}
\label{a1:sec:mda}

\subsection{Meccaniche, dinamiche, estetiche}

Un'altra terminologia fondamentale è quella del framework \textbf{MDA}
(\emph{Mechanics, Dynamics, Aesthetics}; in italiano a volte ``MDE''). Viene da
un \textbf{paper} nato nell'ambito del videogioco\footnote{È il paper di Robin
Hunicke, Marc LeBlanc e Robert Zubek, \emph{MDA: A Formal Approach to Game
Design and Game Research}, del 2004. Daniele aggiunge il paper alla cartella
condivisa durante la lezione.}, che ha ``solidificato'' queste terminologie, e
divide un gioco in tre aree:
\begin{description}
  \item[Meccaniche] le \textbf{regole per come sono scritte}. ``Ogni volta che
        pescate un 2 fate dei punti'': quella è una meccanica.
  \item[Dinamiche] l'\textbf{applicazione in tempo reale} delle meccaniche. La
        dinamica è il giocatore che, pescando una carta, \textbf{spera} che esca
        un 2. Quel processo non è una regola: ``il gioco non vi dice che dovete
        sperare che esca un 2''.
  \item[Estetiche] l'insieme di quello che il gioco vuole comunicare,
        l'\textbf{esperienziale}.
\end{description}

L'esempio di Ruggero mette insieme le tre aree. In un gioco in cui pescare un 2
fa fare punti (meccanica) e in cui quindi si spera che esca un 2 (dinamica),
\textbf{che estetica si ottiene se nel mazzo di 100 carte c'è un solo 2? E se ce
ne sono 50?} L'emozione di aver pescato \emph{l'unico} 2 del mazzo è
completamente diversa da quella di pescarne uno sapendo che ce ne sono
cinquanta e che è molto probabile trovarlo. Le regole sono quasi identiche; cambia
l'esperienza.

\subsection{Il falso amico: ``estetica''}

Daniele avverte di non sottovalutare la parola \textbf{estetica}, che in italiano
è un \textbf{falso amico}: sembra indicare la parte grafica, e invece non lo è.
Si incontreranno designer che hanno studiato MDA e altri che non l'hanno
studiato, e lo si capirà proprio da questo: i secondi diranno che le
\emph{aesthetics} sono la grafica. Le \emph{aesthetics} sono la parte
\textbf{esperienziale}: come nell'esempio di Ruggero, quello che lascia al
giocatore aver pescato l'unico 2 del mazzo.

\subsection{Il DPE: design, play, experience}

Per gli italiani, aggiunge Daniele, è più facile capire l'MDA attraverso un
altro paper, quello del \textbf{DPE}\footnote{È il framework \emph{Design, Play,
Experience} di Brian Winn, proposto come estensione dell'MDA per i giochi
educativi (\emph{serious games}). Anche questo paper viene aggiunto da Daniele
alla cartella condivisa.}, che di fatto fa un parallelismo assegnando
\textbf{design, play, experience}:
\begin{itemize}
  \item le \textbf{meccaniche} sono il lato \textbf{design};
  \item le \textbf{dinamiche} sono il lato \textbf{play};
  \item le \textbf{estetiche} sono l'\textbf{experience}, l'esperienza che riceve
        chi gioca.
\end{itemize}
``Se vivete le \emph{aesthetics} come l'esperienza che vive chi sta giocando,
state a postissimo''. In questo momento Beatrice lascia la lezione (``ciao,
Bea''). I paper di MDA e DPE sono in inglese; chi vuole può dargli una lettura.

\begin{figure}[htbp]
\centering
\begin{tikzpicture}[
    font=\small,
    blocco/.style={draw, thick, rounded corners=3pt, fill=gray!12,
                   minimum width=3.2cm, minimum height=1.0cm, align=center},
    dpe/.style={font=\footnotesize\itshape, align=center},
    esempio/.style={draw, rounded corners=2pt, text width=3.3cm,
                    align=left, font=\footnotesize, inner sep=4pt},
    freccia/.style={-{Stealth[length=2.5mm]}, thick}
  ]
  \node[blocco] (m) at (0,0)   {\textbf{Meccaniche}\\ \footnotesize le regole scritte};
  \node[blocco] (d) at (4.3,0) {\textbf{Dinamiche}\\ \footnotesize le regole in azione};
  \node[blocco] (a) at (8.6,0) {\textbf{Estetiche}\\ \footnotesize l'esperienza};
  \draw[freccia] (m) -- (d);
  \draw[freccia] (d) -- (a);
  \node[dpe, above=0.15cm of m] {DPE: \textbf{design}};
  \node[dpe, above=0.15cm of d] {DPE: \textbf{play}};
  \node[dpe, above=0.15cm of a] {DPE: \textbf{experience}};
  \node[esempio, below=0.45cm of m] (em)
    {``Ogni volta che peschi un 2 fai punti''};
  \node[esempio, below=0.45cm of d] (ed)
    {pescando, il giocatore \emph{spera} che esca un 2 (nessuna regola glielo
     chiede)};
  \node[esempio, below=0.45cm of a] (ea)
    {un solo 2 su 100 carte: pescarlo è un evento;\\
     50 su 100: è quasi scontato};
  \draw[dashed] (m) -- (em);
  \draw[dashed] (d) -- (ed);
  \draw[dashed] (a) -- (ea);
  % prospettive
  \draw[freccia, gray] ([yshift=-0.3cm]em.south west) --
    node[below, font=\footnotesize\itshape, text=black]
    {il designer agisce sulle meccaniche}
    ([yshift=-0.3cm]ea.south east);
  \draw[freccia, gray] ([yshift=-1.2cm]ea.south east) --
    node[below, font=\footnotesize\itshape, text=black]
    {chi gioca incontra prima l'esperienza}
    ([yshift=-1.2cm]em.south west);
\end{tikzpicture}
\caption{Il framework MDA con l'esempio del ``2'' usato da Ruggero, e la
corrispondenza con il DPE suggerita da Daniele. Le due frecce in basso
riassumono la lettura classica del paper: il designer costruisce le meccaniche
per ottenere un'esperienza, mentre il giocatore parte dall'esperienza.}
\label{a1:fig:mda}
\end{figure}

\paragraph{Approfondimento: le due direzioni dell'MDA e le otto estetiche.} Le
frecce della figura~\ref{a1:fig:mda} riprendono un'idea centrale del paper
originale, che la lezione non esplicita: designer e giocatore guardano il gioco
da lati opposti. In \emph{Challenges for Game Designers}, Brenda Brathwaite e
Ian Schreiber la riassumono così: i designer si chiedono quale estetica
vogliono ottenere, definiscono le dinamiche che porterebbero a quella
sensazione, e poi creano le meccaniche che producono quelle dinamiche. Lo stesso
libro conferma l'avvertenza di Daniele sul falso amico: ``le \emph{aesthetics}
non si riferiscono all'aspetto del gioco, ma alla risposta emotiva'' che il
designer vuole suscitare. E offre esempi di dinamiche nel senso di Ruggero: in
uno sparatutto i punti di \emph{spawn} sono una meccanica, mentre appostarsi
accanto a uno di essi per colpire chi riappare (lo \emph{spawn camping}) è una
dinamica; negli scacchi i movimenti dei pezzi sono meccaniche, mentre le
aperture codificate sono dinamiche. Per dare un nome alle estetiche, il paper
MDA propone una tassonomia, poi detta degli ``otto tipi di \emph{fun}'', che
Schell riprende in \emph{The Art of Game Design} come \emph{LeBlanc's Taxonomy of
Game Pleasures}: \textbf{sensazione}, \textbf{fantasia}, \textbf{narrativa},
\textbf{sfida}, \textbf{compagnia} (\emph{fellowship}), \textbf{scoperta},
\textbf{espressione} e \textbf{sottomissione} (il piacere di entrare nel
``cerchio magico'' e lasciarsi andare a un passatempo). È un primo esempio di
quello che i docenti chiedono con le parole vietate della prossima sezione:
invece di dire che un gioco è ``divertente'', dire \emph{quale} esperienza
offre.

% ---------------------------------------------------------------------
\section{Le parole vietate}
\label{a1:sec:parole-vietate}

``Attenzione, adesso qua triggeriamo delle persone''. Ruggero annuncia una regola
che ``tornerà'': nel corso ci sono delle \textbf{parole vietate}. ``Non perché
siamo dei censori''; il perché si vedrà in una lezione apposita, perché
spiegarlo adesso richiederebbe un'ora. Ma Daniele vuole spiegarne ``un
minimo''.

\subsection{La regola}

In \textbf{nessun materiale scritto} consegnato ai docenti, task compresi, sono
accettate queste parole, né in inglese né in italiano:
\begin{center}
\begin{tabular}{@{}ll@{}}
\toprule
\textbf{Inglese} & \textbf{Italiano} \\
\midrule
\emph{fun} & divertimento \\
\emph{immersion} & immersione \\
\emph{entertainment} & intrattenimento \\
\bottomrule
\end{tabular}
\end{center}
``Non li voglio neanche vedere''. Daniele lo ripete con un paragone: ``se mi fate
i disegni con Paint, benissimo, top, fantastico, siamo ancora best friend. Se mi
scrivete \emph{divertimento}, io sbrocco''. È una \textbf{regola dell'accademia}:
negli stessi corsi triennali di narrative design e game design le parole sono
vietate agli studenti per tre anni; ``per voi solo per tre mesi''. Daniele
aggiungerebbe volentieri un quarto termine, \textbf{gamification}.

\subsection{Perché}

Il motivo approfondito lo darà Ruggero nella lezione apposita. Daniele dà il
suo. Ruggero e Diego Ricchiuti sono molto contrari a \emph{fun} e a
\emph{immersion}; lui ha aggiunto \emph{intrattenimento}. Sono \textbf{parole
troppo passe-partout e troppo ambigue}:
\begin{itemize}
  \item \textbf{intrattenimento} in Italia è stata usata per anni come un modo
        per \textbf{sminuire}: ``ma sì, tanto è un gioco solo di
        intrattenimento, è un romanzo solo di intrattenimento''. Ma anche se un
        prodotto è ``solo di intrattenimento'', non vuol dire che non ci si debba
        lavorare con attenzione. E poi intrattenimento ``non significa niente'':
        ci sono tanti modi in cui una cosa può intrattenere;
  \item lo stesso vale per \textbf{fun}, e per \textbf{immersion}, che ha un
        problema in più: è una parola passe-partout per cui si può dire che un
        gioco è ``immersivo'' nel senso generale del termine, ma poi non si sa
        davvero che cosa voglia dire, \emph{immersivo in che senso}.
\end{itemize}

L'obiettivo didattico è che gli studenti \textbf{vadano oltre queste parole}:
``un po' come quando alle medie vi dicono smettetela di dire che tutto è bello,
perché ci sono vari gradi di bello, ci sono altri modi di dire bello''. Bisogna
\textbf{decostruirle}: quando si sta per usarle, chiedersi \textbf{che cosa si
intende} con quella parola e cercarne una \textbf{più precisa}.

Ruggero chiede, in uno dei ``rari casi'', di \textbf{non fare domande adesso} su
questa regola: ogni anno arrivano domande e commenti, e ci saranno lezioni
apposite di tre ore, con strumenti per capire il perché. Ma da subito gli
studenti devono abituarsi a evitare questi termini, sia perché si spera che
abbiano termini più precisi, sia perché queste parole ``portano con sé dei
problemi agganciati'' che si vedranno. ``Non spaventatevi, non verrete sparati in
un vicolo cieco di notte''.

\paragraph{Una piccola gamification.} Il divieto va visto come uno
\textbf{stimolo}, ``un po' come se stessimo gamificando la vostra esperienza'',
una sorta di user experience delle lezioni. Imponendosi di non usare queste
parole, si noterà come migliora la capacità di \textbf{descrivere i giochi} e le
proprie esperienze ludiche, e di essere più precisi. È un invito a ``obbligarsi a
essere più precisi e ad avere una mentalità più critica''. La lezione di
Ruggero sull'\textbf{involvement}, promette Daniele, sarà ``fenomenale'' e
``vi aprirà la testa''.

\paragraph{Approfondimento: anche i libri faticano a definire il ``fun''.} La
difficoltà che il divieto vuole far emergere è visibile anche nei testi di game
design. Schell, in \emph{The Art of Game Design}, prova a definire il \emph{fun}
e propone ``piacere con sorprese'' (\emph{pleasure with surprises}), ammettendo
subito che è strano poter usare una parola per tutta la vita, sapendo con
certezza che cosa significa, senza riuscire a esprimerla chiaramente quando
qualcuno lo chiede; e invita a mettere alla prova la definizione con
controesempi. La sua lente sul \emph{fun} (\emph{The Lens of Fun}) si limita,
non a caso, a chiedere quali parti del gioco siano \emph{fun} e \emph{perché}.
Anche \emph{immersion} ha molte letture: in \emph{21st Century Game Design},
Bateman e Boon la descrivono, tra l'altro, come la capacità di ``perdersi''
nel mondo artificiale presentato dal gioco, legata alla sospensione
dell'incredulità e ai giochi di mimesi. Sono significati diversi sotto la stessa
parola: proprio il motivo per cui, in un documento di design, conviene dire
quale si intende.

% ---------------------------------------------------------------------
\section{Chi crea, chi gioca, chi facilita}
\label{a1:sec:ruoli-esperienza}

Un'altra distinzione importante è quella tra \textbf{chi crea}, \textbf{chi
gioca} e \textbf{chi facilita}. Spesso, quando si fa un gioco, lo si fa provare
di persona a chi gioca. Ma ``la norma dei giochi pubblicati non funziona così'':
quando si pubblica un gioco non si vede chi lo compra e chi lo gioca, e non si
sa se qualcuno, a casa propria con i propri amici, segue le regole alla lettera e
interpreta correttamente quello che il designer aveva descritto. Quindi già in
fase di progettazione bisogna tenere a mente \textbf{chi sta facendo game design
e chi sta giocando}.

E nei giochi a evento, nei giochi pervasivi, nei giochi di ruolo da tavolo e dal
vivo (LARP) c'è un terzo pezzo del puzzle, ``importantissimo'': \textbf{chi
facilita}. La domanda da farsi è: \textbf{il regolamento è scritto per chi gioca
o per chi deve facilitare il gioco?} Bisogna abituarsi a pensare a tutte le
persone coinvolte nell'esperienza.

% ---------------------------------------------------------------------
\section{Playtest e iterazione}
\label{a1:sec:playtest}

\subsection{Nessun gioco nasce perfetto}

L'ultima voce del glossario è ``un'altra cosa importantissima'':
\textbf{playtest} e \textbf{iterazione}. Il corso è \textbf{pragmatico}:
``specialmente da parte mia'', dice Ruggero, non si vedranno teorie ``messe lì
per fare accademia sul game design nelle nostre torri d'avorio''. Ci saranno
teorie e framework, ma tutti estremamente pragmatici, perché l'obiettivo non è
solo capire i giochi, ma \textbf{saperli fare e iterare}.

Playtest e iterazione sono \textbf{inevitabili} nel game design, perché
\textbf{nessun gioco nasce perfetto}, soprattutto perché le persone sono
diverse: un gioco perfetto per una persona può essere incomprensibile per
un'altra. Chi progetta giochi per altre persone deve capire, dopo averli fatti
provare, \textbf{che cosa cambiare, che cosa sistemare, che cosa aggiungere, che
cosa togliere}. Ci saranno lezioni su playtest, prototipazione, raccolta dei
feedback, metodi etnografici. Ma bisogna abituarsi già ora all'idea che
\textbf{non esiste la consegna definitiva}. A un certo punto i task vanno
consegnati per forza di cose, ma, come diceva Daniele all'inizio, alla fine del
corso gli studenti avranno un \textbf{prototipo}, ``che sarà un prototipo'', e
gli strumenti per portare avanti il lavoro anche dopo.

\subsection{L'aiuto dopo il corso}

Non solo gli strumenti: anche l'aiuto dei docenti. Daniele racconta di aver
portato a un editore, a Play\footnote{Probabilmente \emph{Play -- Festival del
Gioco} di Modena; la trascrizione non permette di esserne certi.}, uno dei
progetti realizzati l'anno precedente da due studenti, un gioco di ruolo in
solitario (``un \emph{solo game} RPG molto, molto bello''). Se gli studenti
vorranno trasformare il proprio lavoro in un prototipo vero e proprio, possono
contattarlo: lo fa ``per lavoro, ma anche per passione''. Gestisce anche una
community per aiutare i designer ``un po' depressi'' a rimettersi a progettare.

Quindi: ``non sentitevi persi se durante la fase di iterazione vi sentite
persi, è normale''.

\subsection{Il dizionario comune e il production dictionary}

Daniele riprende un punto detto ``benissimo'' da Ruggero e che secondo lui non è
stato sottolineato abbastanza: il \textbf{dizionario comune} serve anche a
\textbf{non perdersi} mentre si itera. Se da un'iterazione all'altra i
documenti usano termini completamente diversi, e non c'è un dizionario di base,
ci si perde.

Nel videogioco, una delle primissime cose che un designer fa quando entra in
un'azienda è il \textbf{production dictionary}: si annota, in un file privato,
tutte le parole che usano i colleghi, per allinearsi al dizionario comune
dell'azienda, perché lì su un prototipo possono lavorare anche cento persone.
Nel gioco da tavolo non si fa molto, perché il mondo dell'analogico è
``talmente più piccolo''; ma serve comunque imparare un \textbf{dizionario
minimo}, dai docenti e dai professionisti. E il corso ne darà una parte,
entrando nel dettaglio di alcuni termini e framework.

\paragraph{L'esempio del feedback.} Perché serva, lo mostra un esempio. Si
testa un gioco e ci si annota come feedback: ``questo gioco deve essere più
divertente''. Che cosa vuol dire? Come lo si migliora? Quella frase non dà
elementi utili. Se invece si scrive che deve essere \textbf{più adrenalinico},
\textbf{più frenetico}, \textbf{più rilassato}, \textbf{più cozy}, o
\textbf{più bilanciato matematicamente}, si sa dove andare a lavorare.

\paragraph{Approfondimento: la regola del ciclo.} Il principio ``nessun gioco
nasce perfetto'' ha un nome preciso in \emph{The Art of Game Design}: è la
\textbf{Rule of the Loop}, la regola del ciclo. ``Più volte provi e migliori il
tuo design, migliore sarà il tuo gioco''. Schell precisa che non è una lente,
cioè un punto di vista, ma ``una verità assoluta'', senza eccezioni: ci saranno
momenti in cui si cercherà di convincersi che ``questa volta il design è così
buono che non serve testarlo'', e ogni volta se ne pagherà il prezzo. Ed è
interessante la ragione per cui la regola pesa di più sul digitale: secondo
Schell il tempo e il denaro necessari per provare e correggere un sistema sono
molto maggiori nei videogiochi che nei giochi tradizionali. Per un gioco di
carte o da tavolo, con il tempo per provarlo e cambiarlo più e più volte, il
consiglio è semplicemente di farlo. Il gioco analogico è, da questo punto di
vista, la palestra ideale dell'iterazione.

% ---------------------------------------------------------------------
\section{Tabletop Simulator}
\label{a1:sec:tts}

\subsection{Il software di settore}

``Last but not least'': il corso usa un software, \textbf{Tabletop Simulator},
che si compra su \textbf{Steam} e costa tra i \textbf{10 e i 20 euro} (``aspettate
che vada in sconto, va in sconto ogni due per tre''). Daniele spiega perché:
\begin{itemize}
  \item è \textbf{il software di settore} del gioco da tavolo, quello che i
        publisher usano per i prodotti digitali, interni ed esterni;
  \item essendo su Steam, permette di fare \textbf{playtest} e di giocare con gli
        amici di Steam. Daniele stesso sta testando lì il proprio gioco;
  \item funziona anche su Mac: Windows o Mac, nessun problema.
\end{itemize}
Ci sarà una lezione apposita, di Ruggero, sulla creazione del proprio
\textbf{modulo} di gioco.

\paragraph{Che cosa non è.} Daniele chiarisce subito un ``fraintendimento
gigantesco'', soprattutto per chi ha già esperienza: Tabletop Simulator
\textbf{non richiede programmazione} e \textbf{non è un software per
trasformare i giochi da tavolo in videogiochi}. È un software che
\textbf{simula dei componenti fisici su un tavolo}: serve per giocare i board
game in digitale. Simula anche la fisica: c'è perfino il comando per
\textbf{ribaltare il tavolo} (``non usatelo ora'').

\paragraph{La partnership mancata.} Da anni Daniele prova a fare una
partnership con gli sviluppatori: due anni prima ha contattato l'azienda per
acquistare delle chiavi per gli studenti, dato che il software si usa nei corsi.
L'azienda ha risposto che non fa accordi commerciali. ``Credo sia la prima
azienda nella mia vita che dice che non vuole che io le dia dei soldi''. Gli
studenti dovranno quindi provare ad acquistarlo da sé; se i docenti riusciranno a
darglielo, bene.

\paragraph{La domanda di Gianmarco.} Gianmarco segue le lezioni da un computer
aziendale ``super chiuso'', su cui non può installare Steam: deve usare Tabletop
Simulator durante le lezioni? Ruggero risponde che per la lezione dedicata è
meglio poter seguire e rifare in diretta. Ma il software non sarà richiesto
necessariamente per il project work: un gruppo potrebbe fare un gioco pervasivo
o un gioco di ruolo, non un gioco da tavolo. Essendo un software di settore, è
comunque bene conoscerlo. Ruggero aggiunge che le versioni non ufficiali che
circolano non permettono di giocare in molti, e quindi non servono per
playtestare o giocare insieme.

\subsection{Una dimostrazione}

Daniele mostra il software aperto da Steam:
\begin{itemize}
  \item si può creare una partita \textbf{single player} o \textbf{multiplayer};
        per giocare tra compagni si crea un server multiplayer e ci si aggiunge
        come amici su Steam;
  \item si dà al server un nome, un numero di giocatori, eventualmente una
        password; creato il server, si apre la scelta dei giochi;
  \item ci sono giochi gratuiti, giochi da acquistare (DLC, non obbligatori) e
        soprattutto un \textbf{Workshop} ``ampissimo'', con decine di migliaia di
        giochi caricati dagli utenti: ``se c'è un board game che conoscete,
        quel board game è su Tabletop Simulator, giocabile dal Workshop''.
        L'esempio mostrato è l'open beta ufficiale del gioco da tavolo di
        \emph{Slay the Spire}\footnote{La trascrizione dice ``Slade Spire''.}:
        molti publisher, per farsi pubblicità, mettono il proprio gioco in demo o
        giocabile in digitale;
  \item c'è un \textbf{tutorial} iniziale che spiega come muoversi sul tavolo;
        con \emph{Esc} e \emph{Help} si trovano i comandi da tastiera.
\end{itemize}

\subsection{Il primo task}

Il primo (e per ora unico) task è questo: \textbf{entro il 10 ottobre}, giorno
della lezione di Ruggero su Tabletop Simulator, giocare \textbf{tre o quattro
volte in totale}, online, tra compagni di corso o con amici. Non si tratta di
imparare a creare moduli, ma di prendere la mano \textbf{lato giocatore}: aprire
il software, impararne i comandi, giocare. ``Per me potete anche andare a
prendere domino o mahjong'' (``mahjong è bellissimo tra l'altro'', aggiunge
Daniele, che sta facendo un videogioco sul mahjong), o il pachisi o il go. Gli
anni precedenti si chiedeva a ciascuno di proporre ai compagni un gioco da
tavolo che apprezzava. È anche un modo per conoscersi. Ruggero spiega il
perché: usarlo per giocare è facile e intuitivo, usarlo per creare moduli è più
complesso; ma chi arriva al 10 ottobre avendo già aperto il software e giocato,
troverà più intuitivo anche creare un modulo. E l'avvertimento: ``non riducetevi
all'ultimo, non fate una partita il 9 sera''.

\subsection{Tabletopia e altri vantaggi}

Tancredi racconta di aver provato a caricare un suo gioco su
\textbf{Tabletopia}. ``Questo è un vaso di Pandora'', risponde Daniele.
Tabletopia è molto usato, ma secondo lui il suo strumento di creazione dei moduli
è pessimo, ed è il motivo per cui la maggioranza dei moduli ufficiali dei
publisher è su Tabletop Simulator. Tabletopia ha un vantaggio vero, è
\textbf{gratuito}, ma ha scelte di ``qualità della vita'' poco felici.

Tabletop Simulator ha invece un vantaggio ``fenomenale'' che si tende a
sottovalutare: essendo su Steam, \textbf{il modulo si pubblica direttamente sul
Workshop di Steam}; si possono invitare gli amici a iscriversi e giocare con
loro, purché abbiano un account Steam. Rispetto ad altri software da imparare,
costa ``tra i 10 e i 20 euro, ragazzi: si parla proprio di spiccioli''. ``Io con
Tabletop Simulator ci ho playtestato la vita''.

Altre caratteristiche utili:
\begin{itemize}
  \item si possono creare \textbf{moduli non elencati} (\emph{unlisted}),
        perfetti da passare a un publisher prima che il gioco sia pubblicato:
        non sono pubblici, e nessuno può rubare i materiali;
  \item il software ha una grafica datata (``\emph{World Art 1998}'', scherza
        uno dei docenti), ma lo stanno portando su \textbf{Unity} con molti investimenti
        e aggiornamenti; ascolta la comunità e introduce funzioni utili anche per
        chi fa design;
  \item il suo server Discord è molto frequentato, e c'è un canale YouTube
        dedicato solo al design di giochi di carte\footnote{La trascrizione ne
        riporta il nome come ``Distraction Maker''; non è verificabile.} con un
        server Discord dove si fanno playtest: il 99\% è su Tabletop Simulator;
  \item molti editori, quando si propone loro un gioco da tavolo,
        \textbf{chiedono se c'è un modulo per Tabletop Simulator}.
\end{itemize}

% ---------------------------------------------------------------------
\section{Le domande a Daniele: il tema del project work}
\label{a1:sec:domande-daniele}

Daniele ha ``un minuto'' per rispondere alle domande sul corso prima di dover
andare. Uno studente chiede, ``giusto per curiosità'', se il project work avrà
un tema o sarà libero.

\begin{itemize}
  \item C'è un \textbf{mandate}, che i docenti propongono ogni anno e
        aggiornano leggermente; verrà presentato il \textbf{25 settembre}.
        ``Non facciamo spoiler''. Conterrà quello che si chiede di fare e quello
        che non si chiede: dei \textbf{must}, dei \textbf{must not} e dei
        \textbf{nice to have}.
  \item Verrà chiesto di lavorare su un'\textbf{idea originale} che non si è
        portata al corso: originale nel senso di \textbf{creata apposta}, non
        nel senso di ``mai vista prima''.
  \item L'idea va \textbf{partorita insieme}. Altri anni alcuni studenti hanno
        provato a imporre la propria idea al gruppo, ``e l'esperienza è stata
        pessima per tutti''. Quindi: \textbf{non imponete le vostre idee, fate
        brainstorming con gli altri}. Daniele, che fa ``molto spesso
        brainstorming'', darà strumenti per farlo: nessuno arriverà impreparato
        a quella lezione.
\end{itemize}

Ruggero aggiunge che il mandate è \textbf{abbastanza flessibile} perché tutti
possano fare qualcosa che li interessi: ``l'ultima cosa che vogliamo è forzarvi a
lavorare su qualcosa che non vi piace''. Bisogna lavorare su ciò che piace e per
cui si ha uno stimolo personale, anche nella \textbf{tipologia} di gioco: nessuno
sarà obbligato a fare un gioco di ruolo o un board game, e la lezione sui casi
studio mostrerà quante forme diverse di gioco analogico il corso accetta.

L'unica cosa da \textbf{non fare}, ``perché è contraddittorio per questo corso'',
sono i \textbf{giochi grandi}: ``lo chiediamo per voi''. Niente progetti
mastodontici\footnote{La trascrizione dice ``non fate delle cose da gnome'':
probabilmente il nome di un gioco da tavolo molto grande e complesso, non
ricostruibile con certezza.}, niente ``cose alla \emph{Dungeons \& Dragons}''.
Si può fare un gioco di ruolo, ma con un'esperienza ``meno corposa'': ci sono
tanti giochi di ruolo da 15 o 20 pagine, ``bellissimi, fichissimi'', che si
possono leggere online. Il consiglio è farsi una cultura anche su
\textbf{itch.io}, ``uno dei portali più importanti del gioco di ruolo
internazionale'', di cui si parlerà anche nella parte di ricerca e sviluppo.

Poi Daniele saluta (``noi ci vediamo venerdì a lezione'') e lascia la
conduzione a Ruggero.

% ---------------------------------------------------------------------
\section{Le domande finali con Ruggero}
\label{a1:sec:domande}

\subsection{Le regole delle domande}

Ruggero spiega che cosa vuol dire ``stare con lui''. Alla fine di ogni sua
lezione ci saranno \textbf{almeno 20--30 minuti di domande} libere, \emph{free
form}. Vuole che fare domande sia \textbf{normalizzato}: ``io gestisco malissimo
l'\emph{horror vacui} e divento super cringe molto alla svelta se non ci sono
domande in giro''. Sono benvenute le domande sulla lezione appena fatta, ma non
solo: il corso è crossmediale, e anche le domande che vengono da
\textbf{divagazioni}, da parti laterali, ``fatele''. Quello che sembra scontato o
banale a uno studente può non esserlo per il docente, e viceversa; e può essere
utile anche agli altri. ``Se vi sembrano stupide, fatele il doppio''. ``Magari
sembro incazzato, ma vi giuro sono rilassatissimo. Non mordo, magari abbaio un
po'\,''.

\subsection{Come approcciarsi a Steam e Tabletop Simulator}

Rita non ha mai usato né Steam né Tabletop Simulator e chiede come
approcciarsi, e se c'è modo di mettersi d'accordo tra compagni su Discord, magari
con un sondaggio sui giochi, per formare gruppi di gioco. Ruggero conferma: sul
Discord dell'accademia c'è un \textbf{canale apposito per il corso}, ed è degli
studenti. Lo possono usare per organizzarsi a giocare, per chiedere a un
compagno di prestare un gioco, per domande tecniche (come si scarica, come si
installa). ``Sarebbe proprio ideale che vi parlaste e vi organizzaste''.

Anche perché, aggiunge con un ``piccolo spoiler'' (``nessuno vi ha detto
niente''), se fossero i docenti a organizzare i gruppi di gioco, il task lo
starebbero facendo loro. Bisogna \textbf{abituarsi a organizzare le sessioni di
gioco}, perché arriveranno i task sul \textbf{playtest} e su come organizzarlo,
e a quel punto si sarà ``meno disperati''. Sembra scontato, ma saper
\textbf{organizzarsi per giocare in autonomia} ``vi sarà utilissimo come skill in
questo settore, ve lo giuro''.

\subsection{Il bilanciamento}

Angelo confessa ``una paranoia grossissima'': ci saranno lezioni sul
\textbf{bilanciamento}? ``Ma certo''. Sono proprio le lezioni di Marta sul
development: il bilanciamento è di solito un tema che tocca \textbf{i developer
più dei designer}. L'anno precedente c'era una sola lezione con Marta, e i
docenti si sono resi conto che serviva andare ancora più a fondo: per questo
quest'anno le lezioni sono due. ``Non è una paranoia, è una legittima domanda'':
Ruggero insegna da anni, vede altri corsi, parla con altri docenti e studenti, e
``mi incazzo spesso'' perché molti corsi di game design \textbf{non fanno
bilanciamento}. I docenti volevano invece un corso che coprisse tutte le parti
del game design: non solo l'idea del gioco, ma il bilanciamento, la produzione,
i playtest, l'iterazione, la scrittura dei regolamenti.

\subsection{Si può giocare con i docenti?}

Una studentessa chiede se si può giocare con i docenti. Per motivi etici,
durante il corso questo è limitato alle parti di \textbf{valutazione} dei giochi
e di \textbf{tutoraggio}. Dopo il corso, ``nessun problema, almeno con me''.

\subsection{Dove si trovano i giochi su Tabletop Simulator}

Ruggero mostra dove si trovano i giochi: nella propria libreria di Steam, nella
pagina di Tabletop Simulator, alla voce \textbf{Workshop}, dove ci sono ``più o
meno 36.000'' giochi. Una studentessa scopre intanto che l'app di Steam per
telefono è solo una \emph{companion app}: per registrarsi e giocare serve il
PC.

Per provare, chiede agli studenti un titolo:
\begin{itemize}
  \item \emph{Keep the Heroes Out}: non compare, ``non tutti i giochi ci sono,
        perché molti publisher non li mettono''; forse, essendoci già una
        versione del gioco su Steam, il modulo è stato tolto;
  \item \emph{Magic: The Gathering}: ``almeno 50 moduli'', anzi ``123 pagine'';
  \item \emph{Bang!}: tre moduli.
\end{itemize}
Bisogna fare attenzione: alcuni moduli sono in lingue specifiche (``magari il
primo è in russo''), alcuni hanno già l'espansione, altri sono solo il gioco
base. Conviene sfogliare e non fermarsi al primo risultato. E nel Workshop si
trovano anche cose molto particolari: giochi fuori produzione, giochi vecchi,
giochi ancora in prototipo.

\subsection{I ruoli nei team e le dinamiche tra le persone}

Uno studente torna sulla slide delle figure: nelle realtà piccole molti ruoli
hanno ``il doppio o triplo cappello''. Quali ruoli restano distinti, e come si
creano le sinergie tra persone con obiettivi diversi?

Ruggero premette che l'ideale ``platonico'' è un ambiente di lavoro costruttivo,
basato sulla fiducia e ``il più \emph{safe} possibile''; ma nell'industria non
sempre è così, e le strutture variano moltissimo: alcuni publisher sono molto
verticali, altri hanno ruoli già solidificati (se un publisher ha già i suoi
developer e non sta crescendo, entrarci come developer è praticamente
impossibile). ``Se prendiamo due aziende in Italia, vediamo due strutture
interne completamente diverse''. Non esiste una scorciatoia per farsi assumere
da un editore.

Poi prende ``l'esempio più comune di tutti'', la media dei giochi da tavolo
pubblicati (figura~\ref{a1:fig:relazioni}):
\begin{itemize}
  \item il \textbf{game designer} è una figura \textbf{esterna}, con un
        \textbf{contratto a royalty} con il publisher;
  \item la sua figura di riferimento è il \textbf{publisher} stesso, e la
        conversazione designer--publisher è ``molto calda'';
  \item se il publisher ha dei \textbf{developer}, il developer entra in questa
        conversazione: designer e developer collaborano molto, e internamente
        developer e publisher collaborano molto. Spesso, nei publisher piccoli,
        è il publisher a fare le veci del developer;
  \item \textbf{graphic design} e \textbf{illustrazione} tendono a parlare molto
        con il publisher e \textbf{molto poco con il designer}. A meno che non si
        tratti di uno studio indie molto piccolo, è difficile che il designer
        sappia anche solo chi è l'illustratore prima che il gioco esca.
\end{itemize}

\begin{figure}[htbp]
\centering
\begin{tikzpicture}[
    font=\small,
    ruolo/.style={draw, thick, rounded corners=3pt, fill=gray!12,
                  minimum width=2.8cm, minimum height=0.9cm, align=center},
    esterno/.style={draw, thick, dashed, rounded corners=3pt,
                    minimum width=2.8cm, minimum height=0.9cm, align=center},
    calda/.style={{Stealth[length=2.5mm]}-{Stealth[length=2.5mm]}, very thick},
    tiepida/.style={{Stealth[length=2mm]}-{Stealth[length=2mm]}, thin, dashed}
  ]
  \node[esterno] (gd) at (0,0) {Game designer\\ \footnotesize esterno, a royalty};
  \node[ruolo] (dev) at (5,1.4) {Developer\\ \footnotesize se c'è};
  \node[ruolo] (pub) at (5,-1.4) {Publisher};
  \node[ruolo] (art) at (10,-1.4) {Graphic designer\\ e illustratore};
  \draw[calda] (gd) -- (dev);
  \draw[calda] (gd) -- (pub);
  \draw[calda] (dev) -- (pub);
  \draw[calda] (pub) -- (art);
  \draw[tiepida] (gd.south) to[bend right=25]
     node[below, font=\footnotesize\itshape] {quasi nessun contatto} (art.south);
  \node[draw, dotted, rounded corners=4pt, fit=(dev)(pub)(art), inner sep=8pt,
        label={[font=\footnotesize\itshape]above right:interno al publisher (di solito)}] {};
\end{tikzpicture}
\caption{Le relazioni tipiche intorno a un gioco da tavolo pubblicato, secondo la
risposta di Ruggero. Le frecce spesse indicano una collaborazione stretta e
frequente; quella tratteggiata, un contatto quasi assente. Nei publisher piccoli
il publisher fa spesso anche da developer.}
\label{a1:fig:relazioni}
\end{figure}

\subsection{Fare da developer di sé stessi?}

Rita chiede in chat se un designer può fare anche da developer del proprio
gioco, o se Ruggero lo sconsiglia. In generale \textbf{lo sconsiglia}, per lo
stesso motivo per cui si sconsiglia in altri media: i \textbf{darling}. Il
designer ha delle meccaniche a cui è affezionato, e il developer a volte dovrà
\textbf{tagliarle}; è molto difficile, se si è ``gestori di sé stessi'', tagliare
cose che si vogliono nel gioco. ``È un po' come lasciare fare il lavoro di editing
allo scrittore che ha già scritto il libro'': può rivedersi, certo, ma difficilmente
cambierà cose che farebbe bene cambiare. Nelle lezioni con Marta si vedrà perché,
e alla fine del corso ``non avrete dubbi nemmeno voi'' nel separare designer e
developer.

Ma tutto dipende dallo \textbf{scopo del progetto}: se si fa un gioco in tre,
qualcuno quel lavoro deve farlo per forza, e magari il designer fa anche il
developer, compensando con \textbf{più playtest}. ``Il mondo attuale e il mondo
ideale sono due cose molto diverse, e lo vediamo non solo dai giochi ma dalla
politica'', ma ``non entriamo in questo vaso di Pandora''.

\subsection{Il ruolo indispensabile e la pubblicazione}

Lo studente di prima ne trae una conclusione: senza game designer il gioco non
esiste. ``Quello sicuramente''. Ma Ruggero coglie l'occasione per approfondire
un punto che forse non era passato: \textbf{fare design di un gioco per
pubblicarlo non è l'unica opzione possibile}. Un gioco non è ``meno gioco''
se non è pubblicato da un editore. È legittimo mettere un gioco su itch.io in
\textbf{print and play}, fare un gioco a evento senza venderlo, fare un
\emph{murder party}, una escape room, un board game per una persona amica:
``quell'esperienza lì non è meno gioco''. Dipende dall'obiettivo del progetto.

Ruggero ha lavorato molto anche con le \textbf{scuole}, dove il publisher non
esiste: è un gioco per le classi, e la figura più vicina è l'editoria scolastica,
``che però non fa il publisher''. Ci sono compromessi diversi sulla materialità
e su come il gioco si presenta, che si vedranno nella prototipazione e nei casi
studio. Quindi l'unica figura che deve esserci sempre è \textbf{chi fa game
design}, anche se magari non sa di chiamarsi game designer e non ha mai
sentito parlare di game design: ``questo non la rende meno game designer''.

\paragraph{Un libro consigliato.} Ruggero fa anche mentoring e gestisce gruppi
di settore, e spesso gli chiedono che cosa leggere per imparare il game design.
Non è il tipo che obbliga a leggere libri (``non vi obbligherò a leggere
niente''), ma ce n'è uno che consiglia sempre, ed è il primo che ha scelto di
leggere: \textbf{\emph{Rise of the Videogame Zinesters}} di \textbf{Anna
Anthropy}\footnote{La trascrizione dice ``Anna Antrofi'' e ``Rise of the
Videogames in Esters''.}. È un libro un po' datato, ma ancora reperibile; parla
soprattutto di videogiochi, ma ``per l'80\%'' il contenuto si applica al gioco e
al game design in generale. È un libro sulla \textbf{democratizzazione del gioco}:
il gioco come strumento per parlare di sé e come \textbf{medium legittimo}. Come
esistono film drammatici, tristi, intensi, possono esistere giochi su tematiche
delicate, e se ne vedranno nei casi studio. Molti giochi fatti con questa
direttiva non hanno un publisher, che tende a pubblicare ciò che dà un ritorno
commerciale.

\paragraph{Commerciale e personale.} ``Io spero che voi non farete giochi solo
per scopo commerciale''. Certo, servono un tetto sulla testa, il cibo in
tavola, un posto in cui dormire, e così si finisce per fare giochi che non sono
il proprio gioco dei sogni. Spesso è un mix: Ruggero ha fatto giochi commerciali
``per pagarmi l'affitto'' e giochi che voleva fare e su cui ha perso soldi,
``sapendolo proprio che sarebbero stati un fallimento''. I giochi, come i libri e
le canzoni, si fanno per molti motivi, e quelli commerciali di successo non sono
più importanti degli altri. Tra i giochi della sua slide c'erano il gioco da
tavolo ufficiale di \emph{Cyberpunk}\footnote{La trascrizione dice ``Cyberpunk
del 1977'': quasi certamente un gioco legato a \emph{Cyberpunk 2077}.} e
\emph{Marvel United}, su IP grosse e famose, ma anche giochi fatti di ``una
candela e una busta su un tavolo''. E se gli si chiede quale preferisce in
termini di contenuto, non sono quelli sui supereroi.

\subsection{Il gioco di ruolo in Italia}

Nicola, che è interessato soprattutto al gioco di ruolo e conosce già persone
del settore, ha ``una leggera paura'': che chi studia game design finisca per
fare il \textbf{freelance} a vita, ``perché non ci sarà mai un'occupazione per
bene'', o il \textbf{copywriter}. Conosce un ragazzo che ha lavorato in quel
mondo e che ora fa altro.

Ruggero risponde che la maggioranza dei giochi nel suo portfolio sono giochi di
ruolo, quindi il tema lo tocca da vicino, e che in Italia quello che dice Nicola
``è semplicemente vero''. Molti publisher piccoli hanno una persona di
riferimento che copre più ruoli: impagina, fa i playtest, fa un po' di
development. È ``un settore economicamente completamente disastrato'', per
motivi culturali che non apre: oggi in Italia è praticamente impossibile
\textbf{vivere di giochi di ruolo}. Chi ci lavora lo fa nel tempo libero, al di
fuori del lavoro principale, oppure coprendo molte altre cose: molti publisher
italiani di giochi di ruolo pubblicano anche titoli esteri e localizzano, non
fanno solo cose originali. ``È un grosso calderone''. Poi si ferma: ``questo è
un \emph{rant}'', e ``i \emph{rant} tornano sempre'', ma non vuole che gli
studenti escano dalla prima lezione ``con la disperazione nel cuore''. Ci sarà
spazio per parlarne quando avranno gli strumenti per capirlo, ``per non farvi
abbattere''.

\subsection{Che cosa bisogna sapere già}

Uno studente chiede se ci sono nozioni che dovrebbero già avere e che possono
recuperare: per esempio, sente parlare di giochi ``German'' o ``American'' e non
sa che cosa siano. Ruggero spiega l'organizzazione del corso:
\begin{itemize}
  \item le \textbf{prime due lezioni} sono di fondamenta. Questa era ``dall'alto'':
        come è strutturato il corso, le tempistiche;
  \item la \textbf{prossima lezione} è di tre ore di soli \textbf{casi studio},
        ``a manetta, a raffica'', divisi per categorie e tipologie, con
        ragionamenti su categorie e generi di gioco: lì si vedranno German,
        American ``e anche altre più utili'';
  \item nel resto del corso le terminologie specifiche sono inserite nelle
        lezioni dedicate: il \textbf{paragame} nelle lezioni sul paragame, il
        \textbf{modello matematico} nelle lezioni di Marta sul development.
\end{itemize}
Per ora agli studenti si chiede di sapere quello che si è visto oggi:
meccaniche, dinamiche ed estetiche; le figure professionali, designer e
developer; e soprattutto \textbf{game e play}.

\subsection{Game e play: i quattro quadranti}

``Specialmente \emph{game} e \emph{play}'', insiste Ruggero: è più difficile di
quello che sembra, perché l'italiano non abitua a separare le due cose, ed è
possibile averle usate per anni come sinonimi. Nella sua esperienza con chi si
avvicina al design è ``uno dei punti più critici''. Spesso i commenti che si
fanno sui giochi giocati sono commenti sul \emph{play} applicati al
\emph{game}: ``il gioco non mi è piaciuto perché, giocando con Tizio, è successo
questo''; ``il gioco mi sembra sbilanciato perché Tizio al tavolo ha fatto
questo''.

Separare \emph{game} e \emph{play} serve a capire quali giochi sono progettati
bene e quali male, e soprattutto a riconoscere i \textbf{quattro possibili
incroci} (figura~\ref{a1:fig:quadranti}): un gioco progettato bene che piace
giocare; uno progettato male che non piace; uno progettato male che piace; uno
progettato bene che non piace. ``Non è importantissimo in generale nella vita,
ci sono cose più importanti'', ma per chi fa game design lo è: ``a volte vi
chiederanno di lavorare su un gioco che non vi piace giocare, e quel gioco deve
essere fatto bene lo stesso''.

\begin{figure}[htbp]
\centering
\begin{tikzpicture}[
    font=\small,
    cella/.style={draw, thick, minimum width=4.4cm, minimum height=1.9cm,
                  text width=4.0cm, align=center}
  ]
  \node[cella, fill=gray!20] (a) at (0,0)     {\textbf{progettato bene}\\ e \textbf{mi piace} giocarlo};
  \node[cella]               (b) at (4.4,0)   {\textbf{progettato bene}\\ ma \textbf{non mi piace} giocarlo};
  \node[cella]               (c) at (0,-1.9)  {\textbf{progettato male}\\ ma \textbf{mi piace} giocarlo};
  \node[cella, fill=gray!8]  (d) at (4.4,-1.9){\textbf{progettato male}\\ e \textbf{non mi piace} giocarlo};
  \node[font=\footnotesize\itshape, above=0.15cm of a] {mi piace (\emph{play})};
  \node[font=\footnotesize\itshape, above=0.15cm of b] {non mi piace (\emph{play})};
  \node[font=\footnotesize\itshape, rotate=90, left=0.15cm of a, anchor=south]
       {bene (\emph{game})};
  \node[font=\footnotesize\itshape, rotate=90, left=0.15cm of c, anchor=south]
       {male (\emph{game})};
\end{tikzpicture}
\caption{I quattro incroci tra la qualità del \emph{game} (come è progettato) e
l'esperienza di \emph{play} (se piace giocarlo), descritti da Ruggero. Per un
designer i quadranti ``scomodi'' sono i due fuori diagonale: un gioco può
piacere pur essendo progettato male, e va saputo progettare bene anche un gioco
che non piace.}
\label{a1:fig:quadranti}
\end{figure}

\paragraph{La domanda di Tancredi.} Tancredi chiede conferma: il \emph{play}
non è il ``giocare senza regole'', ma può essere anche il giocare a un gioco? Sì:
il \emph{play} è l'\textbf{atto}, e può essere fatto con un \emph{game} o senza.
E un \emph{game} può esistere senza che nessuno lo giochi: gli scacchi non sono
meno scacchi perché in quel momento nessuno ci gioca. E se, giocando, una
persona applica le regole in modo diverso? Allora \textbf{non si sta giocando
quel gioco}.

\subsection{House rule: il caso di Monopoly}

Qui Ruggero fa un po' di \emph{foreshadowing} su una lezione di Daniele: il
discorso degli \textbf{hack}, delle \textbf{house rule}, delle \textbf{mod},
cioè di come modificare un gioco esistente. Il punto è che \textbf{un gioco
modificato da una house rule è un altro gioco}: cambiando le regole cambiano
necessariamente anche le \textbf{dinamiche} che ne derivano, e quindi
l'esperienza. Si immagini \emph{Monopoly} in cui ogni turno si possa, volendo,
prendere gratis dei soldi dalla banca: tutto il resto è identico, eppure il gioco
cambia.

E questo è talmente interiorizzato che ci sono regole che si applicano ``senza
neanche rendercene conto''. Ruggero chiede quanto duri in media una partita a
\emph{Monopoly}. Gianmarco: ``quattro ore''; altri: ``un pomeriggio''; un altro:
``finché uno dei giocatori non aggiunge la regola: ora vince chi\ldots''. In
realtà, \textbf{seguendo le regole scritte}, una partita non dura mai più di
\textbf{un'ora e venti, un'ora e mezza}. La domanda che fa sempre, e la risposta
è sempre la stessa: perché si pensa che \emph{Monopoly} duri tantissimo? Qual è
la regola che nessuno rispetta?

Uno studente indovina: \textbf{le aste}. Secondo il regolamento, se chi finisce
su un terreno non lo compra, il terreno \textbf{va all'asta} tra i giocatori, e
chiunque può aggiudicarselo. Molti invece giocano come se finire su un terreno
desse solo la \emph{facoltà} di comprarlo o di metterlo all'asta\footnote{Più
precisamente, nel regolamento ufficiale, se il giocatore che finisce su una
proprietà libera decide di non comprarla, la banca la mette all'asta e tutti i
giocatori possono fare offerte, compreso chi l'ha rifiutata. L'asta quindi non è
una scelta del giocatore, ma la conseguenza obbligata del suo rifiuto.}. Con la
regola applicata, di aste se ne fanno moltissime: all'ottavo turno ci sono già
otto terreni assegnati, al ventesimo venti, la gente va in bancarotta
velocissimamente e la partita dura un'ora e mezza. Anche le partite competitive
di \emph{Monopoly}, nota Ruggero, durano un'ora e venti, un'ora e mezza.

Uno studente obietta che giocava con le aste e durava comunque quattro ore.
``Avete dei nemici con l'analisi paralisi'', ride Ruggero; e cita un'altra house
rule tipica dell'infanzia: genitori che danno soldi per non far andare in
bancarotta. Ma \textbf{la bancarotta è parte del gioco}: se la si toglie,
\emph{Monopoly} non regge. Sul tema che \emph{Monopoly} vuole trattare, e sul
suo antenato (\emph{The Landlord's Game}), ``non apriamo adesso un altro vaso di
Pandora''.

\paragraph{Approfondimento: Monopoly, la durata e le house rule nei libri.}
L'esempio di Ruggero si ritrova quasi identico in \emph{The Art of Game Design}.
Parlando del bilanciamento tra partite brevi e lunghe (\emph{Short vs. Long}),
Schell scrive che \emph{Monopoly}, giocato secondo le regole ufficiali,
``finisce spesso in circa novanta minuti'', ma molti giocatori trovano quelle
regole troppo dure e le modificano, distribuendo premi in denaro e allentando i
vincoli sull'acquisto delle proprietà: con l'effetto collaterale di far durare la
partita molto di più, ``tipicamente tre ore, o anche di più''. Schell classifica
le \textbf{house rule} come regole create dai giocatori in risposta a un
difetto percepito dopo qualche partita. Anche \emph{Challenges for Game
Designers}, nel capitolo sull'aggiungere e togliere meccaniche, parte dallo
stesso gioco: il \emph{Parcheggio gratuito}, che per le regole originali non fa
succedere nulla, mentre quasi tutti giocano con una house rule che dà un bonus a
chi ci finisce sopra. Le due fonti confermano la lezione: una piccola regola
cambiata cambia le dinamiche, e le dinamiche cambiano l'esperienza.

\subsection{Il fine del game è il play?}

Un altro studente commenta: va bene la distinzione, ma il fine del \emph{game} è
il \emph{play}; quello che si vede è sempre l'effetto. Ruggero è d'accordo: di
solito, se si fa un gioco, lo si fa per farlo giocare, come si scrive un libro
per farlo leggere (anche se esistono, ed ``è già stato esplorato'', libri scritti
per non essere letti e giochi per zero giocatori). Ma la distinzione serve a una
cosa precisa: \textbf{separare, nel playtest, i feedback sul gioco dai feedback
sul gruppo di gioco} e sulle sue dinamiche. Se una persona esce frustratissima da
un playtest, non è necessariamente colpa del gioco; se una persona dice ``questo
è il gioco migliore del mondo'', potrebbe non essere merito del gioco ma del
gruppo. Capire a che cosa si riferisce un feedback serve a capire \textbf{dove
iterare} il prototipo.

Lo studente osserva che l'area è sfumata: alcuni giochi si basano proprio
sull'avere le persone giuste al tavolo. Un gioco di \textbf{negoziazione} con
giocatori abituati ai giochi ``multi-solitario'' diventa ingiocabile; e nei giochi
a \textbf{economia chiusa} (cita \emph{Container}\footnote{La trascrizione cita
anche un secondo titolo non ricostruibile (``Tokyo su Vigimarket'').}) un gruppo
economicamente conservativo, che spende poco e non fa girare l'economia, fa
durare la partita il doppio. Ruggero rilancia: e se il gioco dura il doppio,
questa dinamica cambia la \textbf{percezione di peso} del gioco, le energie spese,
quanto si è stanchi a fine partita. ``Tutto è collegato. \textbf{I giochi sono
ecosistemi di meccaniche}, e cambiandone una, le altre cambiano di
conseguenza''.

\paragraph{La parola vietata, alla prova.} Lo studente aggiunge che ci sono giochi
che ha trovato ``molto sbagliati, ma che mi divertono''. ``Attenzione: non si può
dire \emph{divertimento}!'', lo ferma Ruggero; nelle presentazioni aveva lasciato
correre, ``ma lì non vi avevo ancora letto la regola''. Lo studente ci prova:
``giochi che poi voglio \textbf{rigiocare}'', anche se ``sono evidentemente
sbagliati, rotti''. Ruggero apprezza, ma non vuole anticipare cose che meritano il
loro tempo: se ne parlerà nelle lezioni sui framework e soprattutto in quella di
\textbf{emotion design}, ``forse la mia lezione preferita del corso''. (Lo
studente ammette che ``forse alla fine si gioca per quello, per l'emozione''.)

\subsection{I giochi ``comunicati male''}

L'ultima domanda corposa riprende un'affermazione di inizio lezione: ci sono
giochi belli ma comunicati male. Qualche esempio? ``Mi basta girarmi per
vederne 48'', risponde Ruggero indicando gli scaffali alle sue spalle. Più che
titoli, elenca le \textbf{ragioni}:
\begin{itemize}
  \item \textbf{l'Italia localizza lentamente}. Basta guardare i premi del gioco
        dell'anno in Italia: spesso vanno a giochi usciti due o tre anni prima in
        Germania, Francia, Stati Uniti o Regno Unito, perché localizzare richiede i
        suoi tempi. Il mercato italiano è quindi molto dipendente dall'estero,
        ``molto impaurito'', e spesso aspetta che un gioco vada bene all'estero
        prima di portarlo in Italia; e alcuni giochi ``perdono il treno'';
  \item \textbf{il publisher originale chiude}: come si pubblica un gioco se non
        c'è più l'azienda con cui fare un contratto? A Ruggero è successo più
        volte;
  \item \textbf{il dipartimento commerciale italiano non sa comunicare il gioco}:
        si fa male un'edizione, la stampa vende pochissimo e non viene più
        ristampata;
  \item \textbf{la scala del mercato}: in Italia un gioco da tavolo di peso medio
        (un \emph{eurogame}, un \emph{German}) che vende 1.500 copie ``è tanto'';
        all'estero se ne vendono magari 10.000. Con 10.000 copie in inglese ci
        sono i soldi per una campagna di comunicazione: un trailer, i
        recensori, i tutorial online. In Italia tutto questo è ``ridotto all'osso'',
        e la cosa si riflette su stampa e divulgazione: anche i canali per
        informarsi su un gioco prima di comprarlo sono pochi (se ne parlerà nella
        lezione su ricerca e sviluppo). Un gioco spesso non arriva al mercato non
        perché sia brutto, ma perché l'editore non lo sa comunicare, ne fa poche
        copie e va fuori produzione;
  \item \textbf{il tempismo}: party game progettati nel 2019 e usciti nel 2020, in
        pieno Covid, quando non si poteva giocare al tavolo, sono falliti tutti, e
        gli editori ci hanno perso. ``Non potevano immaginarselo'';
  \item \textbf{i temi caldi}: dopo l'invasione russa dell'Ucraina e la guerra in
        Palestina, alcuni editori non pubblicano o non localizzano giochi
        strategici e di guerra. Ruggero ha lavorato a \textbf{\emph{Status}}, un
        gioco sulla geopolitica che tocca, tra gli altri temi, l'invasione russa
        dell'Ucraina e il genocidio in Palestina, e alcuni editori che di
        solito localizzano i loro giochi (per esempio il partner di riferimento in
        Russia) non lo hanno localizzato per motivi politici.
\end{itemize}
Lo studente voleva capire che cosa voglia dire comunicare male: non metterci
abbastanza energia, o comunicare una cosa per un'altra, ``come il trailer
sbagliato, che poi vai al cinema e ti trovi un altro film''. ``È un mix'', e la
domanda ``era complessa e meritava più spazio''.

\subsection{I giochi di carte collezionabili e i saluti}

L'ultima domanda (``posso rompere le balle un'ultima volta?'') riguarda i
\textbf{giochi di carte collezionabili}: funzionano con le stesse regole dei
giochi da tavolo? Se ne parlerà alla prossima lezione: è proprio il motivo per
cui la lezione sui casi studio dura tre ore, perché ogni genere di gioco ha le
sue dinamiche e le sue regole, e anche il modo di pubblicare cambia molto tra un
TCG, un gioco di ruolo da tavolo, un LARP e un gioco pervasivo.

Ruggero saluta: le email dei tre docenti sono all'inizio delle slide, la
registrazione arriverà, ``ci si vede questo venerdì''.

% ---------------------------------------------------------------------
\section*{Punti chiave del capitolo}
\addcontentsline{toc}{section}{Punti chiave del capitolo}

\begin{itemize}
  \item Il corso di \textbf{Analog Game Design} (seconda edizione) dura circa tre
        mesi, fino a dicembre, ed è tenuto da tre docenti principali:
        \textbf{Daniele} (narrative design, gioco di ruolo, fiction),
        \textbf{Ruggero} (game design, ricerca e sviluppo) e \textbf{Beatrice}
        (giochi pervasivi), con tre ospiti su RNG, development e game production.
  \item Lo scopo è formare \textbf{autori analogici} consapevoli: strumenti da
        usare anche fuori dal corso, non il ``gioco dei sogni''. Il corso è
        \textbf{intersezionale} (nessuna lezione dedicata a un solo tipo di
        gioco), \textbf{cross-generazionale}, \textbf{pratico e adattabile},
        basato sulle \textbf{domande} (nessuna domanda è stupida), sul
        \textbf{feedback} invece che sul voto e sul \emph{fail forward}.
  \item Il \textbf{project work} si fa in tre gruppi da tre, formati dai docenti;
        pitch il 12 ottobre, sviluppo nelle due pause, tutoraggi il sabato,
        consegna il 4 dicembre, presentazione l'11 e giornata finale con gli
        editori il 14. Si valutano documentazione, presentazione,
        autovalutazione, crescita e atteggiamento.
  \item \textbf{Niente intelligenza artificiale} nel corso: i publisher
        rifiutano pitch e prototipi con immagini generate, i regolamenti scritti
        dall'IA descrivono la stessa meccanica in modi diversi, e soprattutto chi
        la usa perde l'occasione di imparare. Un prototipo può essere fatto a
        matita.
  \item Le \textbf{figure dell'industria}: game designer, game developer,
        editor, publisher, sales e marketing, graphic designer, illustratore,
        sensitivity reader (che non è un censore) e producer. Nelle realtà
        piccole una persona copre più ruoli.
  \item Nell'analogico il \textbf{game designer} progetta l'esperienza creando
        le meccaniche; il \textbf{game developer} fa da \textbf{mediatore} con il
        publisher (bilanciamento, modelli matematici, adattamenti). Pretendere
        che il proprio gioco non venga modificato è un \emph{red flag}.
  \item Il \textbf{marketing} non decide che cosa pubblicare ma come; è
        un'azione, ed esiste anche se lo si ignora.
  \item In \textbf{Italia} il gioco è ancora visto come infantile, il mercato
        punta sulla localizzazione e investe poco su giochi e designer interni;
        all'estero (soprattutto dalla Germania in su) ci sono tradizione adulta,
        premi di prestigio e developer interni.
  \item \textbf{Game} è l'artefatto (le regole, il sistema, l'oggetto);
        \textbf{play} è l'atto di giocare. Si fanno giochi, non giocattoli.
        Separarli serve a distinguere i feedback sul gioco da quelli sul gruppo
        e a riconoscere i quattro incroci tra ``progettato bene/male'' e ``mi
        piace/non mi piace''.
  \item \textbf{MDA}: meccaniche (regole scritte), dinamiche (regole in azione),
        estetiche (esperienza, non grafica); nel \textbf{DPE}: design, play,
        experience.
  \item \textbf{Parole vietate} in ogni materiale scritto: \emph{fun},
        \emph{immersion}, \emph{entertainment} (divertimento, immersione,
        intrattenimento). Bisogna trovare parole più precise.
  \item Distinguere \textbf{chi crea, chi gioca e chi facilita}; nessun gioco
        nasce perfetto: \textbf{playtest e iterazione} sono inevitabili, e un
        \textbf{dizionario comune} aiuta a non perdersi.
  \item \textbf{Tabletop Simulator} è il software di settore per giocare (non
        creare videogiochi) i giochi da tavolo in digitale. Primo task: giocarci
        tre o quattro volte entro il 10 ottobre.
  \item Una \textbf{house rule} crea un altro gioco: \emph{Monopoly} con le aste
        applicate come da regolamento dura circa un'ora e mezza.
\end{itemize}

%
%  (Con hyperref conviene aggiungere anche
%  \renewcommand{\theHchapter}{agd.\arabic{chapter}} nella parte, per
%  evitare ancore duplicate con i capitoli 1, 2, ... di Game Design.)
%  Tutte le etichette di questo capitolo hanno il prefisso "a1:"
%  (a1:cap:..., a1:sec:..., a1:fig:..., a1:tab:...). Il capitolo è
%  autonomo: non contiene \ref verso altri capitoli.
%  Le figure sono ricostruzioni schematiche di quanto i docenti spiegano
%  a voce e con le slide, non copie delle slide.
% =====================================================================

\chapter[Introduzione all'Analog Game Design]{Introduzione all'Analog Game
Design: il corso, le figure dell'industria e un primo linguaggio comune}
\label{a1:cap:intro}

Questa è la \textbf{prima lezione} del corso di \textbf{Analog Game Design},
cioè di progettazione di giochi \emph{analogici}: giochi da tavolo, giochi di
carte, giochi di ruolo da tavolo, giochi di ruolo dal vivo (LARP), giochi
\emph{pervasivi} ed esperienze a evento. È un corso \textbf{breve}, di circa tre
mesi (``da qui a dicembre''), separato dal corso di Game Design, e siamo alla sua
\textbf{seconda edizione}. Come avvertono subito i docenti, è ``una giornata di
introduzione'': si spiega come funziona il corso, chi sono i docenti, con quali
metodi e con quale calendario si lavorerà, e ci si presenta a vicenda. ``C'è poco
gioco analogico'', si scusano. Ma nella seconda parte la lezione comincia
comunque a costruire un \textbf{linguaggio comune}: le figure professionali
dell'industria del gioco analogico, la differenza tra \emph{game design} e
\emph{game development}, la situazione del mercato italiano rispetto a quello
estero, e un primo glossario di design (\emph{game} e \emph{play}, il framework
MDA, le parole ``vietate'', il playtest e l'iterazione). Si chiude con lo
strumento di settore per i prototipi digitali dei giochi da tavolo,
\textbf{Tabletop Simulator}, e con una lunga sessione di domande.

La lezione si svolge \textbf{online}, su Zoom, ed è condotta a più voci dai tre
docenti principali del corso: \textbf{Daniele}, narrative designer, che guida la
parte organizzativa; \textbf{Ruggero}, game designer, che guida la parte sui
fondamenti e le domande finali; \textbf{Beatrice}, progettista di giochi
pervasivi, che interviene soprattutto nella prima metà. Gli studenti iscritti
sono \textbf{nove}; uno è assente e un altro arriva a lezione iniziata. Ci sono
una pausa di una decina di minuti (si riprende alle 18:15) e, dopo le 19:15,
un'ultima mezz'ora di domande con il solo Ruggero.

Tre avvertenze di lettura. La prima: la trascrizione è automatica e storpia
molti nomi propri, compresi quelli dei docenti e dell'accademia (``Dive'',
``IVE'', ``Accademia Live'', ``Ai''); dove la ricostruzione è sicura, o molto
probabile, il capitolo usa la forma corretta e lo segnala in nota; dove non lo
è, riporta la forma della trascrizione e lo dice. La seconda: durante la pausa la
trascrizione registra, ripetuta in loop per decine di righe, una frase senza
senso che non ha nulla a che fare con la lezione, e che il capitolo ignora. La
terza: in molti passaggi più docenti parlano uno dopo l'altro senza che la
trascrizione segni il cambio di voce; il capitolo attribuisce le parole a chi le
pronuncia quando il contesto lo rende chiaro (chi parla di sé, chi viene
chiamato per nome, chi ``passa la parola''), e resta generico altrimenti.

Quanto alla data, la lezione si colloca con ogni probabilità \textbf{nella prima
metà di settembre del 2026}\footnote{La data non è mai detta. La si deduce da
alcuni indizi: il form iniziale va compilato ``entro e non oltre il 15
settembre''; i gruppi verranno annunciati il 25 settembre; le date del
calendario (sabato 3 ottobre, lunedì 12 ottobre, sabato 7 e 14 novembre,
venerdì 4 dicembre) cadono nei giorni della settimana indicati soltanto nel
2026; un docente parla di un argomento ``molto complesso nel 2026''; e fa
ancora molto caldo (``le temperature sono ancora indecenti''). È una
deduzione, non un'informazione esplicita.}: i riferimenti a ``oggi'',
``quest'estate'' e ``l'anno scorso'' vanno letti rispetto a quel momento.

Il percorso del capitolo segue quello della lezione:
\begin{enumerate}
  \item l'avvio, i \textbf{docenti} e le \textbf{presentazioni} dei partecipanti
        (sezioni da~\ref{a1:sec:avvio} a~\ref{a1:sec:partecipanti});
  \item i \textbf{materiali} e gli strumenti del corso: il form, la cartella
        condivisa, Discord (sezione~\ref{a1:sec:materiali});
  \item lo \textbf{scopo} e il \textbf{metodo} del corso, le regole di
        comportamento e il carico di lavoro (sezioni~\ref{a1:sec:scopo}
        e~\ref{a1:sec:regole});
  \item il \textbf{project work} e il \textbf{calendario} (sezione~\ref{a1:sec:project});
  \item la regola sull'\textbf{intelligenza artificiale} e la lunga discussione
        che ne segue (sezione~\ref{a1:sec:ia});
  \item le \textbf{figure professionali} dell'industria del gioco analogico,
        con un approfondimento su marketing e sales
        (sezioni~\ref{a1:sec:figure} e~\ref{a1:sec:marketing});
  \item la differenza tra \textbf{game design} e \textbf{game development}
        (sezione~\ref{a1:sec:design-development});
  \item il \textbf{mercato}: Italia ed estero (sezione~\ref{a1:sec:mercato});
  \item i \textbf{fondamenti di design}: perché serve una terminologia comune,
        \emph{game} e \emph{play}, il framework \textbf{MDA} (e il DPE), le
        \textbf{parole vietate}, chi crea, chi gioca e chi facilita, il
        \textbf{playtest} e l'\textbf{iterazione} (sezioni
        da~\ref{a1:sec:terminologia} a~\ref{a1:sec:playtest});
  \item \textbf{Tabletop Simulator} e il primo task (sezione~\ref{a1:sec:tts});
  \item le \textbf{domande finali}: il tema del project work, il bilanciamento,
        i ruoli nei team, la pubblicazione, il gioco di ruolo in Italia, le
        \emph{house rule} di Monopoly, i giochi ``comunicati male''
        (sezioni~\ref{a1:sec:domande-daniele} e~\ref{a1:sec:domande}).
\end{enumerate}

% ---------------------------------------------------------------------
\section{L'avvio della lezione}
\label{a1:sec:avvio}

La lezione comincia con qualche minuto di attesa. ``Siamo in nove in classe'', e
si aspetta che arrivino tutti; uno dei docenti scherza sul fatto che l'unica
persona che arriverà ``sempre in ritardo di venerdì'' è Gianmarco. Mancano
Giordano e Tiziana; Tiziana arriva, mentre la segreteria fa sapere poco dopo che
Giordano non potrà esserci e vedrà la registrazione. Si aspetta anche Beatrice,
la terza docente, che arriva a sua volta dopo pochi minuti.

Nell'attesa, Daniele anticipa il programma della giornata e qualche
informazione pratica:
\begin{itemize}
  \item oggi è una giornata di \textbf{introduzione}: come funziona il corso, chi
        sono i docenti, le presentazioni degli studenti, le \textbf{metodologie}
        e il \textbf{calendario}. Gli studenti avranno ``un fracco di domande'', e
        alla fine della lezione c'è ``proprio uno slot apposta'' per le domande;
  \item le lezioni sono \textbf{registrate}: di solito la registrazione viene
        consegnata dopo circa \textbf{12--24 ore} sul profilo personale della
        piattaforma dell'accademia, che va attivato con la segreteria (Tiziana,
        già studentessa, probabilmente ce l'ha già);
  \item i docenti ``cuore'' del corso sono tre (Daniele, Ruggero e Beatrice),
        ``soprattutto Ruggero''; a loro si aggiungeranno tre docenti ospiti, per
        meno ore, probabilmente una lezione a testa, ``tranne Marta, che avrà due
        lezioni quest'anno''.
\end{itemize}
E una battuta sul caldo: ``spero che voi siate carichi, anche se le temperature
sono ancora indecenti''. ``Carichi moderatamente, per quello che potete''.

Poi Daniele condivide lo schermo (``non solo le slide, ma anche un po' di altre
cose'') e apre ufficialmente il corso: ``benvenuti e benvenute nella seconda
edizione del corso di Analog Game Design. Questa è la prima lezione''. Il
percorso è di circa \textbf{tre mesi}, fino a dicembre, e durante il corso i
docenti cercheranno di far provare agli studenti ``cosa significa fare Analog
Game Design'' con un \textbf{project work}, un progetto di gruppo di cui si
spiegheranno più avanti la struttura e le richieste.

% ---------------------------------------------------------------------
\section{I docenti}
\label{a1:sec:docenti}

``Anzitutto partiamo da delle presentazioni doverose''. Ognuno dei tre docenti
principali si presenta con una slide che riporta i contatti e una selezione dei
lavori fatti.

\subsection{Daniele: narrativa, fiction e gioco di ruolo}

\textbf{Daniele}\footnote{La trascrizione riporta ``Daniel Cusetto''; nel
resto della lezione tutti lo chiamano ``Daniele''. Il cognome non è
ricostruibile con certezza.} ha già conosciuto quasi tutti gli studenti, perché
ha fatto loro il colloquio di ammissione. È \textbf{narrative designer}, o meglio
\emph{game narrative designer}; insegna nell'accademia che ospita il
corso\footnote{La trascrizione storpia il nome dell'accademia in molti modi
(``Dive'', ``IVE'', ``Accademia Live'', ``Ai'', ``AIVA''). Con ogni probabilità
si tratta dell'\textbf{AIV}, l'Accademia Italiana Videogiochi, che ha sede a
Roma: lo fanno pensare i corsi triennali di game design e narrative design citati
dai docenti e il fatto che una studentessa racconti di essersi trasferita a Roma.
È una ricostruzione, non un'informazione esplicita; nel seguito il capitolo
parla semplicemente dell'``accademia''.} ma lavora soprattutto come \textbf{brand
advisor}. Nel corso si occuperà della parte di \textbf{gioco di ruolo},
\textbf{narrazione} e \textbf{fiction}.

I suoi contatti, ``in quest'ordine'', sono email, Telegram e Discord. Ci tiene a
dare alcune indicazioni precise:
\begin{itemize}
  \item se si vuole avere la certezza di una risposta, bisogna scrivergli su
        \textbf{Telegram}, perché lì ha tutti i suoi clienti;
  \item Discord lo guarda di meno; l'\textbf{email} è per le cose ufficiali;
  \item non bisogna contattarlo altrove: non su Facebook, non su LinkedIn, non su
        Instagram, ``perché li vedo molto di meno''.
\end{itemize}
La slide mostra anche alcuni giochi, attività e brand su cui ha lavorato negli
anni, ma non si dilunga: ``nel caso ne parleremo''.

\subsection{Ruggero: game design e ricerca e sviluppo}

\textbf{Ruggero} (spesso chiamato ``Ruger'') ha fatto il \textbf{game designer}
per varie aziende e oggi è game designer e si occupa di \textbf{ricerca e
sviluppo} per un'azienda del settore\footnote{La trascrizione dice ``per MS'';
il nome dell'azienda non è ricostruibile con certezza. Più avanti Ruggero cita
tra i giochi su cui ha lavorato il gioco da tavolo ufficiale di
\emph{Cyberpunk}, \emph{Marvel United} e \emph{Status}, un gioco sulla
geopolitica.}. Ha un background ``misto'': nel suo portfolio ci sono giochi da
tavolo, giochi di ruolo, giochi analogici, LARP, videogiochi, giochi in realtà
virtuale. Si occupa anche di \textbf{divulgazione} sul gioco e di analisi di
\textbf{game studies} e game design. Più avanti, nelle domande finali, dirà di
avere ``50 e più giochi prodotti e progettati'' e di insegnare da vari anni.

Lo si può contattare liberamente via Gmail e Telegram, e cerca di rispondere
\textbf{entro 24 ore}, salvo imprevisti o eventi: ``sentitevi in libertà di
scrivermi se avete dubbi, senza problemi''. Quando arriva a parlare dei propri
giochi, la slide mostra ``un mix di varie cose, dalle più famose alle meno
famose'', ma Ruggero si imbarazza a presentarli, tra le risate dei colleghi.

\subsection{Beatrice: i giochi pervasivi}

\textbf{Beatrice} (``Bea'') si occupa di \textbf{giochi pervasivi}
(\emph{pervasive games}). ``Se non li avete mai sentiti nominare, non
preoccupatevi'': questo corso è ``uno dei pochi, forse l'unico in Italia'' che
li racconta esplicitamente. Li si vedrà nelle prossime lezioni, ma la definizione
di base è questa: i giochi pervasivi \textbf{utilizzano gli spazi e gli oggetti
della vita quotidiana e li trasformano in spazi e oggetti di gioco}.

Beatrice lavora come \textbf{freelance} e si occupa di diversi tipi di
progettazione: lavora sia con le \textbf{aziende} sia con gli \textbf{eventi dal
vivo}. La sua attività diretta si trova sotto l'etichetta \textbf{Minime
Differenze}\footnote{La trascrizione riporta il nome dell'etichetta come
``Minime Differenze'' e l'indirizzo del sito come ``mininedifferenze.com''; la
grafia esatta del sito non è verificabile dalla trascrizione.}, mentre i lavori
per le aziende sono ``sparpagliati'' sui siti delle varie aziende. Molte delle
cose mostrate nella sua slide sono attività proposte al pubblico ``con un
biglietto, con un gioco in formato evento'': un tema che si toccherà nel corso,
per capire come \textbf{comunicare} i propri giochi in modo che siano
comprensibili, accessibili per il pubblico e \textbf{vendibili}.

Ammette un ``grosso limite'' con i messaggi: è un po' lenta a rispondere, ``ma
rispondo sempre a tutto''. La si può contattare per email (cognome punto nome) o
su Telegram. È ``molto curiosa di vedere cosa produrrete''.

\subsection{I docenti ospiti}

Daniele presenta poi i tre docenti che il corso avrà ``solo per alcune
lezioni'', tutte verso \textbf{novembre}:
\begin{itemize}
  \item \textbf{Diego Ricchiuti}, direttore del corso di Game Design
        dell'accademia, farà una lezione sull'\textbf{RNG} (\emph{random number
        generation}, la casualità nei giochi): ``vi parleremo di RNG quando sarà
        necessario. Comunque avrete una lezione con lui, sappiatelo'';
  \item \textbf{Marta}\footnote{La trascrizione riporta il cognome come
        ``Cercasassi''; la grafia non è verificabile.} farà \textbf{due lezioni
        sul development}: quest'anno ``abbiamo deciso di fare doppietta'',
        perché l'anno scorso il suo intervento è stato ``molto, molto utile'';
  \item \textbf{Flavio}\footnote{La trascrizione riporta il cognome come
        ``Mortarino'' e, più avanti, ``Mortalino'', e il nome della sua casa
        editrice come ``Grampibert''; nessuno dei due è ricostruibile con
        certezza.} farà la lezione di \textbf{game production}; è un editore, e
        di solito partecipa anche alla tavola rotonda con gli editori della
        giornata finale.
\end{itemize}
Le slide sono già online e gli studenti avranno accesso alla cartella condivisa:
``non vi preoccupate di queste cose, ne parleremo a breve''. Ma prima di parlare
del corso, ``vogliamo conoscere voi''.

% ---------------------------------------------------------------------
\section{Le presentazioni dei partecipanti}
\label{a1:sec:partecipanti}

Le domande per la presentazione sono semplici, da affrontare ``a turno, senza
fretta, con calma'':
\begin{itemize}
  \item nome e occupazione (``se volete dirla'');
  \item perché si vogliono fare giochi analogici, e a che livello: che cosa ha
        portato lo studente al corso;
  \item che tipo di esperienze o conoscenze si hanno già nell'ambito, se si è già
        fatto qualcosa oppure no.
\end{itemize}
``Chi vuole partire senza nessun tipo di ritegno e di vergogne?''. Per far
parlare gli studenti guardandoli, Daniele interrompe la condivisione dello
schermo e ricorda lui le domande quando serve. Le presentazioni sono interessanti
anche per i commenti dei docenti, che anticipano diversi principi del corso.

\subsection{Tiziana}

Tiziana fa l'\textbf{infermiera}, ma nella vita ha fatto molti lavori:
impiegata, insegnante e ``tanti altri lavoretti''. Gioca ``da sempre'': il gioco
fa parte anche della sua vita familiare, perché suo fratello fa il programmatore
di videogiochi. Ha già \textbf{organizzato} giochi: una caccia al tesoro per far
conoscere il territorio, due scene con un delitto (giochi con omicidio da
risolvere), giochi per intrattenere le persone. Ha riscoperto di recente i giochi
da tavolo, arrivando a \emph{Carcassonne} con famiglia e amici; poiché è
difficile trovare persone con voglia di giocare nella propria cerchia ristretta,
ha scoperto il mondo delle \textbf{associazioni} ludiche, ``e lì è stato
un'esplosione'': prima in Veneto (è di Vicenza), poi a Roma, dove si è
trasferita.

Ha in mente un progetto ambizioso: far conoscere una certa realtà con la modalità
del \textbf{gioco di ruolo}, o comunque proporre un gioco suo.

\paragraph{Il commento: gli strumenti, non il gioco dei sogni.} Ruggero coglie
l'occasione per un'avvertenza che vale per tutti, visto che molti studenti
arrivano con idee già pronte. Nel corso si lavorerà a un progetto, ma
l'obiettivo dei docenti è dare \textbf{strumenti} da applicare anche
\emph{fuori} dal progetto, ai giochi che gli studenti vorranno fare un domani:
``che non siano esercizi per il corso, ma per voi in primis''. Quindi: ``non
rimaneteci male se non riuscirete a lavorare al vostro gioco dei sogni in questo
corso''. Anche perché ``sarebbe impossibile, in due mesi, fare il gioco dei
sogni''. Il corso darà invece la possibilità di fare un'esperienza ludica
insieme e di imparare a replicarla sulle proprie idee: ``le idee non buttatele,
anzi, il contrario: tenetele''.

\subsection{Nicola}

Nicola è uno \textbf{studente}. Sei mesi prima ha fatto un colloquio con Daniele
perché stava valutando di passare all'accademia, e ha deciso, ``nel dubbio'', di
cominciare da questo corso breve per capire come funziona la formazione;
eventualmente si iscriverà al corso di game design vero e proprio. Gli interessa
in particolare il lato dell'esperienza del giocatore, su cui è ``molto critico''
(qui usa la parola ``divertimento'', e Daniele lo ferma con un ``ne parleremo'':
è una delle parole che il corso vieta, come si vedrà nella
sezione~\ref{a1:sec:parole-vietate}).

Ha 21 anni ed è \textbf{giocatore di ruolo da circa otto}. Prima ancora di
conoscere i giochi di ruolo si era creato un sistema tutto suo, ``pessimo,
disgustoso, terribile'', un sistema a percentuale con il d100 ``terrificante'';
poi ha scoperto i giochi di ruolo da tavolo, i classici (\emph{Dungeons \&
Dragons}, \emph{Daggerheart}, \emph{Vampire}\footnote{La trascrizione riporta
``DD'', ``Daggerart'' e ``Vampire'': quasi certamente \emph{Dungeons \&
Dragons}, \emph{Daggerheart} e \emph{Vampire: The Masquerade}.}). Gioca anche a
\textbf{card game} e ad altri giochi sociali, pur avendo sempre avuto difficoltà
a trovare gruppi, ``avendo vissuto in un buco dimenticato dal mondo''. Come
giocatore è ``l'eterno narratore'' del gruppo, quello che prepara la storia e si
costruisce le regole. Come autore ha fatto soprattutto giochi su carta, ``perché
sono incapace a programmare, per adesso''.

Daniele osserva che il tema del giocatore e del suo rapporto con il gioco è ``un
argomento importante'', anzi ``un vaso di Pandora enorme'' (espressione che
tornerà molte volte nella lezione), e che il corso tratterà anche che cos'è
l'analog game design in generale e quanti tipi di prodotti esistono.

\subsection{Rita}

Rita\footnote{Daniele la chiama prima ``Nino'' (probabilmente un errore della
trascrizione) e poi la ringrazia come Rita.} ha 30 anni e si collega da una
regione del Sud, dove ``ancora si muore di caldo''. Ha studiato
all'\textbf{Accademia di Belle Arti di Napoli}, \textbf{design della
comunicazione visiva}: si occupa dell'aspetto più imprenditoriale della
progettazione artistica. Non volendo progettare per un'impresa (``fare loghi,
robe del genere''), vorrebbe applicare la propria creatività alla realizzazione
di giochi analogici.

Non è una giocatrice di ruolo e non lo è mai stata: è una giocatrice
\textbf{tradizionale}, cresciuta con Monopoly, la dama e giochi tradizionali
cinesi. Ha già \textbf{idee, prototipi e progetti} realizzati, e vuole capire
come renderli \textbf{prodotti per la vendita} e se questo è il campo adatto a
lei; magari anche confrontarsi con i colleghi. La sua progettazione è orientata
al \textbf{sociale} e all'\textbf{educativo}: ``responsabilizzare un po' le nuove
generazioni su come sta andando il mondo e su come potrebbe andare, visto che è
tutto nelle nostre mani''.

\paragraph{Il commento: siate curiosi.} ``Pane per i tuoi denti'', le risponde
Daniele, soprattutto per i LARP e i giochi pervasivi: ``vi apriremo un po' la
testa da questo punto di vista''. E aggiunge un consiglio per tutti: ``siate
curiosi'', perché imparare da altri ambiti dà moltissimo. Ognuno è lì per fare
qualcosa di preciso, ``però siate aperti anche ad altre esperienze''.

\subsection{Tancredi}

Tancredi ha 21 anni e non ha ancora grande esperienza lavorativa: qualche
``piccolo progettino'' che non è andato come voleva. Gli interessa il mondo dei
giochi, digitali e non, con un approccio \textbf{educativo} e \textbf{sociale}.
È arrivato soprattutto dai videogiochi, a cui gioca da quando era piccolo; ai
giochi da tavolo si è avvicinato solo nell'ultimo anno, grazie a una realtà
ludica, perché in famiglia ``non siamo molto da giochi''. Ha realizzato qualche
piccolo videogioco con vari engine: prima quelli senza un vero linguaggio di
programmazione, ``più a nodi'', poi qualcosa di più complicato, e ora sta
lavorando a un piccolo prototipo.

\paragraph{Il commento: niente di quello che vedrete è sprecato.} ``Non butterai
via niente di analog game design'', gli risponde Daniele, ``perché il 99 per cento
delle cose che vedremo si ritrovano quasi uno a uno nel mondo dei videogiochi''.
Tanto che nei corsi triennali dell'accademia (anche in quello di narrative
design, dove insegna Ruggero) il primo anno si parte dai \textbf{giochi da tavolo
e di ruolo}, proprio per spiegare meccaniche più complesse che servono poi per i
videogiochi. Ruggero aggiunge che vale anche il contrario: alcuni dei
\textbf{framework} di game design che si vedranno nelle lezioni tecniche vengono
dal videogioco e sono stati riapplicati all'analogico. ``Evitare la
crossmedialità qui è impossibile''. Anzi, tutte le conoscenze ``mediatiche
adiacenti'' (musica, letteratura, cinema) saranno utili.

\subsection{Lorenzo Vergata}

I Lorenzo in classe sono due; il primo a presentarsi è \textbf{Lorenzo
Vergata}, 26 anni, \textbf{dottorando al Politecnico di Milano}. Gioca ai
videogiochi da quando aveva due anni e mezzo, ma ha giocato al suo primo gioco da
tavolo ``credo nove mesi fa''. Gli fa piacere sentire che molti in classe sono
interessati ai giochi \textbf{sociali} ed \textbf{educativi}, un settore che
interessa anche a lui e che ``può essere esplorato molto'', anche con qualche
aiuto dell'intelligenza artificiale per esperienze più ricercate.

Di esperienza nell'analogico dice di avere fatto ``il minimo sindacale''. Ma
racconta che diciotto anni prima, da bambino, era ``più game designer di
adesso'': prendeva un gioco e ne cambiava le regole, anche al biliardino (``no,
ora introduciamo questa regola''). Gli piace sperimentare, creare regole e
vedere che cosa accade in sé e negli altri. Soprattutto vorrebbe ``una chiave di
lettura per poter rimanipolare la realtà'': tre settimane prima è stato a fare la
\textbf{vendemmia in Francia} e ha visto tantissime dinamiche che potrebbero
diventare un gioco da tavolo o un videogioco. Per esempio: finire contro un rovo
pieno di semi che si attaccano ai vestiti lo rallentava nella raccolta
dell'uva, mentre in un altro filare era più rapido. ``Si può fare pure una
dinamica'', e vorrebbe gli strumenti per farne giochi.

Arriva subito la battuta: una ``vendemmia gamificata'' da proporre in Oltrepò
Pavese (``io sono dell'Oltrepò Pavese, quindi devo per forza bere vino'')\footnote{La
trascrizione dice ``nell'altro papavese'' e non chiarisce chi pronunci la
battuta.}. Daniele conferma che quest'anno il corso avrà ``tante cose sul
sociale'': ``ci piace parecchio''.

\subsection{L'altro Lorenzo}

Il secondo Lorenzo (``secondo, ma non meno importante'') è uno \textbf{scrittore
di racconti brevi} e sta scrivendo il suo primo romanzo. Gli è sempre piaciuto
raccontare storie, e da più grande scriverle; gioca ai videogiochi ed è ``un
amante delle storie e della cultura pop''. Negli ultimi anni è entrato nel mondo
del \textbf{gioco di ruolo da tavolo} e vorrebbe unire le due passioni, le
storie e il gioco di ruolo, magari facendone un lavoro; gli interessa anche lo
\emph{strategic writing}. Daniele gli ricorda che ci sarà una lezione proprio
sulla \textbf{fiction}.

\subsection{Angelo}

Angelo si collega senza videocamera dal lavoro: fa il \textbf{birraio} (in un
piccolo birrificio in Brianza, a Bernareggio, come racconterà dopo la pausa) e,
con il caldo, ``vi risparmio l'orrore''. Per molto tempo era convinto che i
giochi da tavolo fossero solo i classici (Monopoly, Risiko) o i party game; poi,
quattro o cinque anni fa, ha ``sbattuto la faccia'' contro giochi davvero
appassionanti, con una specie di narrazione dietro che li rendeva ancora più
coinvolgenti, e non ha più smesso. È uno che ``ha sempre un botto di idee, ma poi
non mette in pratica assolutamente nulla'': cercava un corso che gli desse le
basi e soprattutto lo \textbf{motivasse} a provare a fare qualcosa.

\paragraph{Il commento: bisogna buttarsi.} Daniele racconta che anche l'anno
prima un suo caro amico si era iscritto proprio per questo motivo. Il corso è
breve, ma cerca di portare gli studenti ``da 0 a 100'' e di far realizzare un
\textbf{prodotto} da presentare alla fine. Il punto è che i giochi ``bisogna
farli, bisogna buttarsi''. E, venendo dalla narrazione, Daniele lo paragona allo
scrivere: si scrive e si riscrive, e allo stesso modo si fa \textbf{playtest}, si
rifà, si aggiorna. ``Non è che si fanno così''.

\subsection{Gianmarco (dopo le slide)}

Gianmarco arriva a lezione iniziata e si presenta più tardi, alla fine della
prima parte, dopo che i docenti hanno chiuso le slide organizzative e la
discussione sull'intelligenza artificiale (``Gianmarco, ti giuro, ti facciamo
presentare''). Lavora nel \textbf{team di innovazione} di un'azienda italiana.
Vuole fare giochi perché è una sua passione, che lo gratifica, e in passato ci ha
provato ``in maniera scomposta''. Cerca un \textbf{metodo} e un confronto con
persone più esperte: ``studiare solo dei manuali è molto sterile'', ``stare
chiusi nel pertugio non aiuta'', ed ``essere autodidatta ha estremamente i suoi
limiti''. Come esperienza è un giocatore e ha letto ``un po' di materiale, di
blog, di manuali'', ``niente di rilevante''.

Manca solo Giordano, che si presenterà alla lezione successiva.

\paragraph{Uno sguardo d'insieme.} La classe è molto varia: un'infermiera,
una designer della comunicazione visiva, un birraio, un dottorando, uno
scrittore, un impiegato nell'innovazione, due studenti ventunenni. Alcuni
vengono dal gioco di ruolo, altri dai giochi tradizionali, altri dai
videogiochi; quasi tutti hanno scoperto i giochi da tavolo moderni da pochi anni
o da pochi mesi. E almeno tre studenti dichiarano un interesse per i giochi
\textbf{sociali} ed \textbf{educativi}. Come si vedrà nella
sezione~\ref{a1:sec:scopo}, i docenti considerano questa varietà una risorsa.

% ---------------------------------------------------------------------
\section{I materiali e gli strumenti del corso}
\label{a1:sec:materiali}

Prima delle slide sul corso, Daniele sbriga un po' di ``clean up casalingo''.

\paragraph{Il form iniziale.} C'è un \textbf{form} da compilare ``entro e non
oltre il 15 settembre'': il QR code è nelle slide, e la segreteria lo ricorderà
su Discord. Richiede circa 10--15 minuti; si può fare subito o in giornata, ``e
ve lo togliete di torno''. È un form standardizzato per tutti i corsi (``non
siamo io e Ruggero a farvi passare la tortura del test''), diviso in due parti:
\begin{itemize}
  \item una parte di \textbf{conoscenza del candidato}: quanto tempo si ha a
        disposizione per il corso, come ci si organizza con il lavoro, che cosa
        piace giocare. È la parte che interessa ai docenti, perché li aiuta a
        organizzare le lezioni;
  \item una parte di \textbf{test logico}, che i docenti non hanno scritto (``io
        non sono responsabile del form'', ``no, no, nemmeno io'').
\end{itemize}

\paragraph{La cartella condivisa.} Tutti gli studenti vengono aggiunti alla
cartella del corso su \textbf{Google Drive}, che si trova tra i file
``condivisi con me''. Gli studenti ne sono \textbf{gestori dei contenuti} e
possono caricarci quello che vogliono. Dentro ci sono:
\begin{itemize}
  \item le \textbf{attività in classe};
  \item le \textbf{dispense};
  \item i \textbf{materiali didattici}, con le slide (quelle della lezione sono
        già lì);
  \item le cartelle dei \textbf{project work} futuri;
  \item i \textbf{task}, che ``saranno pochi, ma ci saranno'';
  \item il \textbf{calendario studenti}, ``importantissimo'', con due schede: il
        calendario delle lezioni e quello del project work.
\end{itemize}

\paragraph{Discord.} Il corso ha anche un \textbf{server Discord}, a cui tutti
sono stati invitati; ci sono i docenti (Daniele, Ruggero, Diego e altri), che si
possono taggare. Di solito è lì che la segreteria pubblica le informazioni, per
esempio se una lezione ``dovesse per caso essere cancellata''. Nelle domande
finali Ruggero aggiungerà che c'è un canale apposito per il corso, che gli
studenti possono usare liberamente per organizzarsi tra loro.

\paragraph{Le reaction di Zoom.} Durante le lezioni su Zoom si possono usare le
\textbf{reaction} (il pollice in alto, il cuoricino, la mano alzata): un modo
rapido per far capire ai docenti se si è capito, se si vuole intervenire. Nel
frattempo arrivano i primi pollici.

% ---------------------------------------------------------------------
\section{Lo scopo e il metodo del corso}
\label{a1:sec:scopo}

La prima serie di slide (che Daniele, per un lapsus, presenta come ``cos'è il
corso di pulire il pavimento di casa'') spiega che cosa il corso vuole fare e
come.

\subsection{Diventare autori analogici}

Lo \textbf{scopo del corso} è fornire le competenze per \textbf{diventare autori
analogici} e stimolare la creatività e l'\textbf{autorialità} di ciascuno: ``voi
imparerete a diventare autori di giochi analogici''. La differenza tra
\emph{game designer} e \emph{game developer} verrà chiarita più avanti (sezione~\ref{a1:sec:design-development}),
ma la parola chiave qui è \emph{autorialità}: ``vogliamo voi, vogliamo le vostre
idee''. Magari non si potrà usare l'idea del gioco della propria vita, ma si
potranno portare nelle lezioni le proprie idee e le proprie sensibilità.

Gli obiettivi elencati nella slide sono questi:
\begin{itemize}
  \item dare \textbf{maggiore consapevolezza nel design e nello sviluppo}. In
        Italia, purtroppo, ``si fa design a volte in maniera non consapevole''.
        Consapevole vuol dire avere almeno un'introduzione alle basi del design e
        dello sviluppo, e fare le cose ``con metodo, con raziocinio, con
        consapevolezza'';
  \item fornire gli \textbf{strumenti essenziali} per un design consapevole;
  \item stimolare la \textbf{sperimentazione} e la \textbf{trasversalità
        mediatica}. Non si chiede di fare ``il \emph{Root} della situazione o il
        \emph{Carcassonne} della situazione''\footnote{La trascrizione dice ``il
        lumaven della situazione o il carcassone della situazione''; il primo
        titolo non è ricostruibile con certezza. Il senso è chiaro: non si
        chiede di rifare un grande classico del gioco da tavolo.}, ma un
        project work \textbf{calibrato sulla durata del corso}, e però
        \textbf{sperimentale} e ``assolutamente transmediale'';
  \item portare \textbf{numerosi esempi} di tipologie di gioco, dall'analogico al
        tavolo, al pervasivo ``e oltre''. Lo si vedrà nella lezione sui
        \textbf{casi studio} del venerdì successivo: ``voi non avete idea di
        quante cose esistono, bellissime, sconosciutissime e favolose''. Daniele
        confessa che preparando e aggiornando le slide i colleghi gli hanno
        mostrato cose che lui stesso non sapeva nemmeno che esistessero;
  \item dare gli strumenti per \textbf{analizzare} i giochi altrui. Ruggero lo
        sottolinea: se il corso riesce nel suo intento, gli studenti sapranno
        non solo fare i propri giochi ma anche analizzare giochi già fatti, che è
        ``una delle skill importanti anche nella parte di ricerca e sviluppo''.
        Si vedranno casi studio e giochi pubblicati, sempre con l'obiettivo di
        capire \textbf{che cosa fanno e perché lo fanno}, per poi fare meglio il
        proprio lavoro;
  \item imparare a \textbf{fare giochi insieme}. Esistono gli autori che lavorano
        da soli (i \emph{solo developer}), ``ma sappiate che anche l'autore che
        lavora da solo poi si deve interfacciare con il pubblico, con un editore,
        magari con un editor, con un publisher''. ``Non lavorerete mai da soli'':
        bisogna imparare a lavorare con gli altri.
\end{itemize}

\subsection{Un corso intersezionale}

Il corso è \textbf{intersezionale}: \textbf{non ci sono lezioni dedicate a una
singola tipologia di gioco}. Non ci sarà ``la lezione sul gioco da tavolo'',
``la lezione sul gioco di ruolo'', ``la lezione sul gioco pervasivo''. Ci saranno
lezioni su temi diversi, ciascuna pensata per essere utile a \textbf{tutte e tre
le sfaccettature} del corso; alcune saranno più importanti per il gioco da
tavolo, altre per il pervasivo, altre per il gioco di ruolo, ma nessuna sarà
esclusiva. I docenti vogliono ``contaminazioni e applicazioni assolutamente
trasversali''.

\subsection{Un corso cross-generazionale}

È anche un corso \textbf{cross-generazionale}: i gradi di esperienza di gioco
quest'anno sono abbastanza allineati, ma le esperienze di vita, le età, le
estrazioni sociali sono varie. ``Ed è un pro, non è un contro. Più diversificata
è la classe, più imparate meglio anche voi''. La slide ha una piccola nota per
fasce d'età:
\begin{description}
  \item[Under 30] Concentratevi sulla vostra \textbf{crescita}. Evitate di
        confrontarvi con le persone più grandi che magari ``ce la fanno meglio''
        in alcune cose.
  \item[Over 30] Concentratevi sul \textbf{consolidamento}, ma siate
        \textbf{umili}: i docenti diranno cose che faranno ``storcere il naso''.
        ``Bisogna un po' reimparare per imparare'': bisogna sapersi adattare.
\end{description}
Ruggero contestualizza: lo dicono perché in passato hanno visto studenti che,
essendo più adulti o avendo giocato di più, pensavano che l'esperienza
\emph{di gioco} si trasformasse automaticamente in competenza \emph{di
creazione} di giochi. Non è così. E non bisogna sottovalutare chi gioca meno:
proprio perché gioca cose diverse, per risolvere un problema di design specifico
``è una ricchezza incredibile avere qualcun altro che riesce a uscire dalla
scatola''.

\subsection{Un corso pratico e adattabile}

Il corso è \textbf{pratico}: gli studenti si concentreranno su un progetto di
gruppo, il \textbf{project work}, lavorando in team. Ed è \textbf{adattabile}: il
programma è definito, ma ``si aggiorna sui vostri stimoli, priorità e
necessità''. Se durante il corso i docenti vedono che serve più una cosa
dell'altra, sono ``prontissimi'' a dare più spazio a un argomento invece che a
un altro.

\subsection{Inquiry-based learning: domandate}

Il metodo è quello dell'\textbf{inquiry-based learning}, l'apprendimento basato
sulle domande. Di solito le persone dicono ``ho una domanda, magari sarà
stupida''. ``Non esistono domande stupide. Pensate che ci sia una domanda
stupida? Fatela, alzate la mano''. Vanno bene anche le
\textbf{provocazioni}: se qualcosa non torna, è meglio chiedere. La condivisione
è fondamentale.

I docenti danno \textbf{metodi, dati e conoscenze}, ma lasciano trovare le
risposte agli studenti: ``vi lasciamo un pochettino annegare nel mare prima di
darvi la soluzione, in modo tale che la possiate trovare da soli''. Di
conseguenza \textbf{non ci saranno interrogazioni sulle definizioni}: interessa
vedere \textbf{come} gli studenti applicano ciò che imparano. Se una
definizione non viene in mente, non è un problema; l'importante è saperla usare.

\subsection{Feedback, non valutazione}

``Sì al feedback, no alla valutazione''. I docenti danno \textbf{feedback
verbali}, anche sulle attività da fare a casa (i task, che non saranno tanti).
Il corso è calibrato per \textbf{studenti lavoratori}: c'è un po' da lavorare
fuori dalle lezioni, ma in modo calibrato.

Il principio è \textbf{fallire oggi per imparare per il futuro}: \emph{fail
forward}, \emph{fail to learn}. I docenti lasciano provare e fallire almeno una
volta prima di dare la soluzione: ``il classico metodo Montessori, quello
originale, non quello travisato''.

\subsection{Le soft skill e i criteri di valutazione}

Grande attenzione va alle \textbf{soft skill}: i lavori di gruppo, lo scambio
propositivo, l'atteggiamento positivo, la propositività, la proattività.
``Lavoro di gruppo'' non significa stare ogni giorno in sala studio con gli
altri: va benissimo che un gruppo si organizzi dandosi priorità e tempi in base
alla disponibilità di ciascuno. L'importante è \textbf{collaborare}, essere
\textbf{professionali}, rispettarsi a vicenda, rispettare le idee proprie e
altrui e lasciare spazio agli altri, perché i progetti ``non saranno solo
vostri, saranno di tutte le persone del vostro gruppo''.

La valutazione del project work non riguarda solo la \textbf{documentazione}
prodotta, ma anche:
\begin{itemize}
  \item la \textbf{presentazione} finale;
  \item l'\textbf{autovalutazione};
  \item la \textbf{crescita personale};
  \item l'\textbf{atteggiamento} tenuto nel corso.
\end{itemize}
Daniele cita: ``da ognuno secondo le sue capacità, a ognuno secondo i suoi
bisogni''. Lavorare in gruppo non vuol dire che, in tre, ciascuno lavori
esattamente il 33\% del tempo. Ai docenti serve piuttosto capire \textbf{chi fa
che cosa} nel gruppo, e che il processo interno sia \textbf{trasparente}.

\paragraph{Lavorare in cloud.} Per lo stesso motivo il consiglio è di
\textbf{non lavorare in locale}, ma direttamente nelle cartelle condivise del
project work e dei task. Si è liberi di fare altrimenti, ma lavorando in cloud
si dà ai docenti la possibilità di vedere \emph{come} si lavora, il proprio
``processo mentale''.

% ---------------------------------------------------------------------
\section{Le regole di comportamento e il carico di lavoro}
\label{a1:sec:regole}

Alcune slide sono di quelle che si ripetono in ogni corso (``\emph{repetita
iuvant}'', storpiato per scherzo in ``ripetita Juventus'').

\paragraph{Comportamento.} Prestare attenzione durante le lezioni; interagire;
portare rispetto; essere proattivi; alzare la mano e fare domande.
\textbf{Tolleranza zero} per pregiudizi, maleducazione, mancanza di rispetto,
\emph{cheat} e scorciatoie. Daniele ammette che ``cheat e scorciatoie'' viene da
altri corsi dove il tema è più importante, ma Ruggero lo ferma: ``non
sottovalutate questa cosa'', perché implica anche \textbf{non far fare a ChatGPT
i task}. Daniele ha una slide apposta, ``perché è un discorso molto complesso nel
2026'' (sezione~\ref{a1:sec:ia}).

\paragraph{Le lezioni sono intense.} Le lezioni possono essere \textbf{dense},
``soprattutto le mie, sappiatelo: io sono logorroico'', avverte Daniele. Bisogna
chiedere spiegazioni se serve, prendere appunti, rivedere le registrazioni. Le
lezioni sono ``delle scintille'': modi per far esplorare i temi oltre le slide.
``Andate oltre, se riuscite''.

\paragraph{Il carico di lavoro.} Le regole sui task sono queste:
\begin{itemize}
  \item i task \textbf{sono per voi, non per noi};
  \item se si ha poco tempo, è meglio farli rapidamente piuttosto che non farli:
        ``provateci'';
  \item i docenti cercano di non chiedere più di \textbf{un'ora o due a
        settimana} oltre alle ore di lezione;
  \item non tutte le lezioni avranno un task, e il \textbf{project work} è il più
        importante di tutti.
\end{itemize}

\paragraph{I feedback dei docenti.} I feedback ``si guadagnano'' man mano che si
reagisce bene: non bisogna aspettarsi ``dei papiri'' sui task. Se i docenti
hanno voglia di dare feedback più ampi lo fanno, soprattutto se vedono che lo
studente se li merita. Ad Analog di solito non ci sono task \emph{rejected}; ma
se un docente rifiuta un task, vuol dire che si è fatto \textbf{meno del livello
minimo} richiesto dal corso. I feedback possono non essere immediati: ``siate un
pochettino impietositi dal fatto che siamo dei docenti''. E se un feedback non
arriva, o se ne vogliono di più, bisogna scrivere, ``perché siamo tre esseri
umani e quindi possiamo anche dimenticarci''.

% ---------------------------------------------------------------------
\section{Il project work e il calendario}
\label{a1:sec:project}

\subsection{Come funziona il project work}

Una prima slide con le date del project work è da aggiornare (``perdonatemi il
refuso''); Daniele la corregge al volo e rimanda alla slide successiva, più
precisa. Il funzionamento è questo:
\begin{itemize}
  \item i \textbf{gruppi li formano i docenti}. Gli studenti sono nove, quindi si
        faranno \textbf{tre gruppi da tre}: una dimensione pensata anche per il
        fatto che molti sono studenti lavoratori, così da non ritrovarsi ``in due,
        a esplodere'';
  \item ogni gruppo deve prima preparare un \textbf{pitch} del proprio gioco, da
        presentare in anticipo: un'idea, una bozza del gioco;
  \item poi si lavora al gioco \textbf{nelle pause del corso}: da qui a dicembre
        ci saranno due pause, quindi non ci sarà lezione tutti i lunedì e tutti i
        venerdì;
  \item ci saranno \textbf{review intermedie} con i \textbf{tutoraggi}, per lo più
        il \textbf{sabato mattina}, dalle 11 alle 12, con Ruggero su Google Meet.
        Servono ad aiutare i gruppi a portare avanti il project work.
\end{itemize}

\subsection{Il calendario delle lezioni}

Il \textbf{calendario} è nel calendario studenti della cartella condivisa.
Attenzione: le date \textbf{possono variare}. ``Il 90\% rimarranno così'', ma
può capitare che, per motivi di lavoro, qualcuno stia male o non possa, o che
``scompaia l'elettricità da tutte le città d'Italia''. E l'ordine delle lezioni
può cambiare se i docenti vedono che agli studenti serve prima una lezione
invece di un'altra.

Ci sono \textbf{due pause importantissime}:
\begin{itemize}
  \item dal \textbf{13 ottobre al 5 novembre}: è il periodo delle grandi fiere
        (\textbf{Essen} e \textbf{Lucca}\footnote{La trascrizione dice ``Essen,
        Luca, Cose'': si tratta con ogni probabilità della fiera Spiel di Essen,
        in Germania, e di Lucca Comics \& Games, che si tengono entrambe in
        ottobre.}), ``un periodo un po' frenetico'', e cominciano anche i corsi
        accademici, quindi la pausa permette ai docenti di respirare. Per gli
        studenti è la pausa più importante: quasi un mese in cui sviluppare la
        parte importante del gioco;
  \item dal \textbf{24 novembre all'11 dicembre}, dopo qualche lezione a novembre.
\end{itemize}

E ci sono \textbf{due date da evidenziare}:
\begin{itemize}
  \item il \textbf{25 settembre}: nei primi 5--10 minuti della lezione (che per il
        resto sarà tenuta da Beatrice) i docenti annunceranno i tre gruppi;
  \item il \textbf{12 ottobre}: il \textbf{pitch} dei giochi, con l'idea
        embrionale. Tra il 25 settembre e il 12 ottobre ogni gruppo prepara un
        breve pitch in formato slide; quel giorno i docenti lo rivedono e danno
        feedback: se il gioco va bene, se non va bene, come svilupparlo.
\end{itemize}
Dal 12 ottobre fino all'inizio di dicembre si sviluppa il gioco, con i
tutoraggi: ``sfruttate questi momenti di pausa''.

\subsection{Il calendario del project work}

La slide successiva, ``reale e aggiornata'', riporta il calendario del project
work (tabella~\ref{a1:tab:calendario} e figura~\ref{a1:fig:calendario}).
``Per ora non pensate al project work: ci penserete dal 25 settembre in poi''.

\begin{table}[htbp]
\centering
\begin{tabular}{@{}lp{9.2cm}@{}}
\toprule
\textbf{Data} & \textbf{Che cosa succede} \\
\midrule
ven 25 settembre\footnotemark & lezione: annuncio dei gruppi e del
\emph{mandate}, cioè di che cosa si dovrà fare \\
sab 3 ottobre, ore 11 & prima review e tutoraggio con Ruggero (dopo una
decina di giorni di lavoro) \\
lun 12 ottobre & pitch del gioco durante la lezione \\
13 ottobre -- 5 novembre & prima pausa: sviluppo del gioco \\
sab 7 novembre & tutoraggio \\
sab 14 novembre & tutoraggio \\
24 novembre -- 11 dicembre & seconda pausa: sviluppo del gioco \\
ven 4 dicembre, entro mezzanotte & consegna del gioco sul Drive (non c'è
lezione) \\
ven 11 dicembre & i giochi vengono visti a lezione \\
14 dicembre & giornata finale, con la tavola rotonda con gli editori \\
\bottomrule
\end{tabular}
\caption{Il calendario del project work come lo presenta Daniele. Le date
possono variare.}
\label{a1:tab:calendario}
\end{table}
\footnotetext{Daniele dice ``lunedì 25 settembre'', ma nel 2026 il 25 settembre
è un venerdì; il giorno della settimana nella tabella è corretto di
conseguenza.}

\paragraph{Perché la consegna è il 4 dicembre.} Il 4 dicembre non c'è lezione:
si deve solo consegnare il gioco, entro mezzanotte, sul Drive. La presentazione
a lezione è l'11 dicembre, ma i tre docenti hanno bisogno di \textbf{una
settimana} per rivedere il gioco in privato e valutare anche la
\textbf{documentazione}; per questo la consegna non coincide con la
presentazione. Che cosa voglia dire esattamente ``consegnare il gioco'' verrà
spiegato più avanti.

\paragraph{La giornata finale.} Il 14 dicembre c'è la \textbf{giornata finale},
con una \textbf{tavola rotonda con gli editori}: di solito viene invitato Flavio,
che è anche docente di game production, più qualche altro editore, ``vediamo
quest'anno chi riusciamo a portarvi''.

\begin{figure}[htbp]
\centering
\begin{tikzpicture}[
    font=\footnotesize,
    x=0.1cm,
    evento/.style={draw, thick, rounded corners=2pt, fill=white,
                   text width=1.8cm, align=center, inner sep=2pt},
    pausa/.style={fill=gray!20, draw=gray!60},
    tacca/.style={thick}
  ]
  % asse: giorno 0 = 25 settembre; 1 unità = 1 giorno
  \draw[pausa] (18,-0.35) rectangle (41,0.35);
  \draw[pausa] (60,-0.35) rectangle (77,0.35);
  \node[font=\scriptsize\itshape] at (29.5,0.18) {pausa};
  \node[font=\scriptsize\itshape] at (65,0.18) {pausa};
  \draw[-{Stealth[length=2.5mm]}, thick] (-3,0) -- (85,0);
  % eventi sopra
  \draw[tacca] (0,-0.15) -- (0,0.9);
  \node[evento, above] at (0,0.9) {25 set\\ gruppi e \emph{mandate}};
  \draw[tacca] (17,-0.15) -- (17,2.2);
  \node[evento, above, fill=gray!8] at (17,2.2) {12 ott\\ \textbf{pitch}};
  \draw[tacca] (70,-0.15) -- (70,0.9);
  \node[evento, above, fill=gray!8] at (70,0.9) {4 dic\\ \textbf{consegna}};
  % eventi sotto
  \draw[tacca] (8,0.15) -- (8,-0.9);
  \node[evento, below] at (8,-0.9) {3 ott\\ tutoraggio};
  \draw[tacca] (43,0.15) -- (43,-0.7);
  \draw[tacca] (50,0.15) -- (50,-0.7);
  \draw (43,-0.7) -- (50,-0.7);
  \draw[tacca] (46.5,-0.7) -- (46.5,-0.9);
  \node[evento, below] at (46.5,-0.9) {7 e 14 nov\\ tutoraggi};
  \draw[tacca] (77,0.15) -- (77,-0.9);
  \node[evento, below] at (77,-0.9) {11 dic\\ giochi a lezione};
  \draw[tacca] (80,-0.15) -- (80,2.2);
  \node[evento, above] at (80,2.2) {14 dic\\ giornata finale};
\end{tikzpicture}
\caption{Una linea del tempo del project work, ricostruita dalle slide (scala
approssimativa: dal 25 settembre al 14 dicembre). Le bande grigie sono le due
pause senza lezioni, in cui i gruppi sviluppano il gioco.}
\label{a1:fig:calendario}
\end{figure}

Finita questa parte, Daniele apre alle domande: ``Domande sul calendario? Sulle
slide che ho fatto? Considerazioni? Non abbiate timore, io non ci credo''. E la
domanda arriva, sull'unico tema che aveva lasciato da parte.

% ---------------------------------------------------------------------
\section{L'intelligenza artificiale nel corso}
\label{a1:sec:ia}

\subsection{La slide di Daniele}

Daniele aveva già presentato una slide che definisce ``importantissima''. ``Lo so
che è paradossale, me ne rendo conto. Lo so che magari mi darete del luddista,
dell'anti-tecnologia: non è così. Uso l'intelligenza artificiale per lavoro. Ma
in questo corso, vi prego, \textbf{non usate l'intelligenza artificiale}''. La
richiesta vale per le attività, per i task e \textbf{anche per la
prototipazione}. Le ragioni:
\begin{itemize}
  \item ``voglio che siate voi a fare il gioco, voglio che siate voi a imparare a
        farlo''. Daniele parla da utilizzatore: prima di ChatGPT, che ha smesso
        di usare, oggi soprattutto di Claude, di cui ``si sottovaluta quanto
        possa essere incerto, insicuro, impreciso, non all'altezza'';
  \item nella game industry, \textbf{i publisher rifiutano automaticamente} i
        pitch e i prototipi fatti con l'intelligenza artificiale. Daniele lo ha
        visto di persona in ``tantissime pitch review'' quell'estate, anche al
        \textbf{First Playable} di Firenze\footnote{La trascrizione dice ``First
        Playboy a Firenze'': si tratta del First Playable, l'evento di business
        dell'industria videoludica che si tiene a Firenze.}: anche nel mondo dei
        videogiochi i publisher che vedono pitch o prototipi fatti con l'IA
        ``rigettano automaticamente la build'', perché hanno ancora timore
        dell'IA e delle sue possibili conseguenze;
  \item ``non so come sarà il futuro'': non sa se bisognerà abituarsi a usarla
        tutto il giorno. Non è ``totalmente contrario'', purché sia usata ``con
        raziocinio''; secondo lui oggi, negli Stati Uniti, non lo è;
  \item ma in un corso che si chiama \emph{Analog Game Design} gli studenti devono
        saper fare i giochi da soli. In futuro, se vorranno farsi aiutare
        dall'IA per essere più efficienti, saranno liberissimi di farlo; durante
        il corso no, perché i docenti vogliono valutare \textbf{le loro
        competenze e capacità}.
\end{itemize}

\subsection{La domanda: che cosa vuol dire usare l'IA in un gioco analogico?}

La domanda di uno studente è precisa: che cosa si intende per ``task fatti con
l'IA''? Qual è l'uso dell'IA in un gioco analogico?

\paragraph{Ruggero: l'esperienza di chi valuta i prototipi.} Ruggero risponde
dalla sua esperienza nel settore, dove riceve prototipi da valutare per
un'eventuale pubblicazione. Distingue due usi:
\begin{enumerate}
  \item le \textbf{immagini generate}: illustrazioni, graphic design e così via.
        Per molti publisher è ``un \emph{reject} immediato, proprio immediato'';
  \item l'IA usata per il \textbf{regolamento}, per scrivere le regole e le
        meccaniche. È più difficile da individuare, ma molti publisher se ne
        accorgono comunque, perché le IA ``non sono ancora brave abbastanza da
        creare un regolamento coeso meccanicamente e scritto con un senso per il
        vostro gioco''. Quello che fanno è prendere parti di altri regolamenti,
        elaborare le strutture che vedono nei regolamenti pubblicati e
        riapplicarle. Ma molti giochi si prestano poco a un regolamento
        strutturato come quello di altri, perché \textbf{ogni gioco è un
        ecosistema meccanico a sé}. Il risultato sono regolamenti con
        riferimenti non immediati, e soprattutto con \textbf{la stessa meccanica
        descritta in modo diverso in pagine diverse}: due difetti che fanno
        perdere tempo all'editore, e in questo caso ai docenti.
\end{enumerate}
Poi aggiunge un punto ``personale, soggettivo'': i docenti sono interessati a
vedere il lavoro, la creatività, ``il frutto della vostra creazione''. E
viceversa gli studenti vogliono i commenti dei docenti: ``immaginatevi se voi
fate un gioco e noi vi diamo i commenti scritti con ChatGPT''. Altrimenti, dice
ridendo, gli studenti fanno il gioco con ChatGPT, i docenti lo danno a ChatGPT
per i commenti, ``e facciamo molto alla svelta''.

\paragraph{Daniele: perché gli studenti usano l'IA.} Daniele racconta la sua
esperienza con gli studenti, che secondo lui ricorrono all'IA in due momenti:
\begin{enumerate}
  \item quando vanno in \textbf{crisi}: in ansia da prestazione o per i tempi.
        Cercano di risolvere il problema della propria disorganizzazione
        facendolo risolvere all'IA, senza accorgersi che l'output è ``molto,
        molto scarso'';
  \item quando hanno l'impressione che l'intelligenza dell'IA sia
        \textbf{migliore} di quello che scrivono loro, o di quello che c'è in
        giro. Per Daniele è una ``mancanza di palato critico'': l'output dell'IA
        sembra buono in superficie, ma quando si va nel dettaglio e ci si fa
        delle domande ``di solito fa un po' di \emph{crumbling}'', si sbriciola.
\end{enumerate}
Lui usa Claude ogni giorno e ne conosce i limiti: ``non c'è nemmeno una frase
che io copio direttamente''. Gli fa fare il lavoro più noioso, più ripetitivo e
soprattutto \textbf{strutturato}, ``replicare delle cose come se fosse una
macro''. Lo usa per il testo e per il codice, mai per le immagini: per motivi di
\textbf{diritto d'autore}, e perché ha molti amici artisti ed è stufo di vedere
il loro lavoro ``dilaniato'' dall'IA.

Ammette che i docenti non possono controllare gli studenti; per questo
``ve lo consiglio''. E riprende quella che definisce la cosa più bella detta da
Ruggero (``ti ruberò la frase per il corso di narrative''): i docenti i feedback
li danno di persona. Daniele a volte manda 30 minuti di audio a uno studente, o
fa intere call per rivedere script, pitch e giochi: ``è il mio lavoro, mi piace
farlo e voglio farlo io, quindi mi aspetto che lo studente lo faccia lui''. Come
freelance, poi, è ``stufo marcio'' dei clienti che arrivano con qualcosa fatto
con l'IA (``guarda che bello, pubblica questo''), perché il più delle volte il
risultato è pessimo: non sanno come usarla. ``È anche una questione di
competenza''.

\paragraph{Il prototipo non deve essere bello.} Questo vale anche per il
\textbf{prototipo}, su cui si tornerà nelle lezioni di prototipazione: in un
prototipo non servono ``graficoni allucinanti''. Le grafiche bellissime fatte con
l'IA spesso \textbf{peccano di user experience}: sembrano bellissime, ma quando
il giocatore deve leggere la carta non capisce qual è il costo, qual è il
potere, qual è l'effetto.

\paragraph{Beatrice: si perde l'occasione di imparare.} Beatrice aggiunge due
motivi che valgono ``al di là dei nostri credo''. Il primo: chi usa l'IA in
sostituzione (per esempio per scrivere un testo, per accelerare i tempi o per
trovare strategie) \textbf{perde la parte in cui può imparare}. Nel game design
bisogna ``tenere insieme i pezzi'' di un regolamento, ma anche dell'esperienza
dell'utente e di come trasmetterla; i passaggi necessari (creare il regolamento,
creare le meccaniche, sperimentare, raccogliere i feedback) ``si capiscono solo
vivendoli''. E si sta pagando un corso per imparare. Il secondo, sul piano
editoriale: chi presenta a un editore un prodotto che ha come base l'IA rischia
da un lato di non avere imparato abbastanza (e di ritrovarsi ``a piedi'' sul
campo senza IA), dall'altro che l'editore se ne accorga, ed è un tema
\textbf{reputazionale}. Ci sono editori ``più illuminati di altri'', quindi
anche qui c'è un mercato da capire; ma i docenti preferiscono dire che cosa è
più sensato per il mercato come lo conoscono.

\paragraph{Ruggero: non abbiate paura del giudizio.} Ruggero vuole dare uno
``spotlight'' a un punto detto da Daniele e rimasto ``semi-ignorato'': a volte
l'uso dell'IA nasce dalla \textbf{paura del giudizio}, dal voler consegnare
qualcosa di più pulito, più strutturato possibile. ``Vi giuro, non ve ne
preoccupate''. Ruggero detesta i voti e non ragiona per voti; non valuterà
nessuno per grafiche fatte a matita. Insegna da anni, ha avuto studenti con
dislessia e discalculia: refusi ed errori non gli interessano, ``non è un corso
di editing o di proofreading''. Gli interessa che le \textbf{meccaniche siano
coerenti con lo scopo di design}. Immagini fatte a matita o con Paint ``non sono
un problema, anzi, le preferisco'' rispetto all'``ennesima immagine fatta con
ChatGPT o Midjourney''.

\paragraph{Daniele: il pensiero critico, costo e qualità.} L'ultimo intervento
è di Daniele, sul \textbf{pensiero critico}. Ragionando per assurdo (``un
pensiero di Belzebù''), immagina che un domani l'IA venga usata ovunque. I suoi
limiti, secondo lui, resterebbero due:
\begin{description}
  \item[Il costo] Tutti dicono che sta diminuendo e che diminuirà, ma non se ne
        può essere sicuri al 100\%.
  \item[La qualità] Cioè la capacità di leggere l'output. Se si dà un regolamento
        a Claude o a ChatGPT chiedendo quali problemi ha, il modello risponde con
        i problemi che, in percentuale, è più probabile che gli esseri umani
        abbiano già individuato nei dati su cui è stato addestrato. Ne derivano
        due problemi: se il problema del proprio gioco non è mai stato
        individuato prima, il modello non lo vedrà; e viceversa può segnalare
        problemi che il gioco non ha.
\end{description}
Fa l'esempio di Grammarly, che usa nella versione premium per rivedere i testi
in inglese: spesso non accetta i suoi consigli, perché se vuole scrivere un
dialogo di videogioco in \textbf{dialetto Cockney} londinese, lo deve scrivere
in Cockney. Se Grammarly dice che è sbagliato, ``sei tu che non hai capito il
motivo per cui voglio scrivere questo testo in Cockney''. La conclusione: prima
di usare un'IA per un'attività bisogna sapere \textbf{che cosa si vuole
ottenere}, qual è l'obiettivo, e quali \textbf{competenze} servono per valutare
in modo critico se l'output di un qualunque LLM è davvero valido. Daniele non
vuole dilungarsi oltre: è un argomento ``molto caldo'' per lui, che è anche
scrittore e si sente dire da tutti ``ti sostituiranno''.

Subito dopo si presenta finalmente Gianmarco (lo si è visto nella
sezione~\ref{a1:sec:partecipanti}), e poi si fa una pausa di una decina di
minuti. Daniele avverte che nella terza parte lui non ci sarà (alle sette e un
quarto ha un impegno di lavoro), e lo stesso vale per Beatrice: la parte finale
di domande e risposte sarà con il solo Ruggero, ``in ottime mani''.

% ---------------------------------------------------------------------
\section{Le figure professionali dell'industria analogica}
\label{a1:sec:figure}

Dopo la pausa (in cui si scopre che Angelo lavora in un piccolo birrificio in
Brianza e si scherza sul rischio di un \emph{drinking game}, tra la birra e la
vendemmia) la lezione riprende con le slide di Ruggero. ``Oggi è una lezione
introduttiva'': gli argomenti si vedranno ``molto dall'alto'' e verranno
approfonditi nelle lezioni successive. Ma servono già ora, per creare
\textbf{l'inizio di un linguaggio condiviso}: ci sono termini che torneranno
spesso, a cominciare da quelli del mercato e dell'industria del gioco
analogico, un'industria ``molto varia, molto caotica, anche molto mutevole'',
con tanti tipi diversi di giochi e altrettanti tipi diversi di entità che ci
lavorano.

\subsection{Una decina di etichette}

I docenti hanno raccolto ``una decina di etichette'', cioè di \textbf{ruoli} che
nell'industria sono o possono essere presenti. Due avvertenze:
\begin{itemize}
  \item nelle aziende più piccole \textbf{la stessa persona copre spesso più
        ruoli}, per forza di cose. Sarà anche il caso degli studenti, che
        lavoreranno in gruppi di tre e non di venti: è facile che ciascuno finisca
        a fare più cose. I docenti ne terranno conto, e qui entreranno in gioco
        anche competenze diverse dal game design (programmi di grafica,
        illustrazione, impaginazione, layout, \emph{proofreading}\ldots);
  \item Daniele precisa che \textbf{non si chiede agli studenti di diventare}
        illustratori o altro. Ruggero concorda, ma aggiunge che è utile sapere
        anche \emph{che cosa non è il proprio lavoro}. Se in un gruppo nessuno sa
        illustrare, ai docenti va benissimo un prototipo che mostri dove vanno
        le immagini: \textbf{placeholder}, concept, anche semplici box di testo
        con scritto ``qui vorremmo un'immagine di\ldots''. ``È importante anche
        sapere che cosa sapete fare e che cosa non sapete fare''.
\end{itemize}
La tabella~\ref{a1:tab:figure} riassume le figure; le sottosezioni seguenti le
riprendono una per una, con i commenti dei docenti.

\begin{table}[htbp]
\centering
\small
\begin{tabular}{@{}p{3.3cm}p{9.6cm}@{}}
\toprule
\textbf{Figura} & \textbf{Di che cosa si occupa} \\
\midrule
Game designer & crea le regole alla base dell'esperienza: le meccaniche, come
interagiscono tra loro e come chi gioca interagisce con esse \\
Game developer & fa da ponte tra chi fa game design e il publisher: bilancia il
prototipo, cura i modelli matematici, adatta il gioco (e il paragame) alle
necessità del publisher \\
Editor (proofreader) & revisiona e adatta i testi a partire dalla prima stesura
del designer (correzione, linguaggio inclusivo\ldots) \\
Publisher (editore) & pubblica e produce fisicamente il gioco; il suo marchio è
sulla scatola o sulla copertina \\
Sales e marketing & il lato commerciale: vendite (anche B2B) e comunicazione del
prodotto \\
Graphic designer & icone meccaniche, wireframe, UX: il lato funzionale degli
elementi visivi \\
Illustratore & le illustrazioni: il \emph{look and feel}, il world building, il
``vibe'' del gioco \\
Sensitivity reader (cultural consultant) & verifica l'accuratezza e la
rappresentazione di un tema sensibile che il team non conosce \\
Producer & coordina il progetto (quasi project management), parla con le
tipografie, verifica che i materiali siano sostenibili economicamente \\
\bottomrule
\end{tabular}
\caption{Le figure professionali dell'industria del gioco analogico presentate
nella slide di Ruggero, in sintesi.}
\label{a1:tab:figure}
\end{table}

\subsection{Game designer}

Il \textbf{game designer} (o la game designer) è chi si occupa, ``sembra banale
a dirlo'', di \textbf{creare le regole alla base dell'esperienza} che si vuole
far arrivare. Vale per tutti i tipi di giochi, non solo per quelli analogici, e
che cosa faccia concretamente in un progetto o in una serie specifica ``può
mutare tantissimo''. Nell'analogico non verrà chiesto di programmare o di
prototipare in digitale; verrà chiesto di \textbf{creare le meccaniche di
gioco}, di capire \textbf{come le meccaniche si interfacciano tra loro} e
\textbf{come le persone che giocano si interfacceranno con le meccaniche}.

\subsection{Game developer}

Il \textbf{game developer} (o la game developer) sarà l'argomento delle due
lezioni di Marta. La parola \emph{development} è nota soprattutto
dall'ambito videoludico, ma nell'analogico ha un significato molto specifico
``che non c'entra niente con la programmazione'': il game developer è
\textbf{un ponte tra chi fa game design e il publisher}. Di solito si occupa di
\textbf{bilanciare il prototipo}, di curare i \textbf{modelli matematici}, di
adattare il gioco e il \textbf{paragame}\footnote{La trascrizione dice ``il
Wix, il paragame''. Il \emph{paragame} è, in termini generali, tutto ciò che
accompagna il gioco senza esserne il sistema di regole (regolamento, scatola,
materiali, comunicazione); i docenti lo tratteranno in una lezione dedicata, e
qui il termine è solo citato.} alle necessità del publisher. La distinzione
viene ripresa e approfondita nella sezione~\ref{a1:sec:design-development}.

\subsection{Editor}

L'\textbf{editor}, a volte chiamato \textbf{proofreader} a seconda dell'azienda,
si occupa della revisione dei \textbf{testi}: non della stesura iniziale, perché
di norma, data la scala dell'industria, è chi fa game design a scrivere un
\textbf{primo draft}, e l'editor ci passa sopra. Non solo per correggere i
refusi: spesso il suo lavoro è anche \textbf{adattare} i testi, per esempio
renderli \textbf{gender neutral}. Se ne vedranno casi studio più avanti.

\subsection{Publisher}

Il \textbf{publisher}, o \textbf{editore}, è l'entità che nell'ambito
commerciale \textbf{pubblica} giochi, giochi di ruolo o giochi da tavolo. È una
figura ``più solidificata'' nel gioco da tavolo e nel gioco di ruolo da tavolo;
\textbf{nell'ambito pervasivo non esiste proprio in questa forma}, come
mostrerà Beatrice parlando dei casi studio. Se si compra un gioco da tavolo o un
gioco di ruolo in un negozio, il marchio del publisher è di solito sul retro
della copertina o sul fronte della scatola: è l'entità che, di solito,
\textbf{produce fisicamente} il gioco.

\subsection{Sales e marketing}

L'ambito \textbf{sales e marketing}, a volte affidato a un'unica persona e a
volte separato, riguarda il lato commerciale. In alcune aziende è interno al
publisher, e si occupa di parlare con le fabbriche e le tipografie, di ``far
quadrare i conti''. In altre realtà è completamente diverso, specialmente nel
\textbf{B2B} o nei giochi \textbf{didattici} e \textbf{formativi}, di cui
Daniele e Beatrice ``sapranno parlare meglio''. Il tema viene ripreso più avanti
in dettaglio (sezione~\ref{a1:sec:marketing}).

\subsection{Graphic designer e illustratore}

Ruggero presenta insieme \textbf{graphic designer} e \textbf{illustratore}
proprio perché sono ruoli molto diversi che si tende a far coincidere:
\begin{itemize}
  \item chi fa \textbf{graphic design} si occupa delle \textbf{icone meccaniche},
        del \textbf{wireframing}, della \textbf{UX}: del lato \textbf{meccanico e
        funzionale} degli elementi visivi, cioè di tutto ciò che serve
        meccanicamente nel gioco, o anche alla sua ``qualità della vita'';
  \item chi fa \textbf{illustrazione} si occupa dell'\emph{art}, delle
        illustrazioni: copre il \textbf{look and feel}, la parte di
        \textbf{world building} e il ``vibe'' che si dà al gioco.
\end{itemize}

\paragraph{Daniele: che cos'è la user experience.} Daniele interviene sulla
\textbf{user experience}, ``importantissima per me'': ha un background da
\textbf{layout artist} (impaginava magazine), un mestiere in cui la UX conta
moltissimo. Per user experience non si intende solo l'\textbf{interfaccia
grafica}, la UI, che anche nei videogiochi è disegnata da un artista. La UX è
qualcosa che può fare anche il game designer: \textbf{disegnare l'esperienza del
giocatore}, cioè ``il gesto, l'intenzione, l'intuito'' che stanno dietro
l'aspetto visivo dell'interfaccia. Per fare UX non serve saper disegnare.
L'esempio è quello della \textbf{porta}: se si progetta una porta, bisogna
indicare all'utente come si apre, come si prende la maniglia; ``quella è user
experience''. Il design della maniglia in sé, invece, è
\textbf{interfaccia}.

\paragraph{Un aneddoto sulle icone.} Ruggero aggiunge che molte persone che
fanno \textbf{graphic design ben fatto non sanno illustrare}, e spesso vale il
viceversa. Lo si nota quando le icone di un gioco sono bellissime ma non
funzionali: andando a guardare i crediti, di solito si scopre che le icone le ha
fatte la persona che ha fatto le illustrazioni. ``Comunque non andiamo adesso in
questo, perché è un altro vaso di Pandora''.

\subsection{Sensitivity reader}

Il \textbf{sensitivity reader}, o \textbf{cultural consultant} a seconda del
progetto, fa un lavoro di \emph{sensitivity reading} del gioco. È una figura
molto usata nell'ambito \textbf{formativo e didattico}. L'esempio di Ruggero: se
si fa un gioco per un ospedale su una malattia specifica, chi lo progetta non è
un medico, ma il gioco deve essere \textbf{accurato} e \textbf{rappresentativo}
delle dinamiche che quella malattia comporta. Serve una persona che conosca quel
tema, anche se non sa che cosa sia il game design. Sono figure chiamate di
solito \textbf{a progetto}, per giochi specifici, che compensano competenze che
il team non ha ``per forza di cose, perché noi non sappiamo tutto''; e a volte,
per lavoro, si devono fare giochi su temi che non si conoscono personalmente.

\paragraph{Daniele: non è un censore.} Daniele segnala un problema italiano:
alcune persone, anche online, hanno confuso il sensitivity reader con un
\textbf{censore}, ``una persona che fa censura, soprattutto in materia di
\emph{woke}''. ``Assolutamente non è così''. La parola \emph{sensitive} non
significa che il pubblico sia sensibile a un tema, ma che \textbf{l'argomento di
cui tratta il gioco è sensibile}. In un gioco sulla medicina, realistico o no,
non si possono inserire imprecisioni. Per questo esistono i sensitivity reader
``più disparati'': da quelli per la cultura indiana o afroamericana fino a quello
su ``come si allevano gli gnu in Africa''. È una persona competente su un
argomento che aiuta a essere \textbf{veritieri}. Certo, uno degli ambiti
principali in cui si usano è quello culturale, ma ``non fate l'errore di pensare
che siano dei censori''. Daniele ha lavorato con ``centinaia'' di sensitivity
reader e nessuno gli ha mai detto ``non puoi scrivere questa cosa''. Di solito
sono come degli \textbf{editor}: danno opzioni, mostrano i punti problematici,
anche quelli di cui non ci si era nemmeno accorti.

Ruggero conferma con un esempio di quando lavorava in \textbf{Danimarca} con
l'\textbf{OMS}: giochi sul \textbf{diabete di tipo 2} nei bambini, per un target
di ragazzi delle scuole medie. C'era un team dell'OMS specializzato nella
comunicazione sul diabete, per garantire che il gioco fosse accurato e ne desse
una rappresentazione adatta al target. Il sensitivity reader aiuta a rendere il
gioco più accurato e ``più tematicamente coeso''.

\subsection{Producer}

Il \textbf{producer} è una figura a volte interna al publisher. In un progetto
grande è la figura di riferimento, ``quasi'' un project manager. In un progetto
piccolo si occupa spesso anche di parlare con le \textbf{tipografie} e di fare la
\textbf{lista dei materiali}, per assicurarsi che i materiali non siano solo
funzionali al gioco ma anche \textbf{economicamente sostenibili} da produrre:
``magari è un gioco bellissimo, ma poi costa più di quanto potete permettervi, e
nessuno vi comprerà mai un card game a 200 euro''.

\subsection{Perché conoscere queste figure}

Ruggero chiude la slide ``tornando a ombrello'': perché parlarne, se gli studenti
sono lì per fare game design? Non solo per saper dare un \textbf{nome} alle
figure. Non c'è niente di male a creare un gioco solo per sé; ma chi vuole fare
giochi a livello lavorativo, perché li giochino altre persone, ``magari più
persone possibile'', \textbf{si interfaccerà con queste figure}, e sapere che
cosa fanno e come parlarci è parte importante del percorso. Nel corso se ne
toccheranno diverse: il development con Marta, la produzione con Flavio, e
lezioni sulla UX, sul paragame, sul graphic design. ``Abituatevi mentalmente:
non siete in solitaria contro il mondo''.

% ---------------------------------------------------------------------
\section{Marketing e sales}
\label{a1:sec:marketing}

Arrivato all'angolo della slide con marketing e sales, Ruggero passa la parola
ai colleghi: ``io è meglio che non mi fate parlare di marketing e sales'', anche
per esperienze personali.

\subsection{Beatrice: i giochi trattati come prodotti}

Beatrice premette che le due voci sono in parte autoesplicative. Nel mondo
capitalista in cui viviamo, i giochi sono ``strumenti molto potenti per fare
\emph{n} cose'', ma sono trattati come \textbf{prodotti}, e il corso ragionerà
anche in quest'ottica. Nel \textbf{marketing} delle industrie ludiche e
videoludiche ci sono varie professionalità che gli studenti probabilmente già
conoscono: dal \emph{copy} alla grafica, tutto ciò che riguarda il
\textbf{vendere} il gioco o l'attività di gioco.

Il lato \textbf{sales} è forse la distinzione meno nota. Come in tante aziende, è
la parte che si occupa delle vendite \textbf{B2B}, cioè tra aziende: fare
accordi per vendere prodotti ludici a diverse realtà, per esempio ai
\textbf{distributori}, o a una casa editrice che vende alle librerie, che poi
vendono al consumatore finale. Chi si occupa di sales nel mondo ludico capisce
come vendere \textbf{alle aziende}. I due mondi si parlano spesso, e soprattutto
nelle realtà piccole si trovano figure che fanno ``un po' sales, un po'
marketing''.

\subsection{Daniele: che cos'è (e che cosa non è) il marketing}

Daniele il marketing si è trovato ``purtroppo'' a farlo, e lo fa anche per
l'accademia. ``Purtroppo'' perché in Italia è mal compreso: molti lo
identificano con l'\emph{advertisement} online. Non è così: il marketing è una
\textbf{disciplina molto più antica}, e in Italia ``abbiamo molto dimenticato''
come si fa il marketing tradizionale. Alcune cose da ricordare:
\begin{enumerate}
  \item \textbf{Il marketing non dice che cosa pubblicare, ma come pubblicarlo}.
        Esistono, nei videogiochi e nei giochi da tavolo, realtà che fanno
        prodotti \textbf{marketing driven}: l'esempio sono gli
        \textbf{hyper-casual game} per smartphone, quelli con il personaggio che
        va a destra e a sinistra e l'esercito che si allarga. Possono anche
        guadagnare, ma ``sono un po' la morte della creatività''. Il marketing
        ideale \textbf{prende un prodotto pensato da altri e gli dà le maggiori
        chance di essere venduto}: cerca il pubblico ideale per il prodotto che
        si ha, senza doverlo cambiare. Anzi, il mondo videoludico, e ludico in
        generale, ``è lastricato'' di errori di dipartimenti marketing che hanno
        cambiato completamente non il gioco, ma il modo in cui l'hanno
        pubblicizzato, e il gioco è fallito.
  \item \textbf{Usare strumenti del marketing non significa fare marketing}. Game
        design e game development a volte usano strumenti che vengono dal
        marketing (gli studenti ne vedranno con Beatrice, con Daniele e con
        Marta).
  \item \textbf{Il marketing è un gerundio}: si fa. È un'azione che si compie per
        comunicare; la comunicazione è attiva, e il marketing è l'insieme delle
        attività con cui si comunica il proprio gioco, dal \emph{copy} ai social.
        Può includere il digital marketing, ``ma non è solo quello''.
  \item \textbf{Il marketing esiste anche se si crede che non esista}. In Italia
        molti non vogliono ``fare marketing'' perché lo associano al lucro. Ma
        il marketing è \textbf{lo studio di come gli esseri umani comprano e
        vendono}, ed è innato in noi: ci sono logiche di cui siamo vittime, e
        ``anche se sei il bastian contrario, l'anticapitalista'', magari si
        compra comunque su Amazon perché fa comodo. Sono meccanismi che
        esistono, ``affascinanti anche, a volte''.
\end{enumerate}
Gli studenti non dovranno fare marketing, ma dovranno ricordare alcuni di questi
concetti, soprattutto nelle lezioni con Marta, che tratta argomenti legati al
marketing. E, chiude Daniele, come game designer e game developer ``per lo più
speriamo che il marketing non ci rovini i giochi'', soprattutto nel videoludico.

% ---------------------------------------------------------------------
\section{Game design e game development}
\label{a1:sec:design-development}

Ruggero torna sulla distinzione più importante della slide, ``proseguendo con
questa strada'':
\begin{description}
  \item[Game design] è chi \textbf{progetta l'esperienza} tramite la
        \textbf{creazione delle meccaniche} di gioco.
  \item[Game development] è chi fa \textbf{mediazione} tra la persona che scrive
        il game design e la \textbf{casa editrice}.
\end{description}

\subsection{I red flag dei designer}

Ruggero lo ripete perché lo vede ogni giorno nel suo lavoro: tra giochi
pubblicati e prototipi, ne valuta ``quasi uno al giorno, se faccio un conto
annuale'', e solo una minima percentuale arriva alla pubblicazione. Molte volte
i \textbf{red flag} vengono proprio da designer che non hanno presente questa
definizione, e che cominciano a mettere \textbf{clausole sulla non modifica del
gioco}: il gioco deve restare com'è. ``Questa cosa è impossibile con un
publisher, spesso''. Il publisher, per forza di cose, potrà chiedere di cambiare
la grafica, di reimpaginare il regolamento, di cambiare i materiali perché
troppo costosi: chiederà dei \textbf{compromessi}. E la persona che, all'interno
del publisher, fa questo lavoro è il \textbf{developer}.

Sta al designer capire \textbf{su che cosa del proprio gioco si può trovare un
compromesso e su che cosa no}. Publisher diversi lavorano in modo diverso, e il
designer ha sempre l'autorità di dire ``no, non voglio pubblicare questo gioco
se queste sono le modifiche''. Ma deve sapere chi fa il lavoro di development,
perché se il suo obiettivo è la \textbf{pubblicazione tradizionale} si
interfaccerà con questa figura.

\subsection{Daniele: una distinzione che manca al videogioco}

Daniele sottolinea che questa distinzione è \textbf{tipica del mondo analogico}
ed è ``molto interessante''. In Italia la parola \emph{developer} è diventata
negli anni quasi un sinonimo di programmatore di videogiochi, o di chiunque
lavori nel videoludico. Secondo lui c'entra anche il fatto che qualche anno fa
il sito \textbf{Gamasutra} ha cambiato nome in \textbf{Game Developer}, e la
parola è diventata un ombrello per tutte le professioni del videogioco. ``Mi fa
un po' ribrezzo'': il game developer dovrebbe essere un ruolo di design, un
designer che, come nell'analogico, fa \textbf{mediazione}. Nel videogioco i
publisher hanno in effetti i loro developer interni, che adattano il gioco, per
esempio alle console e ai loro kit di sviluppo. Ma la distinzione dovrebbe
esistere anche nel videogioco, perché è ``difficilissimo'' avere la mentalità
per fare le due cose insieme: chi fa design ha un focus che raramente si può
sovrapporre a quello di chi fa mediazione.

A Daniele piace la parola \emph{develop}, \textbf{sviluppare}, nel senso di
prendere qualcosa che è stato progettato e \textbf{renderlo possibile nel mondo
reale}. Ci sono progetti che hanno bisogno di uno sviluppo che ``polishi''
l'esperienza: la semplifichi, la migliori, la renda comprensibile e la concluda,
ma soprattutto la renda \textbf{stampabile} e \textbf{fruibile} per l'utente. La
metafora crossmediale è quella della \textbf{fotografia}: chi sviluppa le
fotografie non è chi scatta la foto, ma chi la rende visibile, anche
fisicamente. ``È un po' retro come metafora, ma ci capiamo''. ``È vero, però è
giusto'', conferma Ruggero.

% ---------------------------------------------------------------------
\section{Il mercato: Italia ed estero}
\label{a1:sec:mercato}

Il corso è in italiano e si svolge in Italia, ma bisogna conoscere anche il
\textbf{mercato estero}: molti giochi progettati da autori italiani vengono
pubblicati all'estero, e conoscere il mercato estero aiuta a capire come il
proprio gioco può arrivarci, ``o viceversa, come non può arrivarci''. La
tabella~\ref{a1:tab:mercato} riassume il confronto della slide.

\begin{table}[htbp]
\centering
\small
\begin{tabular}{@{}p{6.4cm}p{6.4cm}@{}}
\toprule
\textbf{Italia} & \textbf{Estero (soprattutto Europa)} \\
\midrule
il gioco è ancora visto come un'attività \textbf{infantile}, anche se la
normalizzazione è in crescita &
forte tradizione di gioco anche in un contesto \textbf{familiare adulto},
soprattutto ``dalla Germania in su'' \\
risposta \textbf{conservativa} sul tema ludico &
\textbf{premi di prestigio} riconosciuti in tutto il mondo (Spiel des Jahres,
As d'Or) \\
forte focus sulla \textbf{localizzazione} di giochi stranieri &
si cercano soprattutto giochi \textbf{originali} \\
investimento \textbf{molto basso} sulla creazione di giochi originali interni;
pochissimi designer assunti &
\textbf{developer interni} pagati con un contratto (Francia, Germania,
Danimarca\ldots) \\
\bottomrule
\end{tabular}
\caption{Il mercato del gioco analogico in Italia e all'estero, secondo la
slide di Ruggero e i commenti di Daniele.}
\label{a1:tab:mercato}
\end{table}

\subsection{L'Italia}

\begin{itemize}
  \item In Italia il gioco è ancora visto come un'\textbf{attività infantile},
        anche se la \textbf{normalizzazione} è in crescita. C'è ancora una forte
        risposta \textbf{conservativa} sul tema ludico.
  \item C'è un forte focus sulla \textbf{localizzazione} dei giochi stranieri: in
        una qualunque ludoteca o negozio, molti dei giochi da tavolo e di ruolo
        sono di publisher o designer esteri, pubblicati in italiano da publisher
        italiani. (La localizzazione, cioè portare un titolo dall'estero in
        Italia, esiste come funzione all'interno del publisher; Ruggero non l'ha
        messa nella slide delle figure perché non è il focus del corso.)
  \item C'è, ``purtroppo, preparatevi psicologicamente'', un \textbf{investimento
        molto basso} sulla creazione di giochi originali interni alla casa
        editrice. Pochissime case editrici italiane hanno designer interni pagati
        a tempo pieno con contratto a tempo indeterminato. Daniele lo dice senza
        mezzi termini: ``il mercato è molto piccolo, conosciamo praticamente
        tutto il settore'', e le persone che in Italia fanno game design
        analogico a tempo pieno con un contratto a tempo indeterminato ``sono
        meno delle dita delle mie due mani''. Di solito la formula è quella del
        \textbf{freelance}, del contratto, delle \textbf{royalty}. Chi va da un
        publisher aspettandosi di essere assunto a tempo pieno come game
        designer deve sapere che ``non siamo purtroppo ancora in quel tipo di
        situazione''.
  \item C'è però un contraltare: in Italia, soprattutto nel gioco di ruolo ma
        anche nel gioco da tavolo, c'è stata molta \textbf{imprenditoria}, anche
        grazie a \textbf{Kickstarter}, che ha aperto le porte ad altri sistemi
        di \textbf{crowdfunding}. È un altro argomento, perché richiede un team,
        un investimento e così via.
  \item E gli autori non mancano, anzi: le occasioni di incontro tra autori
        vedono ogni anno un aumento dei possibili autori\footnote{La trascrizione
        dice ``i DAG a Parma''; il riferimento è probabilmente a IDEAG, il
        circuito di incontri tra autori di giochi italiani, ma non è
        ricostruibile con certezza.}. Bisogna ``sgomitare'' tra tanti autori,
        che magari arrivano da esperienze molto diverse.
\end{itemize}
Ruggero precisa: l'investimento basso è ``dall'alto''. Le persone che fanno e
vogliono fare giochi sono tantissime e in aumento; scarseggia l'investimento,
sia dei publisher disposti a pagare stipendi sia dei \textbf{bandi} pubblici. I
bandi esistono (Beatrice ne parlerà per i giochi pervasivi, che spesso si
muovono in quella direzione, e Ruggero cita qualcuno che sta facendo un card
game con un bando italiano, ``una cosa assurda, ma lo sta facendo''), ma sono
una rarità, soprattutto rispetto all'estero.

\subsection{L'estero}

\begin{itemize}
  \item L'estero (dipende dalla nazione, ma in generale il resto d'Europa) ha una
        forte \textbf{tradizione di gioco} anche in un contesto \textbf{familiare
        adulto}, soprattutto ``dalla Germania in su''.
  \item Ha \textbf{premi di prestigio} riconosciuti in tutto il mondo, come lo
        \textbf{Spiel des Jahres} tedesco e l'\textbf{As d'Or} francese\footnote{La
        trascrizione dice ``lo Spield AR, la SDOR''.}.
  \item Cerca soprattutto \textbf{giochi originali}. Lo si vede dal fatto che
        molti autori italiani di giochi da tavolo pubblicano in Germania o in
        Francia e \textbf{non vengono pubblicati in Italia}, ``che è tutto
        dire''.
  \item Ha \textbf{figure di development interne}: non è raro vedere publisher
        francesi, tedeschi, danesi con developer pagati internamente, a
        contratto.
\end{itemize}
``Poi la situazione è ancora più drammatica'', ammette Ruggero, ``ma per ora
vogliamo tenere il morale un po' alto. È l'introduzione. Ci sarà tempo di andare
nella tristezza''.

Prima di proseguire, Ruggero insiste sulle domande: con lui si può interrompere
anche mentre parla, alzando la mano con le reaction o semplicemente togliendo il
microfono e chiedendo. ``Non mi offenderò mai, anzi, sarò felice di vedervelo
fare e mi ricorderò positivamente di voi''.

% ---------------------------------------------------------------------
\section{Perché serve una terminologia comune}
\label{a1:sec:terminologia}

La serie di slide successiva è ``una piccolissima parte di fondamenti di
design'' (``non spaventatevi, vi vedo che vi spaventate un po'\,''). Anche qui
molti temi verranno ripresi in dettaglio, ma l'idea è arrivare alla prossima
lezione con almeno un \textbf{glossario}, un \textbf{linguaggio condiviso}: se
qualcuno dice \emph{game developer}, tutti sanno di che cosa sta parlando.

Ruggero parte da una domanda, e precisa che non è retorica: \textbf{perché ci
serve una terminologia comune?} Gli studenti rispondono:
\begin{itemize}
  \item Gianmarco, ``molto pragmatico'': perché bisogna lavorare in team e
        capirsi. Vero, nel corso e anche fuori;
  \item un'altra studentessa: per la \textbf{comunicazione}, tra di noi e con
        tutti quelli che sono interessati nel settore. Vero anche questo:
        conoscere i termini comuni permette di parlare con chi fa game design o
        development nel dettaglio, nel tecnico, nello specifico.
\end{itemize}
Ruggero aggiunge altre due ragioni.

\paragraph{Saper comunicare il proprio gioco.} La prima riguarda la
\textbf{comunicazione del gioco} verso l'esterno. ``E lo dico non tirando acqua
al mio mulino, ma anzi togliendola'': \textbf{non è spesso il buon game design
a far diffondere un gioco}. ``Vorrei che lo fosse, sarebbe bellissimo, ma
purtroppo non è così''. Ci sono giochi bellissimi che non sanno comunicare, che
non sanno dire che cosa fanno, e che quindi è difficile, da fuori, trovare
invitanti; e ci sono giochi mediocri dal punto di vista del design che si
comunicano molto bene, e che quindi conoscono in molti. Saper parlare di quello
che si fa è utile \textbf{a sé stessi}, in primis.

\paragraph{Documentare per sé stessi.} La seconda è la \textbf{documentazione}:
documenti di design, di brainstorming e così via. ``Sembrerà una battuta, ma
provate a recuperare un documento di brainstorming una settimana dopo, se non
avete scritto bene le cose che avete scritto''. Chi l'ha fatto sa già di che
cosa si parla; ``chi non l'ha ancora vissuto, forse lo vivrà''.

% ---------------------------------------------------------------------
\section{Game e play}
\label{a1:sec:game-play}

La \textbf{prima, importantissima divisione terminologica} è quella tra
\emph{game} e \emph{play}. In italiano c'è un problema: \textbf{gioco} e
\textbf{giocare} hanno la stessa radice, mentre in altre lingue, tra cui
l'inglese, sono due parole strutturalmente diverse:
\begin{description}
  \item[Game] il \textbf{gioco come artefatto}: il gioco come ecosistema di
        meccaniche, come struttura di regole, come prodotto (``tra
        virgolette''), come oggetto fisico.
  \item[Play] l'\textbf{atto}, il verbo: le persone che giocano, l'approccio
        fisico all'attività del gioco.
\end{description}
È importante perché gli studenti faranno \textbf{game design}; esistono anche
ambiti di \emph{play design}, ma qui ``creerete giochi, non giocattoli''.

\subsection{Giochi e giocattoli}

Qual è la differenza? Uno studente risponde: \textbf{la presenza di regole}.
Esatto: un gioco ha delle regole, un'attività di \emph{play} non
necessariamente. Con un giocattolo si può fare \emph{play}: c'è chi gioca con i
Lego, chi gioca a guardia e ladri o a nascondino al parco, e magari alcune di
quelle regole sono ``più o meno semi-solidificate''. Ma il \emph{play} permette
spesso un approccio non strettamente legato alle regole.

Viceversa, \textbf{un game può esistere anche se nessuno lo sta giocando}. Se si
prepara su un tavolo una scacchiera con i pezzi degli scacchi, chi passa e la
vede riconosce il gioco degli scacchi, anche se nessuno sta giocando. Se invece
si vedono dei bambini correre in un parco giochi, quei bambini \emph{stanno
giocando}, ma non stanno giocando \emph{a un gioco} specifico. ``È importante
separare questa cosa''; ci si tornerà, ``ma per ora abituatevi a questa
differenza''.

\paragraph{Approfondimento: giocattoli, giochi e regole.} La stessa distinzione
è il punto di partenza del capitolo che Jesse Schell, in \emph{The Art of Game
Design}, dedica alla definizione di gioco (\emph{The Experience Rises Out of a
Game}). Schell nota che usiamo perfino un linguaggio diverso: ``un gioco è
qualcosa a cui si gioca'' (\emph{a game is something you play}), ``un giocattolo
è qualcosa con cui si gioca'' (\emph{a toy is something you play with}); e
arriva a definire il giocattolo come ``un oggetto con cui si gioca''. Analizzando
poi la definizione di gioco di Avedon e Sutton-Smith, ne ricava che ``i giochi
hanno regole, i giocattoli no'', e che le regole sono uno degli aspetti
\emph{definitori} del gioco, insieme alla volontarietà, agli obiettivi, al
conflitto e alla possibilità di vincere o perdere. Schell osserva anche che
molti giochi sono costruiti \textbf{sopra} un giocattolo (``una palla è un
giocattolo, ma il baseball è un gioco''), e tra i suoi consigli di
prototipazione c'è proprio \emph{Build the Toy First}: costruire prima il
giocattolo, cioè qualcosa che sia già piacevole da manipolare, e poi il gioco
intorno. È una lettura utile della distinzione di Ruggero: il \emph{play} può
fare a meno delle regole, il \emph{game} no.

% ---------------------------------------------------------------------
\section{Il framework MDA (e il DPE)}
\label{a1:sec:mda}

\subsection{Meccaniche, dinamiche, estetiche}

Un'altra terminologia fondamentale è quella del framework \textbf{MDA}
(\emph{Mechanics, Dynamics, Aesthetics}; in italiano a volte ``MDE''). Viene da
un \textbf{paper} nato nell'ambito del videogioco\footnote{È il paper di Robin
Hunicke, Marc LeBlanc e Robert Zubek, \emph{MDA: A Formal Approach to Game
Design and Game Research}, del 2004. Daniele aggiunge il paper alla cartella
condivisa durante la lezione.}, che ha ``solidificato'' queste terminologie, e
divide un gioco in tre aree:
\begin{description}
  \item[Meccaniche] le \textbf{regole per come sono scritte}. ``Ogni volta che
        pescate un 2 fate dei punti'': quella è una meccanica.
  \item[Dinamiche] l'\textbf{applicazione in tempo reale} delle meccaniche. La
        dinamica è il giocatore che, pescando una carta, \textbf{spera} che esca
        un 2. Quel processo non è una regola: ``il gioco non vi dice che dovete
        sperare che esca un 2''.
  \item[Estetiche] l'insieme di quello che il gioco vuole comunicare,
        l'\textbf{esperienziale}.
\end{description}

L'esempio di Ruggero mette insieme le tre aree. In un gioco in cui pescare un 2
fa fare punti (meccanica) e in cui quindi si spera che esca un 2 (dinamica),
\textbf{che estetica si ottiene se nel mazzo di 100 carte c'è un solo 2? E se ce
ne sono 50?} L'emozione di aver pescato \emph{l'unico} 2 del mazzo è
completamente diversa da quella di pescarne uno sapendo che ce ne sono
cinquanta e che è molto probabile trovarlo. Le regole sono quasi identiche; cambia
l'esperienza.

\subsection{Il falso amico: ``estetica''}

Daniele avverte di non sottovalutare la parola \textbf{estetica}, che in italiano
è un \textbf{falso amico}: sembra indicare la parte grafica, e invece non lo è.
Si incontreranno designer che hanno studiato MDA e altri che non l'hanno
studiato, e lo si capirà proprio da questo: i secondi diranno che le
\emph{aesthetics} sono la grafica. Le \emph{aesthetics} sono la parte
\textbf{esperienziale}: come nell'esempio di Ruggero, quello che lascia al
giocatore aver pescato l'unico 2 del mazzo.

\subsection{Il DPE: design, play, experience}

Per gli italiani, aggiunge Daniele, è più facile capire l'MDA attraverso un
altro paper, quello del \textbf{DPE}\footnote{È il framework \emph{Design, Play,
Experience} di Brian Winn, proposto come estensione dell'MDA per i giochi
educativi (\emph{serious games}). Anche questo paper viene aggiunto da Daniele
alla cartella condivisa.}, che di fatto fa un parallelismo assegnando
\textbf{design, play, experience}:
\begin{itemize}
  \item le \textbf{meccaniche} sono il lato \textbf{design};
  \item le \textbf{dinamiche} sono il lato \textbf{play};
  \item le \textbf{estetiche} sono l'\textbf{experience}, l'esperienza che riceve
        chi gioca.
\end{itemize}
``Se vivete le \emph{aesthetics} come l'esperienza che vive chi sta giocando,
state a postissimo''. In questo momento Beatrice lascia la lezione (``ciao,
Bea''). I paper di MDA e DPE sono in inglese; chi vuole può dargli una lettura.

\begin{figure}[htbp]
\centering
\begin{tikzpicture}[
    font=\small,
    blocco/.style={draw, thick, rounded corners=3pt, fill=gray!12,
                   minimum width=3.2cm, minimum height=1.0cm, align=center},
    dpe/.style={font=\footnotesize\itshape, align=center},
    esempio/.style={draw, rounded corners=2pt, text width=3.3cm,
                    align=left, font=\footnotesize, inner sep=4pt},
    freccia/.style={-{Stealth[length=2.5mm]}, thick}
  ]
  \node[blocco] (m) at (0,0)   {\textbf{Meccaniche}\\ \footnotesize le regole scritte};
  \node[blocco] (d) at (4.3,0) {\textbf{Dinamiche}\\ \footnotesize le regole in azione};
  \node[blocco] (a) at (8.6,0) {\textbf{Estetiche}\\ \footnotesize l'esperienza};
  \draw[freccia] (m) -- (d);
  \draw[freccia] (d) -- (a);
  \node[dpe, above=0.15cm of m] {DPE: \textbf{design}};
  \node[dpe, above=0.15cm of d] {DPE: \textbf{play}};
  \node[dpe, above=0.15cm of a] {DPE: \textbf{experience}};
  \node[esempio, below=0.45cm of m] (em)
    {``Ogni volta che peschi un 2 fai punti''};
  \node[esempio, below=0.45cm of d] (ed)
    {pescando, il giocatore \emph{spera} che esca un 2 (nessuna regola glielo
     chiede)};
  \node[esempio, below=0.45cm of a] (ea)
    {un solo 2 su 100 carte: pescarlo è un evento;\\
     50 su 100: è quasi scontato};
  \draw[dashed] (m) -- (em);
  \draw[dashed] (d) -- (ed);
  \draw[dashed] (a) -- (ea);
  % prospettive
  \draw[freccia, gray] ([yshift=-0.3cm]em.south west) --
    node[below, font=\footnotesize\itshape, text=black]
    {il designer agisce sulle meccaniche}
    ([yshift=-0.3cm]ea.south east);
  \draw[freccia, gray] ([yshift=-1.2cm]ea.south east) --
    node[below, font=\footnotesize\itshape, text=black]
    {chi gioca incontra prima l'esperienza}
    ([yshift=-1.2cm]em.south west);
\end{tikzpicture}
\caption{Il framework MDA con l'esempio del ``2'' usato da Ruggero, e la
corrispondenza con il DPE suggerita da Daniele. Le due frecce in basso
riassumono la lettura classica del paper: il designer costruisce le meccaniche
per ottenere un'esperienza, mentre il giocatore parte dall'esperienza.}
\label{a1:fig:mda}
\end{figure}

\paragraph{Approfondimento: le due direzioni dell'MDA e le otto estetiche.} Le
frecce della figura~\ref{a1:fig:mda} riprendono un'idea centrale del paper
originale, che la lezione non esplicita: designer e giocatore guardano il gioco
da lati opposti. In \emph{Challenges for Game Designers}, Brenda Brathwaite e
Ian Schreiber la riassumono così: i designer si chiedono quale estetica
vogliono ottenere, definiscono le dinamiche che porterebbero a quella
sensazione, e poi creano le meccaniche che producono quelle dinamiche. Lo stesso
libro conferma l'avvertenza di Daniele sul falso amico: ``le \emph{aesthetics}
non si riferiscono all'aspetto del gioco, ma alla risposta emotiva'' che il
designer vuole suscitare. E offre esempi di dinamiche nel senso di Ruggero: in
uno sparatutto i punti di \emph{spawn} sono una meccanica, mentre appostarsi
accanto a uno di essi per colpire chi riappare (lo \emph{spawn camping}) è una
dinamica; negli scacchi i movimenti dei pezzi sono meccaniche, mentre le
aperture codificate sono dinamiche. Per dare un nome alle estetiche, il paper
MDA propone una tassonomia, poi detta degli ``otto tipi di \emph{fun}'', che
Schell riprende in \emph{The Art of Game Design} come \emph{LeBlanc's Taxonomy of
Game Pleasures}: \textbf{sensazione}, \textbf{fantasia}, \textbf{narrativa},
\textbf{sfida}, \textbf{compagnia} (\emph{fellowship}), \textbf{scoperta},
\textbf{espressione} e \textbf{sottomissione} (il piacere di entrare nel
``cerchio magico'' e lasciarsi andare a un passatempo). È un primo esempio di
quello che i docenti chiedono con le parole vietate della prossima sezione:
invece di dire che un gioco è ``divertente'', dire \emph{quale} esperienza
offre.

% ---------------------------------------------------------------------
\section{Le parole vietate}
\label{a1:sec:parole-vietate}

``Attenzione, adesso qua triggeriamo delle persone''. Ruggero annuncia una regola
che ``tornerà'': nel corso ci sono delle \textbf{parole vietate}. ``Non perché
siamo dei censori''; il perché si vedrà in una lezione apposita, perché
spiegarlo adesso richiederebbe un'ora. Ma Daniele vuole spiegarne ``un
minimo''.

\subsection{La regola}

In \textbf{nessun materiale scritto} consegnato ai docenti, task compresi, sono
accettate queste parole, né in inglese né in italiano:
\begin{center}
\begin{tabular}{@{}ll@{}}
\toprule
\textbf{Inglese} & \textbf{Italiano} \\
\midrule
\emph{fun} & divertimento \\
\emph{immersion} & immersione \\
\emph{entertainment} & intrattenimento \\
\bottomrule
\end{tabular}
\end{center}
``Non li voglio neanche vedere''. Daniele lo ripete con un paragone: ``se mi fate
i disegni con Paint, benissimo, top, fantastico, siamo ancora best friend. Se mi
scrivete \emph{divertimento}, io sbrocco''. È una \textbf{regola dell'accademia}:
negli stessi corsi triennali di narrative design e game design le parole sono
vietate agli studenti per tre anni; ``per voi solo per tre mesi''. Daniele
aggiungerebbe volentieri un quarto termine, \textbf{gamification}.

\subsection{Perché}

Il motivo approfondito lo darà Ruggero nella lezione apposita. Daniele dà il
suo. Ruggero e Diego Ricchiuti sono molto contrari a \emph{fun} e a
\emph{immersion}; lui ha aggiunto \emph{intrattenimento}. Sono \textbf{parole
troppo passe-partout e troppo ambigue}:
\begin{itemize}
  \item \textbf{intrattenimento} in Italia è stata usata per anni come un modo
        per \textbf{sminuire}: ``ma sì, tanto è un gioco solo di
        intrattenimento, è un romanzo solo di intrattenimento''. Ma anche se un
        prodotto è ``solo di intrattenimento'', non vuol dire che non ci si debba
        lavorare con attenzione. E poi intrattenimento ``non significa niente'':
        ci sono tanti modi in cui una cosa può intrattenere;
  \item lo stesso vale per \textbf{fun}, e per \textbf{immersion}, che ha un
        problema in più: è una parola passe-partout per cui si può dire che un
        gioco è ``immersivo'' nel senso generale del termine, ma poi non si sa
        davvero che cosa voglia dire, \emph{immersivo in che senso}.
\end{itemize}

L'obiettivo didattico è che gli studenti \textbf{vadano oltre queste parole}:
``un po' come quando alle medie vi dicono smettetela di dire che tutto è bello,
perché ci sono vari gradi di bello, ci sono altri modi di dire bello''. Bisogna
\textbf{decostruirle}: quando si sta per usarle, chiedersi \textbf{che cosa si
intende} con quella parola e cercarne una \textbf{più precisa}.

Ruggero chiede, in uno dei ``rari casi'', di \textbf{non fare domande adesso} su
questa regola: ogni anno arrivano domande e commenti, e ci saranno lezioni
apposite di tre ore, con strumenti per capire il perché. Ma da subito gli
studenti devono abituarsi a evitare questi termini, sia perché si spera che
abbiano termini più precisi, sia perché queste parole ``portano con sé dei
problemi agganciati'' che si vedranno. ``Non spaventatevi, non verrete sparati in
un vicolo cieco di notte''.

\paragraph{Una piccola gamification.} Il divieto va visto come uno
\textbf{stimolo}, ``un po' come se stessimo gamificando la vostra esperienza'',
una sorta di user experience delle lezioni. Imponendosi di non usare queste
parole, si noterà come migliora la capacità di \textbf{descrivere i giochi} e le
proprie esperienze ludiche, e di essere più precisi. È un invito a ``obbligarsi a
essere più precisi e ad avere una mentalità più critica''. La lezione di
Ruggero sull'\textbf{involvement}, promette Daniele, sarà ``fenomenale'' e
``vi aprirà la testa''.

\paragraph{Approfondimento: anche i libri faticano a definire il ``fun''.} La
difficoltà che il divieto vuole far emergere è visibile anche nei testi di game
design. Schell, in \emph{The Art of Game Design}, prova a definire il \emph{fun}
e propone ``piacere con sorprese'' (\emph{pleasure with surprises}), ammettendo
subito che è strano poter usare una parola per tutta la vita, sapendo con
certezza che cosa significa, senza riuscire a esprimerla chiaramente quando
qualcuno lo chiede; e invita a mettere alla prova la definizione con
controesempi. La sua lente sul \emph{fun} (\emph{The Lens of Fun}) si limita,
non a caso, a chiedere quali parti del gioco siano \emph{fun} e \emph{perché}.
Anche \emph{immersion} ha molte letture: in \emph{21st Century Game Design},
Bateman e Boon la descrivono, tra l'altro, come la capacità di ``perdersi''
nel mondo artificiale presentato dal gioco, legata alla sospensione
dell'incredulità e ai giochi di mimesi. Sono significati diversi sotto la stessa
parola: proprio il motivo per cui, in un documento di design, conviene dire
quale si intende.

% ---------------------------------------------------------------------
\section{Chi crea, chi gioca, chi facilita}
\label{a1:sec:ruoli-esperienza}

Un'altra distinzione importante è quella tra \textbf{chi crea}, \textbf{chi
gioca} e \textbf{chi facilita}. Spesso, quando si fa un gioco, lo si fa provare
di persona a chi gioca. Ma ``la norma dei giochi pubblicati non funziona così'':
quando si pubblica un gioco non si vede chi lo compra e chi lo gioca, e non si
sa se qualcuno, a casa propria con i propri amici, segue le regole alla lettera e
interpreta correttamente quello che il designer aveva descritto. Quindi già in
fase di progettazione bisogna tenere a mente \textbf{chi sta facendo game design
e chi sta giocando}.

E nei giochi a evento, nei giochi pervasivi, nei giochi di ruolo da tavolo e dal
vivo (LARP) c'è un terzo pezzo del puzzle, ``importantissimo'': \textbf{chi
facilita}. La domanda da farsi è: \textbf{il regolamento è scritto per chi gioca
o per chi deve facilitare il gioco?} Bisogna abituarsi a pensare a tutte le
persone coinvolte nell'esperienza.

% ---------------------------------------------------------------------
\section{Playtest e iterazione}
\label{a1:sec:playtest}

\subsection{Nessun gioco nasce perfetto}

L'ultima voce del glossario è ``un'altra cosa importantissima'':
\textbf{playtest} e \textbf{iterazione}. Il corso è \textbf{pragmatico}:
``specialmente da parte mia'', dice Ruggero, non si vedranno teorie ``messe lì
per fare accademia sul game design nelle nostre torri d'avorio''. Ci saranno
teorie e framework, ma tutti estremamente pragmatici, perché l'obiettivo non è
solo capire i giochi, ma \textbf{saperli fare e iterare}.

Playtest e iterazione sono \textbf{inevitabili} nel game design, perché
\textbf{nessun gioco nasce perfetto}, soprattutto perché le persone sono
diverse: un gioco perfetto per una persona può essere incomprensibile per
un'altra. Chi progetta giochi per altre persone deve capire, dopo averli fatti
provare, \textbf{che cosa cambiare, che cosa sistemare, che cosa aggiungere, che
cosa togliere}. Ci saranno lezioni su playtest, prototipazione, raccolta dei
feedback, metodi etnografici. Ma bisogna abituarsi già ora all'idea che
\textbf{non esiste la consegna definitiva}. A un certo punto i task vanno
consegnati per forza di cose, ma, come diceva Daniele all'inizio, alla fine del
corso gli studenti avranno un \textbf{prototipo}, ``che sarà un prototipo'', e
gli strumenti per portare avanti il lavoro anche dopo.

\subsection{L'aiuto dopo il corso}

Non solo gli strumenti: anche l'aiuto dei docenti. Daniele racconta di aver
portato a un editore, a Play\footnote{Probabilmente \emph{Play -- Festival del
Gioco} di Modena; la trascrizione non permette di esserne certi.}, uno dei
progetti realizzati l'anno precedente da due studenti, un gioco di ruolo in
solitario (``un \emph{solo game} RPG molto, molto bello''). Se gli studenti
vorranno trasformare il proprio lavoro in un prototipo vero e proprio, possono
contattarlo: lo fa ``per lavoro, ma anche per passione''. Gestisce anche una
community per aiutare i designer ``un po' depressi'' a rimettersi a progettare.

Quindi: ``non sentitevi persi se durante la fase di iterazione vi sentite
persi, è normale''.

\subsection{Il dizionario comune e il production dictionary}

Daniele riprende un punto detto ``benissimo'' da Ruggero e che secondo lui non è
stato sottolineato abbastanza: il \textbf{dizionario comune} serve anche a
\textbf{non perdersi} mentre si itera. Se da un'iterazione all'altra i
documenti usano termini completamente diversi, e non c'è un dizionario di base,
ci si perde.

Nel videogioco, una delle primissime cose che un designer fa quando entra in
un'azienda è il \textbf{production dictionary}: si annota, in un file privato,
tutte le parole che usano i colleghi, per allinearsi al dizionario comune
dell'azienda, perché lì su un prototipo possono lavorare anche cento persone.
Nel gioco da tavolo non si fa molto, perché il mondo dell'analogico è
``talmente più piccolo''; ma serve comunque imparare un \textbf{dizionario
minimo}, dai docenti e dai professionisti. E il corso ne darà una parte,
entrando nel dettaglio di alcuni termini e framework.

\paragraph{L'esempio del feedback.} Perché serva, lo mostra un esempio. Si
testa un gioco e ci si annota come feedback: ``questo gioco deve essere più
divertente''. Che cosa vuol dire? Come lo si migliora? Quella frase non dà
elementi utili. Se invece si scrive che deve essere \textbf{più adrenalinico},
\textbf{più frenetico}, \textbf{più rilassato}, \textbf{più cozy}, o
\textbf{più bilanciato matematicamente}, si sa dove andare a lavorare.

\paragraph{Approfondimento: la regola del ciclo.} Il principio ``nessun gioco
nasce perfetto'' ha un nome preciso in \emph{The Art of Game Design}: è la
\textbf{Rule of the Loop}, la regola del ciclo. ``Più volte provi e migliori il
tuo design, migliore sarà il tuo gioco''. Schell precisa che non è una lente,
cioè un punto di vista, ma ``una verità assoluta'', senza eccezioni: ci saranno
momenti in cui si cercherà di convincersi che ``questa volta il design è così
buono che non serve testarlo'', e ogni volta se ne pagherà il prezzo. Ed è
interessante la ragione per cui la regola pesa di più sul digitale: secondo
Schell il tempo e il denaro necessari per provare e correggere un sistema sono
molto maggiori nei videogiochi che nei giochi tradizionali. Per un gioco di
carte o da tavolo, con il tempo per provarlo e cambiarlo più e più volte, il
consiglio è semplicemente di farlo. Il gioco analogico è, da questo punto di
vista, la palestra ideale dell'iterazione.

% ---------------------------------------------------------------------
\section{Tabletop Simulator}
\label{a1:sec:tts}

\subsection{Il software di settore}

``Last but not least'': il corso usa un software, \textbf{Tabletop Simulator},
che si compra su \textbf{Steam} e costa tra i \textbf{10 e i 20 euro} (``aspettate
che vada in sconto, va in sconto ogni due per tre''). Daniele spiega perché:
\begin{itemize}
  \item è \textbf{il software di settore} del gioco da tavolo, quello che i
        publisher usano per i prodotti digitali, interni ed esterni;
  \item essendo su Steam, permette di fare \textbf{playtest} e di giocare con gli
        amici di Steam. Daniele stesso sta testando lì il proprio gioco;
  \item funziona anche su Mac: Windows o Mac, nessun problema.
\end{itemize}
Ci sarà una lezione apposita, di Ruggero, sulla creazione del proprio
\textbf{modulo} di gioco.

\paragraph{Che cosa non è.} Daniele chiarisce subito un ``fraintendimento
gigantesco'', soprattutto per chi ha già esperienza: Tabletop Simulator
\textbf{non richiede programmazione} e \textbf{non è un software per
trasformare i giochi da tavolo in videogiochi}. È un software che
\textbf{simula dei componenti fisici su un tavolo}: serve per giocare i board
game in digitale. Simula anche la fisica: c'è perfino il comando per
\textbf{ribaltare il tavolo} (``non usatelo ora'').

\paragraph{La partnership mancata.} Da anni Daniele prova a fare una
partnership con gli sviluppatori: due anni prima ha contattato l'azienda per
acquistare delle chiavi per gli studenti, dato che il software si usa nei corsi.
L'azienda ha risposto che non fa accordi commerciali. ``Credo sia la prima
azienda nella mia vita che dice che non vuole che io le dia dei soldi''. Gli
studenti dovranno quindi provare ad acquistarlo da sé; se i docenti riusciranno a
darglielo, bene.

\paragraph{La domanda di Gianmarco.} Gianmarco segue le lezioni da un computer
aziendale ``super chiuso'', su cui non può installare Steam: deve usare Tabletop
Simulator durante le lezioni? Ruggero risponde che per la lezione dedicata è
meglio poter seguire e rifare in diretta. Ma il software non sarà richiesto
necessariamente per il project work: un gruppo potrebbe fare un gioco pervasivo
o un gioco di ruolo, non un gioco da tavolo. Essendo un software di settore, è
comunque bene conoscerlo. Ruggero aggiunge che le versioni non ufficiali che
circolano non permettono di giocare in molti, e quindi non servono per
playtestare o giocare insieme.

\subsection{Una dimostrazione}

Daniele mostra il software aperto da Steam:
\begin{itemize}
  \item si può creare una partita \textbf{single player} o \textbf{multiplayer};
        per giocare tra compagni si crea un server multiplayer e ci si aggiunge
        come amici su Steam;
  \item si dà al server un nome, un numero di giocatori, eventualmente una
        password; creato il server, si apre la scelta dei giochi;
  \item ci sono giochi gratuiti, giochi da acquistare (DLC, non obbligatori) e
        soprattutto un \textbf{Workshop} ``ampissimo'', con decine di migliaia di
        giochi caricati dagli utenti: ``se c'è un board game che conoscete,
        quel board game è su Tabletop Simulator, giocabile dal Workshop''.
        L'esempio mostrato è l'open beta ufficiale del gioco da tavolo di
        \emph{Slay the Spire}\footnote{La trascrizione dice ``Slade Spire''.}:
        molti publisher, per farsi pubblicità, mettono il proprio gioco in demo o
        giocabile in digitale;
  \item c'è un \textbf{tutorial} iniziale che spiega come muoversi sul tavolo;
        con \emph{Esc} e \emph{Help} si trovano i comandi da tastiera.
\end{itemize}

\subsection{Il primo task}

Il primo (e per ora unico) task è questo: \textbf{entro il 10 ottobre}, giorno
della lezione di Ruggero su Tabletop Simulator, giocare \textbf{tre o quattro
volte in totale}, online, tra compagni di corso o con amici. Non si tratta di
imparare a creare moduli, ma di prendere la mano \textbf{lato giocatore}: aprire
il software, impararne i comandi, giocare. ``Per me potete anche andare a
prendere domino o mahjong'' (``mahjong è bellissimo tra l'altro'', aggiunge
Daniele, che sta facendo un videogioco sul mahjong), o il pachisi o il go. Gli
anni precedenti si chiedeva a ciascuno di proporre ai compagni un gioco da
tavolo che apprezzava. È anche un modo per conoscersi. Ruggero spiega il
perché: usarlo per giocare è facile e intuitivo, usarlo per creare moduli è più
complesso; ma chi arriva al 10 ottobre avendo già aperto il software e giocato,
troverà più intuitivo anche creare un modulo. E l'avvertimento: ``non riducetevi
all'ultimo, non fate una partita il 9 sera''.

\subsection{Tabletopia e altri vantaggi}

Tancredi racconta di aver provato a caricare un suo gioco su
\textbf{Tabletopia}. ``Questo è un vaso di Pandora'', risponde Daniele.
Tabletopia è molto usato, ma secondo lui il suo strumento di creazione dei moduli
è pessimo, ed è il motivo per cui la maggioranza dei moduli ufficiali dei
publisher è su Tabletop Simulator. Tabletopia ha un vantaggio vero, è
\textbf{gratuito}, ma ha scelte di ``qualità della vita'' poco felici.

Tabletop Simulator ha invece un vantaggio ``fenomenale'' che si tende a
sottovalutare: essendo su Steam, \textbf{il modulo si pubblica direttamente sul
Workshop di Steam}; si possono invitare gli amici a iscriversi e giocare con
loro, purché abbiano un account Steam. Rispetto ad altri software da imparare,
costa ``tra i 10 e i 20 euro, ragazzi: si parla proprio di spiccioli''. ``Io con
Tabletop Simulator ci ho playtestato la vita''.

Altre caratteristiche utili:
\begin{itemize}
  \item si possono creare \textbf{moduli non elencati} (\emph{unlisted}),
        perfetti da passare a un publisher prima che il gioco sia pubblicato:
        non sono pubblici, e nessuno può rubare i materiali;
  \item il software ha una grafica datata (``\emph{World Art 1998}'', scherza
        uno dei docenti), ma lo stanno portando su \textbf{Unity} con molti investimenti
        e aggiornamenti; ascolta la comunità e introduce funzioni utili anche per
        chi fa design;
  \item il suo server Discord è molto frequentato, e c'è un canale YouTube
        dedicato solo al design di giochi di carte\footnote{La trascrizione ne
        riporta il nome come ``Distraction Maker''; non è verificabile.} con un
        server Discord dove si fanno playtest: il 99\% è su Tabletop Simulator;
  \item molti editori, quando si propone loro un gioco da tavolo,
        \textbf{chiedono se c'è un modulo per Tabletop Simulator}.
\end{itemize}

% ---------------------------------------------------------------------
\section{Le domande a Daniele: il tema del project work}
\label{a1:sec:domande-daniele}

Daniele ha ``un minuto'' per rispondere alle domande sul corso prima di dover
andare. Uno studente chiede, ``giusto per curiosità'', se il project work avrà
un tema o sarà libero.

\begin{itemize}
  \item C'è un \textbf{mandate}, che i docenti propongono ogni anno e
        aggiornano leggermente; verrà presentato il \textbf{25 settembre}.
        ``Non facciamo spoiler''. Conterrà quello che si chiede di fare e quello
        che non si chiede: dei \textbf{must}, dei \textbf{must not} e dei
        \textbf{nice to have}.
  \item Verrà chiesto di lavorare su un'\textbf{idea originale} che non si è
        portata al corso: originale nel senso di \textbf{creata apposta}, non
        nel senso di ``mai vista prima''.
  \item L'idea va \textbf{partorita insieme}. Altri anni alcuni studenti hanno
        provato a imporre la propria idea al gruppo, ``e l'esperienza è stata
        pessima per tutti''. Quindi: \textbf{non imponete le vostre idee, fate
        brainstorming con gli altri}. Daniele, che fa ``molto spesso
        brainstorming'', darà strumenti per farlo: nessuno arriverà impreparato
        a quella lezione.
\end{itemize}

Ruggero aggiunge che il mandate è \textbf{abbastanza flessibile} perché tutti
possano fare qualcosa che li interessi: ``l'ultima cosa che vogliamo è forzarvi a
lavorare su qualcosa che non vi piace''. Bisogna lavorare su ciò che piace e per
cui si ha uno stimolo personale, anche nella \textbf{tipologia} di gioco: nessuno
sarà obbligato a fare un gioco di ruolo o un board game, e la lezione sui casi
studio mostrerà quante forme diverse di gioco analogico il corso accetta.

L'unica cosa da \textbf{non fare}, ``perché è contraddittorio per questo corso'',
sono i \textbf{giochi grandi}: ``lo chiediamo per voi''. Niente progetti
mastodontici\footnote{La trascrizione dice ``non fate delle cose da gnome'':
probabilmente il nome di un gioco da tavolo molto grande e complesso, non
ricostruibile con certezza.}, niente ``cose alla \emph{Dungeons \& Dragons}''.
Si può fare un gioco di ruolo, ma con un'esperienza ``meno corposa'': ci sono
tanti giochi di ruolo da 15 o 20 pagine, ``bellissimi, fichissimi'', che si
possono leggere online. Il consiglio è farsi una cultura anche su
\textbf{itch.io}, ``uno dei portali più importanti del gioco di ruolo
internazionale'', di cui si parlerà anche nella parte di ricerca e sviluppo.

Poi Daniele saluta (``noi ci vediamo venerdì a lezione'') e lascia la
conduzione a Ruggero.

% ---------------------------------------------------------------------
\section{Le domande finali con Ruggero}
\label{a1:sec:domande}

\subsection{Le regole delle domande}

Ruggero spiega che cosa vuol dire ``stare con lui''. Alla fine di ogni sua
lezione ci saranno \textbf{almeno 20--30 minuti di domande} libere, \emph{free
form}. Vuole che fare domande sia \textbf{normalizzato}: ``io gestisco malissimo
l'\emph{horror vacui} e divento super cringe molto alla svelta se non ci sono
domande in giro''. Sono benvenute le domande sulla lezione appena fatta, ma non
solo: il corso è crossmediale, e anche le domande che vengono da
\textbf{divagazioni}, da parti laterali, ``fatele''. Quello che sembra scontato o
banale a uno studente può non esserlo per il docente, e viceversa; e può essere
utile anche agli altri. ``Se vi sembrano stupide, fatele il doppio''. ``Magari
sembro incazzato, ma vi giuro sono rilassatissimo. Non mordo, magari abbaio un
po'\,''.

\subsection{Come approcciarsi a Steam e Tabletop Simulator}

Rita non ha mai usato né Steam né Tabletop Simulator e chiede come
approcciarsi, e se c'è modo di mettersi d'accordo tra compagni su Discord, magari
con un sondaggio sui giochi, per formare gruppi di gioco. Ruggero conferma: sul
Discord dell'accademia c'è un \textbf{canale apposito per il corso}, ed è degli
studenti. Lo possono usare per organizzarsi a giocare, per chiedere a un
compagno di prestare un gioco, per domande tecniche (come si scarica, come si
installa). ``Sarebbe proprio ideale che vi parlaste e vi organizzaste''.

Anche perché, aggiunge con un ``piccolo spoiler'' (``nessuno vi ha detto
niente''), se fossero i docenti a organizzare i gruppi di gioco, il task lo
starebbero facendo loro. Bisogna \textbf{abituarsi a organizzare le sessioni di
gioco}, perché arriveranno i task sul \textbf{playtest} e su come organizzarlo,
e a quel punto si sarà ``meno disperati''. Sembra scontato, ma saper
\textbf{organizzarsi per giocare in autonomia} ``vi sarà utilissimo come skill in
questo settore, ve lo giuro''.

\subsection{Il bilanciamento}

Angelo confessa ``una paranoia grossissima'': ci saranno lezioni sul
\textbf{bilanciamento}? ``Ma certo''. Sono proprio le lezioni di Marta sul
development: il bilanciamento è di solito un tema che tocca \textbf{i developer
più dei designer}. L'anno precedente c'era una sola lezione con Marta, e i
docenti si sono resi conto che serviva andare ancora più a fondo: per questo
quest'anno le lezioni sono due. ``Non è una paranoia, è una legittima domanda'':
Ruggero insegna da anni, vede altri corsi, parla con altri docenti e studenti, e
``mi incazzo spesso'' perché molti corsi di game design \textbf{non fanno
bilanciamento}. I docenti volevano invece un corso che coprisse tutte le parti
del game design: non solo l'idea del gioco, ma il bilanciamento, la produzione,
i playtest, l'iterazione, la scrittura dei regolamenti.

\subsection{Si può giocare con i docenti?}

Una studentessa chiede se si può giocare con i docenti. Per motivi etici,
durante il corso questo è limitato alle parti di \textbf{valutazione} dei giochi
e di \textbf{tutoraggio}. Dopo il corso, ``nessun problema, almeno con me''.

\subsection{Dove si trovano i giochi su Tabletop Simulator}

Ruggero mostra dove si trovano i giochi: nella propria libreria di Steam, nella
pagina di Tabletop Simulator, alla voce \textbf{Workshop}, dove ci sono ``più o
meno 36.000'' giochi. Una studentessa scopre intanto che l'app di Steam per
telefono è solo una \emph{companion app}: per registrarsi e giocare serve il
PC.

Per provare, chiede agli studenti un titolo:
\begin{itemize}
  \item \emph{Keep the Heroes Out}: non compare, ``non tutti i giochi ci sono,
        perché molti publisher non li mettono''; forse, essendoci già una
        versione del gioco su Steam, il modulo è stato tolto;
  \item \emph{Magic: The Gathering}: ``almeno 50 moduli'', anzi ``123 pagine'';
  \item \emph{Bang!}: tre moduli.
\end{itemize}
Bisogna fare attenzione: alcuni moduli sono in lingue specifiche (``magari il
primo è in russo''), alcuni hanno già l'espansione, altri sono solo il gioco
base. Conviene sfogliare e non fermarsi al primo risultato. E nel Workshop si
trovano anche cose molto particolari: giochi fuori produzione, giochi vecchi,
giochi ancora in prototipo.

\subsection{I ruoli nei team e le dinamiche tra le persone}

Uno studente torna sulla slide delle figure: nelle realtà piccole molti ruoli
hanno ``il doppio o triplo cappello''. Quali ruoli restano distinti, e come si
creano le sinergie tra persone con obiettivi diversi?

Ruggero premette che l'ideale ``platonico'' è un ambiente di lavoro costruttivo,
basato sulla fiducia e ``il più \emph{safe} possibile''; ma nell'industria non
sempre è così, e le strutture variano moltissimo: alcuni publisher sono molto
verticali, altri hanno ruoli già solidificati (se un publisher ha già i suoi
developer e non sta crescendo, entrarci come developer è praticamente
impossibile). ``Se prendiamo due aziende in Italia, vediamo due strutture
interne completamente diverse''. Non esiste una scorciatoia per farsi assumere
da un editore.

Poi prende ``l'esempio più comune di tutti'', la media dei giochi da tavolo
pubblicati (figura~\ref{a1:fig:relazioni}):
\begin{itemize}
  \item il \textbf{game designer} è una figura \textbf{esterna}, con un
        \textbf{contratto a royalty} con il publisher;
  \item la sua figura di riferimento è il \textbf{publisher} stesso, e la
        conversazione designer--publisher è ``molto calda'';
  \item se il publisher ha dei \textbf{developer}, il developer entra in questa
        conversazione: designer e developer collaborano molto, e internamente
        developer e publisher collaborano molto. Spesso, nei publisher piccoli,
        è il publisher a fare le veci del developer;
  \item \textbf{graphic design} e \textbf{illustrazione} tendono a parlare molto
        con il publisher e \textbf{molto poco con il designer}. A meno che non si
        tratti di uno studio indie molto piccolo, è difficile che il designer
        sappia anche solo chi è l'illustratore prima che il gioco esca.
\end{itemize}

\begin{figure}[htbp]
\centering
\begin{tikzpicture}[
    font=\small,
    ruolo/.style={draw, thick, rounded corners=3pt, fill=gray!12,
                  minimum width=2.8cm, minimum height=0.9cm, align=center},
    esterno/.style={draw, thick, dashed, rounded corners=3pt,
                    minimum width=2.8cm, minimum height=0.9cm, align=center},
    calda/.style={{Stealth[length=2.5mm]}-{Stealth[length=2.5mm]}, very thick},
    tiepida/.style={{Stealth[length=2mm]}-{Stealth[length=2mm]}, thin, dashed}
  ]
  \node[esterno] (gd) at (0,0) {Game designer\\ \footnotesize esterno, a royalty};
  \node[ruolo] (dev) at (5,1.4) {Developer\\ \footnotesize se c'è};
  \node[ruolo] (pub) at (5,-1.4) {Publisher};
  \node[ruolo] (art) at (10,-1.4) {Graphic designer\\ e illustratore};
  \draw[calda] (gd) -- (dev);
  \draw[calda] (gd) -- (pub);
  \draw[calda] (dev) -- (pub);
  \draw[calda] (pub) -- (art);
  \draw[tiepida] (gd.south) to[bend right=25]
     node[below, font=\footnotesize\itshape] {quasi nessun contatto} (art.south);
  \node[draw, dotted, rounded corners=4pt, fit=(dev)(pub)(art), inner sep=8pt,
        label={[font=\footnotesize\itshape]above right:interno al publisher (di solito)}] {};
\end{tikzpicture}
\caption{Le relazioni tipiche intorno a un gioco da tavolo pubblicato, secondo la
risposta di Ruggero. Le frecce spesse indicano una collaborazione stretta e
frequente; quella tratteggiata, un contatto quasi assente. Nei publisher piccoli
il publisher fa spesso anche da developer.}
\label{a1:fig:relazioni}
\end{figure}

\subsection{Fare da developer di sé stessi?}

Rita chiede in chat se un designer può fare anche da developer del proprio
gioco, o se Ruggero lo sconsiglia. In generale \textbf{lo sconsiglia}, per lo
stesso motivo per cui si sconsiglia in altri media: i \textbf{darling}. Il
designer ha delle meccaniche a cui è affezionato, e il developer a volte dovrà
\textbf{tagliarle}; è molto difficile, se si è ``gestori di sé stessi'', tagliare
cose che si vogliono nel gioco. ``È un po' come lasciare fare il lavoro di editing
allo scrittore che ha già scritto il libro'': può rivedersi, certo, ma difficilmente
cambierà cose che farebbe bene cambiare. Nelle lezioni con Marta si vedrà perché,
e alla fine del corso ``non avrete dubbi nemmeno voi'' nel separare designer e
developer.

Ma tutto dipende dallo \textbf{scopo del progetto}: se si fa un gioco in tre,
qualcuno quel lavoro deve farlo per forza, e magari il designer fa anche il
developer, compensando con \textbf{più playtest}. ``Il mondo attuale e il mondo
ideale sono due cose molto diverse, e lo vediamo non solo dai giochi ma dalla
politica'', ma ``non entriamo in questo vaso di Pandora''.

\subsection{Il ruolo indispensabile e la pubblicazione}

Lo studente di prima ne trae una conclusione: senza game designer il gioco non
esiste. ``Quello sicuramente''. Ma Ruggero coglie l'occasione per approfondire
un punto che forse non era passato: \textbf{fare design di un gioco per
pubblicarlo non è l'unica opzione possibile}. Un gioco non è ``meno gioco''
se non è pubblicato da un editore. È legittimo mettere un gioco su itch.io in
\textbf{print and play}, fare un gioco a evento senza venderlo, fare un
\emph{murder party}, una escape room, un board game per una persona amica:
``quell'esperienza lì non è meno gioco''. Dipende dall'obiettivo del progetto.

Ruggero ha lavorato molto anche con le \textbf{scuole}, dove il publisher non
esiste: è un gioco per le classi, e la figura più vicina è l'editoria scolastica,
``che però non fa il publisher''. Ci sono compromessi diversi sulla materialità
e su come il gioco si presenta, che si vedranno nella prototipazione e nei casi
studio. Quindi l'unica figura che deve esserci sempre è \textbf{chi fa game
design}, anche se magari non sa di chiamarsi game designer e non ha mai
sentito parlare di game design: ``questo non la rende meno game designer''.

\paragraph{Un libro consigliato.} Ruggero fa anche mentoring e gestisce gruppi
di settore, e spesso gli chiedono che cosa leggere per imparare il game design.
Non è il tipo che obbliga a leggere libri (``non vi obbligherò a leggere
niente''), ma ce n'è uno che consiglia sempre, ed è il primo che ha scelto di
leggere: \textbf{\emph{Rise of the Videogame Zinesters}} di \textbf{Anna
Anthropy}\footnote{La trascrizione dice ``Anna Antrofi'' e ``Rise of the
Videogames in Esters''.}. È un libro un po' datato, ma ancora reperibile; parla
soprattutto di videogiochi, ma ``per l'80\%'' il contenuto si applica al gioco e
al game design in generale. È un libro sulla \textbf{democratizzazione del gioco}:
il gioco come strumento per parlare di sé e come \textbf{medium legittimo}. Come
esistono film drammatici, tristi, intensi, possono esistere giochi su tematiche
delicate, e se ne vedranno nei casi studio. Molti giochi fatti con questa
direttiva non hanno un publisher, che tende a pubblicare ciò che dà un ritorno
commerciale.

\paragraph{Commerciale e personale.} ``Io spero che voi non farete giochi solo
per scopo commerciale''. Certo, servono un tetto sulla testa, il cibo in
tavola, un posto in cui dormire, e così si finisce per fare giochi che non sono
il proprio gioco dei sogni. Spesso è un mix: Ruggero ha fatto giochi commerciali
``per pagarmi l'affitto'' e giochi che voleva fare e su cui ha perso soldi,
``sapendolo proprio che sarebbero stati un fallimento''. I giochi, come i libri e
le canzoni, si fanno per molti motivi, e quelli commerciali di successo non sono
più importanti degli altri. Tra i giochi della sua slide c'erano il gioco da
tavolo ufficiale di \emph{Cyberpunk}\footnote{La trascrizione dice ``Cyberpunk
del 1977'': quasi certamente un gioco legato a \emph{Cyberpunk 2077}.} e
\emph{Marvel United}, su IP grosse e famose, ma anche giochi fatti di ``una
candela e una busta su un tavolo''. E se gli si chiede quale preferisce in
termini di contenuto, non sono quelli sui supereroi.

\subsection{Il gioco di ruolo in Italia}

Nicola, che è interessato soprattutto al gioco di ruolo e conosce già persone
del settore, ha ``una leggera paura'': che chi studia game design finisca per
fare il \textbf{freelance} a vita, ``perché non ci sarà mai un'occupazione per
bene'', o il \textbf{copywriter}. Conosce un ragazzo che ha lavorato in quel
mondo e che ora fa altro.

Ruggero risponde che la maggioranza dei giochi nel suo portfolio sono giochi di
ruolo, quindi il tema lo tocca da vicino, e che in Italia quello che dice Nicola
``è semplicemente vero''. Molti publisher piccoli hanno una persona di
riferimento che copre più ruoli: impagina, fa i playtest, fa un po' di
development. È ``un settore economicamente completamente disastrato'', per
motivi culturali che non apre: oggi in Italia è praticamente impossibile
\textbf{vivere di giochi di ruolo}. Chi ci lavora lo fa nel tempo libero, al di
fuori del lavoro principale, oppure coprendo molte altre cose: molti publisher
italiani di giochi di ruolo pubblicano anche titoli esteri e localizzano, non
fanno solo cose originali. ``È un grosso calderone''. Poi si ferma: ``questo è
un \emph{rant}'', e ``i \emph{rant} tornano sempre'', ma non vuole che gli
studenti escano dalla prima lezione ``con la disperazione nel cuore''. Ci sarà
spazio per parlarne quando avranno gli strumenti per capirlo, ``per non farvi
abbattere''.

\subsection{Che cosa bisogna sapere già}

Uno studente chiede se ci sono nozioni che dovrebbero già avere e che possono
recuperare: per esempio, sente parlare di giochi ``German'' o ``American'' e non
sa che cosa siano. Ruggero spiega l'organizzazione del corso:
\begin{itemize}
  \item le \textbf{prime due lezioni} sono di fondamenta. Questa era ``dall'alto'':
        come è strutturato il corso, le tempistiche;
  \item la \textbf{prossima lezione} è di tre ore di soli \textbf{casi studio},
        ``a manetta, a raffica'', divisi per categorie e tipologie, con
        ragionamenti su categorie e generi di gioco: lì si vedranno German,
        American ``e anche altre più utili'';
  \item nel resto del corso le terminologie specifiche sono inserite nelle
        lezioni dedicate: il \textbf{paragame} nelle lezioni sul paragame, il
        \textbf{modello matematico} nelle lezioni di Marta sul development.
\end{itemize}
Per ora agli studenti si chiede di sapere quello che si è visto oggi:
meccaniche, dinamiche ed estetiche; le figure professionali, designer e
developer; e soprattutto \textbf{game e play}.

\subsection{Game e play: i quattro quadranti}

``Specialmente \emph{game} e \emph{play}'', insiste Ruggero: è più difficile di
quello che sembra, perché l'italiano non abitua a separare le due cose, ed è
possibile averle usate per anni come sinonimi. Nella sua esperienza con chi si
avvicina al design è ``uno dei punti più critici''. Spesso i commenti che si
fanno sui giochi giocati sono commenti sul \emph{play} applicati al
\emph{game}: ``il gioco non mi è piaciuto perché, giocando con Tizio, è successo
questo''; ``il gioco mi sembra sbilanciato perché Tizio al tavolo ha fatto
questo''.

Separare \emph{game} e \emph{play} serve a capire quali giochi sono progettati
bene e quali male, e soprattutto a riconoscere i \textbf{quattro possibili
incroci} (figura~\ref{a1:fig:quadranti}): un gioco progettato bene che piace
giocare; uno progettato male che non piace; uno progettato male che piace; uno
progettato bene che non piace. ``Non è importantissimo in generale nella vita,
ci sono cose più importanti'', ma per chi fa game design lo è: ``a volte vi
chiederanno di lavorare su un gioco che non vi piace giocare, e quel gioco deve
essere fatto bene lo stesso''.

\begin{figure}[htbp]
\centering
\begin{tikzpicture}[
    font=\small,
    cella/.style={draw, thick, minimum width=4.4cm, minimum height=1.9cm,
                  text width=4.0cm, align=center}
  ]
  \node[cella, fill=gray!20] (a) at (0,0)     {\textbf{progettato bene}\\ e \textbf{mi piace} giocarlo};
  \node[cella]               (b) at (4.4,0)   {\textbf{progettato bene}\\ ma \textbf{non mi piace} giocarlo};
  \node[cella]               (c) at (0,-1.9)  {\textbf{progettato male}\\ ma \textbf{mi piace} giocarlo};
  \node[cella, fill=gray!8]  (d) at (4.4,-1.9){\textbf{progettato male}\\ e \textbf{non mi piace} giocarlo};
  \node[font=\footnotesize\itshape, above=0.15cm of a] {mi piace (\emph{play})};
  \node[font=\footnotesize\itshape, above=0.15cm of b] {non mi piace (\emph{play})};
  \node[font=\footnotesize\itshape, rotate=90, left=0.15cm of a, anchor=south]
       {bene (\emph{game})};
  \node[font=\footnotesize\itshape, rotate=90, left=0.15cm of c, anchor=south]
       {male (\emph{game})};
\end{tikzpicture}
\caption{I quattro incroci tra la qualità del \emph{game} (come è progettato) e
l'esperienza di \emph{play} (se piace giocarlo), descritti da Ruggero. Per un
designer i quadranti ``scomodi'' sono i due fuori diagonale: un gioco può
piacere pur essendo progettato male, e va saputo progettare bene anche un gioco
che non piace.}
\label{a1:fig:quadranti}
\end{figure}

\paragraph{La domanda di Tancredi.} Tancredi chiede conferma: il \emph{play}
non è il ``giocare senza regole'', ma può essere anche il giocare a un gioco? Sì:
il \emph{play} è l'\textbf{atto}, e può essere fatto con un \emph{game} o senza.
E un \emph{game} può esistere senza che nessuno lo giochi: gli scacchi non sono
meno scacchi perché in quel momento nessuno ci gioca. E se, giocando, una
persona applica le regole in modo diverso? Allora \textbf{non si sta giocando
quel gioco}.

\subsection{House rule: il caso di Monopoly}

Qui Ruggero fa un po' di \emph{foreshadowing} su una lezione di Daniele: il
discorso degli \textbf{hack}, delle \textbf{house rule}, delle \textbf{mod},
cioè di come modificare un gioco esistente. Il punto è che \textbf{un gioco
modificato da una house rule è un altro gioco}: cambiando le regole cambiano
necessariamente anche le \textbf{dinamiche} che ne derivano, e quindi
l'esperienza. Si immagini \emph{Monopoly} in cui ogni turno si possa, volendo,
prendere gratis dei soldi dalla banca: tutto il resto è identico, eppure il gioco
cambia.

E questo è talmente interiorizzato che ci sono regole che si applicano ``senza
neanche rendercene conto''. Ruggero chiede quanto duri in media una partita a
\emph{Monopoly}. Gianmarco: ``quattro ore''; altri: ``un pomeriggio''; un altro:
``finché uno dei giocatori non aggiunge la regola: ora vince chi\ldots''. In
realtà, \textbf{seguendo le regole scritte}, una partita non dura mai più di
\textbf{un'ora e venti, un'ora e mezza}. La domanda che fa sempre, e la risposta
è sempre la stessa: perché si pensa che \emph{Monopoly} duri tantissimo? Qual è
la regola che nessuno rispetta?

Uno studente indovina: \textbf{le aste}. Secondo il regolamento, se chi finisce
su un terreno non lo compra, il terreno \textbf{va all'asta} tra i giocatori, e
chiunque può aggiudicarselo. Molti invece giocano come se finire su un terreno
desse solo la \emph{facoltà} di comprarlo o di metterlo all'asta\footnote{Più
precisamente, nel regolamento ufficiale, se il giocatore che finisce su una
proprietà libera decide di non comprarla, la banca la mette all'asta e tutti i
giocatori possono fare offerte, compreso chi l'ha rifiutata. L'asta quindi non è
una scelta del giocatore, ma la conseguenza obbligata del suo rifiuto.}. Con la
regola applicata, di aste se ne fanno moltissime: all'ottavo turno ci sono già
otto terreni assegnati, al ventesimo venti, la gente va in bancarotta
velocissimamente e la partita dura un'ora e mezza. Anche le partite competitive
di \emph{Monopoly}, nota Ruggero, durano un'ora e venti, un'ora e mezza.

Uno studente obietta che giocava con le aste e durava comunque quattro ore.
``Avete dei nemici con l'analisi paralisi'', ride Ruggero; e cita un'altra house
rule tipica dell'infanzia: genitori che danno soldi per non far andare in
bancarotta. Ma \textbf{la bancarotta è parte del gioco}: se la si toglie,
\emph{Monopoly} non regge. Sul tema che \emph{Monopoly} vuole trattare, e sul
suo antenato (\emph{The Landlord's Game}), ``non apriamo adesso un altro vaso di
Pandora''.

\paragraph{Approfondimento: Monopoly, la durata e le house rule nei libri.}
L'esempio di Ruggero si ritrova quasi identico in \emph{The Art of Game Design}.
Parlando del bilanciamento tra partite brevi e lunghe (\emph{Short vs. Long}),
Schell scrive che \emph{Monopoly}, giocato secondo le regole ufficiali,
``finisce spesso in circa novanta minuti'', ma molti giocatori trovano quelle
regole troppo dure e le modificano, distribuendo premi in denaro e allentando i
vincoli sull'acquisto delle proprietà: con l'effetto collaterale di far durare la
partita molto di più, ``tipicamente tre ore, o anche di più''. Schell classifica
le \textbf{house rule} come regole create dai giocatori in risposta a un
difetto percepito dopo qualche partita. Anche \emph{Challenges for Game
Designers}, nel capitolo sull'aggiungere e togliere meccaniche, parte dallo
stesso gioco: il \emph{Parcheggio gratuito}, che per le regole originali non fa
succedere nulla, mentre quasi tutti giocano con una house rule che dà un bonus a
chi ci finisce sopra. Le due fonti confermano la lezione: una piccola regola
cambiata cambia le dinamiche, e le dinamiche cambiano l'esperienza.

\subsection{Il fine del game è il play?}

Un altro studente commenta: va bene la distinzione, ma il fine del \emph{game} è
il \emph{play}; quello che si vede è sempre l'effetto. Ruggero è d'accordo: di
solito, se si fa un gioco, lo si fa per farlo giocare, come si scrive un libro
per farlo leggere (anche se esistono, ed ``è già stato esplorato'', libri scritti
per non essere letti e giochi per zero giocatori). Ma la distinzione serve a una
cosa precisa: \textbf{separare, nel playtest, i feedback sul gioco dai feedback
sul gruppo di gioco} e sulle sue dinamiche. Se una persona esce frustratissima da
un playtest, non è necessariamente colpa del gioco; se una persona dice ``questo
è il gioco migliore del mondo'', potrebbe non essere merito del gioco ma del
gruppo. Capire a che cosa si riferisce un feedback serve a capire \textbf{dove
iterare} il prototipo.

Lo studente osserva che l'area è sfumata: alcuni giochi si basano proprio
sull'avere le persone giuste al tavolo. Un gioco di \textbf{negoziazione} con
giocatori abituati ai giochi ``multi-solitario'' diventa ingiocabile; e nei giochi
a \textbf{economia chiusa} (cita \emph{Container}\footnote{La trascrizione cita
anche un secondo titolo non ricostruibile (``Tokyo su Vigimarket'').}) un gruppo
economicamente conservativo, che spende poco e non fa girare l'economia, fa
durare la partita il doppio. Ruggero rilancia: e se il gioco dura il doppio,
questa dinamica cambia la \textbf{percezione di peso} del gioco, le energie spese,
quanto si è stanchi a fine partita. ``Tutto è collegato. \textbf{I giochi sono
ecosistemi di meccaniche}, e cambiandone una, le altre cambiano di
conseguenza''.

\paragraph{La parola vietata, alla prova.} Lo studente aggiunge che ci sono giochi
che ha trovato ``molto sbagliati, ma che mi divertono''. ``Attenzione: non si può
dire \emph{divertimento}!'', lo ferma Ruggero; nelle presentazioni aveva lasciato
correre, ``ma lì non vi avevo ancora letto la regola''. Lo studente ci prova:
``giochi che poi voglio \textbf{rigiocare}'', anche se ``sono evidentemente
sbagliati, rotti''. Ruggero apprezza, ma non vuole anticipare cose che meritano il
loro tempo: se ne parlerà nelle lezioni sui framework e soprattutto in quella di
\textbf{emotion design}, ``forse la mia lezione preferita del corso''. (Lo
studente ammette che ``forse alla fine si gioca per quello, per l'emozione''.)

\subsection{I giochi ``comunicati male''}

L'ultima domanda corposa riprende un'affermazione di inizio lezione: ci sono
giochi belli ma comunicati male. Qualche esempio? ``Mi basta girarmi per
vederne 48'', risponde Ruggero indicando gli scaffali alle sue spalle. Più che
titoli, elenca le \textbf{ragioni}:
\begin{itemize}
  \item \textbf{l'Italia localizza lentamente}. Basta guardare i premi del gioco
        dell'anno in Italia: spesso vanno a giochi usciti due o tre anni prima in
        Germania, Francia, Stati Uniti o Regno Unito, perché localizzare richiede i
        suoi tempi. Il mercato italiano è quindi molto dipendente dall'estero,
        ``molto impaurito'', e spesso aspetta che un gioco vada bene all'estero
        prima di portarlo in Italia; e alcuni giochi ``perdono il treno'';
  \item \textbf{il publisher originale chiude}: come si pubblica un gioco se non
        c'è più l'azienda con cui fare un contratto? A Ruggero è successo più
        volte;
  \item \textbf{il dipartimento commerciale italiano non sa comunicare il gioco}:
        si fa male un'edizione, la stampa vende pochissimo e non viene più
        ristampata;
  \item \textbf{la scala del mercato}: in Italia un gioco da tavolo di peso medio
        (un \emph{eurogame}, un \emph{German}) che vende 1.500 copie ``è tanto'';
        all'estero se ne vendono magari 10.000. Con 10.000 copie in inglese ci
        sono i soldi per una campagna di comunicazione: un trailer, i
        recensori, i tutorial online. In Italia tutto questo è ``ridotto all'osso'',
        e la cosa si riflette su stampa e divulgazione: anche i canali per
        informarsi su un gioco prima di comprarlo sono pochi (se ne parlerà nella
        lezione su ricerca e sviluppo). Un gioco spesso non arriva al mercato non
        perché sia brutto, ma perché l'editore non lo sa comunicare, ne fa poche
        copie e va fuori produzione;
  \item \textbf{il tempismo}: party game progettati nel 2019 e usciti nel 2020, in
        pieno Covid, quando non si poteva giocare al tavolo, sono falliti tutti, e
        gli editori ci hanno perso. ``Non potevano immaginarselo'';
  \item \textbf{i temi caldi}: dopo l'invasione russa dell'Ucraina e la guerra in
        Palestina, alcuni editori non pubblicano o non localizzano giochi
        strategici e di guerra. Ruggero ha lavorato a \textbf{\emph{Status}}, un
        gioco sulla geopolitica che tocca, tra gli altri temi, l'invasione russa
        dell'Ucraina e il genocidio in Palestina, e alcuni editori che di
        solito localizzano i loro giochi (per esempio il partner di riferimento in
        Russia) non lo hanno localizzato per motivi politici.
\end{itemize}
Lo studente voleva capire che cosa voglia dire comunicare male: non metterci
abbastanza energia, o comunicare una cosa per un'altra, ``come il trailer
sbagliato, che poi vai al cinema e ti trovi un altro film''. ``È un mix'', e la
domanda ``era complessa e meritava più spazio''.

\subsection{I giochi di carte collezionabili e i saluti}

L'ultima domanda (``posso rompere le balle un'ultima volta?'') riguarda i
\textbf{giochi di carte collezionabili}: funzionano con le stesse regole dei
giochi da tavolo? Se ne parlerà alla prossima lezione: è proprio il motivo per
cui la lezione sui casi studio dura tre ore, perché ogni genere di gioco ha le
sue dinamiche e le sue regole, e anche il modo di pubblicare cambia molto tra un
TCG, un gioco di ruolo da tavolo, un LARP e un gioco pervasivo.

Ruggero saluta: le email dei tre docenti sono all'inizio delle slide, la
registrazione arriverà, ``ci si vede questo venerdì''.

% ---------------------------------------------------------------------
\section*{Punti chiave del capitolo}
\addcontentsline{toc}{section}{Punti chiave del capitolo}

\begin{itemize}
  \item Il corso di \textbf{Analog Game Design} (seconda edizione) dura circa tre
        mesi, fino a dicembre, ed è tenuto da tre docenti principali:
        \textbf{Daniele} (narrative design, gioco di ruolo, fiction),
        \textbf{Ruggero} (game design, ricerca e sviluppo) e \textbf{Beatrice}
        (giochi pervasivi), con tre ospiti su RNG, development e game production.
  \item Lo scopo è formare \textbf{autori analogici} consapevoli: strumenti da
        usare anche fuori dal corso, non il ``gioco dei sogni''. Il corso è
        \textbf{intersezionale} (nessuna lezione dedicata a un solo tipo di
        gioco), \textbf{cross-generazionale}, \textbf{pratico e adattabile},
        basato sulle \textbf{domande} (nessuna domanda è stupida), sul
        \textbf{feedback} invece che sul voto e sul \emph{fail forward}.
  \item Il \textbf{project work} si fa in tre gruppi da tre, formati dai docenti;
        pitch il 12 ottobre, sviluppo nelle due pause, tutoraggi il sabato,
        consegna il 4 dicembre, presentazione l'11 e giornata finale con gli
        editori il 14. Si valutano documentazione, presentazione,
        autovalutazione, crescita e atteggiamento.
  \item \textbf{Niente intelligenza artificiale} nel corso: i publisher
        rifiutano pitch e prototipi con immagini generate, i regolamenti scritti
        dall'IA descrivono la stessa meccanica in modi diversi, e soprattutto chi
        la usa perde l'occasione di imparare. Un prototipo può essere fatto a
        matita.
  \item Le \textbf{figure dell'industria}: game designer, game developer,
        editor, publisher, sales e marketing, graphic designer, illustratore,
        sensitivity reader (che non è un censore) e producer. Nelle realtà
        piccole una persona copre più ruoli.
  \item Nell'analogico il \textbf{game designer} progetta l'esperienza creando
        le meccaniche; il \textbf{game developer} fa da \textbf{mediatore} con il
        publisher (bilanciamento, modelli matematici, adattamenti). Pretendere
        che il proprio gioco non venga modificato è un \emph{red flag}.
  \item Il \textbf{marketing} non decide che cosa pubblicare ma come; è
        un'azione, ed esiste anche se lo si ignora.
  \item In \textbf{Italia} il gioco è ancora visto come infantile, il mercato
        punta sulla localizzazione e investe poco su giochi e designer interni;
        all'estero (soprattutto dalla Germania in su) ci sono tradizione adulta,
        premi di prestigio e developer interni.
  \item \textbf{Game} è l'artefatto (le regole, il sistema, l'oggetto);
        \textbf{play} è l'atto di giocare. Si fanno giochi, non giocattoli.
        Separarli serve a distinguere i feedback sul gioco da quelli sul gruppo
        e a riconoscere i quattro incroci tra ``progettato bene/male'' e ``mi
        piace/non mi piace''.
  \item \textbf{MDA}: meccaniche (regole scritte), dinamiche (regole in azione),
        estetiche (esperienza, non grafica); nel \textbf{DPE}: design, play,
        experience.
  \item \textbf{Parole vietate} in ogni materiale scritto: \emph{fun},
        \emph{immersion}, \emph{entertainment} (divertimento, immersione,
        intrattenimento). Bisogna trovare parole più precise.
  \item Distinguere \textbf{chi crea, chi gioca e chi facilita}; nessun gioco
        nasce perfetto: \textbf{playtest e iterazione} sono inevitabili, e un
        \textbf{dizionario comune} aiuta a non perdersi.
  \item \textbf{Tabletop Simulator} è il software di settore per giocare (non
        creare videogiochi) i giochi da tavolo in digitale. Primo task: giocarci
        tre o quattro volte entro il 10 ottobre.
  \item Una \textbf{house rule} crea un altro gioco: \emph{Monopoly} con le aste
        applicate come da regolamento dura circa un'ora e mezza.
\end{itemize}

%
%  (Con hyperref conviene aggiungere anche
%  \renewcommand{\theHchapter}{agd.\arabic{chapter}} nella parte, per
%  evitare ancore duplicate con i capitoli 1, 2, ... di Game Design.)
%  Tutte le etichette di questo capitolo hanno il prefisso "a1:"
%  (a1:cap:..., a1:sec:..., a1:fig:..., a1:tab:...). Il capitolo è
%  autonomo: non contiene \ref verso altri capitoli.
%  Le figure sono ricostruzioni schematiche di quanto i docenti spiegano
%  a voce e con le slide, non copie delle slide.
% =====================================================================

\chapter[Introduzione all'Analog Game Design]{Introduzione all'Analog Game
Design: il corso, le figure dell'industria e un primo linguaggio comune}
\label{a1:cap:intro}

Questa è la \textbf{prima lezione} del corso di \textbf{Analog Game Design},
cioè di progettazione di giochi \emph{analogici}: giochi da tavolo, giochi di
carte, giochi di ruolo da tavolo, giochi di ruolo dal vivo (LARP), giochi
\emph{pervasivi} ed esperienze a evento. È un corso \textbf{breve}, di circa tre
mesi (``da qui a dicembre''), separato dal corso di Game Design, e siamo alla sua
\textbf{seconda edizione}. Come avvertono subito i docenti, è ``una giornata di
introduzione'': si spiega come funziona il corso, chi sono i docenti, con quali
metodi e con quale calendario si lavorerà, e ci si presenta a vicenda. ``C'è poco
gioco analogico'', si scusano. Ma nella seconda parte la lezione comincia
comunque a costruire un \textbf{linguaggio comune}: le figure professionali
dell'industria del gioco analogico, la differenza tra \emph{game design} e
\emph{game development}, la situazione del mercato italiano rispetto a quello
estero, e un primo glossario di design (\emph{game} e \emph{play}, il framework
MDA, le parole ``vietate'', il playtest e l'iterazione). Si chiude con lo
strumento di settore per i prototipi digitali dei giochi da tavolo,
\textbf{Tabletop Simulator}, e con una lunga sessione di domande.

La lezione si svolge \textbf{online}, su Zoom, ed è condotta a più voci dai tre
docenti principali del corso: \textbf{Daniele}, narrative designer, che guida la
parte organizzativa; \textbf{Ruggero}, game designer, che guida la parte sui
fondamenti e le domande finali; \textbf{Beatrice}, progettista di giochi
pervasivi, che interviene soprattutto nella prima metà. Gli studenti iscritti
sono \textbf{nove}; uno è assente e un altro arriva a lezione iniziata. Ci sono
una pausa di una decina di minuti (si riprende alle 18:15) e, dopo le 19:15,
un'ultima mezz'ora di domande con il solo Ruggero.

Tre avvertenze di lettura. La prima: la trascrizione è automatica e storpia
molti nomi propri, compresi quelli dei docenti e dell'accademia (``Dive'',
``IVE'', ``Accademia Live'', ``Ai''); dove la ricostruzione è sicura, o molto
probabile, il capitolo usa la forma corretta e lo segnala in nota; dove non lo
è, riporta la forma della trascrizione e lo dice. La seconda: durante la pausa la
trascrizione registra, ripetuta in loop per decine di righe, una frase senza
senso che non ha nulla a che fare con la lezione, e che il capitolo ignora. La
terza: in molti passaggi più docenti parlano uno dopo l'altro senza che la
trascrizione segni il cambio di voce; il capitolo attribuisce le parole a chi le
pronuncia quando il contesto lo rende chiaro (chi parla di sé, chi viene
chiamato per nome, chi ``passa la parola''), e resta generico altrimenti.

Quanto alla data, la lezione si colloca con ogni probabilità \textbf{nella prima
metà di settembre del 2026}\footnote{La data non è mai detta. La si deduce da
alcuni indizi: il form iniziale va compilato ``entro e non oltre il 15
settembre''; i gruppi verranno annunciati il 25 settembre; le date del
calendario (sabato 3 ottobre, lunedì 12 ottobre, sabato 7 e 14 novembre,
venerdì 4 dicembre) cadono nei giorni della settimana indicati soltanto nel
2026; un docente parla di un argomento ``molto complesso nel 2026''; e fa
ancora molto caldo (``le temperature sono ancora indecenti''). È una
deduzione, non un'informazione esplicita.}: i riferimenti a ``oggi'',
``quest'estate'' e ``l'anno scorso'' vanno letti rispetto a quel momento.

Il percorso del capitolo segue quello della lezione:
\begin{enumerate}
  \item l'avvio, i \textbf{docenti} e le \textbf{presentazioni} dei partecipanti
        (sezioni da~\ref{a1:sec:avvio} a~\ref{a1:sec:partecipanti});
  \item i \textbf{materiali} e gli strumenti del corso: il form, la cartella
        condivisa, Discord (sezione~\ref{a1:sec:materiali});
  \item lo \textbf{scopo} e il \textbf{metodo} del corso, le regole di
        comportamento e il carico di lavoro (sezioni~\ref{a1:sec:scopo}
        e~\ref{a1:sec:regole});
  \item il \textbf{project work} e il \textbf{calendario} (sezione~\ref{a1:sec:project});
  \item la regola sull'\textbf{intelligenza artificiale} e la lunga discussione
        che ne segue (sezione~\ref{a1:sec:ia});
  \item le \textbf{figure professionali} dell'industria del gioco analogico,
        con un approfondimento su marketing e sales
        (sezioni~\ref{a1:sec:figure} e~\ref{a1:sec:marketing});
  \item la differenza tra \textbf{game design} e \textbf{game development}
        (sezione~\ref{a1:sec:design-development});
  \item il \textbf{mercato}: Italia ed estero (sezione~\ref{a1:sec:mercato});
  \item i \textbf{fondamenti di design}: perché serve una terminologia comune,
        \emph{game} e \emph{play}, il framework \textbf{MDA} (e il DPE), le
        \textbf{parole vietate}, chi crea, chi gioca e chi facilita, il
        \textbf{playtest} e l'\textbf{iterazione} (sezioni
        da~\ref{a1:sec:terminologia} a~\ref{a1:sec:playtest});
  \item \textbf{Tabletop Simulator} e il primo task (sezione~\ref{a1:sec:tts});
  \item le \textbf{domande finali}: il tema del project work, il bilanciamento,
        i ruoli nei team, la pubblicazione, il gioco di ruolo in Italia, le
        \emph{house rule} di Monopoly, i giochi ``comunicati male''
        (sezioni~\ref{a1:sec:domande-daniele} e~\ref{a1:sec:domande}).
\end{enumerate}

% ---------------------------------------------------------------------
\section{L'avvio della lezione}
\label{a1:sec:avvio}

La lezione comincia con qualche minuto di attesa. ``Siamo in nove in classe'', e
si aspetta che arrivino tutti; uno dei docenti scherza sul fatto che l'unica
persona che arriverà ``sempre in ritardo di venerdì'' è Gianmarco. Mancano
Giordano e Tiziana; Tiziana arriva, mentre la segreteria fa sapere poco dopo che
Giordano non potrà esserci e vedrà la registrazione. Si aspetta anche Beatrice,
la terza docente, che arriva a sua volta dopo pochi minuti.

Nell'attesa, Daniele anticipa il programma della giornata e qualche
informazione pratica:
\begin{itemize}
  \item oggi è una giornata di \textbf{introduzione}: come funziona il corso, chi
        sono i docenti, le presentazioni degli studenti, le \textbf{metodologie}
        e il \textbf{calendario}. Gli studenti avranno ``un fracco di domande'', e
        alla fine della lezione c'è ``proprio uno slot apposta'' per le domande;
  \item le lezioni sono \textbf{registrate}: di solito la registrazione viene
        consegnata dopo circa \textbf{12--24 ore} sul profilo personale della
        piattaforma dell'accademia, che va attivato con la segreteria (Tiziana,
        già studentessa, probabilmente ce l'ha già);
  \item i docenti ``cuore'' del corso sono tre (Daniele, Ruggero e Beatrice),
        ``soprattutto Ruggero''; a loro si aggiungeranno tre docenti ospiti, per
        meno ore, probabilmente una lezione a testa, ``tranne Marta, che avrà due
        lezioni quest'anno''.
\end{itemize}
E una battuta sul caldo: ``spero che voi siate carichi, anche se le temperature
sono ancora indecenti''. ``Carichi moderatamente, per quello che potete''.

Poi Daniele condivide lo schermo (``non solo le slide, ma anche un po' di altre
cose'') e apre ufficialmente il corso: ``benvenuti e benvenute nella seconda
edizione del corso di Analog Game Design. Questa è la prima lezione''. Il
percorso è di circa \textbf{tre mesi}, fino a dicembre, e durante il corso i
docenti cercheranno di far provare agli studenti ``cosa significa fare Analog
Game Design'' con un \textbf{project work}, un progetto di gruppo di cui si
spiegheranno più avanti la struttura e le richieste.

% ---------------------------------------------------------------------
\section{I docenti}
\label{a1:sec:docenti}

``Anzitutto partiamo da delle presentazioni doverose''. Ognuno dei tre docenti
principali si presenta con una slide che riporta i contatti e una selezione dei
lavori fatti.

\subsection{Daniele: narrativa, fiction e gioco di ruolo}

\textbf{Daniele}\footnote{La trascrizione riporta ``Daniel Cusetto''; nel
resto della lezione tutti lo chiamano ``Daniele''. Il cognome non è
ricostruibile con certezza.} ha già conosciuto quasi tutti gli studenti, perché
ha fatto loro il colloquio di ammissione. È \textbf{narrative designer}, o meglio
\emph{game narrative designer}; insegna nell'accademia che ospita il
corso\footnote{La trascrizione storpia il nome dell'accademia in molti modi
(``Dive'', ``IVE'', ``Accademia Live'', ``Ai'', ``AIVA''). Con ogni probabilità
si tratta dell'\textbf{AIV}, l'Accademia Italiana Videogiochi, che ha sede a
Roma: lo fanno pensare i corsi triennali di game design e narrative design citati
dai docenti e il fatto che una studentessa racconti di essersi trasferita a Roma.
È una ricostruzione, non un'informazione esplicita; nel seguito il capitolo
parla semplicemente dell'``accademia''.} ma lavora soprattutto come \textbf{brand
advisor}. Nel corso si occuperà della parte di \textbf{gioco di ruolo},
\textbf{narrazione} e \textbf{fiction}.

I suoi contatti, ``in quest'ordine'', sono email, Telegram e Discord. Ci tiene a
dare alcune indicazioni precise:
\begin{itemize}
  \item se si vuole avere la certezza di una risposta, bisogna scrivergli su
        \textbf{Telegram}, perché lì ha tutti i suoi clienti;
  \item Discord lo guarda di meno; l'\textbf{email} è per le cose ufficiali;
  \item non bisogna contattarlo altrove: non su Facebook, non su LinkedIn, non su
        Instagram, ``perché li vedo molto di meno''.
\end{itemize}
La slide mostra anche alcuni giochi, attività e brand su cui ha lavorato negli
anni, ma non si dilunga: ``nel caso ne parleremo''.

\subsection{Ruggero: game design e ricerca e sviluppo}

\textbf{Ruggero} (spesso chiamato ``Ruger'') ha fatto il \textbf{game designer}
per varie aziende e oggi è game designer e si occupa di \textbf{ricerca e
sviluppo} per un'azienda del settore\footnote{La trascrizione dice ``per MS'';
il nome dell'azienda non è ricostruibile con certezza. Più avanti Ruggero cita
tra i giochi su cui ha lavorato il gioco da tavolo ufficiale di
\emph{Cyberpunk}, \emph{Marvel United} e \emph{Status}, un gioco sulla
geopolitica.}. Ha un background ``misto'': nel suo portfolio ci sono giochi da
tavolo, giochi di ruolo, giochi analogici, LARP, videogiochi, giochi in realtà
virtuale. Si occupa anche di \textbf{divulgazione} sul gioco e di analisi di
\textbf{game studies} e game design. Più avanti, nelle domande finali, dirà di
avere ``50 e più giochi prodotti e progettati'' e di insegnare da vari anni.

Lo si può contattare liberamente via Gmail e Telegram, e cerca di rispondere
\textbf{entro 24 ore}, salvo imprevisti o eventi: ``sentitevi in libertà di
scrivermi se avete dubbi, senza problemi''. Quando arriva a parlare dei propri
giochi, la slide mostra ``un mix di varie cose, dalle più famose alle meno
famose'', ma Ruggero si imbarazza a presentarli, tra le risate dei colleghi.

\subsection{Beatrice: i giochi pervasivi}

\textbf{Beatrice} (``Bea'') si occupa di \textbf{giochi pervasivi}
(\emph{pervasive games}). ``Se non li avete mai sentiti nominare, non
preoccupatevi'': questo corso è ``uno dei pochi, forse l'unico in Italia'' che
li racconta esplicitamente. Li si vedrà nelle prossime lezioni, ma la definizione
di base è questa: i giochi pervasivi \textbf{utilizzano gli spazi e gli oggetti
della vita quotidiana e li trasformano in spazi e oggetti di gioco}.

Beatrice lavora come \textbf{freelance} e si occupa di diversi tipi di
progettazione: lavora sia con le \textbf{aziende} sia con gli \textbf{eventi dal
vivo}. La sua attività diretta si trova sotto l'etichetta \textbf{Minime
Differenze}\footnote{La trascrizione riporta il nome dell'etichetta come
``Minime Differenze'' e l'indirizzo del sito come ``mininedifferenze.com''; la
grafia esatta del sito non è verificabile dalla trascrizione.}, mentre i lavori
per le aziende sono ``sparpagliati'' sui siti delle varie aziende. Molte delle
cose mostrate nella sua slide sono attività proposte al pubblico ``con un
biglietto, con un gioco in formato evento'': un tema che si toccherà nel corso,
per capire come \textbf{comunicare} i propri giochi in modo che siano
comprensibili, accessibili per il pubblico e \textbf{vendibili}.

Ammette un ``grosso limite'' con i messaggi: è un po' lenta a rispondere, ``ma
rispondo sempre a tutto''. La si può contattare per email (cognome punto nome) o
su Telegram. È ``molto curiosa di vedere cosa produrrete''.

\subsection{I docenti ospiti}

Daniele presenta poi i tre docenti che il corso avrà ``solo per alcune
lezioni'', tutte verso \textbf{novembre}:
\begin{itemize}
  \item \textbf{Diego Ricchiuti}, direttore del corso di Game Design
        dell'accademia, farà una lezione sull'\textbf{RNG} (\emph{random number
        generation}, la casualità nei giochi): ``vi parleremo di RNG quando sarà
        necessario. Comunque avrete una lezione con lui, sappiatelo'';
  \item \textbf{Marta}\footnote{La trascrizione riporta il cognome come
        ``Cercasassi''; la grafia non è verificabile.} farà \textbf{due lezioni
        sul development}: quest'anno ``abbiamo deciso di fare doppietta'',
        perché l'anno scorso il suo intervento è stato ``molto, molto utile'';
  \item \textbf{Flavio}\footnote{La trascrizione riporta il cognome come
        ``Mortarino'' e, più avanti, ``Mortalino'', e il nome della sua casa
        editrice come ``Grampibert''; nessuno dei due è ricostruibile con
        certezza.} farà la lezione di \textbf{game production}; è un editore, e
        di solito partecipa anche alla tavola rotonda con gli editori della
        giornata finale.
\end{itemize}
Le slide sono già online e gli studenti avranno accesso alla cartella condivisa:
``non vi preoccupate di queste cose, ne parleremo a breve''. Ma prima di parlare
del corso, ``vogliamo conoscere voi''.

% ---------------------------------------------------------------------
\section{Le presentazioni dei partecipanti}
\label{a1:sec:partecipanti}

Le domande per la presentazione sono semplici, da affrontare ``a turno, senza
fretta, con calma'':
\begin{itemize}
  \item nome e occupazione (``se volete dirla'');
  \item perché si vogliono fare giochi analogici, e a che livello: che cosa ha
        portato lo studente al corso;
  \item che tipo di esperienze o conoscenze si hanno già nell'ambito, se si è già
        fatto qualcosa oppure no.
\end{itemize}
``Chi vuole partire senza nessun tipo di ritegno e di vergogne?''. Per far
parlare gli studenti guardandoli, Daniele interrompe la condivisione dello
schermo e ricorda lui le domande quando serve. Le presentazioni sono interessanti
anche per i commenti dei docenti, che anticipano diversi principi del corso.

\subsection{Tiziana}

Tiziana fa l'\textbf{infermiera}, ma nella vita ha fatto molti lavori:
impiegata, insegnante e ``tanti altri lavoretti''. Gioca ``da sempre'': il gioco
fa parte anche della sua vita familiare, perché suo fratello fa il programmatore
di videogiochi. Ha già \textbf{organizzato} giochi: una caccia al tesoro per far
conoscere il territorio, due scene con un delitto (giochi con omicidio da
risolvere), giochi per intrattenere le persone. Ha riscoperto di recente i giochi
da tavolo, arrivando a \emph{Carcassonne} con famiglia e amici; poiché è
difficile trovare persone con voglia di giocare nella propria cerchia ristretta,
ha scoperto il mondo delle \textbf{associazioni} ludiche, ``e lì è stato
un'esplosione'': prima in Veneto (è di Vicenza), poi a Roma, dove si è
trasferita.

Ha in mente un progetto ambizioso: far conoscere una certa realtà con la modalità
del \textbf{gioco di ruolo}, o comunque proporre un gioco suo.

\paragraph{Il commento: gli strumenti, non il gioco dei sogni.} Ruggero coglie
l'occasione per un'avvertenza che vale per tutti, visto che molti studenti
arrivano con idee già pronte. Nel corso si lavorerà a un progetto, ma
l'obiettivo dei docenti è dare \textbf{strumenti} da applicare anche
\emph{fuori} dal progetto, ai giochi che gli studenti vorranno fare un domani:
``che non siano esercizi per il corso, ma per voi in primis''. Quindi: ``non
rimaneteci male se non riuscirete a lavorare al vostro gioco dei sogni in questo
corso''. Anche perché ``sarebbe impossibile, in due mesi, fare il gioco dei
sogni''. Il corso darà invece la possibilità di fare un'esperienza ludica
insieme e di imparare a replicarla sulle proprie idee: ``le idee non buttatele,
anzi, il contrario: tenetele''.

\subsection{Nicola}

Nicola è uno \textbf{studente}. Sei mesi prima ha fatto un colloquio con Daniele
perché stava valutando di passare all'accademia, e ha deciso, ``nel dubbio'', di
cominciare da questo corso breve per capire come funziona la formazione;
eventualmente si iscriverà al corso di game design vero e proprio. Gli interessa
in particolare il lato dell'esperienza del giocatore, su cui è ``molto critico''
(qui usa la parola ``divertimento'', e Daniele lo ferma con un ``ne parleremo'':
è una delle parole che il corso vieta, come si vedrà nella
sezione~\ref{a1:sec:parole-vietate}).

Ha 21 anni ed è \textbf{giocatore di ruolo da circa otto}. Prima ancora di
conoscere i giochi di ruolo si era creato un sistema tutto suo, ``pessimo,
disgustoso, terribile'', un sistema a percentuale con il d100 ``terrificante'';
poi ha scoperto i giochi di ruolo da tavolo, i classici (\emph{Dungeons \&
Dragons}, \emph{Daggerheart}, \emph{Vampire}\footnote{La trascrizione riporta
``DD'', ``Daggerart'' e ``Vampire'': quasi certamente \emph{Dungeons \&
Dragons}, \emph{Daggerheart} e \emph{Vampire: The Masquerade}.}). Gioca anche a
\textbf{card game} e ad altri giochi sociali, pur avendo sempre avuto difficoltà
a trovare gruppi, ``avendo vissuto in un buco dimenticato dal mondo''. Come
giocatore è ``l'eterno narratore'' del gruppo, quello che prepara la storia e si
costruisce le regole. Come autore ha fatto soprattutto giochi su carta, ``perché
sono incapace a programmare, per adesso''.

Daniele osserva che il tema del giocatore e del suo rapporto con il gioco è ``un
argomento importante'', anzi ``un vaso di Pandora enorme'' (espressione che
tornerà molte volte nella lezione), e che il corso tratterà anche che cos'è
l'analog game design in generale e quanti tipi di prodotti esistono.

\subsection{Rita}

Rita\footnote{Daniele la chiama prima ``Nino'' (probabilmente un errore della
trascrizione) e poi la ringrazia come Rita.} ha 30 anni e si collega da una
regione del Sud, dove ``ancora si muore di caldo''. Ha studiato
all'\textbf{Accademia di Belle Arti di Napoli}, \textbf{design della
comunicazione visiva}: si occupa dell'aspetto più imprenditoriale della
progettazione artistica. Non volendo progettare per un'impresa (``fare loghi,
robe del genere''), vorrebbe applicare la propria creatività alla realizzazione
di giochi analogici.

Non è una giocatrice di ruolo e non lo è mai stata: è una giocatrice
\textbf{tradizionale}, cresciuta con Monopoly, la dama e giochi tradizionali
cinesi. Ha già \textbf{idee, prototipi e progetti} realizzati, e vuole capire
come renderli \textbf{prodotti per la vendita} e se questo è il campo adatto a
lei; magari anche confrontarsi con i colleghi. La sua progettazione è orientata
al \textbf{sociale} e all'\textbf{educativo}: ``responsabilizzare un po' le nuove
generazioni su come sta andando il mondo e su come potrebbe andare, visto che è
tutto nelle nostre mani''.

\paragraph{Il commento: siate curiosi.} ``Pane per i tuoi denti'', le risponde
Daniele, soprattutto per i LARP e i giochi pervasivi: ``vi apriremo un po' la
testa da questo punto di vista''. E aggiunge un consiglio per tutti: ``siate
curiosi'', perché imparare da altri ambiti dà moltissimo. Ognuno è lì per fare
qualcosa di preciso, ``però siate aperti anche ad altre esperienze''.

\subsection{Tancredi}

Tancredi ha 21 anni e non ha ancora grande esperienza lavorativa: qualche
``piccolo progettino'' che non è andato come voleva. Gli interessa il mondo dei
giochi, digitali e non, con un approccio \textbf{educativo} e \textbf{sociale}.
È arrivato soprattutto dai videogiochi, a cui gioca da quando era piccolo; ai
giochi da tavolo si è avvicinato solo nell'ultimo anno, grazie a una realtà
ludica, perché in famiglia ``non siamo molto da giochi''. Ha realizzato qualche
piccolo videogioco con vari engine: prima quelli senza un vero linguaggio di
programmazione, ``più a nodi'', poi qualcosa di più complicato, e ora sta
lavorando a un piccolo prototipo.

\paragraph{Il commento: niente di quello che vedrete è sprecato.} ``Non butterai
via niente di analog game design'', gli risponde Daniele, ``perché il 99 per cento
delle cose che vedremo si ritrovano quasi uno a uno nel mondo dei videogiochi''.
Tanto che nei corsi triennali dell'accademia (anche in quello di narrative
design, dove insegna Ruggero) il primo anno si parte dai \textbf{giochi da tavolo
e di ruolo}, proprio per spiegare meccaniche più complesse che servono poi per i
videogiochi. Ruggero aggiunge che vale anche il contrario: alcuni dei
\textbf{framework} di game design che si vedranno nelle lezioni tecniche vengono
dal videogioco e sono stati riapplicati all'analogico. ``Evitare la
crossmedialità qui è impossibile''. Anzi, tutte le conoscenze ``mediatiche
adiacenti'' (musica, letteratura, cinema) saranno utili.

\subsection{Lorenzo Vergata}

I Lorenzo in classe sono due; il primo a presentarsi è \textbf{Lorenzo
Vergata}, 26 anni, \textbf{dottorando al Politecnico di Milano}. Gioca ai
videogiochi da quando aveva due anni e mezzo, ma ha giocato al suo primo gioco da
tavolo ``credo nove mesi fa''. Gli fa piacere sentire che molti in classe sono
interessati ai giochi \textbf{sociali} ed \textbf{educativi}, un settore che
interessa anche a lui e che ``può essere esplorato molto'', anche con qualche
aiuto dell'intelligenza artificiale per esperienze più ricercate.

Di esperienza nell'analogico dice di avere fatto ``il minimo sindacale''. Ma
racconta che diciotto anni prima, da bambino, era ``più game designer di
adesso'': prendeva un gioco e ne cambiava le regole, anche al biliardino (``no,
ora introduciamo questa regola''). Gli piace sperimentare, creare regole e
vedere che cosa accade in sé e negli altri. Soprattutto vorrebbe ``una chiave di
lettura per poter rimanipolare la realtà'': tre settimane prima è stato a fare la
\textbf{vendemmia in Francia} e ha visto tantissime dinamiche che potrebbero
diventare un gioco da tavolo o un videogioco. Per esempio: finire contro un rovo
pieno di semi che si attaccano ai vestiti lo rallentava nella raccolta
dell'uva, mentre in un altro filare era più rapido. ``Si può fare pure una
dinamica'', e vorrebbe gli strumenti per farne giochi.

Arriva subito la battuta: una ``vendemmia gamificata'' da proporre in Oltrepò
Pavese (``io sono dell'Oltrepò Pavese, quindi devo per forza bere vino'')\footnote{La
trascrizione dice ``nell'altro papavese'' e non chiarisce chi pronunci la
battuta.}. Daniele conferma che quest'anno il corso avrà ``tante cose sul
sociale'': ``ci piace parecchio''.

\subsection{L'altro Lorenzo}

Il secondo Lorenzo (``secondo, ma non meno importante'') è uno \textbf{scrittore
di racconti brevi} e sta scrivendo il suo primo romanzo. Gli è sempre piaciuto
raccontare storie, e da più grande scriverle; gioca ai videogiochi ed è ``un
amante delle storie e della cultura pop''. Negli ultimi anni è entrato nel mondo
del \textbf{gioco di ruolo da tavolo} e vorrebbe unire le due passioni, le
storie e il gioco di ruolo, magari facendone un lavoro; gli interessa anche lo
\emph{strategic writing}. Daniele gli ricorda che ci sarà una lezione proprio
sulla \textbf{fiction}.

\subsection{Angelo}

Angelo si collega senza videocamera dal lavoro: fa il \textbf{birraio} (in un
piccolo birrificio in Brianza, a Bernareggio, come racconterà dopo la pausa) e,
con il caldo, ``vi risparmio l'orrore''. Per molto tempo era convinto che i
giochi da tavolo fossero solo i classici (Monopoly, Risiko) o i party game; poi,
quattro o cinque anni fa, ha ``sbattuto la faccia'' contro giochi davvero
appassionanti, con una specie di narrazione dietro che li rendeva ancora più
coinvolgenti, e non ha più smesso. È uno che ``ha sempre un botto di idee, ma poi
non mette in pratica assolutamente nulla'': cercava un corso che gli desse le
basi e soprattutto lo \textbf{motivasse} a provare a fare qualcosa.

\paragraph{Il commento: bisogna buttarsi.} Daniele racconta che anche l'anno
prima un suo caro amico si era iscritto proprio per questo motivo. Il corso è
breve, ma cerca di portare gli studenti ``da 0 a 100'' e di far realizzare un
\textbf{prodotto} da presentare alla fine. Il punto è che i giochi ``bisogna
farli, bisogna buttarsi''. E, venendo dalla narrazione, Daniele lo paragona allo
scrivere: si scrive e si riscrive, e allo stesso modo si fa \textbf{playtest}, si
rifà, si aggiorna. ``Non è che si fanno così''.

\subsection{Gianmarco (dopo le slide)}

Gianmarco arriva a lezione iniziata e si presenta più tardi, alla fine della
prima parte, dopo che i docenti hanno chiuso le slide organizzative e la
discussione sull'intelligenza artificiale (``Gianmarco, ti giuro, ti facciamo
presentare''). Lavora nel \textbf{team di innovazione} di un'azienda italiana.
Vuole fare giochi perché è una sua passione, che lo gratifica, e in passato ci ha
provato ``in maniera scomposta''. Cerca un \textbf{metodo} e un confronto con
persone più esperte: ``studiare solo dei manuali è molto sterile'', ``stare
chiusi nel pertugio non aiuta'', ed ``essere autodidatta ha estremamente i suoi
limiti''. Come esperienza è un giocatore e ha letto ``un po' di materiale, di
blog, di manuali'', ``niente di rilevante''.

Manca solo Giordano, che si presenterà alla lezione successiva.

\paragraph{Uno sguardo d'insieme.} La classe è molto varia: un'infermiera,
una designer della comunicazione visiva, un birraio, un dottorando, uno
scrittore, un impiegato nell'innovazione, due studenti ventunenni. Alcuni
vengono dal gioco di ruolo, altri dai giochi tradizionali, altri dai
videogiochi; quasi tutti hanno scoperto i giochi da tavolo moderni da pochi anni
o da pochi mesi. E almeno tre studenti dichiarano un interesse per i giochi
\textbf{sociali} ed \textbf{educativi}. Come si vedrà nella
sezione~\ref{a1:sec:scopo}, i docenti considerano questa varietà una risorsa.

% ---------------------------------------------------------------------
\section{I materiali e gli strumenti del corso}
\label{a1:sec:materiali}

Prima delle slide sul corso, Daniele sbriga un po' di ``clean up casalingo''.

\paragraph{Il form iniziale.} C'è un \textbf{form} da compilare ``entro e non
oltre il 15 settembre'': il QR code è nelle slide, e la segreteria lo ricorderà
su Discord. Richiede circa 10--15 minuti; si può fare subito o in giornata, ``e
ve lo togliete di torno''. È un form standardizzato per tutti i corsi (``non
siamo io e Ruggero a farvi passare la tortura del test''), diviso in due parti:
\begin{itemize}
  \item una parte di \textbf{conoscenza del candidato}: quanto tempo si ha a
        disposizione per il corso, come ci si organizza con il lavoro, che cosa
        piace giocare. È la parte che interessa ai docenti, perché li aiuta a
        organizzare le lezioni;
  \item una parte di \textbf{test logico}, che i docenti non hanno scritto (``io
        non sono responsabile del form'', ``no, no, nemmeno io'').
\end{itemize}

\paragraph{La cartella condivisa.} Tutti gli studenti vengono aggiunti alla
cartella del corso su \textbf{Google Drive}, che si trova tra i file
``condivisi con me''. Gli studenti ne sono \textbf{gestori dei contenuti} e
possono caricarci quello che vogliono. Dentro ci sono:
\begin{itemize}
  \item le \textbf{attività in classe};
  \item le \textbf{dispense};
  \item i \textbf{materiali didattici}, con le slide (quelle della lezione sono
        già lì);
  \item le cartelle dei \textbf{project work} futuri;
  \item i \textbf{task}, che ``saranno pochi, ma ci saranno'';
  \item il \textbf{calendario studenti}, ``importantissimo'', con due schede: il
        calendario delle lezioni e quello del project work.
\end{itemize}

\paragraph{Discord.} Il corso ha anche un \textbf{server Discord}, a cui tutti
sono stati invitati; ci sono i docenti (Daniele, Ruggero, Diego e altri), che si
possono taggare. Di solito è lì che la segreteria pubblica le informazioni, per
esempio se una lezione ``dovesse per caso essere cancellata''. Nelle domande
finali Ruggero aggiungerà che c'è un canale apposito per il corso, che gli
studenti possono usare liberamente per organizzarsi tra loro.

\paragraph{Le reaction di Zoom.} Durante le lezioni su Zoom si possono usare le
\textbf{reaction} (il pollice in alto, il cuoricino, la mano alzata): un modo
rapido per far capire ai docenti se si è capito, se si vuole intervenire. Nel
frattempo arrivano i primi pollici.

% ---------------------------------------------------------------------
\section{Lo scopo e il metodo del corso}
\label{a1:sec:scopo}

La prima serie di slide (che Daniele, per un lapsus, presenta come ``cos'è il
corso di pulire il pavimento di casa'') spiega che cosa il corso vuole fare e
come.

\subsection{Diventare autori analogici}

Lo \textbf{scopo del corso} è fornire le competenze per \textbf{diventare autori
analogici} e stimolare la creatività e l'\textbf{autorialità} di ciascuno: ``voi
imparerete a diventare autori di giochi analogici''. La differenza tra
\emph{game designer} e \emph{game developer} verrà chiarita più avanti (sezione~\ref{a1:sec:design-development}),
ma la parola chiave qui è \emph{autorialità}: ``vogliamo voi, vogliamo le vostre
idee''. Magari non si potrà usare l'idea del gioco della propria vita, ma si
potranno portare nelle lezioni le proprie idee e le proprie sensibilità.

Gli obiettivi elencati nella slide sono questi:
\begin{itemize}
  \item dare \textbf{maggiore consapevolezza nel design e nello sviluppo}. In
        Italia, purtroppo, ``si fa design a volte in maniera non consapevole''.
        Consapevole vuol dire avere almeno un'introduzione alle basi del design e
        dello sviluppo, e fare le cose ``con metodo, con raziocinio, con
        consapevolezza'';
  \item fornire gli \textbf{strumenti essenziali} per un design consapevole;
  \item stimolare la \textbf{sperimentazione} e la \textbf{trasversalità
        mediatica}. Non si chiede di fare ``il \emph{Root} della situazione o il
        \emph{Carcassonne} della situazione''\footnote{La trascrizione dice ``il
        lumaven della situazione o il carcassone della situazione''; il primo
        titolo non è ricostruibile con certezza. Il senso è chiaro: non si
        chiede di rifare un grande classico del gioco da tavolo.}, ma un
        project work \textbf{calibrato sulla durata del corso}, e però
        \textbf{sperimentale} e ``assolutamente transmediale'';
  \item portare \textbf{numerosi esempi} di tipologie di gioco, dall'analogico al
        tavolo, al pervasivo ``e oltre''. Lo si vedrà nella lezione sui
        \textbf{casi studio} del venerdì successivo: ``voi non avete idea di
        quante cose esistono, bellissime, sconosciutissime e favolose''. Daniele
        confessa che preparando e aggiornando le slide i colleghi gli hanno
        mostrato cose che lui stesso non sapeva nemmeno che esistessero;
  \item dare gli strumenti per \textbf{analizzare} i giochi altrui. Ruggero lo
        sottolinea: se il corso riesce nel suo intento, gli studenti sapranno
        non solo fare i propri giochi ma anche analizzare giochi già fatti, che è
        ``una delle skill importanti anche nella parte di ricerca e sviluppo''.
        Si vedranno casi studio e giochi pubblicati, sempre con l'obiettivo di
        capire \textbf{che cosa fanno e perché lo fanno}, per poi fare meglio il
        proprio lavoro;
  \item imparare a \textbf{fare giochi insieme}. Esistono gli autori che lavorano
        da soli (i \emph{solo developer}), ``ma sappiate che anche l'autore che
        lavora da solo poi si deve interfacciare con il pubblico, con un editore,
        magari con un editor, con un publisher''. ``Non lavorerete mai da soli'':
        bisogna imparare a lavorare con gli altri.
\end{itemize}

\subsection{Un corso intersezionale}

Il corso è \textbf{intersezionale}: \textbf{non ci sono lezioni dedicate a una
singola tipologia di gioco}. Non ci sarà ``la lezione sul gioco da tavolo'',
``la lezione sul gioco di ruolo'', ``la lezione sul gioco pervasivo''. Ci saranno
lezioni su temi diversi, ciascuna pensata per essere utile a \textbf{tutte e tre
le sfaccettature} del corso; alcune saranno più importanti per il gioco da
tavolo, altre per il pervasivo, altre per il gioco di ruolo, ma nessuna sarà
esclusiva. I docenti vogliono ``contaminazioni e applicazioni assolutamente
trasversali''.

\subsection{Un corso cross-generazionale}

È anche un corso \textbf{cross-generazionale}: i gradi di esperienza di gioco
quest'anno sono abbastanza allineati, ma le esperienze di vita, le età, le
estrazioni sociali sono varie. ``Ed è un pro, non è un contro. Più diversificata
è la classe, più imparate meglio anche voi''. La slide ha una piccola nota per
fasce d'età:
\begin{description}
  \item[Under 30] Concentratevi sulla vostra \textbf{crescita}. Evitate di
        confrontarvi con le persone più grandi che magari ``ce la fanno meglio''
        in alcune cose.
  \item[Over 30] Concentratevi sul \textbf{consolidamento}, ma siate
        \textbf{umili}: i docenti diranno cose che faranno ``storcere il naso''.
        ``Bisogna un po' reimparare per imparare'': bisogna sapersi adattare.
\end{description}
Ruggero contestualizza: lo dicono perché in passato hanno visto studenti che,
essendo più adulti o avendo giocato di più, pensavano che l'esperienza
\emph{di gioco} si trasformasse automaticamente in competenza \emph{di
creazione} di giochi. Non è così. E non bisogna sottovalutare chi gioca meno:
proprio perché gioca cose diverse, per risolvere un problema di design specifico
``è una ricchezza incredibile avere qualcun altro che riesce a uscire dalla
scatola''.

\subsection{Un corso pratico e adattabile}

Il corso è \textbf{pratico}: gli studenti si concentreranno su un progetto di
gruppo, il \textbf{project work}, lavorando in team. Ed è \textbf{adattabile}: il
programma è definito, ma ``si aggiorna sui vostri stimoli, priorità e
necessità''. Se durante il corso i docenti vedono che serve più una cosa
dell'altra, sono ``prontissimi'' a dare più spazio a un argomento invece che a
un altro.

\subsection{Inquiry-based learning: domandate}

Il metodo è quello dell'\textbf{inquiry-based learning}, l'apprendimento basato
sulle domande. Di solito le persone dicono ``ho una domanda, magari sarà
stupida''. ``Non esistono domande stupide. Pensate che ci sia una domanda
stupida? Fatela, alzate la mano''. Vanno bene anche le
\textbf{provocazioni}: se qualcosa non torna, è meglio chiedere. La condivisione
è fondamentale.

I docenti danno \textbf{metodi, dati e conoscenze}, ma lasciano trovare le
risposte agli studenti: ``vi lasciamo un pochettino annegare nel mare prima di
darvi la soluzione, in modo tale che la possiate trovare da soli''. Di
conseguenza \textbf{non ci saranno interrogazioni sulle definizioni}: interessa
vedere \textbf{come} gli studenti applicano ciò che imparano. Se una
definizione non viene in mente, non è un problema; l'importante è saperla usare.

\subsection{Feedback, non valutazione}

``Sì al feedback, no alla valutazione''. I docenti danno \textbf{feedback
verbali}, anche sulle attività da fare a casa (i task, che non saranno tanti).
Il corso è calibrato per \textbf{studenti lavoratori}: c'è un po' da lavorare
fuori dalle lezioni, ma in modo calibrato.

Il principio è \textbf{fallire oggi per imparare per il futuro}: \emph{fail
forward}, \emph{fail to learn}. I docenti lasciano provare e fallire almeno una
volta prima di dare la soluzione: ``il classico metodo Montessori, quello
originale, non quello travisato''.

\subsection{Le soft skill e i criteri di valutazione}

Grande attenzione va alle \textbf{soft skill}: i lavori di gruppo, lo scambio
propositivo, l'atteggiamento positivo, la propositività, la proattività.
``Lavoro di gruppo'' non significa stare ogni giorno in sala studio con gli
altri: va benissimo che un gruppo si organizzi dandosi priorità e tempi in base
alla disponibilità di ciascuno. L'importante è \textbf{collaborare}, essere
\textbf{professionali}, rispettarsi a vicenda, rispettare le idee proprie e
altrui e lasciare spazio agli altri, perché i progetti ``non saranno solo
vostri, saranno di tutte le persone del vostro gruppo''.

La valutazione del project work non riguarda solo la \textbf{documentazione}
prodotta, ma anche:
\begin{itemize}
  \item la \textbf{presentazione} finale;
  \item l'\textbf{autovalutazione};
  \item la \textbf{crescita personale};
  \item l'\textbf{atteggiamento} tenuto nel corso.
\end{itemize}
Daniele cita: ``da ognuno secondo le sue capacità, a ognuno secondo i suoi
bisogni''. Lavorare in gruppo non vuol dire che, in tre, ciascuno lavori
esattamente il 33\% del tempo. Ai docenti serve piuttosto capire \textbf{chi fa
che cosa} nel gruppo, e che il processo interno sia \textbf{trasparente}.

\paragraph{Lavorare in cloud.} Per lo stesso motivo il consiglio è di
\textbf{non lavorare in locale}, ma direttamente nelle cartelle condivise del
project work e dei task. Si è liberi di fare altrimenti, ma lavorando in cloud
si dà ai docenti la possibilità di vedere \emph{come} si lavora, il proprio
``processo mentale''.

% ---------------------------------------------------------------------
\section{Le regole di comportamento e il carico di lavoro}
\label{a1:sec:regole}

Alcune slide sono di quelle che si ripetono in ogni corso (``\emph{repetita
iuvant}'', storpiato per scherzo in ``ripetita Juventus'').

\paragraph{Comportamento.} Prestare attenzione durante le lezioni; interagire;
portare rispetto; essere proattivi; alzare la mano e fare domande.
\textbf{Tolleranza zero} per pregiudizi, maleducazione, mancanza di rispetto,
\emph{cheat} e scorciatoie. Daniele ammette che ``cheat e scorciatoie'' viene da
altri corsi dove il tema è più importante, ma Ruggero lo ferma: ``non
sottovalutate questa cosa'', perché implica anche \textbf{non far fare a ChatGPT
i task}. Daniele ha una slide apposta, ``perché è un discorso molto complesso nel
2026'' (sezione~\ref{a1:sec:ia}).

\paragraph{Le lezioni sono intense.} Le lezioni possono essere \textbf{dense},
``soprattutto le mie, sappiatelo: io sono logorroico'', avverte Daniele. Bisogna
chiedere spiegazioni se serve, prendere appunti, rivedere le registrazioni. Le
lezioni sono ``delle scintille'': modi per far esplorare i temi oltre le slide.
``Andate oltre, se riuscite''.

\paragraph{Il carico di lavoro.} Le regole sui task sono queste:
\begin{itemize}
  \item i task \textbf{sono per voi, non per noi};
  \item se si ha poco tempo, è meglio farli rapidamente piuttosto che non farli:
        ``provateci'';
  \item i docenti cercano di non chiedere più di \textbf{un'ora o due a
        settimana} oltre alle ore di lezione;
  \item non tutte le lezioni avranno un task, e il \textbf{project work} è il più
        importante di tutti.
\end{itemize}

\paragraph{I feedback dei docenti.} I feedback ``si guadagnano'' man mano che si
reagisce bene: non bisogna aspettarsi ``dei papiri'' sui task. Se i docenti
hanno voglia di dare feedback più ampi lo fanno, soprattutto se vedono che lo
studente se li merita. Ad Analog di solito non ci sono task \emph{rejected}; ma
se un docente rifiuta un task, vuol dire che si è fatto \textbf{meno del livello
minimo} richiesto dal corso. I feedback possono non essere immediati: ``siate un
pochettino impietositi dal fatto che siamo dei docenti''. E se un feedback non
arriva, o se ne vogliono di più, bisogna scrivere, ``perché siamo tre esseri
umani e quindi possiamo anche dimenticarci''.

% ---------------------------------------------------------------------
\section{Il project work e il calendario}
\label{a1:sec:project}

\subsection{Come funziona il project work}

Una prima slide con le date del project work è da aggiornare (``perdonatemi il
refuso''); Daniele la corregge al volo e rimanda alla slide successiva, più
precisa. Il funzionamento è questo:
\begin{itemize}
  \item i \textbf{gruppi li formano i docenti}. Gli studenti sono nove, quindi si
        faranno \textbf{tre gruppi da tre}: una dimensione pensata anche per il
        fatto che molti sono studenti lavoratori, così da non ritrovarsi ``in due,
        a esplodere'';
  \item ogni gruppo deve prima preparare un \textbf{pitch} del proprio gioco, da
        presentare in anticipo: un'idea, una bozza del gioco;
  \item poi si lavora al gioco \textbf{nelle pause del corso}: da qui a dicembre
        ci saranno due pause, quindi non ci sarà lezione tutti i lunedì e tutti i
        venerdì;
  \item ci saranno \textbf{review intermedie} con i \textbf{tutoraggi}, per lo più
        il \textbf{sabato mattina}, dalle 11 alle 12, con Ruggero su Google Meet.
        Servono ad aiutare i gruppi a portare avanti il project work.
\end{itemize}

\subsection{Il calendario delle lezioni}

Il \textbf{calendario} è nel calendario studenti della cartella condivisa.
Attenzione: le date \textbf{possono variare}. ``Il 90\% rimarranno così'', ma
può capitare che, per motivi di lavoro, qualcuno stia male o non possa, o che
``scompaia l'elettricità da tutte le città d'Italia''. E l'ordine delle lezioni
può cambiare se i docenti vedono che agli studenti serve prima una lezione
invece di un'altra.

Ci sono \textbf{due pause importantissime}:
\begin{itemize}
  \item dal \textbf{13 ottobre al 5 novembre}: è il periodo delle grandi fiere
        (\textbf{Essen} e \textbf{Lucca}\footnote{La trascrizione dice ``Essen,
        Luca, Cose'': si tratta con ogni probabilità della fiera Spiel di Essen,
        in Germania, e di Lucca Comics \& Games, che si tengono entrambe in
        ottobre.}), ``un periodo un po' frenetico'', e cominciano anche i corsi
        accademici, quindi la pausa permette ai docenti di respirare. Per gli
        studenti è la pausa più importante: quasi un mese in cui sviluppare la
        parte importante del gioco;
  \item dal \textbf{24 novembre all'11 dicembre}, dopo qualche lezione a novembre.
\end{itemize}

E ci sono \textbf{due date da evidenziare}:
\begin{itemize}
  \item il \textbf{25 settembre}: nei primi 5--10 minuti della lezione (che per il
        resto sarà tenuta da Beatrice) i docenti annunceranno i tre gruppi;
  \item il \textbf{12 ottobre}: il \textbf{pitch} dei giochi, con l'idea
        embrionale. Tra il 25 settembre e il 12 ottobre ogni gruppo prepara un
        breve pitch in formato slide; quel giorno i docenti lo rivedono e danno
        feedback: se il gioco va bene, se non va bene, come svilupparlo.
\end{itemize}
Dal 12 ottobre fino all'inizio di dicembre si sviluppa il gioco, con i
tutoraggi: ``sfruttate questi momenti di pausa''.

\subsection{Il calendario del project work}

La slide successiva, ``reale e aggiornata'', riporta il calendario del project
work (tabella~\ref{a1:tab:calendario} e figura~\ref{a1:fig:calendario}).
``Per ora non pensate al project work: ci penserete dal 25 settembre in poi''.

\begin{table}[htbp]
\centering
\begin{tabular}{@{}lp{9.2cm}@{}}
\toprule
\textbf{Data} & \textbf{Che cosa succede} \\
\midrule
ven 25 settembre\footnotemark & lezione: annuncio dei gruppi e del
\emph{mandate}, cioè di che cosa si dovrà fare \\
sab 3 ottobre, ore 11 & prima review e tutoraggio con Ruggero (dopo una
decina di giorni di lavoro) \\
lun 12 ottobre & pitch del gioco durante la lezione \\
13 ottobre -- 5 novembre & prima pausa: sviluppo del gioco \\
sab 7 novembre & tutoraggio \\
sab 14 novembre & tutoraggio \\
24 novembre -- 11 dicembre & seconda pausa: sviluppo del gioco \\
ven 4 dicembre, entro mezzanotte & consegna del gioco sul Drive (non c'è
lezione) \\
ven 11 dicembre & i giochi vengono visti a lezione \\
14 dicembre & giornata finale, con la tavola rotonda con gli editori \\
\bottomrule
\end{tabular}
\caption{Il calendario del project work come lo presenta Daniele. Le date
possono variare.}
\label{a1:tab:calendario}
\end{table}
\footnotetext{Daniele dice ``lunedì 25 settembre'', ma nel 2026 il 25 settembre
è un venerdì; il giorno della settimana nella tabella è corretto di
conseguenza.}

\paragraph{Perché la consegna è il 4 dicembre.} Il 4 dicembre non c'è lezione:
si deve solo consegnare il gioco, entro mezzanotte, sul Drive. La presentazione
a lezione è l'11 dicembre, ma i tre docenti hanno bisogno di \textbf{una
settimana} per rivedere il gioco in privato e valutare anche la
\textbf{documentazione}; per questo la consegna non coincide con la
presentazione. Che cosa voglia dire esattamente ``consegnare il gioco'' verrà
spiegato più avanti.

\paragraph{La giornata finale.} Il 14 dicembre c'è la \textbf{giornata finale},
con una \textbf{tavola rotonda con gli editori}: di solito viene invitato Flavio,
che è anche docente di game production, più qualche altro editore, ``vediamo
quest'anno chi riusciamo a portarvi''.

\begin{figure}[htbp]
\centering
\begin{tikzpicture}[
    font=\footnotesize,
    x=0.1cm,
    evento/.style={draw, thick, rounded corners=2pt, fill=white,
                   text width=1.8cm, align=center, inner sep=2pt},
    pausa/.style={fill=gray!20, draw=gray!60},
    tacca/.style={thick}
  ]
  % asse: giorno 0 = 25 settembre; 1 unità = 1 giorno
  \draw[pausa] (18,-0.35) rectangle (41,0.35);
  \draw[pausa] (60,-0.35) rectangle (77,0.35);
  \node[font=\scriptsize\itshape] at (29.5,0.18) {pausa};
  \node[font=\scriptsize\itshape] at (65,0.18) {pausa};
  \draw[-{Stealth[length=2.5mm]}, thick] (-3,0) -- (85,0);
  % eventi sopra
  \draw[tacca] (0,-0.15) -- (0,0.9);
  \node[evento, above] at (0,0.9) {25 set\\ gruppi e \emph{mandate}};
  \draw[tacca] (17,-0.15) -- (17,2.2);
  \node[evento, above, fill=gray!8] at (17,2.2) {12 ott\\ \textbf{pitch}};
  \draw[tacca] (70,-0.15) -- (70,0.9);
  \node[evento, above, fill=gray!8] at (70,0.9) {4 dic\\ \textbf{consegna}};
  % eventi sotto
  \draw[tacca] (8,0.15) -- (8,-0.9);
  \node[evento, below] at (8,-0.9) {3 ott\\ tutoraggio};
  \draw[tacca] (43,0.15) -- (43,-0.7);
  \draw[tacca] (50,0.15) -- (50,-0.7);
  \draw (43,-0.7) -- (50,-0.7);
  \draw[tacca] (46.5,-0.7) -- (46.5,-0.9);
  \node[evento, below] at (46.5,-0.9) {7 e 14 nov\\ tutoraggi};
  \draw[tacca] (77,0.15) -- (77,-0.9);
  \node[evento, below] at (77,-0.9) {11 dic\\ giochi a lezione};
  \draw[tacca] (80,-0.15) -- (80,2.2);
  \node[evento, above] at (80,2.2) {14 dic\\ giornata finale};
\end{tikzpicture}
\caption{Una linea del tempo del project work, ricostruita dalle slide (scala
approssimativa: dal 25 settembre al 14 dicembre). Le bande grigie sono le due
pause senza lezioni, in cui i gruppi sviluppano il gioco.}
\label{a1:fig:calendario}
\end{figure}

Finita questa parte, Daniele apre alle domande: ``Domande sul calendario? Sulle
slide che ho fatto? Considerazioni? Non abbiate timore, io non ci credo''. E la
domanda arriva, sull'unico tema che aveva lasciato da parte.

% ---------------------------------------------------------------------
\section{L'intelligenza artificiale nel corso}
\label{a1:sec:ia}

\subsection{La slide di Daniele}

Daniele aveva già presentato una slide che definisce ``importantissima''. ``Lo so
che è paradossale, me ne rendo conto. Lo so che magari mi darete del luddista,
dell'anti-tecnologia: non è così. Uso l'intelligenza artificiale per lavoro. Ma
in questo corso, vi prego, \textbf{non usate l'intelligenza artificiale}''. La
richiesta vale per le attività, per i task e \textbf{anche per la
prototipazione}. Le ragioni:
\begin{itemize}
  \item ``voglio che siate voi a fare il gioco, voglio che siate voi a imparare a
        farlo''. Daniele parla da utilizzatore: prima di ChatGPT, che ha smesso
        di usare, oggi soprattutto di Claude, di cui ``si sottovaluta quanto
        possa essere incerto, insicuro, impreciso, non all'altezza'';
  \item nella game industry, \textbf{i publisher rifiutano automaticamente} i
        pitch e i prototipi fatti con l'intelligenza artificiale. Daniele lo ha
        visto di persona in ``tantissime pitch review'' quell'estate, anche al
        \textbf{First Playable} di Firenze\footnote{La trascrizione dice ``First
        Playboy a Firenze'': si tratta del First Playable, l'evento di business
        dell'industria videoludica che si tiene a Firenze.}: anche nel mondo dei
        videogiochi i publisher che vedono pitch o prototipi fatti con l'IA
        ``rigettano automaticamente la build'', perché hanno ancora timore
        dell'IA e delle sue possibili conseguenze;
  \item ``non so come sarà il futuro'': non sa se bisognerà abituarsi a usarla
        tutto il giorno. Non è ``totalmente contrario'', purché sia usata ``con
        raziocinio''; secondo lui oggi, negli Stati Uniti, non lo è;
  \item ma in un corso che si chiama \emph{Analog Game Design} gli studenti devono
        saper fare i giochi da soli. In futuro, se vorranno farsi aiutare
        dall'IA per essere più efficienti, saranno liberissimi di farlo; durante
        il corso no, perché i docenti vogliono valutare \textbf{le loro
        competenze e capacità}.
\end{itemize}

\subsection{La domanda: che cosa vuol dire usare l'IA in un gioco analogico?}

La domanda di uno studente è precisa: che cosa si intende per ``task fatti con
l'IA''? Qual è l'uso dell'IA in un gioco analogico?

\paragraph{Ruggero: l'esperienza di chi valuta i prototipi.} Ruggero risponde
dalla sua esperienza nel settore, dove riceve prototipi da valutare per
un'eventuale pubblicazione. Distingue due usi:
\begin{enumerate}
  \item le \textbf{immagini generate}: illustrazioni, graphic design e così via.
        Per molti publisher è ``un \emph{reject} immediato, proprio immediato'';
  \item l'IA usata per il \textbf{regolamento}, per scrivere le regole e le
        meccaniche. È più difficile da individuare, ma molti publisher se ne
        accorgono comunque, perché le IA ``non sono ancora brave abbastanza da
        creare un regolamento coeso meccanicamente e scritto con un senso per il
        vostro gioco''. Quello che fanno è prendere parti di altri regolamenti,
        elaborare le strutture che vedono nei regolamenti pubblicati e
        riapplicarle. Ma molti giochi si prestano poco a un regolamento
        strutturato come quello di altri, perché \textbf{ogni gioco è un
        ecosistema meccanico a sé}. Il risultato sono regolamenti con
        riferimenti non immediati, e soprattutto con \textbf{la stessa meccanica
        descritta in modo diverso in pagine diverse}: due difetti che fanno
        perdere tempo all'editore, e in questo caso ai docenti.
\end{enumerate}
Poi aggiunge un punto ``personale, soggettivo'': i docenti sono interessati a
vedere il lavoro, la creatività, ``il frutto della vostra creazione''. E
viceversa gli studenti vogliono i commenti dei docenti: ``immaginatevi se voi
fate un gioco e noi vi diamo i commenti scritti con ChatGPT''. Altrimenti, dice
ridendo, gli studenti fanno il gioco con ChatGPT, i docenti lo danno a ChatGPT
per i commenti, ``e facciamo molto alla svelta''.

\paragraph{Daniele: perché gli studenti usano l'IA.} Daniele racconta la sua
esperienza con gli studenti, che secondo lui ricorrono all'IA in due momenti:
\begin{enumerate}
  \item quando vanno in \textbf{crisi}: in ansia da prestazione o per i tempi.
        Cercano di risolvere il problema della propria disorganizzazione
        facendolo risolvere all'IA, senza accorgersi che l'output è ``molto,
        molto scarso'';
  \item quando hanno l'impressione che l'intelligenza dell'IA sia
        \textbf{migliore} di quello che scrivono loro, o di quello che c'è in
        giro. Per Daniele è una ``mancanza di palato critico'': l'output dell'IA
        sembra buono in superficie, ma quando si va nel dettaglio e ci si fa
        delle domande ``di solito fa un po' di \emph{crumbling}'', si sbriciola.
\end{enumerate}
Lui usa Claude ogni giorno e ne conosce i limiti: ``non c'è nemmeno una frase
che io copio direttamente''. Gli fa fare il lavoro più noioso, più ripetitivo e
soprattutto \textbf{strutturato}, ``replicare delle cose come se fosse una
macro''. Lo usa per il testo e per il codice, mai per le immagini: per motivi di
\textbf{diritto d'autore}, e perché ha molti amici artisti ed è stufo di vedere
il loro lavoro ``dilaniato'' dall'IA.

Ammette che i docenti non possono controllare gli studenti; per questo
``ve lo consiglio''. E riprende quella che definisce la cosa più bella detta da
Ruggero (``ti ruberò la frase per il corso di narrative''): i docenti i feedback
li danno di persona. Daniele a volte manda 30 minuti di audio a uno studente, o
fa intere call per rivedere script, pitch e giochi: ``è il mio lavoro, mi piace
farlo e voglio farlo io, quindi mi aspetto che lo studente lo faccia lui''. Come
freelance, poi, è ``stufo marcio'' dei clienti che arrivano con qualcosa fatto
con l'IA (``guarda che bello, pubblica questo''), perché il più delle volte il
risultato è pessimo: non sanno come usarla. ``È anche una questione di
competenza''.

\paragraph{Il prototipo non deve essere bello.} Questo vale anche per il
\textbf{prototipo}, su cui si tornerà nelle lezioni di prototipazione: in un
prototipo non servono ``graficoni allucinanti''. Le grafiche bellissime fatte con
l'IA spesso \textbf{peccano di user experience}: sembrano bellissime, ma quando
il giocatore deve leggere la carta non capisce qual è il costo, qual è il
potere, qual è l'effetto.

\paragraph{Beatrice: si perde l'occasione di imparare.} Beatrice aggiunge due
motivi che valgono ``al di là dei nostri credo''. Il primo: chi usa l'IA in
sostituzione (per esempio per scrivere un testo, per accelerare i tempi o per
trovare strategie) \textbf{perde la parte in cui può imparare}. Nel game design
bisogna ``tenere insieme i pezzi'' di un regolamento, ma anche dell'esperienza
dell'utente e di come trasmetterla; i passaggi necessari (creare il regolamento,
creare le meccaniche, sperimentare, raccogliere i feedback) ``si capiscono solo
vivendoli''. E si sta pagando un corso per imparare. Il secondo, sul piano
editoriale: chi presenta a un editore un prodotto che ha come base l'IA rischia
da un lato di non avere imparato abbastanza (e di ritrovarsi ``a piedi'' sul
campo senza IA), dall'altro che l'editore se ne accorga, ed è un tema
\textbf{reputazionale}. Ci sono editori ``più illuminati di altri'', quindi
anche qui c'è un mercato da capire; ma i docenti preferiscono dire che cosa è
più sensato per il mercato come lo conoscono.

\paragraph{Ruggero: non abbiate paura del giudizio.} Ruggero vuole dare uno
``spotlight'' a un punto detto da Daniele e rimasto ``semi-ignorato'': a volte
l'uso dell'IA nasce dalla \textbf{paura del giudizio}, dal voler consegnare
qualcosa di più pulito, più strutturato possibile. ``Vi giuro, non ve ne
preoccupate''. Ruggero detesta i voti e non ragiona per voti; non valuterà
nessuno per grafiche fatte a matita. Insegna da anni, ha avuto studenti con
dislessia e discalculia: refusi ed errori non gli interessano, ``non è un corso
di editing o di proofreading''. Gli interessa che le \textbf{meccaniche siano
coerenti con lo scopo di design}. Immagini fatte a matita o con Paint ``non sono
un problema, anzi, le preferisco'' rispetto all'``ennesima immagine fatta con
ChatGPT o Midjourney''.

\paragraph{Daniele: il pensiero critico, costo e qualità.} L'ultimo intervento
è di Daniele, sul \textbf{pensiero critico}. Ragionando per assurdo (``un
pensiero di Belzebù''), immagina che un domani l'IA venga usata ovunque. I suoi
limiti, secondo lui, resterebbero due:
\begin{description}
  \item[Il costo] Tutti dicono che sta diminuendo e che diminuirà, ma non se ne
        può essere sicuri al 100\%.
  \item[La qualità] Cioè la capacità di leggere l'output. Se si dà un regolamento
        a Claude o a ChatGPT chiedendo quali problemi ha, il modello risponde con
        i problemi che, in percentuale, è più probabile che gli esseri umani
        abbiano già individuato nei dati su cui è stato addestrato. Ne derivano
        due problemi: se il problema del proprio gioco non è mai stato
        individuato prima, il modello non lo vedrà; e viceversa può segnalare
        problemi che il gioco non ha.
\end{description}
Fa l'esempio di Grammarly, che usa nella versione premium per rivedere i testi
in inglese: spesso non accetta i suoi consigli, perché se vuole scrivere un
dialogo di videogioco in \textbf{dialetto Cockney} londinese, lo deve scrivere
in Cockney. Se Grammarly dice che è sbagliato, ``sei tu che non hai capito il
motivo per cui voglio scrivere questo testo in Cockney''. La conclusione: prima
di usare un'IA per un'attività bisogna sapere \textbf{che cosa si vuole
ottenere}, qual è l'obiettivo, e quali \textbf{competenze} servono per valutare
in modo critico se l'output di un qualunque LLM è davvero valido. Daniele non
vuole dilungarsi oltre: è un argomento ``molto caldo'' per lui, che è anche
scrittore e si sente dire da tutti ``ti sostituiranno''.

Subito dopo si presenta finalmente Gianmarco (lo si è visto nella
sezione~\ref{a1:sec:partecipanti}), e poi si fa una pausa di una decina di
minuti. Daniele avverte che nella terza parte lui non ci sarà (alle sette e un
quarto ha un impegno di lavoro), e lo stesso vale per Beatrice: la parte finale
di domande e risposte sarà con il solo Ruggero, ``in ottime mani''.

% ---------------------------------------------------------------------
\section{Le figure professionali dell'industria analogica}
\label{a1:sec:figure}

Dopo la pausa (in cui si scopre che Angelo lavora in un piccolo birrificio in
Brianza e si scherza sul rischio di un \emph{drinking game}, tra la birra e la
vendemmia) la lezione riprende con le slide di Ruggero. ``Oggi è una lezione
introduttiva'': gli argomenti si vedranno ``molto dall'alto'' e verranno
approfonditi nelle lezioni successive. Ma servono già ora, per creare
\textbf{l'inizio di un linguaggio condiviso}: ci sono termini che torneranno
spesso, a cominciare da quelli del mercato e dell'industria del gioco
analogico, un'industria ``molto varia, molto caotica, anche molto mutevole'',
con tanti tipi diversi di giochi e altrettanti tipi diversi di entità che ci
lavorano.

\subsection{Una decina di etichette}

I docenti hanno raccolto ``una decina di etichette'', cioè di \textbf{ruoli} che
nell'industria sono o possono essere presenti. Due avvertenze:
\begin{itemize}
  \item nelle aziende più piccole \textbf{la stessa persona copre spesso più
        ruoli}, per forza di cose. Sarà anche il caso degli studenti, che
        lavoreranno in gruppi di tre e non di venti: è facile che ciascuno finisca
        a fare più cose. I docenti ne terranno conto, e qui entreranno in gioco
        anche competenze diverse dal game design (programmi di grafica,
        illustrazione, impaginazione, layout, \emph{proofreading}\ldots);
  \item Daniele precisa che \textbf{non si chiede agli studenti di diventare}
        illustratori o altro. Ruggero concorda, ma aggiunge che è utile sapere
        anche \emph{che cosa non è il proprio lavoro}. Se in un gruppo nessuno sa
        illustrare, ai docenti va benissimo un prototipo che mostri dove vanno
        le immagini: \textbf{placeholder}, concept, anche semplici box di testo
        con scritto ``qui vorremmo un'immagine di\ldots''. ``È importante anche
        sapere che cosa sapete fare e che cosa non sapete fare''.
\end{itemize}
La tabella~\ref{a1:tab:figure} riassume le figure; le sottosezioni seguenti le
riprendono una per una, con i commenti dei docenti.

\begin{table}[htbp]
\centering
\small
\begin{tabular}{@{}p{3.3cm}p{9.6cm}@{}}
\toprule
\textbf{Figura} & \textbf{Di che cosa si occupa} \\
\midrule
Game designer & crea le regole alla base dell'esperienza: le meccaniche, come
interagiscono tra loro e come chi gioca interagisce con esse \\
Game developer & fa da ponte tra chi fa game design e il publisher: bilancia il
prototipo, cura i modelli matematici, adatta il gioco (e il paragame) alle
necessità del publisher \\
Editor (proofreader) & revisiona e adatta i testi a partire dalla prima stesura
del designer (correzione, linguaggio inclusivo\ldots) \\
Publisher (editore) & pubblica e produce fisicamente il gioco; il suo marchio è
sulla scatola o sulla copertina \\
Sales e marketing & il lato commerciale: vendite (anche B2B) e comunicazione del
prodotto \\
Graphic designer & icone meccaniche, wireframe, UX: il lato funzionale degli
elementi visivi \\
Illustratore & le illustrazioni: il \emph{look and feel}, il world building, il
``vibe'' del gioco \\
Sensitivity reader (cultural consultant) & verifica l'accuratezza e la
rappresentazione di un tema sensibile che il team non conosce \\
Producer & coordina il progetto (quasi project management), parla con le
tipografie, verifica che i materiali siano sostenibili economicamente \\
\bottomrule
\end{tabular}
\caption{Le figure professionali dell'industria del gioco analogico presentate
nella slide di Ruggero, in sintesi.}
\label{a1:tab:figure}
\end{table}

\subsection{Game designer}

Il \textbf{game designer} (o la game designer) è chi si occupa, ``sembra banale
a dirlo'', di \textbf{creare le regole alla base dell'esperienza} che si vuole
far arrivare. Vale per tutti i tipi di giochi, non solo per quelli analogici, e
che cosa faccia concretamente in un progetto o in una serie specifica ``può
mutare tantissimo''. Nell'analogico non verrà chiesto di programmare o di
prototipare in digitale; verrà chiesto di \textbf{creare le meccaniche di
gioco}, di capire \textbf{come le meccaniche si interfacciano tra loro} e
\textbf{come le persone che giocano si interfacceranno con le meccaniche}.

\subsection{Game developer}

Il \textbf{game developer} (o la game developer) sarà l'argomento delle due
lezioni di Marta. La parola \emph{development} è nota soprattutto
dall'ambito videoludico, ma nell'analogico ha un significato molto specifico
``che non c'entra niente con la programmazione'': il game developer è
\textbf{un ponte tra chi fa game design e il publisher}. Di solito si occupa di
\textbf{bilanciare il prototipo}, di curare i \textbf{modelli matematici}, di
adattare il gioco e il \textbf{paragame}\footnote{La trascrizione dice ``il
Wix, il paragame''. Il \emph{paragame} è, in termini generali, tutto ciò che
accompagna il gioco senza esserne il sistema di regole (regolamento, scatola,
materiali, comunicazione); i docenti lo tratteranno in una lezione dedicata, e
qui il termine è solo citato.} alle necessità del publisher. La distinzione
viene ripresa e approfondita nella sezione~\ref{a1:sec:design-development}.

\subsection{Editor}

L'\textbf{editor}, a volte chiamato \textbf{proofreader} a seconda dell'azienda,
si occupa della revisione dei \textbf{testi}: non della stesura iniziale, perché
di norma, data la scala dell'industria, è chi fa game design a scrivere un
\textbf{primo draft}, e l'editor ci passa sopra. Non solo per correggere i
refusi: spesso il suo lavoro è anche \textbf{adattare} i testi, per esempio
renderli \textbf{gender neutral}. Se ne vedranno casi studio più avanti.

\subsection{Publisher}

Il \textbf{publisher}, o \textbf{editore}, è l'entità che nell'ambito
commerciale \textbf{pubblica} giochi, giochi di ruolo o giochi da tavolo. È una
figura ``più solidificata'' nel gioco da tavolo e nel gioco di ruolo da tavolo;
\textbf{nell'ambito pervasivo non esiste proprio in questa forma}, come
mostrerà Beatrice parlando dei casi studio. Se si compra un gioco da tavolo o un
gioco di ruolo in un negozio, il marchio del publisher è di solito sul retro
della copertina o sul fronte della scatola: è l'entità che, di solito,
\textbf{produce fisicamente} il gioco.

\subsection{Sales e marketing}

L'ambito \textbf{sales e marketing}, a volte affidato a un'unica persona e a
volte separato, riguarda il lato commerciale. In alcune aziende è interno al
publisher, e si occupa di parlare con le fabbriche e le tipografie, di ``far
quadrare i conti''. In altre realtà è completamente diverso, specialmente nel
\textbf{B2B} o nei giochi \textbf{didattici} e \textbf{formativi}, di cui
Daniele e Beatrice ``sapranno parlare meglio''. Il tema viene ripreso più avanti
in dettaglio (sezione~\ref{a1:sec:marketing}).

\subsection{Graphic designer e illustratore}

Ruggero presenta insieme \textbf{graphic designer} e \textbf{illustratore}
proprio perché sono ruoli molto diversi che si tende a far coincidere:
\begin{itemize}
  \item chi fa \textbf{graphic design} si occupa delle \textbf{icone meccaniche},
        del \textbf{wireframing}, della \textbf{UX}: del lato \textbf{meccanico e
        funzionale} degli elementi visivi, cioè di tutto ciò che serve
        meccanicamente nel gioco, o anche alla sua ``qualità della vita'';
  \item chi fa \textbf{illustrazione} si occupa dell'\emph{art}, delle
        illustrazioni: copre il \textbf{look and feel}, la parte di
        \textbf{world building} e il ``vibe'' che si dà al gioco.
\end{itemize}

\paragraph{Daniele: che cos'è la user experience.} Daniele interviene sulla
\textbf{user experience}, ``importantissima per me'': ha un background da
\textbf{layout artist} (impaginava magazine), un mestiere in cui la UX conta
moltissimo. Per user experience non si intende solo l'\textbf{interfaccia
grafica}, la UI, che anche nei videogiochi è disegnata da un artista. La UX è
qualcosa che può fare anche il game designer: \textbf{disegnare l'esperienza del
giocatore}, cioè ``il gesto, l'intenzione, l'intuito'' che stanno dietro
l'aspetto visivo dell'interfaccia. Per fare UX non serve saper disegnare.
L'esempio è quello della \textbf{porta}: se si progetta una porta, bisogna
indicare all'utente come si apre, come si prende la maniglia; ``quella è user
experience''. Il design della maniglia in sé, invece, è
\textbf{interfaccia}.

\paragraph{Un aneddoto sulle icone.} Ruggero aggiunge che molte persone che
fanno \textbf{graphic design ben fatto non sanno illustrare}, e spesso vale il
viceversa. Lo si nota quando le icone di un gioco sono bellissime ma non
funzionali: andando a guardare i crediti, di solito si scopre che le icone le ha
fatte la persona che ha fatto le illustrazioni. ``Comunque non andiamo adesso in
questo, perché è un altro vaso di Pandora''.

\subsection{Sensitivity reader}

Il \textbf{sensitivity reader}, o \textbf{cultural consultant} a seconda del
progetto, fa un lavoro di \emph{sensitivity reading} del gioco. È una figura
molto usata nell'ambito \textbf{formativo e didattico}. L'esempio di Ruggero: se
si fa un gioco per un ospedale su una malattia specifica, chi lo progetta non è
un medico, ma il gioco deve essere \textbf{accurato} e \textbf{rappresentativo}
delle dinamiche che quella malattia comporta. Serve una persona che conosca quel
tema, anche se non sa che cosa sia il game design. Sono figure chiamate di
solito \textbf{a progetto}, per giochi specifici, che compensano competenze che
il team non ha ``per forza di cose, perché noi non sappiamo tutto''; e a volte,
per lavoro, si devono fare giochi su temi che non si conoscono personalmente.

\paragraph{Daniele: non è un censore.} Daniele segnala un problema italiano:
alcune persone, anche online, hanno confuso il sensitivity reader con un
\textbf{censore}, ``una persona che fa censura, soprattutto in materia di
\emph{woke}''. ``Assolutamente non è così''. La parola \emph{sensitive} non
significa che il pubblico sia sensibile a un tema, ma che \textbf{l'argomento di
cui tratta il gioco è sensibile}. In un gioco sulla medicina, realistico o no,
non si possono inserire imprecisioni. Per questo esistono i sensitivity reader
``più disparati'': da quelli per la cultura indiana o afroamericana fino a quello
su ``come si allevano gli gnu in Africa''. È una persona competente su un
argomento che aiuta a essere \textbf{veritieri}. Certo, uno degli ambiti
principali in cui si usano è quello culturale, ma ``non fate l'errore di pensare
che siano dei censori''. Daniele ha lavorato con ``centinaia'' di sensitivity
reader e nessuno gli ha mai detto ``non puoi scrivere questa cosa''. Di solito
sono come degli \textbf{editor}: danno opzioni, mostrano i punti problematici,
anche quelli di cui non ci si era nemmeno accorti.

Ruggero conferma con un esempio di quando lavorava in \textbf{Danimarca} con
l'\textbf{OMS}: giochi sul \textbf{diabete di tipo 2} nei bambini, per un target
di ragazzi delle scuole medie. C'era un team dell'OMS specializzato nella
comunicazione sul diabete, per garantire che il gioco fosse accurato e ne desse
una rappresentazione adatta al target. Il sensitivity reader aiuta a rendere il
gioco più accurato e ``più tematicamente coeso''.

\subsection{Producer}

Il \textbf{producer} è una figura a volte interna al publisher. In un progetto
grande è la figura di riferimento, ``quasi'' un project manager. In un progetto
piccolo si occupa spesso anche di parlare con le \textbf{tipografie} e di fare la
\textbf{lista dei materiali}, per assicurarsi che i materiali non siano solo
funzionali al gioco ma anche \textbf{economicamente sostenibili} da produrre:
``magari è un gioco bellissimo, ma poi costa più di quanto potete permettervi, e
nessuno vi comprerà mai un card game a 200 euro''.

\subsection{Perché conoscere queste figure}

Ruggero chiude la slide ``tornando a ombrello'': perché parlarne, se gli studenti
sono lì per fare game design? Non solo per saper dare un \textbf{nome} alle
figure. Non c'è niente di male a creare un gioco solo per sé; ma chi vuole fare
giochi a livello lavorativo, perché li giochino altre persone, ``magari più
persone possibile'', \textbf{si interfaccerà con queste figure}, e sapere che
cosa fanno e come parlarci è parte importante del percorso. Nel corso se ne
toccheranno diverse: il development con Marta, la produzione con Flavio, e
lezioni sulla UX, sul paragame, sul graphic design. ``Abituatevi mentalmente:
non siete in solitaria contro il mondo''.

% ---------------------------------------------------------------------
\section{Marketing e sales}
\label{a1:sec:marketing}

Arrivato all'angolo della slide con marketing e sales, Ruggero passa la parola
ai colleghi: ``io è meglio che non mi fate parlare di marketing e sales'', anche
per esperienze personali.

\subsection{Beatrice: i giochi trattati come prodotti}

Beatrice premette che le due voci sono in parte autoesplicative. Nel mondo
capitalista in cui viviamo, i giochi sono ``strumenti molto potenti per fare
\emph{n} cose'', ma sono trattati come \textbf{prodotti}, e il corso ragionerà
anche in quest'ottica. Nel \textbf{marketing} delle industrie ludiche e
videoludiche ci sono varie professionalità che gli studenti probabilmente già
conoscono: dal \emph{copy} alla grafica, tutto ciò che riguarda il
\textbf{vendere} il gioco o l'attività di gioco.

Il lato \textbf{sales} è forse la distinzione meno nota. Come in tante aziende, è
la parte che si occupa delle vendite \textbf{B2B}, cioè tra aziende: fare
accordi per vendere prodotti ludici a diverse realtà, per esempio ai
\textbf{distributori}, o a una casa editrice che vende alle librerie, che poi
vendono al consumatore finale. Chi si occupa di sales nel mondo ludico capisce
come vendere \textbf{alle aziende}. I due mondi si parlano spesso, e soprattutto
nelle realtà piccole si trovano figure che fanno ``un po' sales, un po'
marketing''.

\subsection{Daniele: che cos'è (e che cosa non è) il marketing}

Daniele il marketing si è trovato ``purtroppo'' a farlo, e lo fa anche per
l'accademia. ``Purtroppo'' perché in Italia è mal compreso: molti lo
identificano con l'\emph{advertisement} online. Non è così: il marketing è una
\textbf{disciplina molto più antica}, e in Italia ``abbiamo molto dimenticato''
come si fa il marketing tradizionale. Alcune cose da ricordare:
\begin{enumerate}
  \item \textbf{Il marketing non dice che cosa pubblicare, ma come pubblicarlo}.
        Esistono, nei videogiochi e nei giochi da tavolo, realtà che fanno
        prodotti \textbf{marketing driven}: l'esempio sono gli
        \textbf{hyper-casual game} per smartphone, quelli con il personaggio che
        va a destra e a sinistra e l'esercito che si allarga. Possono anche
        guadagnare, ma ``sono un po' la morte della creatività''. Il marketing
        ideale \textbf{prende un prodotto pensato da altri e gli dà le maggiori
        chance di essere venduto}: cerca il pubblico ideale per il prodotto che
        si ha, senza doverlo cambiare. Anzi, il mondo videoludico, e ludico in
        generale, ``è lastricato'' di errori di dipartimenti marketing che hanno
        cambiato completamente non il gioco, ma il modo in cui l'hanno
        pubblicizzato, e il gioco è fallito.
  \item \textbf{Usare strumenti del marketing non significa fare marketing}. Game
        design e game development a volte usano strumenti che vengono dal
        marketing (gli studenti ne vedranno con Beatrice, con Daniele e con
        Marta).
  \item \textbf{Il marketing è un gerundio}: si fa. È un'azione che si compie per
        comunicare; la comunicazione è attiva, e il marketing è l'insieme delle
        attività con cui si comunica il proprio gioco, dal \emph{copy} ai social.
        Può includere il digital marketing, ``ma non è solo quello''.
  \item \textbf{Il marketing esiste anche se si crede che non esista}. In Italia
        molti non vogliono ``fare marketing'' perché lo associano al lucro. Ma
        il marketing è \textbf{lo studio di come gli esseri umani comprano e
        vendono}, ed è innato in noi: ci sono logiche di cui siamo vittime, e
        ``anche se sei il bastian contrario, l'anticapitalista'', magari si
        compra comunque su Amazon perché fa comodo. Sono meccanismi che
        esistono, ``affascinanti anche, a volte''.
\end{enumerate}
Gli studenti non dovranno fare marketing, ma dovranno ricordare alcuni di questi
concetti, soprattutto nelle lezioni con Marta, che tratta argomenti legati al
marketing. E, chiude Daniele, come game designer e game developer ``per lo più
speriamo che il marketing non ci rovini i giochi'', soprattutto nel videoludico.

% ---------------------------------------------------------------------
\section{Game design e game development}
\label{a1:sec:design-development}

Ruggero torna sulla distinzione più importante della slide, ``proseguendo con
questa strada'':
\begin{description}
  \item[Game design] è chi \textbf{progetta l'esperienza} tramite la
        \textbf{creazione delle meccaniche} di gioco.
  \item[Game development] è chi fa \textbf{mediazione} tra la persona che scrive
        il game design e la \textbf{casa editrice}.
\end{description}

\subsection{I red flag dei designer}

Ruggero lo ripete perché lo vede ogni giorno nel suo lavoro: tra giochi
pubblicati e prototipi, ne valuta ``quasi uno al giorno, se faccio un conto
annuale'', e solo una minima percentuale arriva alla pubblicazione. Molte volte
i \textbf{red flag} vengono proprio da designer che non hanno presente questa
definizione, e che cominciano a mettere \textbf{clausole sulla non modifica del
gioco}: il gioco deve restare com'è. ``Questa cosa è impossibile con un
publisher, spesso''. Il publisher, per forza di cose, potrà chiedere di cambiare
la grafica, di reimpaginare il regolamento, di cambiare i materiali perché
troppo costosi: chiederà dei \textbf{compromessi}. E la persona che, all'interno
del publisher, fa questo lavoro è il \textbf{developer}.

Sta al designer capire \textbf{su che cosa del proprio gioco si può trovare un
compromesso e su che cosa no}. Publisher diversi lavorano in modo diverso, e il
designer ha sempre l'autorità di dire ``no, non voglio pubblicare questo gioco
se queste sono le modifiche''. Ma deve sapere chi fa il lavoro di development,
perché se il suo obiettivo è la \textbf{pubblicazione tradizionale} si
interfaccerà con questa figura.

\subsection{Daniele: una distinzione che manca al videogioco}

Daniele sottolinea che questa distinzione è \textbf{tipica del mondo analogico}
ed è ``molto interessante''. In Italia la parola \emph{developer} è diventata
negli anni quasi un sinonimo di programmatore di videogiochi, o di chiunque
lavori nel videoludico. Secondo lui c'entra anche il fatto che qualche anno fa
il sito \textbf{Gamasutra} ha cambiato nome in \textbf{Game Developer}, e la
parola è diventata un ombrello per tutte le professioni del videogioco. ``Mi fa
un po' ribrezzo'': il game developer dovrebbe essere un ruolo di design, un
designer che, come nell'analogico, fa \textbf{mediazione}. Nel videogioco i
publisher hanno in effetti i loro developer interni, che adattano il gioco, per
esempio alle console e ai loro kit di sviluppo. Ma la distinzione dovrebbe
esistere anche nel videogioco, perché è ``difficilissimo'' avere la mentalità
per fare le due cose insieme: chi fa design ha un focus che raramente si può
sovrapporre a quello di chi fa mediazione.

A Daniele piace la parola \emph{develop}, \textbf{sviluppare}, nel senso di
prendere qualcosa che è stato progettato e \textbf{renderlo possibile nel mondo
reale}. Ci sono progetti che hanno bisogno di uno sviluppo che ``polishi''
l'esperienza: la semplifichi, la migliori, la renda comprensibile e la concluda,
ma soprattutto la renda \textbf{stampabile} e \textbf{fruibile} per l'utente. La
metafora crossmediale è quella della \textbf{fotografia}: chi sviluppa le
fotografie non è chi scatta la foto, ma chi la rende visibile, anche
fisicamente. ``È un po' retro come metafora, ma ci capiamo''. ``È vero, però è
giusto'', conferma Ruggero.

% ---------------------------------------------------------------------
\section{Il mercato: Italia ed estero}
\label{a1:sec:mercato}

Il corso è in italiano e si svolge in Italia, ma bisogna conoscere anche il
\textbf{mercato estero}: molti giochi progettati da autori italiani vengono
pubblicati all'estero, e conoscere il mercato estero aiuta a capire come il
proprio gioco può arrivarci, ``o viceversa, come non può arrivarci''. La
tabella~\ref{a1:tab:mercato} riassume il confronto della slide.

\begin{table}[htbp]
\centering
\small
\begin{tabular}{@{}p{6.4cm}p{6.4cm}@{}}
\toprule
\textbf{Italia} & \textbf{Estero (soprattutto Europa)} \\
\midrule
il gioco è ancora visto come un'attività \textbf{infantile}, anche se la
normalizzazione è in crescita &
forte tradizione di gioco anche in un contesto \textbf{familiare adulto},
soprattutto ``dalla Germania in su'' \\
risposta \textbf{conservativa} sul tema ludico &
\textbf{premi di prestigio} riconosciuti in tutto il mondo (Spiel des Jahres,
As d'Or) \\
forte focus sulla \textbf{localizzazione} di giochi stranieri &
si cercano soprattutto giochi \textbf{originali} \\
investimento \textbf{molto basso} sulla creazione di giochi originali interni;
pochissimi designer assunti &
\textbf{developer interni} pagati con un contratto (Francia, Germania,
Danimarca\ldots) \\
\bottomrule
\end{tabular}
\caption{Il mercato del gioco analogico in Italia e all'estero, secondo la
slide di Ruggero e i commenti di Daniele.}
\label{a1:tab:mercato}
\end{table}

\subsection{L'Italia}

\begin{itemize}
  \item In Italia il gioco è ancora visto come un'\textbf{attività infantile},
        anche se la \textbf{normalizzazione} è in crescita. C'è ancora una forte
        risposta \textbf{conservativa} sul tema ludico.
  \item C'è un forte focus sulla \textbf{localizzazione} dei giochi stranieri: in
        una qualunque ludoteca o negozio, molti dei giochi da tavolo e di ruolo
        sono di publisher o designer esteri, pubblicati in italiano da publisher
        italiani. (La localizzazione, cioè portare un titolo dall'estero in
        Italia, esiste come funzione all'interno del publisher; Ruggero non l'ha
        messa nella slide delle figure perché non è il focus del corso.)
  \item C'è, ``purtroppo, preparatevi psicologicamente'', un \textbf{investimento
        molto basso} sulla creazione di giochi originali interni alla casa
        editrice. Pochissime case editrici italiane hanno designer interni pagati
        a tempo pieno con contratto a tempo indeterminato. Daniele lo dice senza
        mezzi termini: ``il mercato è molto piccolo, conosciamo praticamente
        tutto il settore'', e le persone che in Italia fanno game design
        analogico a tempo pieno con un contratto a tempo indeterminato ``sono
        meno delle dita delle mie due mani''. Di solito la formula è quella del
        \textbf{freelance}, del contratto, delle \textbf{royalty}. Chi va da un
        publisher aspettandosi di essere assunto a tempo pieno come game
        designer deve sapere che ``non siamo purtroppo ancora in quel tipo di
        situazione''.
  \item C'è però un contraltare: in Italia, soprattutto nel gioco di ruolo ma
        anche nel gioco da tavolo, c'è stata molta \textbf{imprenditoria}, anche
        grazie a \textbf{Kickstarter}, che ha aperto le porte ad altri sistemi
        di \textbf{crowdfunding}. È un altro argomento, perché richiede un team,
        un investimento e così via.
  \item E gli autori non mancano, anzi: le occasioni di incontro tra autori
        vedono ogni anno un aumento dei possibili autori\footnote{La trascrizione
        dice ``i DAG a Parma''; il riferimento è probabilmente a IDEAG, il
        circuito di incontri tra autori di giochi italiani, ma non è
        ricostruibile con certezza.}. Bisogna ``sgomitare'' tra tanti autori,
        che magari arrivano da esperienze molto diverse.
\end{itemize}
Ruggero precisa: l'investimento basso è ``dall'alto''. Le persone che fanno e
vogliono fare giochi sono tantissime e in aumento; scarseggia l'investimento,
sia dei publisher disposti a pagare stipendi sia dei \textbf{bandi} pubblici. I
bandi esistono (Beatrice ne parlerà per i giochi pervasivi, che spesso si
muovono in quella direzione, e Ruggero cita qualcuno che sta facendo un card
game con un bando italiano, ``una cosa assurda, ma lo sta facendo''), ma sono
una rarità, soprattutto rispetto all'estero.

\subsection{L'estero}

\begin{itemize}
  \item L'estero (dipende dalla nazione, ma in generale il resto d'Europa) ha una
        forte \textbf{tradizione di gioco} anche in un contesto \textbf{familiare
        adulto}, soprattutto ``dalla Germania in su''.
  \item Ha \textbf{premi di prestigio} riconosciuti in tutto il mondo, come lo
        \textbf{Spiel des Jahres} tedesco e l'\textbf{As d'Or} francese\footnote{La
        trascrizione dice ``lo Spield AR, la SDOR''.}.
  \item Cerca soprattutto \textbf{giochi originali}. Lo si vede dal fatto che
        molti autori italiani di giochi da tavolo pubblicano in Germania o in
        Francia e \textbf{non vengono pubblicati in Italia}, ``che è tutto
        dire''.
  \item Ha \textbf{figure di development interne}: non è raro vedere publisher
        francesi, tedeschi, danesi con developer pagati internamente, a
        contratto.
\end{itemize}
``Poi la situazione è ancora più drammatica'', ammette Ruggero, ``ma per ora
vogliamo tenere il morale un po' alto. È l'introduzione. Ci sarà tempo di andare
nella tristezza''.

Prima di proseguire, Ruggero insiste sulle domande: con lui si può interrompere
anche mentre parla, alzando la mano con le reaction o semplicemente togliendo il
microfono e chiedendo. ``Non mi offenderò mai, anzi, sarò felice di vedervelo
fare e mi ricorderò positivamente di voi''.

% ---------------------------------------------------------------------
\section{Perché serve una terminologia comune}
\label{a1:sec:terminologia}

La serie di slide successiva è ``una piccolissima parte di fondamenti di
design'' (``non spaventatevi, vi vedo che vi spaventate un po'\,''). Anche qui
molti temi verranno ripresi in dettaglio, ma l'idea è arrivare alla prossima
lezione con almeno un \textbf{glossario}, un \textbf{linguaggio condiviso}: se
qualcuno dice \emph{game developer}, tutti sanno di che cosa sta parlando.

Ruggero parte da una domanda, e precisa che non è retorica: \textbf{perché ci
serve una terminologia comune?} Gli studenti rispondono:
\begin{itemize}
  \item Gianmarco, ``molto pragmatico'': perché bisogna lavorare in team e
        capirsi. Vero, nel corso e anche fuori;
  \item un'altra studentessa: per la \textbf{comunicazione}, tra di noi e con
        tutti quelli che sono interessati nel settore. Vero anche questo:
        conoscere i termini comuni permette di parlare con chi fa game design o
        development nel dettaglio, nel tecnico, nello specifico.
\end{itemize}
Ruggero aggiunge altre due ragioni.

\paragraph{Saper comunicare il proprio gioco.} La prima riguarda la
\textbf{comunicazione del gioco} verso l'esterno. ``E lo dico non tirando acqua
al mio mulino, ma anzi togliendola'': \textbf{non è spesso il buon game design
a far diffondere un gioco}. ``Vorrei che lo fosse, sarebbe bellissimo, ma
purtroppo non è così''. Ci sono giochi bellissimi che non sanno comunicare, che
non sanno dire che cosa fanno, e che quindi è difficile, da fuori, trovare
invitanti; e ci sono giochi mediocri dal punto di vista del design che si
comunicano molto bene, e che quindi conoscono in molti. Saper parlare di quello
che si fa è utile \textbf{a sé stessi}, in primis.

\paragraph{Documentare per sé stessi.} La seconda è la \textbf{documentazione}:
documenti di design, di brainstorming e così via. ``Sembrerà una battuta, ma
provate a recuperare un documento di brainstorming una settimana dopo, se non
avete scritto bene le cose che avete scritto''. Chi l'ha fatto sa già di che
cosa si parla; ``chi non l'ha ancora vissuto, forse lo vivrà''.

% ---------------------------------------------------------------------
\section{Game e play}
\label{a1:sec:game-play}

La \textbf{prima, importantissima divisione terminologica} è quella tra
\emph{game} e \emph{play}. In italiano c'è un problema: \textbf{gioco} e
\textbf{giocare} hanno la stessa radice, mentre in altre lingue, tra cui
l'inglese, sono due parole strutturalmente diverse:
\begin{description}
  \item[Game] il \textbf{gioco come artefatto}: il gioco come ecosistema di
        meccaniche, come struttura di regole, come prodotto (``tra
        virgolette''), come oggetto fisico.
  \item[Play] l'\textbf{atto}, il verbo: le persone che giocano, l'approccio
        fisico all'attività del gioco.
\end{description}
È importante perché gli studenti faranno \textbf{game design}; esistono anche
ambiti di \emph{play design}, ma qui ``creerete giochi, non giocattoli''.

\subsection{Giochi e giocattoli}

Qual è la differenza? Uno studente risponde: \textbf{la presenza di regole}.
Esatto: un gioco ha delle regole, un'attività di \emph{play} non
necessariamente. Con un giocattolo si può fare \emph{play}: c'è chi gioca con i
Lego, chi gioca a guardia e ladri o a nascondino al parco, e magari alcune di
quelle regole sono ``più o meno semi-solidificate''. Ma il \emph{play} permette
spesso un approccio non strettamente legato alle regole.

Viceversa, \textbf{un game può esistere anche se nessuno lo sta giocando}. Se si
prepara su un tavolo una scacchiera con i pezzi degli scacchi, chi passa e la
vede riconosce il gioco degli scacchi, anche se nessuno sta giocando. Se invece
si vedono dei bambini correre in un parco giochi, quei bambini \emph{stanno
giocando}, ma non stanno giocando \emph{a un gioco} specifico. ``È importante
separare questa cosa''; ci si tornerà, ``ma per ora abituatevi a questa
differenza''.

\paragraph{Approfondimento: giocattoli, giochi e regole.} La stessa distinzione
è il punto di partenza del capitolo che Jesse Schell, in \emph{The Art of Game
Design}, dedica alla definizione di gioco (\emph{The Experience Rises Out of a
Game}). Schell nota che usiamo perfino un linguaggio diverso: ``un gioco è
qualcosa a cui si gioca'' (\emph{a game is something you play}), ``un giocattolo
è qualcosa con cui si gioca'' (\emph{a toy is something you play with}); e
arriva a definire il giocattolo come ``un oggetto con cui si gioca''. Analizzando
poi la definizione di gioco di Avedon e Sutton-Smith, ne ricava che ``i giochi
hanno regole, i giocattoli no'', e che le regole sono uno degli aspetti
\emph{definitori} del gioco, insieme alla volontarietà, agli obiettivi, al
conflitto e alla possibilità di vincere o perdere. Schell osserva anche che
molti giochi sono costruiti \textbf{sopra} un giocattolo (``una palla è un
giocattolo, ma il baseball è un gioco''), e tra i suoi consigli di
prototipazione c'è proprio \emph{Build the Toy First}: costruire prima il
giocattolo, cioè qualcosa che sia già piacevole da manipolare, e poi il gioco
intorno. È una lettura utile della distinzione di Ruggero: il \emph{play} può
fare a meno delle regole, il \emph{game} no.

% ---------------------------------------------------------------------
\section{Il framework MDA (e il DPE)}
\label{a1:sec:mda}

\subsection{Meccaniche, dinamiche, estetiche}

Un'altra terminologia fondamentale è quella del framework \textbf{MDA}
(\emph{Mechanics, Dynamics, Aesthetics}; in italiano a volte ``MDE''). Viene da
un \textbf{paper} nato nell'ambito del videogioco\footnote{È il paper di Robin
Hunicke, Marc LeBlanc e Robert Zubek, \emph{MDA: A Formal Approach to Game
Design and Game Research}, del 2004. Daniele aggiunge il paper alla cartella
condivisa durante la lezione.}, che ha ``solidificato'' queste terminologie, e
divide un gioco in tre aree:
\begin{description}
  \item[Meccaniche] le \textbf{regole per come sono scritte}. ``Ogni volta che
        pescate un 2 fate dei punti'': quella è una meccanica.
  \item[Dinamiche] l'\textbf{applicazione in tempo reale} delle meccaniche. La
        dinamica è il giocatore che, pescando una carta, \textbf{spera} che esca
        un 2. Quel processo non è una regola: ``il gioco non vi dice che dovete
        sperare che esca un 2''.
  \item[Estetiche] l'insieme di quello che il gioco vuole comunicare,
        l'\textbf{esperienziale}.
\end{description}

L'esempio di Ruggero mette insieme le tre aree. In un gioco in cui pescare un 2
fa fare punti (meccanica) e in cui quindi si spera che esca un 2 (dinamica),
\textbf{che estetica si ottiene se nel mazzo di 100 carte c'è un solo 2? E se ce
ne sono 50?} L'emozione di aver pescato \emph{l'unico} 2 del mazzo è
completamente diversa da quella di pescarne uno sapendo che ce ne sono
cinquanta e che è molto probabile trovarlo. Le regole sono quasi identiche; cambia
l'esperienza.

\subsection{Il falso amico: ``estetica''}

Daniele avverte di non sottovalutare la parola \textbf{estetica}, che in italiano
è un \textbf{falso amico}: sembra indicare la parte grafica, e invece non lo è.
Si incontreranno designer che hanno studiato MDA e altri che non l'hanno
studiato, e lo si capirà proprio da questo: i secondi diranno che le
\emph{aesthetics} sono la grafica. Le \emph{aesthetics} sono la parte
\textbf{esperienziale}: come nell'esempio di Ruggero, quello che lascia al
giocatore aver pescato l'unico 2 del mazzo.

\subsection{Il DPE: design, play, experience}

Per gli italiani, aggiunge Daniele, è più facile capire l'MDA attraverso un
altro paper, quello del \textbf{DPE}\footnote{È il framework \emph{Design, Play,
Experience} di Brian Winn, proposto come estensione dell'MDA per i giochi
educativi (\emph{serious games}). Anche questo paper viene aggiunto da Daniele
alla cartella condivisa.}, che di fatto fa un parallelismo assegnando
\textbf{design, play, experience}:
\begin{itemize}
  \item le \textbf{meccaniche} sono il lato \textbf{design};
  \item le \textbf{dinamiche} sono il lato \textbf{play};
  \item le \textbf{estetiche} sono l'\textbf{experience}, l'esperienza che riceve
        chi gioca.
\end{itemize}
``Se vivete le \emph{aesthetics} come l'esperienza che vive chi sta giocando,
state a postissimo''. In questo momento Beatrice lascia la lezione (``ciao,
Bea''). I paper di MDA e DPE sono in inglese; chi vuole può dargli una lettura.

\begin{figure}[htbp]
\centering
\begin{tikzpicture}[
    font=\small,
    blocco/.style={draw, thick, rounded corners=3pt, fill=gray!12,
                   minimum width=3.2cm, minimum height=1.0cm, align=center},
    dpe/.style={font=\footnotesize\itshape, align=center},
    esempio/.style={draw, rounded corners=2pt, text width=3.3cm,
                    align=left, font=\footnotesize, inner sep=4pt},
    freccia/.style={-{Stealth[length=2.5mm]}, thick}
  ]
  \node[blocco] (m) at (0,0)   {\textbf{Meccaniche}\\ \footnotesize le regole scritte};
  \node[blocco] (d) at (4.3,0) {\textbf{Dinamiche}\\ \footnotesize le regole in azione};
  \node[blocco] (a) at (8.6,0) {\textbf{Estetiche}\\ \footnotesize l'esperienza};
  \draw[freccia] (m) -- (d);
  \draw[freccia] (d) -- (a);
  \node[dpe, above=0.15cm of m] {DPE: \textbf{design}};
  \node[dpe, above=0.15cm of d] {DPE: \textbf{play}};
  \node[dpe, above=0.15cm of a] {DPE: \textbf{experience}};
  \node[esempio, below=0.45cm of m] (em)
    {``Ogni volta che peschi un 2 fai punti''};
  \node[esempio, below=0.45cm of d] (ed)
    {pescando, il giocatore \emph{spera} che esca un 2 (nessuna regola glielo
     chiede)};
  \node[esempio, below=0.45cm of a] (ea)
    {un solo 2 su 100 carte: pescarlo è un evento;\\
     50 su 100: è quasi scontato};
  \draw[dashed] (m) -- (em);
  \draw[dashed] (d) -- (ed);
  \draw[dashed] (a) -- (ea);
  % prospettive
  \draw[freccia, gray] ([yshift=-0.3cm]em.south west) --
    node[below, font=\footnotesize\itshape, text=black]
    {il designer agisce sulle meccaniche}
    ([yshift=-0.3cm]ea.south east);
  \draw[freccia, gray] ([yshift=-1.2cm]ea.south east) --
    node[below, font=\footnotesize\itshape, text=black]
    {chi gioca incontra prima l'esperienza}
    ([yshift=-1.2cm]em.south west);
\end{tikzpicture}
\caption{Il framework MDA con l'esempio del ``2'' usato da Ruggero, e la
corrispondenza con il DPE suggerita da Daniele. Le due frecce in basso
riassumono la lettura classica del paper: il designer costruisce le meccaniche
per ottenere un'esperienza, mentre il giocatore parte dall'esperienza.}
\label{a1:fig:mda}
\end{figure}

\paragraph{Approfondimento: le due direzioni dell'MDA e le otto estetiche.} Le
frecce della figura~\ref{a1:fig:mda} riprendono un'idea centrale del paper
originale, che la lezione non esplicita: designer e giocatore guardano il gioco
da lati opposti. In \emph{Challenges for Game Designers}, Brenda Brathwaite e
Ian Schreiber la riassumono così: i designer si chiedono quale estetica
vogliono ottenere, definiscono le dinamiche che porterebbero a quella
sensazione, e poi creano le meccaniche che producono quelle dinamiche. Lo stesso
libro conferma l'avvertenza di Daniele sul falso amico: ``le \emph{aesthetics}
non si riferiscono all'aspetto del gioco, ma alla risposta emotiva'' che il
designer vuole suscitare. E offre esempi di dinamiche nel senso di Ruggero: in
uno sparatutto i punti di \emph{spawn} sono una meccanica, mentre appostarsi
accanto a uno di essi per colpire chi riappare (lo \emph{spawn camping}) è una
dinamica; negli scacchi i movimenti dei pezzi sono meccaniche, mentre le
aperture codificate sono dinamiche. Per dare un nome alle estetiche, il paper
MDA propone una tassonomia, poi detta degli ``otto tipi di \emph{fun}'', che
Schell riprende in \emph{The Art of Game Design} come \emph{LeBlanc's Taxonomy of
Game Pleasures}: \textbf{sensazione}, \textbf{fantasia}, \textbf{narrativa},
\textbf{sfida}, \textbf{compagnia} (\emph{fellowship}), \textbf{scoperta},
\textbf{espressione} e \textbf{sottomissione} (il piacere di entrare nel
``cerchio magico'' e lasciarsi andare a un passatempo). È un primo esempio di
quello che i docenti chiedono con le parole vietate della prossima sezione:
invece di dire che un gioco è ``divertente'', dire \emph{quale} esperienza
offre.

% ---------------------------------------------------------------------
\section{Le parole vietate}
\label{a1:sec:parole-vietate}

``Attenzione, adesso qua triggeriamo delle persone''. Ruggero annuncia una regola
che ``tornerà'': nel corso ci sono delle \textbf{parole vietate}. ``Non perché
siamo dei censori''; il perché si vedrà in una lezione apposita, perché
spiegarlo adesso richiederebbe un'ora. Ma Daniele vuole spiegarne ``un
minimo''.

\subsection{La regola}

In \textbf{nessun materiale scritto} consegnato ai docenti, task compresi, sono
accettate queste parole, né in inglese né in italiano:
\begin{center}
\begin{tabular}{@{}ll@{}}
\toprule
\textbf{Inglese} & \textbf{Italiano} \\
\midrule
\emph{fun} & divertimento \\
\emph{immersion} & immersione \\
\emph{entertainment} & intrattenimento \\
\bottomrule
\end{tabular}
\end{center}
``Non li voglio neanche vedere''. Daniele lo ripete con un paragone: ``se mi fate
i disegni con Paint, benissimo, top, fantastico, siamo ancora best friend. Se mi
scrivete \emph{divertimento}, io sbrocco''. È una \textbf{regola dell'accademia}:
negli stessi corsi triennali di narrative design e game design le parole sono
vietate agli studenti per tre anni; ``per voi solo per tre mesi''. Daniele
aggiungerebbe volentieri un quarto termine, \textbf{gamification}.

\subsection{Perché}

Il motivo approfondito lo darà Ruggero nella lezione apposita. Daniele dà il
suo. Ruggero e Diego Ricchiuti sono molto contrari a \emph{fun} e a
\emph{immersion}; lui ha aggiunto \emph{intrattenimento}. Sono \textbf{parole
troppo passe-partout e troppo ambigue}:
\begin{itemize}
  \item \textbf{intrattenimento} in Italia è stata usata per anni come un modo
        per \textbf{sminuire}: ``ma sì, tanto è un gioco solo di
        intrattenimento, è un romanzo solo di intrattenimento''. Ma anche se un
        prodotto è ``solo di intrattenimento'', non vuol dire che non ci si debba
        lavorare con attenzione. E poi intrattenimento ``non significa niente'':
        ci sono tanti modi in cui una cosa può intrattenere;
  \item lo stesso vale per \textbf{fun}, e per \textbf{immersion}, che ha un
        problema in più: è una parola passe-partout per cui si può dire che un
        gioco è ``immersivo'' nel senso generale del termine, ma poi non si sa
        davvero che cosa voglia dire, \emph{immersivo in che senso}.
\end{itemize}

L'obiettivo didattico è che gli studenti \textbf{vadano oltre queste parole}:
``un po' come quando alle medie vi dicono smettetela di dire che tutto è bello,
perché ci sono vari gradi di bello, ci sono altri modi di dire bello''. Bisogna
\textbf{decostruirle}: quando si sta per usarle, chiedersi \textbf{che cosa si
intende} con quella parola e cercarne una \textbf{più precisa}.

Ruggero chiede, in uno dei ``rari casi'', di \textbf{non fare domande adesso} su
questa regola: ogni anno arrivano domande e commenti, e ci saranno lezioni
apposite di tre ore, con strumenti per capire il perché. Ma da subito gli
studenti devono abituarsi a evitare questi termini, sia perché si spera che
abbiano termini più precisi, sia perché queste parole ``portano con sé dei
problemi agganciati'' che si vedranno. ``Non spaventatevi, non verrete sparati in
un vicolo cieco di notte''.

\paragraph{Una piccola gamification.} Il divieto va visto come uno
\textbf{stimolo}, ``un po' come se stessimo gamificando la vostra esperienza'',
una sorta di user experience delle lezioni. Imponendosi di non usare queste
parole, si noterà come migliora la capacità di \textbf{descrivere i giochi} e le
proprie esperienze ludiche, e di essere più precisi. È un invito a ``obbligarsi a
essere più precisi e ad avere una mentalità più critica''. La lezione di
Ruggero sull'\textbf{involvement}, promette Daniele, sarà ``fenomenale'' e
``vi aprirà la testa''.

\paragraph{Approfondimento: anche i libri faticano a definire il ``fun''.} La
difficoltà che il divieto vuole far emergere è visibile anche nei testi di game
design. Schell, in \emph{The Art of Game Design}, prova a definire il \emph{fun}
e propone ``piacere con sorprese'' (\emph{pleasure with surprises}), ammettendo
subito che è strano poter usare una parola per tutta la vita, sapendo con
certezza che cosa significa, senza riuscire a esprimerla chiaramente quando
qualcuno lo chiede; e invita a mettere alla prova la definizione con
controesempi. La sua lente sul \emph{fun} (\emph{The Lens of Fun}) si limita,
non a caso, a chiedere quali parti del gioco siano \emph{fun} e \emph{perché}.
Anche \emph{immersion} ha molte letture: in \emph{21st Century Game Design},
Bateman e Boon la descrivono, tra l'altro, come la capacità di ``perdersi''
nel mondo artificiale presentato dal gioco, legata alla sospensione
dell'incredulità e ai giochi di mimesi. Sono significati diversi sotto la stessa
parola: proprio il motivo per cui, in un documento di design, conviene dire
quale si intende.

% ---------------------------------------------------------------------
\section{Chi crea, chi gioca, chi facilita}
\label{a1:sec:ruoli-esperienza}

Un'altra distinzione importante è quella tra \textbf{chi crea}, \textbf{chi
gioca} e \textbf{chi facilita}. Spesso, quando si fa un gioco, lo si fa provare
di persona a chi gioca. Ma ``la norma dei giochi pubblicati non funziona così'':
quando si pubblica un gioco non si vede chi lo compra e chi lo gioca, e non si
sa se qualcuno, a casa propria con i propri amici, segue le regole alla lettera e
interpreta correttamente quello che il designer aveva descritto. Quindi già in
fase di progettazione bisogna tenere a mente \textbf{chi sta facendo game design
e chi sta giocando}.

E nei giochi a evento, nei giochi pervasivi, nei giochi di ruolo da tavolo e dal
vivo (LARP) c'è un terzo pezzo del puzzle, ``importantissimo'': \textbf{chi
facilita}. La domanda da farsi è: \textbf{il regolamento è scritto per chi gioca
o per chi deve facilitare il gioco?} Bisogna abituarsi a pensare a tutte le
persone coinvolte nell'esperienza.

% ---------------------------------------------------------------------
\section{Playtest e iterazione}
\label{a1:sec:playtest}

\subsection{Nessun gioco nasce perfetto}

L'ultima voce del glossario è ``un'altra cosa importantissima'':
\textbf{playtest} e \textbf{iterazione}. Il corso è \textbf{pragmatico}:
``specialmente da parte mia'', dice Ruggero, non si vedranno teorie ``messe lì
per fare accademia sul game design nelle nostre torri d'avorio''. Ci saranno
teorie e framework, ma tutti estremamente pragmatici, perché l'obiettivo non è
solo capire i giochi, ma \textbf{saperli fare e iterare}.

Playtest e iterazione sono \textbf{inevitabili} nel game design, perché
\textbf{nessun gioco nasce perfetto}, soprattutto perché le persone sono
diverse: un gioco perfetto per una persona può essere incomprensibile per
un'altra. Chi progetta giochi per altre persone deve capire, dopo averli fatti
provare, \textbf{che cosa cambiare, che cosa sistemare, che cosa aggiungere, che
cosa togliere}. Ci saranno lezioni su playtest, prototipazione, raccolta dei
feedback, metodi etnografici. Ma bisogna abituarsi già ora all'idea che
\textbf{non esiste la consegna definitiva}. A un certo punto i task vanno
consegnati per forza di cose, ma, come diceva Daniele all'inizio, alla fine del
corso gli studenti avranno un \textbf{prototipo}, ``che sarà un prototipo'', e
gli strumenti per portare avanti il lavoro anche dopo.

\subsection{L'aiuto dopo il corso}

Non solo gli strumenti: anche l'aiuto dei docenti. Daniele racconta di aver
portato a un editore, a Play\footnote{Probabilmente \emph{Play -- Festival del
Gioco} di Modena; la trascrizione non permette di esserne certi.}, uno dei
progetti realizzati l'anno precedente da due studenti, un gioco di ruolo in
solitario (``un \emph{solo game} RPG molto, molto bello''). Se gli studenti
vorranno trasformare il proprio lavoro in un prototipo vero e proprio, possono
contattarlo: lo fa ``per lavoro, ma anche per passione''. Gestisce anche una
community per aiutare i designer ``un po' depressi'' a rimettersi a progettare.

Quindi: ``non sentitevi persi se durante la fase di iterazione vi sentite
persi, è normale''.

\subsection{Il dizionario comune e il production dictionary}

Daniele riprende un punto detto ``benissimo'' da Ruggero e che secondo lui non è
stato sottolineato abbastanza: il \textbf{dizionario comune} serve anche a
\textbf{non perdersi} mentre si itera. Se da un'iterazione all'altra i
documenti usano termini completamente diversi, e non c'è un dizionario di base,
ci si perde.

Nel videogioco, una delle primissime cose che un designer fa quando entra in
un'azienda è il \textbf{production dictionary}: si annota, in un file privato,
tutte le parole che usano i colleghi, per allinearsi al dizionario comune
dell'azienda, perché lì su un prototipo possono lavorare anche cento persone.
Nel gioco da tavolo non si fa molto, perché il mondo dell'analogico è
``talmente più piccolo''; ma serve comunque imparare un \textbf{dizionario
minimo}, dai docenti e dai professionisti. E il corso ne darà una parte,
entrando nel dettaglio di alcuni termini e framework.

\paragraph{L'esempio del feedback.} Perché serva, lo mostra un esempio. Si
testa un gioco e ci si annota come feedback: ``questo gioco deve essere più
divertente''. Che cosa vuol dire? Come lo si migliora? Quella frase non dà
elementi utili. Se invece si scrive che deve essere \textbf{più adrenalinico},
\textbf{più frenetico}, \textbf{più rilassato}, \textbf{più cozy}, o
\textbf{più bilanciato matematicamente}, si sa dove andare a lavorare.

\paragraph{Approfondimento: la regola del ciclo.} Il principio ``nessun gioco
nasce perfetto'' ha un nome preciso in \emph{The Art of Game Design}: è la
\textbf{Rule of the Loop}, la regola del ciclo. ``Più volte provi e migliori il
tuo design, migliore sarà il tuo gioco''. Schell precisa che non è una lente,
cioè un punto di vista, ma ``una verità assoluta'', senza eccezioni: ci saranno
momenti in cui si cercherà di convincersi che ``questa volta il design è così
buono che non serve testarlo'', e ogni volta se ne pagherà il prezzo. Ed è
interessante la ragione per cui la regola pesa di più sul digitale: secondo
Schell il tempo e il denaro necessari per provare e correggere un sistema sono
molto maggiori nei videogiochi che nei giochi tradizionali. Per un gioco di
carte o da tavolo, con il tempo per provarlo e cambiarlo più e più volte, il
consiglio è semplicemente di farlo. Il gioco analogico è, da questo punto di
vista, la palestra ideale dell'iterazione.

% ---------------------------------------------------------------------
\section{Tabletop Simulator}
\label{a1:sec:tts}

\subsection{Il software di settore}

``Last but not least'': il corso usa un software, \textbf{Tabletop Simulator},
che si compra su \textbf{Steam} e costa tra i \textbf{10 e i 20 euro} (``aspettate
che vada in sconto, va in sconto ogni due per tre''). Daniele spiega perché:
\begin{itemize}
  \item è \textbf{il software di settore} del gioco da tavolo, quello che i
        publisher usano per i prodotti digitali, interni ed esterni;
  \item essendo su Steam, permette di fare \textbf{playtest} e di giocare con gli
        amici di Steam. Daniele stesso sta testando lì il proprio gioco;
  \item funziona anche su Mac: Windows o Mac, nessun problema.
\end{itemize}
Ci sarà una lezione apposita, di Ruggero, sulla creazione del proprio
\textbf{modulo} di gioco.

\paragraph{Che cosa non è.} Daniele chiarisce subito un ``fraintendimento
gigantesco'', soprattutto per chi ha già esperienza: Tabletop Simulator
\textbf{non richiede programmazione} e \textbf{non è un software per
trasformare i giochi da tavolo in videogiochi}. È un software che
\textbf{simula dei componenti fisici su un tavolo}: serve per giocare i board
game in digitale. Simula anche la fisica: c'è perfino il comando per
\textbf{ribaltare il tavolo} (``non usatelo ora'').

\paragraph{La partnership mancata.} Da anni Daniele prova a fare una
partnership con gli sviluppatori: due anni prima ha contattato l'azienda per
acquistare delle chiavi per gli studenti, dato che il software si usa nei corsi.
L'azienda ha risposto che non fa accordi commerciali. ``Credo sia la prima
azienda nella mia vita che dice che non vuole che io le dia dei soldi''. Gli
studenti dovranno quindi provare ad acquistarlo da sé; se i docenti riusciranno a
darglielo, bene.

\paragraph{La domanda di Gianmarco.} Gianmarco segue le lezioni da un computer
aziendale ``super chiuso'', su cui non può installare Steam: deve usare Tabletop
Simulator durante le lezioni? Ruggero risponde che per la lezione dedicata è
meglio poter seguire e rifare in diretta. Ma il software non sarà richiesto
necessariamente per il project work: un gruppo potrebbe fare un gioco pervasivo
o un gioco di ruolo, non un gioco da tavolo. Essendo un software di settore, è
comunque bene conoscerlo. Ruggero aggiunge che le versioni non ufficiali che
circolano non permettono di giocare in molti, e quindi non servono per
playtestare o giocare insieme.

\subsection{Una dimostrazione}

Daniele mostra il software aperto da Steam:
\begin{itemize}
  \item si può creare una partita \textbf{single player} o \textbf{multiplayer};
        per giocare tra compagni si crea un server multiplayer e ci si aggiunge
        come amici su Steam;
  \item si dà al server un nome, un numero di giocatori, eventualmente una
        password; creato il server, si apre la scelta dei giochi;
  \item ci sono giochi gratuiti, giochi da acquistare (DLC, non obbligatori) e
        soprattutto un \textbf{Workshop} ``ampissimo'', con decine di migliaia di
        giochi caricati dagli utenti: ``se c'è un board game che conoscete,
        quel board game è su Tabletop Simulator, giocabile dal Workshop''.
        L'esempio mostrato è l'open beta ufficiale del gioco da tavolo di
        \emph{Slay the Spire}\footnote{La trascrizione dice ``Slade Spire''.}:
        molti publisher, per farsi pubblicità, mettono il proprio gioco in demo o
        giocabile in digitale;
  \item c'è un \textbf{tutorial} iniziale che spiega come muoversi sul tavolo;
        con \emph{Esc} e \emph{Help} si trovano i comandi da tastiera.
\end{itemize}

\subsection{Il primo task}

Il primo (e per ora unico) task è questo: \textbf{entro il 10 ottobre}, giorno
della lezione di Ruggero su Tabletop Simulator, giocare \textbf{tre o quattro
volte in totale}, online, tra compagni di corso o con amici. Non si tratta di
imparare a creare moduli, ma di prendere la mano \textbf{lato giocatore}: aprire
il software, impararne i comandi, giocare. ``Per me potete anche andare a
prendere domino o mahjong'' (``mahjong è bellissimo tra l'altro'', aggiunge
Daniele, che sta facendo un videogioco sul mahjong), o il pachisi o il go. Gli
anni precedenti si chiedeva a ciascuno di proporre ai compagni un gioco da
tavolo che apprezzava. È anche un modo per conoscersi. Ruggero spiega il
perché: usarlo per giocare è facile e intuitivo, usarlo per creare moduli è più
complesso; ma chi arriva al 10 ottobre avendo già aperto il software e giocato,
troverà più intuitivo anche creare un modulo. E l'avvertimento: ``non riducetevi
all'ultimo, non fate una partita il 9 sera''.

\subsection{Tabletopia e altri vantaggi}

Tancredi racconta di aver provato a caricare un suo gioco su
\textbf{Tabletopia}. ``Questo è un vaso di Pandora'', risponde Daniele.
Tabletopia è molto usato, ma secondo lui il suo strumento di creazione dei moduli
è pessimo, ed è il motivo per cui la maggioranza dei moduli ufficiali dei
publisher è su Tabletop Simulator. Tabletopia ha un vantaggio vero, è
\textbf{gratuito}, ma ha scelte di ``qualità della vita'' poco felici.

Tabletop Simulator ha invece un vantaggio ``fenomenale'' che si tende a
sottovalutare: essendo su Steam, \textbf{il modulo si pubblica direttamente sul
Workshop di Steam}; si possono invitare gli amici a iscriversi e giocare con
loro, purché abbiano un account Steam. Rispetto ad altri software da imparare,
costa ``tra i 10 e i 20 euro, ragazzi: si parla proprio di spiccioli''. ``Io con
Tabletop Simulator ci ho playtestato la vita''.

Altre caratteristiche utili:
\begin{itemize}
  \item si possono creare \textbf{moduli non elencati} (\emph{unlisted}),
        perfetti da passare a un publisher prima che il gioco sia pubblicato:
        non sono pubblici, e nessuno può rubare i materiali;
  \item il software ha una grafica datata (``\emph{World Art 1998}'', scherza
        uno dei docenti), ma lo stanno portando su \textbf{Unity} con molti investimenti
        e aggiornamenti; ascolta la comunità e introduce funzioni utili anche per
        chi fa design;
  \item il suo server Discord è molto frequentato, e c'è un canale YouTube
        dedicato solo al design di giochi di carte\footnote{La trascrizione ne
        riporta il nome come ``Distraction Maker''; non è verificabile.} con un
        server Discord dove si fanno playtest: il 99\% è su Tabletop Simulator;
  \item molti editori, quando si propone loro un gioco da tavolo,
        \textbf{chiedono se c'è un modulo per Tabletop Simulator}.
\end{itemize}

% ---------------------------------------------------------------------
\section{Le domande a Daniele: il tema del project work}
\label{a1:sec:domande-daniele}

Daniele ha ``un minuto'' per rispondere alle domande sul corso prima di dover
andare. Uno studente chiede, ``giusto per curiosità'', se il project work avrà
un tema o sarà libero.

\begin{itemize}
  \item C'è un \textbf{mandate}, che i docenti propongono ogni anno e
        aggiornano leggermente; verrà presentato il \textbf{25 settembre}.
        ``Non facciamo spoiler''. Conterrà quello che si chiede di fare e quello
        che non si chiede: dei \textbf{must}, dei \textbf{must not} e dei
        \textbf{nice to have}.
  \item Verrà chiesto di lavorare su un'\textbf{idea originale} che non si è
        portata al corso: originale nel senso di \textbf{creata apposta}, non
        nel senso di ``mai vista prima''.
  \item L'idea va \textbf{partorita insieme}. Altri anni alcuni studenti hanno
        provato a imporre la propria idea al gruppo, ``e l'esperienza è stata
        pessima per tutti''. Quindi: \textbf{non imponete le vostre idee, fate
        brainstorming con gli altri}. Daniele, che fa ``molto spesso
        brainstorming'', darà strumenti per farlo: nessuno arriverà impreparato
        a quella lezione.
\end{itemize}

Ruggero aggiunge che il mandate è \textbf{abbastanza flessibile} perché tutti
possano fare qualcosa che li interessi: ``l'ultima cosa che vogliamo è forzarvi a
lavorare su qualcosa che non vi piace''. Bisogna lavorare su ciò che piace e per
cui si ha uno stimolo personale, anche nella \textbf{tipologia} di gioco: nessuno
sarà obbligato a fare un gioco di ruolo o un board game, e la lezione sui casi
studio mostrerà quante forme diverse di gioco analogico il corso accetta.

L'unica cosa da \textbf{non fare}, ``perché è contraddittorio per questo corso'',
sono i \textbf{giochi grandi}: ``lo chiediamo per voi''. Niente progetti
mastodontici\footnote{La trascrizione dice ``non fate delle cose da gnome'':
probabilmente il nome di un gioco da tavolo molto grande e complesso, non
ricostruibile con certezza.}, niente ``cose alla \emph{Dungeons \& Dragons}''.
Si può fare un gioco di ruolo, ma con un'esperienza ``meno corposa'': ci sono
tanti giochi di ruolo da 15 o 20 pagine, ``bellissimi, fichissimi'', che si
possono leggere online. Il consiglio è farsi una cultura anche su
\textbf{itch.io}, ``uno dei portali più importanti del gioco di ruolo
internazionale'', di cui si parlerà anche nella parte di ricerca e sviluppo.

Poi Daniele saluta (``noi ci vediamo venerdì a lezione'') e lascia la
conduzione a Ruggero.

% ---------------------------------------------------------------------
\section{Le domande finali con Ruggero}
\label{a1:sec:domande}

\subsection{Le regole delle domande}

Ruggero spiega che cosa vuol dire ``stare con lui''. Alla fine di ogni sua
lezione ci saranno \textbf{almeno 20--30 minuti di domande} libere, \emph{free
form}. Vuole che fare domande sia \textbf{normalizzato}: ``io gestisco malissimo
l'\emph{horror vacui} e divento super cringe molto alla svelta se non ci sono
domande in giro''. Sono benvenute le domande sulla lezione appena fatta, ma non
solo: il corso è crossmediale, e anche le domande che vengono da
\textbf{divagazioni}, da parti laterali, ``fatele''. Quello che sembra scontato o
banale a uno studente può non esserlo per il docente, e viceversa; e può essere
utile anche agli altri. ``Se vi sembrano stupide, fatele il doppio''. ``Magari
sembro incazzato, ma vi giuro sono rilassatissimo. Non mordo, magari abbaio un
po'\,''.

\subsection{Come approcciarsi a Steam e Tabletop Simulator}

Rita non ha mai usato né Steam né Tabletop Simulator e chiede come
approcciarsi, e se c'è modo di mettersi d'accordo tra compagni su Discord, magari
con un sondaggio sui giochi, per formare gruppi di gioco. Ruggero conferma: sul
Discord dell'accademia c'è un \textbf{canale apposito per il corso}, ed è degli
studenti. Lo possono usare per organizzarsi a giocare, per chiedere a un
compagno di prestare un gioco, per domande tecniche (come si scarica, come si
installa). ``Sarebbe proprio ideale che vi parlaste e vi organizzaste''.

Anche perché, aggiunge con un ``piccolo spoiler'' (``nessuno vi ha detto
niente''), se fossero i docenti a organizzare i gruppi di gioco, il task lo
starebbero facendo loro. Bisogna \textbf{abituarsi a organizzare le sessioni di
gioco}, perché arriveranno i task sul \textbf{playtest} e su come organizzarlo,
e a quel punto si sarà ``meno disperati''. Sembra scontato, ma saper
\textbf{organizzarsi per giocare in autonomia} ``vi sarà utilissimo come skill in
questo settore, ve lo giuro''.

\subsection{Il bilanciamento}

Angelo confessa ``una paranoia grossissima'': ci saranno lezioni sul
\textbf{bilanciamento}? ``Ma certo''. Sono proprio le lezioni di Marta sul
development: il bilanciamento è di solito un tema che tocca \textbf{i developer
più dei designer}. L'anno precedente c'era una sola lezione con Marta, e i
docenti si sono resi conto che serviva andare ancora più a fondo: per questo
quest'anno le lezioni sono due. ``Non è una paranoia, è una legittima domanda'':
Ruggero insegna da anni, vede altri corsi, parla con altri docenti e studenti, e
``mi incazzo spesso'' perché molti corsi di game design \textbf{non fanno
bilanciamento}. I docenti volevano invece un corso che coprisse tutte le parti
del game design: non solo l'idea del gioco, ma il bilanciamento, la produzione,
i playtest, l'iterazione, la scrittura dei regolamenti.

\subsection{Si può giocare con i docenti?}

Una studentessa chiede se si può giocare con i docenti. Per motivi etici,
durante il corso questo è limitato alle parti di \textbf{valutazione} dei giochi
e di \textbf{tutoraggio}. Dopo il corso, ``nessun problema, almeno con me''.

\subsection{Dove si trovano i giochi su Tabletop Simulator}

Ruggero mostra dove si trovano i giochi: nella propria libreria di Steam, nella
pagina di Tabletop Simulator, alla voce \textbf{Workshop}, dove ci sono ``più o
meno 36.000'' giochi. Una studentessa scopre intanto che l'app di Steam per
telefono è solo una \emph{companion app}: per registrarsi e giocare serve il
PC.

Per provare, chiede agli studenti un titolo:
\begin{itemize}
  \item \emph{Keep the Heroes Out}: non compare, ``non tutti i giochi ci sono,
        perché molti publisher non li mettono''; forse, essendoci già una
        versione del gioco su Steam, il modulo è stato tolto;
  \item \emph{Magic: The Gathering}: ``almeno 50 moduli'', anzi ``123 pagine'';
  \item \emph{Bang!}: tre moduli.
\end{itemize}
Bisogna fare attenzione: alcuni moduli sono in lingue specifiche (``magari il
primo è in russo''), alcuni hanno già l'espansione, altri sono solo il gioco
base. Conviene sfogliare e non fermarsi al primo risultato. E nel Workshop si
trovano anche cose molto particolari: giochi fuori produzione, giochi vecchi,
giochi ancora in prototipo.

\subsection{I ruoli nei team e le dinamiche tra le persone}

Uno studente torna sulla slide delle figure: nelle realtà piccole molti ruoli
hanno ``il doppio o triplo cappello''. Quali ruoli restano distinti, e come si
creano le sinergie tra persone con obiettivi diversi?

Ruggero premette che l'ideale ``platonico'' è un ambiente di lavoro costruttivo,
basato sulla fiducia e ``il più \emph{safe} possibile''; ma nell'industria non
sempre è così, e le strutture variano moltissimo: alcuni publisher sono molto
verticali, altri hanno ruoli già solidificati (se un publisher ha già i suoi
developer e non sta crescendo, entrarci come developer è praticamente
impossibile). ``Se prendiamo due aziende in Italia, vediamo due strutture
interne completamente diverse''. Non esiste una scorciatoia per farsi assumere
da un editore.

Poi prende ``l'esempio più comune di tutti'', la media dei giochi da tavolo
pubblicati (figura~\ref{a1:fig:relazioni}):
\begin{itemize}
  \item il \textbf{game designer} è una figura \textbf{esterna}, con un
        \textbf{contratto a royalty} con il publisher;
  \item la sua figura di riferimento è il \textbf{publisher} stesso, e la
        conversazione designer--publisher è ``molto calda'';
  \item se il publisher ha dei \textbf{developer}, il developer entra in questa
        conversazione: designer e developer collaborano molto, e internamente
        developer e publisher collaborano molto. Spesso, nei publisher piccoli,
        è il publisher a fare le veci del developer;
  \item \textbf{graphic design} e \textbf{illustrazione} tendono a parlare molto
        con il publisher e \textbf{molto poco con il designer}. A meno che non si
        tratti di uno studio indie molto piccolo, è difficile che il designer
        sappia anche solo chi è l'illustratore prima che il gioco esca.
\end{itemize}

\begin{figure}[htbp]
\centering
\begin{tikzpicture}[
    font=\small,
    ruolo/.style={draw, thick, rounded corners=3pt, fill=gray!12,
                  minimum width=2.8cm, minimum height=0.9cm, align=center},
    esterno/.style={draw, thick, dashed, rounded corners=3pt,
                    minimum width=2.8cm, minimum height=0.9cm, align=center},
    calda/.style={{Stealth[length=2.5mm]}-{Stealth[length=2.5mm]}, very thick},
    tiepida/.style={{Stealth[length=2mm]}-{Stealth[length=2mm]}, thin, dashed}
  ]
  \node[esterno] (gd) at (0,0) {Game designer\\ \footnotesize esterno, a royalty};
  \node[ruolo] (dev) at (5,1.4) {Developer\\ \footnotesize se c'è};
  \node[ruolo] (pub) at (5,-1.4) {Publisher};
  \node[ruolo] (art) at (10,-1.4) {Graphic designer\\ e illustratore};
  \draw[calda] (gd) -- (dev);
  \draw[calda] (gd) -- (pub);
  \draw[calda] (dev) -- (pub);
  \draw[calda] (pub) -- (art);
  \draw[tiepida] (gd.south) to[bend right=25]
     node[below, font=\footnotesize\itshape] {quasi nessun contatto} (art.south);
  \node[draw, dotted, rounded corners=4pt, fit=(dev)(pub)(art), inner sep=8pt,
        label={[font=\footnotesize\itshape]above right:interno al publisher (di solito)}] {};
\end{tikzpicture}
\caption{Le relazioni tipiche intorno a un gioco da tavolo pubblicato, secondo la
risposta di Ruggero. Le frecce spesse indicano una collaborazione stretta e
frequente; quella tratteggiata, un contatto quasi assente. Nei publisher piccoli
il publisher fa spesso anche da developer.}
\label{a1:fig:relazioni}
\end{figure}

\subsection{Fare da developer di sé stessi?}

Rita chiede in chat se un designer può fare anche da developer del proprio
gioco, o se Ruggero lo sconsiglia. In generale \textbf{lo sconsiglia}, per lo
stesso motivo per cui si sconsiglia in altri media: i \textbf{darling}. Il
designer ha delle meccaniche a cui è affezionato, e il developer a volte dovrà
\textbf{tagliarle}; è molto difficile, se si è ``gestori di sé stessi'', tagliare
cose che si vogliono nel gioco. ``È un po' come lasciare fare il lavoro di editing
allo scrittore che ha già scritto il libro'': può rivedersi, certo, ma difficilmente
cambierà cose che farebbe bene cambiare. Nelle lezioni con Marta si vedrà perché,
e alla fine del corso ``non avrete dubbi nemmeno voi'' nel separare designer e
developer.

Ma tutto dipende dallo \textbf{scopo del progetto}: se si fa un gioco in tre,
qualcuno quel lavoro deve farlo per forza, e magari il designer fa anche il
developer, compensando con \textbf{più playtest}. ``Il mondo attuale e il mondo
ideale sono due cose molto diverse, e lo vediamo non solo dai giochi ma dalla
politica'', ma ``non entriamo in questo vaso di Pandora''.

\subsection{Il ruolo indispensabile e la pubblicazione}

Lo studente di prima ne trae una conclusione: senza game designer il gioco non
esiste. ``Quello sicuramente''. Ma Ruggero coglie l'occasione per approfondire
un punto che forse non era passato: \textbf{fare design di un gioco per
pubblicarlo non è l'unica opzione possibile}. Un gioco non è ``meno gioco''
se non è pubblicato da un editore. È legittimo mettere un gioco su itch.io in
\textbf{print and play}, fare un gioco a evento senza venderlo, fare un
\emph{murder party}, una escape room, un board game per una persona amica:
``quell'esperienza lì non è meno gioco''. Dipende dall'obiettivo del progetto.

Ruggero ha lavorato molto anche con le \textbf{scuole}, dove il publisher non
esiste: è un gioco per le classi, e la figura più vicina è l'editoria scolastica,
``che però non fa il publisher''. Ci sono compromessi diversi sulla materialità
e su come il gioco si presenta, che si vedranno nella prototipazione e nei casi
studio. Quindi l'unica figura che deve esserci sempre è \textbf{chi fa game
design}, anche se magari non sa di chiamarsi game designer e non ha mai
sentito parlare di game design: ``questo non la rende meno game designer''.

\paragraph{Un libro consigliato.} Ruggero fa anche mentoring e gestisce gruppi
di settore, e spesso gli chiedono che cosa leggere per imparare il game design.
Non è il tipo che obbliga a leggere libri (``non vi obbligherò a leggere
niente''), ma ce n'è uno che consiglia sempre, ed è il primo che ha scelto di
leggere: \textbf{\emph{Rise of the Videogame Zinesters}} di \textbf{Anna
Anthropy}\footnote{La trascrizione dice ``Anna Antrofi'' e ``Rise of the
Videogames in Esters''.}. È un libro un po' datato, ma ancora reperibile; parla
soprattutto di videogiochi, ma ``per l'80\%'' il contenuto si applica al gioco e
al game design in generale. È un libro sulla \textbf{democratizzazione del gioco}:
il gioco come strumento per parlare di sé e come \textbf{medium legittimo}. Come
esistono film drammatici, tristi, intensi, possono esistere giochi su tematiche
delicate, e se ne vedranno nei casi studio. Molti giochi fatti con questa
direttiva non hanno un publisher, che tende a pubblicare ciò che dà un ritorno
commerciale.

\paragraph{Commerciale e personale.} ``Io spero che voi non farete giochi solo
per scopo commerciale''. Certo, servono un tetto sulla testa, il cibo in
tavola, un posto in cui dormire, e così si finisce per fare giochi che non sono
il proprio gioco dei sogni. Spesso è un mix: Ruggero ha fatto giochi commerciali
``per pagarmi l'affitto'' e giochi che voleva fare e su cui ha perso soldi,
``sapendolo proprio che sarebbero stati un fallimento''. I giochi, come i libri e
le canzoni, si fanno per molti motivi, e quelli commerciali di successo non sono
più importanti degli altri. Tra i giochi della sua slide c'erano il gioco da
tavolo ufficiale di \emph{Cyberpunk}\footnote{La trascrizione dice ``Cyberpunk
del 1977'': quasi certamente un gioco legato a \emph{Cyberpunk 2077}.} e
\emph{Marvel United}, su IP grosse e famose, ma anche giochi fatti di ``una
candela e una busta su un tavolo''. E se gli si chiede quale preferisce in
termini di contenuto, non sono quelli sui supereroi.

\subsection{Il gioco di ruolo in Italia}

Nicola, che è interessato soprattutto al gioco di ruolo e conosce già persone
del settore, ha ``una leggera paura'': che chi studia game design finisca per
fare il \textbf{freelance} a vita, ``perché non ci sarà mai un'occupazione per
bene'', o il \textbf{copywriter}. Conosce un ragazzo che ha lavorato in quel
mondo e che ora fa altro.

Ruggero risponde che la maggioranza dei giochi nel suo portfolio sono giochi di
ruolo, quindi il tema lo tocca da vicino, e che in Italia quello che dice Nicola
``è semplicemente vero''. Molti publisher piccoli hanno una persona di
riferimento che copre più ruoli: impagina, fa i playtest, fa un po' di
development. È ``un settore economicamente completamente disastrato'', per
motivi culturali che non apre: oggi in Italia è praticamente impossibile
\textbf{vivere di giochi di ruolo}. Chi ci lavora lo fa nel tempo libero, al di
fuori del lavoro principale, oppure coprendo molte altre cose: molti publisher
italiani di giochi di ruolo pubblicano anche titoli esteri e localizzano, non
fanno solo cose originali. ``È un grosso calderone''. Poi si ferma: ``questo è
un \emph{rant}'', e ``i \emph{rant} tornano sempre'', ma non vuole che gli
studenti escano dalla prima lezione ``con la disperazione nel cuore''. Ci sarà
spazio per parlarne quando avranno gli strumenti per capirlo, ``per non farvi
abbattere''.

\subsection{Che cosa bisogna sapere già}

Uno studente chiede se ci sono nozioni che dovrebbero già avere e che possono
recuperare: per esempio, sente parlare di giochi ``German'' o ``American'' e non
sa che cosa siano. Ruggero spiega l'organizzazione del corso:
\begin{itemize}
  \item le \textbf{prime due lezioni} sono di fondamenta. Questa era ``dall'alto'':
        come è strutturato il corso, le tempistiche;
  \item la \textbf{prossima lezione} è di tre ore di soli \textbf{casi studio},
        ``a manetta, a raffica'', divisi per categorie e tipologie, con
        ragionamenti su categorie e generi di gioco: lì si vedranno German,
        American ``e anche altre più utili'';
  \item nel resto del corso le terminologie specifiche sono inserite nelle
        lezioni dedicate: il \textbf{paragame} nelle lezioni sul paragame, il
        \textbf{modello matematico} nelle lezioni di Marta sul development.
\end{itemize}
Per ora agli studenti si chiede di sapere quello che si è visto oggi:
meccaniche, dinamiche ed estetiche; le figure professionali, designer e
developer; e soprattutto \textbf{game e play}.

\subsection{Game e play: i quattro quadranti}

``Specialmente \emph{game} e \emph{play}'', insiste Ruggero: è più difficile di
quello che sembra, perché l'italiano non abitua a separare le due cose, ed è
possibile averle usate per anni come sinonimi. Nella sua esperienza con chi si
avvicina al design è ``uno dei punti più critici''. Spesso i commenti che si
fanno sui giochi giocati sono commenti sul \emph{play} applicati al
\emph{game}: ``il gioco non mi è piaciuto perché, giocando con Tizio, è successo
questo''; ``il gioco mi sembra sbilanciato perché Tizio al tavolo ha fatto
questo''.

Separare \emph{game} e \emph{play} serve a capire quali giochi sono progettati
bene e quali male, e soprattutto a riconoscere i \textbf{quattro possibili
incroci} (figura~\ref{a1:fig:quadranti}): un gioco progettato bene che piace
giocare; uno progettato male che non piace; uno progettato male che piace; uno
progettato bene che non piace. ``Non è importantissimo in generale nella vita,
ci sono cose più importanti'', ma per chi fa game design lo è: ``a volte vi
chiederanno di lavorare su un gioco che non vi piace giocare, e quel gioco deve
essere fatto bene lo stesso''.

\begin{figure}[htbp]
\centering
\begin{tikzpicture}[
    font=\small,
    cella/.style={draw, thick, minimum width=4.4cm, minimum height=1.9cm,
                  text width=4.0cm, align=center}
  ]
  \node[cella, fill=gray!20] (a) at (0,0)     {\textbf{progettato bene}\\ e \textbf{mi piace} giocarlo};
  \node[cella]               (b) at (4.4,0)   {\textbf{progettato bene}\\ ma \textbf{non mi piace} giocarlo};
  \node[cella]               (c) at (0,-1.9)  {\textbf{progettato male}\\ ma \textbf{mi piace} giocarlo};
  \node[cella, fill=gray!8]  (d) at (4.4,-1.9){\textbf{progettato male}\\ e \textbf{non mi piace} giocarlo};
  \node[font=\footnotesize\itshape, above=0.15cm of a] {mi piace (\emph{play})};
  \node[font=\footnotesize\itshape, above=0.15cm of b] {non mi piace (\emph{play})};
  \node[font=\footnotesize\itshape, rotate=90, left=0.15cm of a, anchor=south]
       {bene (\emph{game})};
  \node[font=\footnotesize\itshape, rotate=90, left=0.15cm of c, anchor=south]
       {male (\emph{game})};
\end{tikzpicture}
\caption{I quattro incroci tra la qualità del \emph{game} (come è progettato) e
l'esperienza di \emph{play} (se piace giocarlo), descritti da Ruggero. Per un
designer i quadranti ``scomodi'' sono i due fuori diagonale: un gioco può
piacere pur essendo progettato male, e va saputo progettare bene anche un gioco
che non piace.}
\label{a1:fig:quadranti}
\end{figure}

\paragraph{La domanda di Tancredi.} Tancredi chiede conferma: il \emph{play}
non è il ``giocare senza regole'', ma può essere anche il giocare a un gioco? Sì:
il \emph{play} è l'\textbf{atto}, e può essere fatto con un \emph{game} o senza.
E un \emph{game} può esistere senza che nessuno lo giochi: gli scacchi non sono
meno scacchi perché in quel momento nessuno ci gioca. E se, giocando, una
persona applica le regole in modo diverso? Allora \textbf{non si sta giocando
quel gioco}.

\subsection{House rule: il caso di Monopoly}

Qui Ruggero fa un po' di \emph{foreshadowing} su una lezione di Daniele: il
discorso degli \textbf{hack}, delle \textbf{house rule}, delle \textbf{mod},
cioè di come modificare un gioco esistente. Il punto è che \textbf{un gioco
modificato da una house rule è un altro gioco}: cambiando le regole cambiano
necessariamente anche le \textbf{dinamiche} che ne derivano, e quindi
l'esperienza. Si immagini \emph{Monopoly} in cui ogni turno si possa, volendo,
prendere gratis dei soldi dalla banca: tutto il resto è identico, eppure il gioco
cambia.

E questo è talmente interiorizzato che ci sono regole che si applicano ``senza
neanche rendercene conto''. Ruggero chiede quanto duri in media una partita a
\emph{Monopoly}. Gianmarco: ``quattro ore''; altri: ``un pomeriggio''; un altro:
``finché uno dei giocatori non aggiunge la regola: ora vince chi\ldots''. In
realtà, \textbf{seguendo le regole scritte}, una partita non dura mai più di
\textbf{un'ora e venti, un'ora e mezza}. La domanda che fa sempre, e la risposta
è sempre la stessa: perché si pensa che \emph{Monopoly} duri tantissimo? Qual è
la regola che nessuno rispetta?

Uno studente indovina: \textbf{le aste}. Secondo il regolamento, se chi finisce
su un terreno non lo compra, il terreno \textbf{va all'asta} tra i giocatori, e
chiunque può aggiudicarselo. Molti invece giocano come se finire su un terreno
desse solo la \emph{facoltà} di comprarlo o di metterlo all'asta\footnote{Più
precisamente, nel regolamento ufficiale, se il giocatore che finisce su una
proprietà libera decide di non comprarla, la banca la mette all'asta e tutti i
giocatori possono fare offerte, compreso chi l'ha rifiutata. L'asta quindi non è
una scelta del giocatore, ma la conseguenza obbligata del suo rifiuto.}. Con la
regola applicata, di aste se ne fanno moltissime: all'ottavo turno ci sono già
otto terreni assegnati, al ventesimo venti, la gente va in bancarotta
velocissimamente e la partita dura un'ora e mezza. Anche le partite competitive
di \emph{Monopoly}, nota Ruggero, durano un'ora e venti, un'ora e mezza.

Uno studente obietta che giocava con le aste e durava comunque quattro ore.
``Avete dei nemici con l'analisi paralisi'', ride Ruggero; e cita un'altra house
rule tipica dell'infanzia: genitori che danno soldi per non far andare in
bancarotta. Ma \textbf{la bancarotta è parte del gioco}: se la si toglie,
\emph{Monopoly} non regge. Sul tema che \emph{Monopoly} vuole trattare, e sul
suo antenato (\emph{The Landlord's Game}), ``non apriamo adesso un altro vaso di
Pandora''.

\paragraph{Approfondimento: Monopoly, la durata e le house rule nei libri.}
L'esempio di Ruggero si ritrova quasi identico in \emph{The Art of Game Design}.
Parlando del bilanciamento tra partite brevi e lunghe (\emph{Short vs. Long}),
Schell scrive che \emph{Monopoly}, giocato secondo le regole ufficiali,
``finisce spesso in circa novanta minuti'', ma molti giocatori trovano quelle
regole troppo dure e le modificano, distribuendo premi in denaro e allentando i
vincoli sull'acquisto delle proprietà: con l'effetto collaterale di far durare la
partita molto di più, ``tipicamente tre ore, o anche di più''. Schell classifica
le \textbf{house rule} come regole create dai giocatori in risposta a un
difetto percepito dopo qualche partita. Anche \emph{Challenges for Game
Designers}, nel capitolo sull'aggiungere e togliere meccaniche, parte dallo
stesso gioco: il \emph{Parcheggio gratuito}, che per le regole originali non fa
succedere nulla, mentre quasi tutti giocano con una house rule che dà un bonus a
chi ci finisce sopra. Le due fonti confermano la lezione: una piccola regola
cambiata cambia le dinamiche, e le dinamiche cambiano l'esperienza.

\subsection{Il fine del game è il play?}

Un altro studente commenta: va bene la distinzione, ma il fine del \emph{game} è
il \emph{play}; quello che si vede è sempre l'effetto. Ruggero è d'accordo: di
solito, se si fa un gioco, lo si fa per farlo giocare, come si scrive un libro
per farlo leggere (anche se esistono, ed ``è già stato esplorato'', libri scritti
per non essere letti e giochi per zero giocatori). Ma la distinzione serve a una
cosa precisa: \textbf{separare, nel playtest, i feedback sul gioco dai feedback
sul gruppo di gioco} e sulle sue dinamiche. Se una persona esce frustratissima da
un playtest, non è necessariamente colpa del gioco; se una persona dice ``questo
è il gioco migliore del mondo'', potrebbe non essere merito del gioco ma del
gruppo. Capire a che cosa si riferisce un feedback serve a capire \textbf{dove
iterare} il prototipo.

Lo studente osserva che l'area è sfumata: alcuni giochi si basano proprio
sull'avere le persone giuste al tavolo. Un gioco di \textbf{negoziazione} con
giocatori abituati ai giochi ``multi-solitario'' diventa ingiocabile; e nei giochi
a \textbf{economia chiusa} (cita \emph{Container}\footnote{La trascrizione cita
anche un secondo titolo non ricostruibile (``Tokyo su Vigimarket'').}) un gruppo
economicamente conservativo, che spende poco e non fa girare l'economia, fa
durare la partita il doppio. Ruggero rilancia: e se il gioco dura il doppio,
questa dinamica cambia la \textbf{percezione di peso} del gioco, le energie spese,
quanto si è stanchi a fine partita. ``Tutto è collegato. \textbf{I giochi sono
ecosistemi di meccaniche}, e cambiandone una, le altre cambiano di
conseguenza''.

\paragraph{La parola vietata, alla prova.} Lo studente aggiunge che ci sono giochi
che ha trovato ``molto sbagliati, ma che mi divertono''. ``Attenzione: non si può
dire \emph{divertimento}!'', lo ferma Ruggero; nelle presentazioni aveva lasciato
correre, ``ma lì non vi avevo ancora letto la regola''. Lo studente ci prova:
``giochi che poi voglio \textbf{rigiocare}'', anche se ``sono evidentemente
sbagliati, rotti''. Ruggero apprezza, ma non vuole anticipare cose che meritano il
loro tempo: se ne parlerà nelle lezioni sui framework e soprattutto in quella di
\textbf{emotion design}, ``forse la mia lezione preferita del corso''. (Lo
studente ammette che ``forse alla fine si gioca per quello, per l'emozione''.)

\subsection{I giochi ``comunicati male''}

L'ultima domanda corposa riprende un'affermazione di inizio lezione: ci sono
giochi belli ma comunicati male. Qualche esempio? ``Mi basta girarmi per
vederne 48'', risponde Ruggero indicando gli scaffali alle sue spalle. Più che
titoli, elenca le \textbf{ragioni}:
\begin{itemize}
  \item \textbf{l'Italia localizza lentamente}. Basta guardare i premi del gioco
        dell'anno in Italia: spesso vanno a giochi usciti due o tre anni prima in
        Germania, Francia, Stati Uniti o Regno Unito, perché localizzare richiede i
        suoi tempi. Il mercato italiano è quindi molto dipendente dall'estero,
        ``molto impaurito'', e spesso aspetta che un gioco vada bene all'estero
        prima di portarlo in Italia; e alcuni giochi ``perdono il treno'';
  \item \textbf{il publisher originale chiude}: come si pubblica un gioco se non
        c'è più l'azienda con cui fare un contratto? A Ruggero è successo più
        volte;
  \item \textbf{il dipartimento commerciale italiano non sa comunicare il gioco}:
        si fa male un'edizione, la stampa vende pochissimo e non viene più
        ristampata;
  \item \textbf{la scala del mercato}: in Italia un gioco da tavolo di peso medio
        (un \emph{eurogame}, un \emph{German}) che vende 1.500 copie ``è tanto'';
        all'estero se ne vendono magari 10.000. Con 10.000 copie in inglese ci
        sono i soldi per una campagna di comunicazione: un trailer, i
        recensori, i tutorial online. In Italia tutto questo è ``ridotto all'osso'',
        e la cosa si riflette su stampa e divulgazione: anche i canali per
        informarsi su un gioco prima di comprarlo sono pochi (se ne parlerà nella
        lezione su ricerca e sviluppo). Un gioco spesso non arriva al mercato non
        perché sia brutto, ma perché l'editore non lo sa comunicare, ne fa poche
        copie e va fuori produzione;
  \item \textbf{il tempismo}: party game progettati nel 2019 e usciti nel 2020, in
        pieno Covid, quando non si poteva giocare al tavolo, sono falliti tutti, e
        gli editori ci hanno perso. ``Non potevano immaginarselo'';
  \item \textbf{i temi caldi}: dopo l'invasione russa dell'Ucraina e la guerra in
        Palestina, alcuni editori non pubblicano o non localizzano giochi
        strategici e di guerra. Ruggero ha lavorato a \textbf{\emph{Status}}, un
        gioco sulla geopolitica che tocca, tra gli altri temi, l'invasione russa
        dell'Ucraina e il genocidio in Palestina, e alcuni editori che di
        solito localizzano i loro giochi (per esempio il partner di riferimento in
        Russia) non lo hanno localizzato per motivi politici.
\end{itemize}
Lo studente voleva capire che cosa voglia dire comunicare male: non metterci
abbastanza energia, o comunicare una cosa per un'altra, ``come il trailer
sbagliato, che poi vai al cinema e ti trovi un altro film''. ``È un mix'', e la
domanda ``era complessa e meritava più spazio''.

\subsection{I giochi di carte collezionabili e i saluti}

L'ultima domanda (``posso rompere le balle un'ultima volta?'') riguarda i
\textbf{giochi di carte collezionabili}: funzionano con le stesse regole dei
giochi da tavolo? Se ne parlerà alla prossima lezione: è proprio il motivo per
cui la lezione sui casi studio dura tre ore, perché ogni genere di gioco ha le
sue dinamiche e le sue regole, e anche il modo di pubblicare cambia molto tra un
TCG, un gioco di ruolo da tavolo, un LARP e un gioco pervasivo.

Ruggero saluta: le email dei tre docenti sono all'inizio delle slide, la
registrazione arriverà, ``ci si vede questo venerdì''.

% ---------------------------------------------------------------------
\section*{Punti chiave del capitolo}
\addcontentsline{toc}{section}{Punti chiave del capitolo}

\begin{itemize}
  \item Il corso di \textbf{Analog Game Design} (seconda edizione) dura circa tre
        mesi, fino a dicembre, ed è tenuto da tre docenti principali:
        \textbf{Daniele} (narrative design, gioco di ruolo, fiction),
        \textbf{Ruggero} (game design, ricerca e sviluppo) e \textbf{Beatrice}
        (giochi pervasivi), con tre ospiti su RNG, development e game production.
  \item Lo scopo è formare \textbf{autori analogici} consapevoli: strumenti da
        usare anche fuori dal corso, non il ``gioco dei sogni''. Il corso è
        \textbf{intersezionale} (nessuna lezione dedicata a un solo tipo di
        gioco), \textbf{cross-generazionale}, \textbf{pratico e adattabile},
        basato sulle \textbf{domande} (nessuna domanda è stupida), sul
        \textbf{feedback} invece che sul voto e sul \emph{fail forward}.
  \item Il \textbf{project work} si fa in tre gruppi da tre, formati dai docenti;
        pitch il 12 ottobre, sviluppo nelle due pause, tutoraggi il sabato,
        consegna il 4 dicembre, presentazione l'11 e giornata finale con gli
        editori il 14. Si valutano documentazione, presentazione,
        autovalutazione, crescita e atteggiamento.
  \item \textbf{Niente intelligenza artificiale} nel corso: i publisher
        rifiutano pitch e prototipi con immagini generate, i regolamenti scritti
        dall'IA descrivono la stessa meccanica in modi diversi, e soprattutto chi
        la usa perde l'occasione di imparare. Un prototipo può essere fatto a
        matita.
  \item Le \textbf{figure dell'industria}: game designer, game developer,
        editor, publisher, sales e marketing, graphic designer, illustratore,
        sensitivity reader (che non è un censore) e producer. Nelle realtà
        piccole una persona copre più ruoli.
  \item Nell'analogico il \textbf{game designer} progetta l'esperienza creando
        le meccaniche; il \textbf{game developer} fa da \textbf{mediatore} con il
        publisher (bilanciamento, modelli matematici, adattamenti). Pretendere
        che il proprio gioco non venga modificato è un \emph{red flag}.
  \item Il \textbf{marketing} non decide che cosa pubblicare ma come; è
        un'azione, ed esiste anche se lo si ignora.
  \item In \textbf{Italia} il gioco è ancora visto come infantile, il mercato
        punta sulla localizzazione e investe poco su giochi e designer interni;
        all'estero (soprattutto dalla Germania in su) ci sono tradizione adulta,
        premi di prestigio e developer interni.
  \item \textbf{Game} è l'artefatto (le regole, il sistema, l'oggetto);
        \textbf{play} è l'atto di giocare. Si fanno giochi, non giocattoli.
        Separarli serve a distinguere i feedback sul gioco da quelli sul gruppo
        e a riconoscere i quattro incroci tra ``progettato bene/male'' e ``mi
        piace/non mi piace''.
  \item \textbf{MDA}: meccaniche (regole scritte), dinamiche (regole in azione),
        estetiche (esperienza, non grafica); nel \textbf{DPE}: design, play,
        experience.
  \item \textbf{Parole vietate} in ogni materiale scritto: \emph{fun},
        \emph{immersion}, \emph{entertainment} (divertimento, immersione,
        intrattenimento). Bisogna trovare parole più precise.
  \item Distinguere \textbf{chi crea, chi gioca e chi facilita}; nessun gioco
        nasce perfetto: \textbf{playtest e iterazione} sono inevitabili, e un
        \textbf{dizionario comune} aiuta a non perdersi.
  \item \textbf{Tabletop Simulator} è il software di settore per giocare (non
        creare videogiochi) i giochi da tavolo in digitale. Primo task: giocarci
        tre o quattro volte entro il 10 ottobre.
  \item Una \textbf{house rule} crea un altro gioco: \emph{Monopoly} con le aste
        applicate come da regolamento dura circa un'ora e mezza.
\end{itemize}
